\documentclass[11pt,a4paper]{article}
\usepackage[T1]{fontenc}
\usepackage[utf8]{inputenc}
\usepackage{lmodern}
\usepackage[margin=28mm]{geometry}
\usepackage{amsmath,amssymb,amsthm,mathtools,bm}
\usepackage{microtype,booktabs,longtable,enumitem,xcolor,needspace}
\usepackage{listings}
\usepackage[hidelinks]{hyperref}
\hypersetup{pdftitle={Microscopic jamming in the negative spherical perceptron},pdfauthor={OpenAI},pdfsubject={Critical exponents and stability under Gaussian scale mixtures},pdfkeywords={spherical perceptron, jamming, critical exponents, universality, computer-assisted proof}}
\usepackage[nameinlink,noabbrev]{cleveref}
\allowdisplaybreaks[2]
\setlength{\emergencystretch}{1em}
\numberwithin{equation}{section}
\newtheorem{theorem}{Theorem}[section]
\newtheorem{lemma}[theorem]{Lemma}
\newtheorem{proposition}[theorem]{Proposition}
\newtheorem{corollary}[theorem]{Corollary}
\theoremstyle{definition}
\newtheorem{definition}[theorem]{Definition}
\newtheorem{assumption}[theorem]{Assumption}
\theoremstyle{remark}
\newtheorem{remark}[theorem]{Remark}
\newcommand{\E}{\mathbb E}
\newcommand{\Prob}{\mathbb P}
\newcommand{\R}{\mathbb R}
\newcommand{\N}{\mathbb N}
\newcommand{\dd}{\mathrm d}
\newcommand{\ind}{\mathbf 1}
\newcommand{\cN}{\mathcal N}
\newcommand{\eps}{\varepsilon}
\providecommand{\mathscr}[1]{\mathcal{#1}}
\DeclareMathOperator*{\esssup}{ess\,sup}
\DeclareMathOperator{\supp}{supp}
\DeclareMathOperator{\Var}{Var}
\DeclareMathOperator{\Cov}{Cov}
\DeclareMathOperator{\Lip}{Lip}
\DeclareMathOperator{\diag}{diag}
\DeclareMathOperator*{\argmin}{arg\,min}
\DeclareMathOperator*{\argmax}{arg\,max}
\DeclareMathOperator{\Tr}{Tr}
\lstset{language=Python,basicstyle=\ttfamily\scriptsize,breaklines=true,breakatwhitespace=false,breakautoindent=true,breakindent=1.5em,postbreak=\mbox{\textcolor{gray}{$\hookrightarrow$}\space},columns=fullflexible,keepspaces=true,showstringspaces=false,numbers=left,numberstyle=\tiny\color{gray},numbersep=5pt,frame=single,rulecolor=\color{gray!40},tabsize=2,literate={²}{{\textsuperscript{2}}}1}
\title{Microscopic jamming in the negative spherical perceptron}
\author{OpenAI}
\date{September 24, 2026}
\begin{document}
\maketitle
\begin{abstract}
We prove a sharp feasibility threshold and limiting gap and force laws for
the spherical perceptron with margin $-1$ and quadratic penalty. The limits
are taken successively in system size, inverse temperature, and density
approaching the threshold from above. They agree for Gaussian coordinates
and for the equal mixture of centered Gaussian coordinates with variances
$1-\eps$ and $1+\eps$, for every sufficiently small fixed $\eps$.
The contact-removed gap cumulative law and the mean-one force cumulative
law satisfy
\[
 G_J(u)=u^{1-\gamma+o(1)},\qquad F_J(s)=s^{1+\theta+o(1)},
\]
as $u\downarrow0$ and $s\downarrow0$, with $\gamma=(2+\theta)^{-1}$,
$0.4126930<\gamma<0.4126934$, and
$0.4231063<\theta<0.4231088$.
A finite numerical certificate for these exponent intervals, together with
its mathematical error bounds, is included.
\end{abstract}

\section{The problem and the proof architecture}\label{sec:introduction}

\subsection{Background and contribution}
Jamming links the loss of feasibility in a system of random constraints to
singular distributions of small positive gaps and weak forces.  The
negative-margin spherical perceptron is a particularly simple setting in
which this connection can be studied: its variables lie on a sphere, its
constraints are linear inequalities, and its feasible set is nonconvex.
Franz and Parisi~\cite{FranzParisi2016} identified this model with the
jamming universality class of high-dimensional spheres.  Their analysis
predicted isostaticity, the balance between contact constraints and degrees
of freedom, together with the gap and force exponents.  Their equations
use a continuum of replica-overlap levels, called full replica-symmetry
breaking.

The scaling mechanism and numerical values have an earlier origin in the
full replica-symmetry-breaking solution for hard spheres developed by
Charbonneau, Kurchan, Parisi, Urbani and Zamponi~\cite{CKPUZ2014}.
Franz, Parisi, Sevelev, Urbani and Zamponi~\cite{FPSUZ2017} subsequently
analyzed the spherical perceptron on both sides of jamming.  In particular,
they derived the zero-temperature scaling from the unsatisfiable side,
the contact atom and normalized force distribution, and the relation
between the critical exponents by matching the scaling regimes.
Parisi and Zamponi~\cite{ParisiZamponi2026} gave a direct analytic proof
of that relation within the scaling equations.  Their argument assumes the existence of the full replica-symmetry-breaking
profile and the convergence and matching of its scaling expansion.
Constructing and selecting the microscopic limit is a separate problem.

The relation between gap and force exponents also has a mechanical
stability antecedent. Wyart~\cite{Wyart2012} derived a stability inequality
for sphere packings and predicted its saturation at marginal stability.
Lerner, D\"uring and Wyart~\cite{LernerDuringWyart2013} distinguished the
extended rearrangements governed by that inequality from localized
buckling, which obeys a different bound. The force exponent in the
full replica-symmetry-breaking hard-sphere solution is the extended-mode
exponent~\cite{CKPUZ2014}. These mechanical arguments provide a conceptual
comparison; the present proof derives its endpoint mass estimates from
the microscopic equilibrium model.

The feasibility question also belongs to the older spherical
storage-capacity problem, which asks how many random linear constraints
admit a vector on a sphere.  Gardner~\cite{Gardner1988} and Gardner and
Derrida~\cite{GardnerDerrida1988} developed replica calculations for
binary-pattern neural-network models; the latter identified instability
of the replica-symmetric continuation at negative margins.  For Gaussian
constraints at a fixed negative margin, Stojnic~\cite[Section 2.3]{Stojnic2013}
proved a one-sided infeasibility bound.  Montanari, Zhong and
Zhou~\cite[Theorems 4.1--4.2]{MontanariZhongZhou2024} subsequently proved
feasible and infeasible density regions whose bounds agree to leading
order as the margin tends to $-\infty$.  These bounds do not establish a
sharp threshold at margin $-1$ or determine the microscopic gap and force
laws studied here.

Here we establish the corresponding equilibrium statement directly for
the random spherical model with quadratic penalty.  The result includes
a sharp feasibility threshold, the existence and identification of the
ordered limiting gap and force laws, and their invariance under a fixed
small independent Gaussian scale mixture in each pattern coordinate.
It also identifies the critical susceptibility power and encloses the exponents
with a finite certificate.  The exponent values, relation, normalized
force convention, and approach from the unsatisfiable side are thus
consistent with the preceding physical predictions.  The mathematical
content is the passage from the microscopic model through these limits,
including the selection and error bounds needed for the stated
conclusions.

There are three linked difficulties.  First, the spherical constraint
prevents a direct convex-feasibility argument, while the quadratic penalty
is unbounded.  Second, vanishing energy alone does not imply exact
satisfiability, and a weak critical limit need not preserve contact mass
or force normalization.  Third, writing down a stationary scaling profile
does not show that the microscopic critical layer approaches it.  The
proof addresses these issues through a variational pressure formula and
a separate feasibility repair, a convex Brownian martingale problem
with endpoint control, and a contraction for all possible rescaled
critical limits.

The pressure proof uses the interpolation scheme of
Guerra~\cite{Guerra2003} and the block-cavity organization of
Aizenman, Sims and Starr~\cite{AizenmanSimsStarr2003}, including its
spherical implementation by Chen~\cite{Chen2013SphericalASS}.
For the present row Hamiltonian, concavity of the row functional supplies
the interpolation sign, while the lower bound requires a separate row
expansion and identification of the empirical cavity covariance.
Coordinate comparison uses the smooth-function form of Lindeberg
replacement~\cite{Chatterjee2006Lindeberg}; controlling spins concentrated
on a few coordinates requires the additional entropy argument described
below.

The hierarchical change-of-measure and cavity viewpoint follows the
Ruelle probability-cascade framework of Bolthausen and
Sznitman~\cite{BolthausenSznitman1998}.  The Brownian formulation is
related to the variational representation of Bou\'e and
Dupuis~\cite{BoueDupuis1998} and the stochastic-control treatment of the
Parisi equation by Auffinger and Chen~\cite{AuffingerChen2015}.
The passage from stochastic controls to a variational problem over Brownian
martingales has a close precedent in Mourrat's dual formulation of the
Parisi functional~\cite{MourratUninverting2025}.  The martingale
representation and uniqueness results of Chen, Issa and
Mourrat~\cite{ChenIssaMourrat2025} further develop this approach for
multi-species spin glasses.  These works provide methodological precedents
for the convex martingale problem used here.  The present formulation
requires a concave quadratic terminal penalty, nonnegative martingales
with square-integrable terminal values, and control of a singular critical
endpoint.  The arguments below establish the corresponding identities,
critical uniqueness, and scaling selection for the microscopic limits
constructed here.

\subsection{Model and observables}
Fix $\alpha>0$, let $M=\lfloor\alpha N\rfloor$, and let $A$ be an $M\times N$ matrix with independent entries. In empirical observables containing $M^{-1}$, take $N$ large enough that $M\ge1$. Initially the entries are standard normal. On
\[
 \mathbb S_N=\{x\in\R^N:\|x\|^2=N\}
\]
use normalized surface measure $\sigma_N$. For $y=x/\sqrt N$, define
\begin{equation}\label{eq:intro-model}
 g_\mu(x)=A_\mu y,\qquad h_\mu(x)=1+g_\mu(x),\qquad
 V(z)=\frac12(-1-z)_+^2,\qquad H_N(x)=\sum_{\mu=1}^M V(g_\mu(x)).
\end{equation}
Here $z_+=\max(z,0)$ and $z_-=\max(-z,0)$. Since $\|y\|=1$, the condition $H_N(x)=0$ is exactly $Ay\ge-\mathbf 1$ componentwise, the spherical storage-capacity problem at margin $\kappa=-1$. The realization is satisfiable (SAT) if $H_N$ vanishes somewhere on $\mathbb S_N$ and is unsatisfiable (UNSAT) otherwise. At inverse temperature $\beta>0$, write $\langle\cdot\rangle_{N,\beta,\alpha}$ for integration against
\[
 \frac{e^{-\beta H_N(x)}\,\sigma_N(\dd x)}
 {\int e^{-\beta H_N(x)}\,\sigma_N(\dd x)}.
\]
Disorder expectation is $\E$; subscripts on the Gibbs brackets are omitted when unambiguous.

For $0<\delta<u$, put
\begin{equation}\label{eq:intro-gap-observable}
 G_{N,\beta,\alpha}(\delta,u)
 =\E\left\langle\frac1M\sum_{\mu=1}^M
       \ind\{\delta<h_\mu(x)\le u\}\right\rangle.
\end{equation}
Forces are defined on the violated constraints
$C(x)=\{\mu:h_\mu(x)<0\}$. When this set is nonempty, let
\[
 \overline r(x)=\frac1{|C(x)|}\sum_{\nu\in C(x)}(-h_\nu(x)),
 \qquad f_\mu(x)=\frac{-h_\mu(x)}{\overline r(x)}\quad(\mu\in C(x)),
\]
and set
\begin{equation}\label{eq:intro-force-observable}
 F_{N,\beta,\alpha}(s)=
 \E\left\langle\frac1{|C(x)|}\sum_{\mu\in C(x)}
             \ind\{f_\mu(x)\le s\}\right\rangle,\qquad s>0.
\end{equation}
The integrand is defined to be zero when $C(x)=\varnothing$.

All limits at jamming are taken from UNSAT in the order
\begin{equation}\label{eq:intro-order}
 \lim_{\alpha\downarrow\alpha_c}\ \lim_{\beta\to\infty}\ \lim_{N\to\infty}.
\end{equation}
For gaps the lower cutoff is removed subsequently:
\begin{align}
 G_J(u)&=\lim_{\delta\downarrow0}
          \lim_{\alpha\downarrow\alpha_c}
          \lim_{\beta\to\infty}\lim_{N\to\infty}
          G_{N,\beta,\alpha}(\delta,u),\label{eq:intro-GJ}\\
 F_J(s)&=\lim_{\alpha\downarrow\alpha_c}
          \lim_{\beta\to\infty}\lim_{N\to\infty}
          F_{N,\beta,\alpha}(s).\label{eq:intro-FJ}
\end{align}
The limits $u\downarrow0$ and $s\downarrow0$ are taken last. This convention distinguishes positive gaps from contacts without assuming a limiting density.

The perturbations considered here have coordinate law
\begin{equation}\label{eq:intro-mixture}
 \nu_\eps=\tfrac12\cN(0,1-\eps)+\tfrac12\cN(0,1+\eps),\qquad 0\le\eps<1,
\end{equation}
where the second argument is the variance. Coordinates and rows remain independent.
This law is centered with variance one; for $\eps>0$ it is non-Gaussian,
since its fourth moment is $3(1+\eps^2)$.

\subsection{Main result}
\begin{theorem}[Microscopic criticality and mixture universality]\label{thm:main}
There exist $\alpha_c\in(0,\infty)$ and $\eps_0\in(0,1)$ with the following properties for every fixed $\eps\in[0,\eps_0)$ and coordinate law $\nu_\eps$.
\begin{enumerate}[label=\textup{(\roman*)}]
\item The SAT probability tends to one at every fixed $\alpha<\alpha_c$ and to zero at every fixed $\alpha>\alpha_c$.
\item All the ordered limits in \eqref{eq:intro-GJ}--\eqref{eq:intro-FJ} exist for their indicated positive arguments. They are independent of $\eps$.
\item $F_J$ is the cumulative distribution function on $(0,\infty)$ of a probability law of mean one. It has no atom at zero. There are exponents $\gamma,\theta$ such that
\[
 G_J(u)=u^{1-\gamma+o(1)},\qquad
 F_J(s)=s^{1+\theta+o(1)},
\]
as the corresponding argument decreases to zero, and
\[
 \gamma=\frac1{2+\theta},\qquad
 0.4126930<\gamma<0.4126934,\qquad
 0.4231063<\theta<0.4231088.
\]
\end{enumerate}
\end{theorem}

The limiting full gap law has a contact atom of mass $1/\alpha_c$
(Lemma~\ref{crit:contacts}). Multiplying this contact fraction by the
constraint density gives one contact per degree of freedom in the ordered
critical equilibrium limit. This is the isostaticity predicted for the
model; it also identifies the conditioning mass used to normalize forces.

\subsection{The endpoint mechanism and the spectral characterization}
\label{subsec:intro-endpoint}
The two exponents arise from different fluctuation scales with the same
small probability mass. The proof constructs a positive, nonincreasing
susceptibility $\chi:[0,1)\to(0,1]$, normalized by $\chi(0)=1$ and
vanishing at the critical endpoint $t=1$. Here $t$ is the variance time
of the Brownian representation of the critical limit, rather than a
microscopic dynamical time. Put
\[
 \tau=1-t,\qquad w(t)=\int_t^1\chi(v)^{-2}\,\dd v.
\]
The remaining Brownian scales for gaps and residual forces are,
respectively, $\sqrt\tau$ and $\sqrt{w(t)}$. Section~\ref{sec:scales}
proves that $w(t)\asymp\tau/\chi(t)^2$ and that the strictly positive
masses up to these two scales are both comparable to $\chi(t)$.
Equivalently, for a fixed $c_f>0$ accounting for mean force normalization,
\[
 G_J(\sqrt\tau)\asymp\chi(t),\qquad
 F_J(c_f\sqrt{w(t)})\asymp\chi(t).
\]
Here $\asymp$ denotes bounds in both directions by fixed positive
constants, uniformly as $t\uparrow1$. These statements concern cumulative
masses and require no assumption of a limiting density.

If $\chi(1-\tau)=\tau^{a+o(1)}$ with $0<a<1/2$, the gap and force
cutoffs have logarithmic powers $1/2$ and $(1-2a)/2$, while their masses
have power $a$. Their cumulative powers are therefore $2a$ and
$2a/(1-2a)$. The remaining task is to prove that the susceptibility has
such a power and to select $a$; a stationary scaling ansatz alone would
not establish either assertion.

The selected value is characterized by a scalar boundary-value and
spectral problem. For $a\in[0,0.41]$, Section~\ref{sec:contraction} constructs the unique stationary reference $M_a:\R\to\R$ in the class $-1\le M_a'\le0$, $M_a''>0$, satisfying
\begin{equation}\label{eq:intro-reference}
 \frac12M_a''+\left(\frac z2+aM_a\right)M_a'
                  +\left(a-\frac12\right)M_a=0,
\end{equation}
with the bounds
\[
 \phi(z)-z\Phi(-z)\le M_a(z)\le\frac{\phi(z)}{\Phi(z)}.
\]
Here $\phi$ and $\Phi$ are the standard normal density and distribution function. Put $B_a(z)=z/2+aM_a(z)$ and let $\lambda(a)$ be the lowest eigenvalue of the Friedrichs realization on $L^2(\R)$ of
\begin{equation}\label{eq:intro-spectral}
 \mathcal H_a=-\frac12\frac{\dd^2}{\dd z^2}
                  +\frac12\bigl(B_a^2+B_a'\bigr).
\end{equation}
There is a unique root $a_*$ of $\lambda(a)=a$ in
\begin{equation}\label{eq:intro-root-bracket}
 0.2936533<a_*<0.2936535.
\end{equation}
The exponents in Theorem~\ref{thm:main} are
\begin{equation}\label{eq:intro-exponents}
 \gamma=1-2a_*,\qquad
 \theta=\frac{2a_*}{1-2a_*}-1.
\end{equation}
The root enclosure and the contraction selecting it are proved with the finite certificate in Section~\ref{sec:certificate} and Appendix~\ref{app:code}.

\subsection{Architecture and conventions}
Sections~\ref{sec:hierarchy}--\ref{sec:row-concavity} construct the
hierarchical functionals and the concavity estimates used in the Gaussian
pressure formula, where pressure means the expected logarithm of the
partition function divided by $N$.  The hierarchy describes the possible
laws of the overlap $N^{-1}x\cdot x^{\prime}$ between replicas, meaning
independent samples from a fixed-disorder Gibbs measure.
The row functional records how adding one constraint changes the pressure;
its concavity supplies the sign needed in the interpolation and cavity
bounds of Section~\ref{sec:gaussian-pressure}.

At fixed positive temperature, small row perturbations turn pressure
information into an empirical row law.  Adding $d\psi$ to the single-row exponent $-\beta V$
couples the pressure to the empirical average of a smooth compactly
supported test $\psi$.  An upper bound that touches the pressure at $d=0$
for both signs of $d$ pins down this average; a determining family of tests
then identifies the row law.  Section~\ref{sec:universality} compares that
law with the coordinate-mixture model, including configurations whose mass
is concentrated on a few coordinates.  On a fixed density neighborhood of
the threshold, the same sufficiently small fixed mixture is admissible at
every fixed positive temperature, so the resulting identical functions
pass through the prescribed later limits.  For
feasibility, Section~\ref{sec:threshold} starts from subextensive minimum
energy at a slightly higher density, deletes a fixed fraction of the rows,
and repairs the remaining constraints exactly at the target density.

Sections~\ref{sec:zero-temperature} and~\ref{sec:critical} pass to zero
temperature and then to the threshold.  A susceptibility profile measures
the response along the hierarchy; a nonnegative Brownian martingale
represents the residual force.  Joint convexity and a uniqueness argument
on one Wiener space identify the critical profile.  Strong endpoint
convergence and small-cutoff estimates preserve the contact mass and
normalized force law.

The final two stages address the critical singularity.  Section~\ref{sec:scales}
relates small-gap and small-force masses to the susceptibility and proves
that its associated increasing order parameter has no constant interval
sufficiently near the endpoint.  Section~\ref{sec:contraction} takes limits after translating
logarithmic endpoint time.  Every resulting solution is defined for all
translated times and satisfies the same integral identities.  Quantitative
barriers first confine these solutions; a strict contraction then forces
their susceptibility slope to be the selected scalar root.  Section~\ref{sec:certificate}
certifies the finite inequalities used in that step.  The principal outputs
and their precise locations are:
\begin{center}
\small
\begin{tabular}{@{}p{.69\textwidth}p{.27\textwidth}@{}}
\toprule
Output & Proof location \\
\midrule
Gaussian variational pressure and touching upper bounds for both signs of a
compact row perturbation
& Proposition~\ref{prop:press-formula}; \eqref{eq:src11} \\
Equality of the thermodynamic row laws, and of the gap and normalized-force
observables, at each fixed positive temperature for sufficiently small mixtures
& Proposition~\ref{prop:univ-observables} \\
Exact satisfiability below the common threshold, obtained by repairing a
system with subextensive energy
& Theorem~\ref{thresh:sharp} \\
A unique normalized physical critical limit, strong convergence through the endpoint,
and the prescribed ordered gap and force laws
& Proposition~\ref{crit:unique-critical}, Proposition~\ref{crit:strong},
Theorem~\ref{crit:ordered-laws} \\
Two-sided endpoint mass estimates and full increase support near the endpoint
& Proposition~\ref{prop:scale-endpoint-masses};
Proposition~\ref{prop:scale-full-support} \\
The selected susceptibility slope, from entire-limit rigidity and the finite
certificate
& Proposition~\ref{prop:contr-rigidity}; Section~\ref{sec:certificate} \\
The cumulative exponents and their exact relation
& Section~\ref{sec:conclusion} \\
\bottomrule
\end{tabular}
\end{center}

The mixture argument identifies the fixed-temperature functions;
Corollary~\ref{thresh:ordered-limits} transfers each subsequent Gaussian
limit in the prescribed order. Sections~\ref{sec:contraction}
and~\ref{sec:certificate} form one computer-assisted argument: the former
proves the analytic implications of the quantitative inequalities,
including the required tail bounds, and the latter verifies their finite
enclosures. The short completion in Section~\ref{sec:conclusion} then
combines the endpoint mass estimates with the selected slope.

We use $C,c$ for finite positive constants whose permitted dependencies are stated locally. The notation $f\asymp g$ means two-sided bounds by such constants, uniform in the limit under discussion. Profiles are identified up to equality almost everywhere. A reference to a support concerns the support of the associated increasing measure, not the set where its cumulative distribution is positive.

The analytic prerequisites are Gaussian concentration and integration by parts, Brownian martingale representation, Brownian time change and reflection, It\^o calculus, elementary parabolic comparison and interior regularity, the Dovbysh--Sudakov representation~\cite{Panchenko2010DS}, the support-radius consequence of the Ghirlanda--Guerra identities~\cite{Panchenko2010Sphere}, Panchenko's ultrametricity theorem~\cite{Panchenko2013}, and the Marchenko--Pastur limit~\cite{MarcenkoPastur1967}. Their precise uses are specified where they enter. All model-specific compactness, comparison, limiting, and numerical arguments are supplied below.


\tableofcontents
\section{Hierarchical Gaussian functionals and spherical entropy}
\label{sec:hierarchy}

The pressure calculation separates the contribution of a new constraint
from that of new spin coordinates.  We construct the corresponding row
and spherical functionals on a common Gaussian hierarchy.  For the row
functional we identify its derivative with respect to the covariance
profile.  For the spherical functional we compute its limit and the
entropy conjugate $E(q)$ used in the pressure formula.

The common hierarchy is a random family of leaf weights, together with
Gaussian marks shared by leaves having a common ancestor.  We first prove
its change-of-measure and sampling rules.  These rules identify which
marks remain independent after a row is inserted.  The continuity lemma
then passes these conclusions through limits of replica arrays under
explicit moment bounds.  Throughout, a profile specifies only
off-diagonal covariances; the diagonal is a separate parameter.  This
distinction determines the independent terminal Gaussian noise in a row
insertion.

The finite cascade and its reweighting are the objects of
Bolthausen--Sznitman's abstract cavity method~\cite{BolthausenSznitman1998}.
We give the marking calculation explicitly, including the genealogy and
independent terminal noise used in the row insertions below.

\subsection{Finite cascades and their change of measure}

Let $\mathcal Q_Q$ be the set of nonnegative nondecreasing functions
$q:[0,1)\to[0,Q]$, identified up to Lebesgue-null sets.  Write
$q_*:=\esssup q$.  For a step profile, fix
\[
 0=m_{-1}<m_0<\cdots<m_k=1,\qquad
 q(s)=q_j\quad(m_{j-1}\le s<m_j),
 \qquad w_j:=m_j-m_{j-1}.
\]
Repeated values $q_j$ are allowed.  A cascade of depth $k$ has leaves
$\tau=(\tau_1,\ldots,\tau_k)\in\N^k$.  At every vertex of depth $j-1$, independently,
place a Poisson point process on $(0,\infty)$ with intensity
$m_{j-1}s^{-1-m_{j-1}}\,\dd s$; its points are the weights of the outgoing edges.
Let $W_\tau$ be the product of the edge weights on the path to $\tau$, and set
\[
 T=\sum_\tau W_\tau,\qquad v_\tau=W_\tau/T.
\]
A depth-zero cascade consists of one leaf of weight one.  The following
facts also prove that the normalizing sum is finite and positive.

\begin{lemma}[Stable sums and the cascade change of measure]
\label{lem:hier-cascade-tilt}
For a fixed finite cascade, $0<T<\infty$ almost surely and
$\E|\log T|^2<\infty$.  Attach marks to its vertices, independently of the
Poisson processes, with conditionally independent and identically distributed
child marks given the ancestral marks.  A terminal function $X_\tau$ of the
path marks determines recursive functions by
\begin{equation}
 X_v=\frac1{m_{j-1}}\log
       \E\bigl[\exp(m_{j-1}X_{vi})\mid\text{marks up to }v\bigr]
 \quad\text{at depth }j-1.
 \label{eq:hier-mark-recursion}
\end{equation}
Suppose these moments are finite and $\E|X_\varnothing|<\infty$.
Reweighting the leaf probabilities by $e^{X_\tau}$ preserves the law of their
genealogy.  Conditional on the transformed genealogy, the child-mark
transition at depth $j$ is multiplied by
\begin{equation}
 \exp\{m_{j-1}(X_{vi}-X_v)\}.
 \label{eq:hier-child-tilt}
\end{equation}
Root marks retain their original distribution.  Moreover,
\begin{equation}
 \E\log\sum_\tau v_\tau e^{X_\tau}=\E X_\varnothing.
 \label{eq:hier-root-expectation}
\end{equation}
If the marks are independent families and $X_\tau$ is the sum of one function
of each family, both the recursion and the tilted path law factor over the
families, including conditional on the genealogy of finitely many replicas.
\end{lemma}

\begin{proof}
For $0<m<1$, a Poisson sum $S=\sum_i s_iY_i$, with intensity
$m s^{-1-m}\,\dd s$ and iid nonnegative marks satisfying
$0<\E Y^m<\infty$, has Laplace transform
\[
 \E e^{-tS}
 =\exp\{-\Gamma(1-m)\E(Y^m)t^m\}.
\]
Indeed, the Poisson Laplace functional reduces the exponent to
$m\E\int_0^\infty(1-e^{-tsY})s^{-1-m}\,\dd s$; integration by parts evaluates
this integral.  The resulting positive stable random variable has finite
positive moments of every order less than $m$ and finite negative moments
of every order.  The latter follows directly from
$x^{-a}=\Gamma(a)^{-1}\int_0^\infty t^{a-1}e^{-tx}\,\dd t$; the former follows
by applying the corresponding integral formula to $x^a$, $0<a<m$.
Consequently $\E|\log S|^2<\infty$.  Starting at the last generation and
using $m_0<\cdots<m_{k-1}$ proves these assertions for every subtree total,
and hence for $T$.

Here is the change of measure at one vertex.  Write $\mu(\dd\omega)$ for
the conditional law of a child mark and put $c(\omega)=e^{X_{vi}-X_v}$.
Under the map $(s,\omega)\mapsto(sc(\omega),\omega)$, the marked intensity
becomes
\[
 m_{j-1}s^{-1-m_{j-1}}\,\dd s\,
 c(\omega)^{m_{j-1}}\mu(\dd\omega).
\]
The second factor is a probability measure by
\eqref{eq:hier-mark-recursion}.  Thus the transformed weights have their
original Poisson law and are independent of marks with transition
\eqref{eq:hier-child-tilt}.

We spell out why this yields independence of the \emph{entire} transformed
weight array from the marked tree.  At depth two, write the root mark as
$B$, a child mark as $A_i$, and its descendant marks as $C_{ij}$.  Let $s_i$
and $t_{ij}$ be the edge processes of exponents $m_0$ and $m_1$.  Set
\[
 X_1(B,A)=m_1^{-1}\log\E_C e^{m_1X_2(B,A,C)},\qquad
 X_0(B)=m_0^{-1}\log\E_A e^{m_0X_1(B,A)}.
\]
The lower transformation
$t'_{ij}=t_{ij}e^{X_2(B,A_i,C_{ij})-X_1(B,A_i)}$ has, conditional on
$(B,A_i)$, the original unmarked Poisson law, denoted $\mathcal P_{m_1}$,
independently of child marks with tilted transition $K_1(\dd C\mid B,A_i)$.
In particular this law of the complete unmarked lower process is
independent of $A_i$ and $B$.  Before the upper transformation a decorated
child therefore has law
\[
 \mu_0(\dd A\mid B)\,\mathcal P_{m_1}(\dd t')
              \prod_jK_1(\dd C_j\mid B,A).
\]
The upper map $s'_i=s_i e^{X_1(B,A_i)-X_0(B)}$ replaces only $\mu_0$ by
its tilted transition $K_0$; it leaves $\mathcal P_{m_1}$ unchanged.
The marked-Poisson theorem makes the ordered $s'_i$ independent of these
iid decorated children.  Thus, after both reorderings, all the transformed
weights are independent of the hierarchical marked tree, conditional on
$B$.  The same argument inducts upward for any depth: replace
$\mathcal P_{m_1}$ by the already transformed descendant weight-array law,
which is independent of its ancestor mark.  A sampled genealogy depends
only on the complete weight array and independent sampling uniforms.  This
includes the descendant total masses used in selecting a child.
Conditioning on that genealogy therefore does not tilt its marks again.
Replicas sharing a vertex reuse the one mark already drawn there; no
additional power depending on the number of replicas is introduced.

Reordering preserves a bijection of leaves.  Under that correspondence,
products of the edge multipliers telescope:
\[
 W'_\tau=W_\tau e^{X_\tau-X_\varnothing},\qquad
 \log\sum_\tau v_\tau e^{X_\tau}
   =X_\varnothing+\log T'-\log T.
\]
Conditionally on root marks, $T'$ and $T$ have the same marginal law; they
need not be independent.  Their integrable logarithms therefore cancel in
expectation.  The root factor cancels in normalized leaf probabilities, so
the root marks are not tilted.  Finally, exponential moments of a sum of
independent families factor at each recursion step, and so do the
Radon--Nikodym factors in \eqref{eq:hier-child-tilt}.  This proves the last
assertion.
\end{proof}

For clarity, the genealogy used here can be constructed without appealing
to any unproved identification of an overlap array.  We record its sampling
rule.

\begin{lemma}[Sampling a cascade and the continuous rank array]
\label{lem:hier-rank-array}
For two independent samples from a depth-$k$ cascade, the probability that
their leaves share exactly $j$ edges is $w_j$, $0\le j\le k$; sharing $k$
edges includes selecting the same leaf.  There exists a weakly exchangeable
ultrametric array $(U_{ab})_{a,b\ge1}$ with $U_{aa}=1$ and $U_{12}$ uniform
on $[0,1]$, such that every step-profile cascade covariance array has law
\[
 C_{aa}=Q,\qquad C_{ab}=q(U_{ab})\quad(a\ne b).
\]
All the finite-dimensional laws of this rank array are determined by its
nested partition sampling rule.
\end{lemma}

\begin{proof}
We give the sampling computation, including the change caused by an
ancestor.  Let a Poisson process have intensity
$c\beta s^{-1-\beta}\,\dd s$, $0<\beta<1$, and total $S$.  Change its law by
$S^\eta/\E S^\eta$, where $0\le\eta<\beta$.  For a specified partition of
$n$ sampled indices into $b$ blocks of sizes $n_1,\ldots,n_b$, normalized
Poisson weights give probability
\begin{equation}
 \frac{\prod_{a=1}^{b-1}(a\beta-\eta)}
      {\prod_{a=1}^{n-1}(a-\eta)}
 \prod_{i=1}^b\prod_{a=1}^{n_i-1}(a-\beta).
 \label{eq:hier-eppf}
\end{equation}
Empty products are one.  To verify this, insert
\[
 S^{\eta-n}=\frac1{\Gamma(n-\eta)}
       \int_0^\infty t^{n-\eta-1}e^{-tS}\,\dd t
\]
into the expectation of the sum over distinct Poisson points.  The Poisson
factorial-moment formula gives one factor
$c\beta\Gamma(n_i-\beta)t^{\beta-n_i}$ for each selected point and the
Laplace factor $\exp\{-c\Gamma(1-\beta)t^\beta\}$.  The remaining integral,
after $v=c\Gamma(1-\beta)t^\beta$, is a gamma integral.  Dividing by
\[
 \E S^\eta=
 [c\Gamma(1-\beta)]^{\eta/\beta}
 \frac{\Gamma(1-\eta/\beta)}{\Gamma(1-\eta)}
\]
(with value one at $\eta=0$) gives \eqref{eq:hier-eppf}.
Its ratios show that an existing block of size $l$ receives the next sample
with probability $(l-\beta)/(n-\eta)$; a new block has probability
$(b\beta-\eta)/(n-\eta)$.

At a cascade vertex, multiply each outgoing edge by the total mass of its
child subtree.  The Poisson mapping calculation in the preceding proof
makes the resulting child masses a stable Poisson process and tilts the
child subtree law by its total mass raised to the parent parameter.
More explicitly, if the outgoing parameter is $\beta$, the masses
$r_i=s_iT_i$ are independent of child-subtree marks whose law is multiplied
by $T_i^\beta/\E T_i^\beta$.  An incoming bias
$(\sum_i r_i)^\eta$ changes only the mass process, so those child marks
remain independent, also after conditioning on the choices of children.
Thus \eqref{eq:hier-eppf} applies at the first level with
$(\eta,\beta)=(0,m_0)$ and, inside a selected parent at later levels, with
$(\eta,\beta)=(m_{j-2},m_{j-1})$.  For two samples the conditional joining
probabilities are
\[
 1-m_0,\quad \frac{1-m_1}{1-m_0},\quad\ldots,
 \quad\frac{1-m_{k-1}}{1-m_{k-2}}.
\]
Their products give probability $1-m_{j-1}$ of sharing the first $j$
edges.  Taking differences proves the first assertion.

More generally, multiply the conditional joining probabilities along the
ancestral path to a cluster containing $l$ of the first $n$ samples at a
level of parameter $m$.  The denominators telescope, giving
$(l-m)/n$.  These probabilities, and differences between probabilities of
joining a parent and joining its existing children, determine all ways
the next sample can enter a finite nested partition.  They also show that
deleting intermediate levels leaves the same rule at the retained levels.
Construct nested partitions at all dyadic parameters by these consistent
finite-dimensional laws, taking the universal partition at parameter zero
and the singleton partition at parameter one, and put
$U_{ab}=\sup\{m:a,b\text{ belong to the same cluster at level }m\}$.
The rule gives $\Prob(U_{12}\ge m)=1-m$ at dyadic $m$, hence uniformly on
$[0,1]$.  Nesting gives ultrametricity.  The predictive probabilities for a
finite nested partition are differences of $(l-m)/n$, and hence its
probability is continuous in the finitely many level parameters.
Approximating arbitrary real levels by dyadic levels, and using that no
$U_{ab}$ has an atom at a prescribed level, therefore recovers the same
sampling law at every real level.  Applying a step function $q$ now recovers
the finite cascade covariance law.  Set $q(1)=q_*$; nonnegative
nondecreasing transforms of this nested array are positive semidefinite:
for a step transform they are sums, with nonnegative coefficients, of
block-indicator matrices, and the general case follows by limits.  An
additional nonnegative diagonal correction supplies any $Q\ge q_*$.
\end{proof}

\subsection{The row functional and tilted derivatives}

Attach independent centered Gaussian increments of variances
$q_0,q_1-q_0,\ldots,q_k-q_{k-1}$ to the common root and subsequent vertices,
and denote their path sum by $z_\tau$.  Initially suppose $q_k=Q$ and that
$u\in C^2(\R)$ has bounded first and second derivatives.  Define
\[
 \Phi_u(q;Q)=\E\log\sum_\tau v_\tau e^{u(z_\tau)}.
\]
When $q_k<Q$, replace $e^{u(z)}$ at a leaf by
$\E_G e^{u(z+\sqrt{Q-q_k}\,G)}$, where $G\sim\cN(0,1)$.
This is the \emph{terminal convolution}; a residual Gaussian is drawn
independently for each sampled replica, even if two replicas select the
same leaf.  We omit $Q$ when the diagonal has been fixed.

For a finite profile put
\[
 f_k(z)=\log\E_G e^{u(z+\sqrt{Q-q_k}\,G)},\qquad
 f_{j-1}(z)=\frac1{m_{j-1}}\log\E_G
 e^{m_{j-1}f_j(z+\sqrt{q_j-q_{j-1}}\,G)}.
\]
The terminal formula means $f_k=u$ when $q_k=Q$.  Lemma
\ref{lem:hier-cascade-tilt} gives
$\Phi_u(q;Q)=\E_G f_0(\sqrt{q_0}\,G)$.
Let $Z_j$ be the partial path sum under the tilted transitions
\eqref{eq:hier-child-tilt}, with the original Gaussian root law, and put
$A_j=f'_j(Z_j)$.  Differentiation under the Gaussian integrals gives
\[
 f'_{j-1}(z)=
 \E^{\mathrm{tilt}}[f'_j(Z_j)\mid Z_{j-1}=z].
\]
Thus $(A_j)$ is a martingale.  At the terminal convolution there is one
further tilted transition, of parameter one, from $f'_k$ to $u'$.

To pass from finite cascades to arbitrary covariance profiles, we use the
following continuity principle.  Its hypotheses concern replica arrays,
so no convergence of random measures in a chosen representation is needed.

\begin{lemma}[Replica-array continuity under a Gaussian tilt]
\label{lem:hier-array-continuity}
Let $G_n$ be random probability measures.  Suppose that the jointly sampled
arrays of a finite list of uniformly bounded positive-semidefinite kernels,
and any other bounded observables being retained, converge in law on every
finite set of replicas.  Conditional on these arrays, adjoin independent
centered Gaussian fields with the prescribed kernels.  Let $X_n$ be the
evaluation of one fixed continuous function on the retained single-replica
variables and those fields, and assume, for every fixed $a>0$,
\begin{equation}
 \sup_n\E G_n(e^{a|X_n|})<\infty.
 \label{eq:hier-exponential-moments}
\end{equation}
Then the expectations of $\log G_n(e^{X_n})$, and the expectations of bounded
continuous functions of finitely many replicas sampled from
$e^{X_n}G_n/G_n(e^{X_n})$, converge to the corresponding functionals of the
limiting array.  Independent auxiliary variables of a fixed law may be
adjoined to each replica.  The conclusion also holds for unbounded
observables whose truncation errors vanish uniformly under the tilted law;
in particular this follows for polynomial functions of the Gaussian fields
under \eqref{eq:hier-exponential-moments}.
\end{lemma}

\begin{proof}
For a fixed number of replicas, the conditional characteristic function of
the adjoined fields is $\exp(-t^{\mathsf T}Ct/2)$.  It is a bounded continuous
function of the covariance matrix, including at singular matrices.
Consequently the fields converge jointly with the arrays.  The same
argument with a product law adjoins the independent auxiliary variables.

Put $W=e^X$, $W_L=(W\vee L^{-1})\wedge L$, $Z=G(W)$, and $Z_L=G(W_L)$.
For every fixed $r\ge1$, Jensen's inequality gives
\[
 Z^{-r}\le G(e^{-rX}),\qquad
 Z_L^{-r}\le G(W_L^{-r})\le G(1+e^{-rX}).
\]
All these inverse moments are uniformly bounded in expectation.  For
arbitrary fixed $r\ge1$ and $A>0$, \eqref{eq:hier-exponential-moments} and another
application of Jensen give
\[
 \sup_n\E|Z-Z_L|^r
 \le\sup_n\E G(|W-W_L|^r)=O(L^{-A}),
\]
after using a sufficiently large moment in
\eqref{eq:hier-exponential-moments}.  The elementary inequality
$|\log a-\log b|\le |a-b|(a^{-1}+b^{-1})$ and H\"older's inequality show that
replacing $Z$ by $Z_L$ changes its expected logarithm by a quantity tending
to zero uniformly in $n$.

At fixed $L$, uniformly approximate $\log z$ on $[L^{-1},L]$ by polynomials.
Every expected polynomial of $G(W_L)$ is a finite linear combination of
replica expectations, which converge by the first paragraph.  This proves
the assertion for logarithms.  For an $r$-replica observable $B$, its tilted
expectation is
\[
 \E\frac{G_n^{\otimes r}(B\prod_{i=1}^rW_i)}{Z^r}.
\]
Replace the $W_i$ and $Z$ by their clipped versions.  H\"older's inequality,
the inverse-moment bounds, and the displayed tail estimate make this error
uniformly small for bounded $B$.  At fixed $L$, approximate $z^{-r}$ on
$[L^{-1},L]$ by polynomials and again use finitely many replicas.  A further
truncation proves the unbounded version whenever the stated uniform
integrability holds.  Polynomial Gaussian observables have moments of all
orders uniformly, because the covariance kernels are bounded; H\"older's
inequality with \eqref{eq:hier-exponential-moments} and the inverse-moment
bound then supplies this uniform integrability.
\end{proof}

\begin{lemma}[Derivative profile and Gaussian interpolation]
\label{lem:hier-row-interpolation}
For a step profile define
$D_q(s)=\E A_j^2$ on $[m_{j-1},m_j)$.  This is nonnegative and
nondecreasing, satisfies $D_q\le\|u'\|_\infty^2$, and is constant across
adjacent equal values of $q$.  The diagonal quantity
$D_{q,d}:=\E^{\mathrm{tilt}}u'(Z)^2$ is at least the final off-diagonal
value.  For two profiles with common diagonal $Q$,
\begin{equation}
 \frac{\dd}{\dd t}\Phi_u((1-t)q+tr;Q)
 =-\frac12\int_0^1(r-q)(s)D_{(1-t)q+tr}(s)\,\dd s.
 \label{eq:src1}
\end{equation}
For arbitrary $q,r\in\mathcal Q_Q$ this identity holds for $0<t<1$ and
as the appropriate one-sided derivative at the endpoints.  Moreover,
\begin{equation}
 |\Phi_u(q;Q)-\Phi_u(r;Q)|
 \le\frac12\|u'\|_\infty^2\|q-r\|_1,
 \qquad q\longmapsto D_q\text{ is continuous into }L^1[0,1].
 \label{eq:hier-row-continuity}
\end{equation}
For a general profile, $D_q$ is constant almost everywhere on every
interval on which $q$ is constant almost everywhere.
\end{lemma}

\begin{proof}
We first prove the finite-profile identity, then use replica-array
continuity to construct its extension and pass the derivative formula to
the limit.
The square of a bounded martingale is a submartingale.  A zero-variance
increment has identical incoming and outgoing derivative, proving all the
monotonicity and plateau assertions.  Given that two replicas share exactly
$j$ levels, their later marks are independent under the tilted transitions.
Their terminal derivatives therefore have product expectation $\E A_j^2$.
The shared path itself has the single-path tilted law, because transformed
weights and marks are independent, also conditional on the genealogy.

Put $q$ and $r$ on a common finite partition, first with no residual
convolution.  Gaussian covariance differentiation gives
\[
 \frac{\dd}{\dd t}\Phi_u(q_t;Q)
 =\frac12\E\Big\langle
 [u''(z_1)+u'(z_1)^2]\dot C_{11}
       -u'(z_1)u'(z_2)\dot C_{12}\Big\rangle_t.
\]
The diagonal is fixed, so $\dot C_{11}=0$; off the diagonal
$\dot C_{12}=(r-q)(U_{12})$.  Lemmas
\ref{lem:hier-cascade-tilt} and \ref{lem:hier-rank-array}, followed by the
conditional product calculation, give \eqref{eq:src1}.  Bounded derivatives
justify Gaussian integration by parts; one may truncate the number of
leaves and then pass by dominated convergence and the exponential moments
of the Gaussian marks.  Integrating the derivative proves the Lipschitz
bound for finite profiles.

Use the common rank array of Lemma \ref{lem:hier-rank-array}.  If bounded
profiles $q_n\to q$ in $L^1$, then
$q_n(U_{ab})\to q(U_{ab})$ in probability for every $a\ne b$; a union bound
handles any finite subarray.  Keep diagonal $Q$ fixed.  Since $u$ has at
most linear growth, \eqref{eq:hier-exponential-moments} holds uniformly for
$X=u(z)$, as the base single-replica field has law $\cN(0,Q)$.  Lemma
\ref{lem:hier-array-continuity} therefore constructs the limiting functional
and tilted replica expectations independently of the step approximation.
If an approximating profile is raised to $Q$ on a last interval of length
$\varepsilon$, its final recursion is
\[
 \frac1{1-\varepsilon}\log
 \E e^{(1-\varepsilon)u(z+\sqrt{Q-q_k}\,G)}.
\]
As $\varepsilon\downarrow0$ this is exactly the terminal convolution.
The limiting covariance array has an independent residual Gaussian for
each replica.  This verifies both descriptions of the limit.

The $D_{q_n}$ are bounded monotone functions and hence precompact in $L^1$.
The two-replica identity can be tested against the full rank $U_{12}$:
refine a step profile at the endpoints of any step test function, using
zero Gaussian increments on the added levels inside each plateau.
Lemma \ref{lem:hier-cascade-tilt} preserves the entire refined genealogy,
while the derivative martingale is unchanged across those zero increments.
The finite shared-level calculation thus proves the identity for step
test functions, and uniform approximation proves it for continuous ones.
For each continuous $\varphi:[0,1]\to\R$, this identity and Lemma
\ref{lem:hier-array-continuity} give convergence of
\[
 \int_0^1\varphi(s)D_{q_n}(s)\,\dd s
 =\E\langle\varphi(U_{12})u'(z_1)u'(z_2)\rangle_{q_n}.
\]
These limits identify every $L^1$ subsequential limit, proving convergence
of $D_{q_n}$ and defining $D_q$.  The same argument proves continuity for
any convergent sequence of general profiles.  On an interval where $q$ is
constant, choose step approximations that have that value on every compact
subinterval.  Their derivative profiles are constant there, so the limit
is constant almost everywhere there.

Pass to the limit in the integrated finite-profile version of
\eqref{eq:src1}, using boundedness and $L^1$ continuity of $D$.  For fixed
$q,r$, the resulting integrand is continuous in $t$.  The fundamental
theorem of calculus gives the claimed derivative, including one-sided
endpoint derivatives, and completes the proof.
\end{proof}

The moment assumption in Lemma \ref{lem:hier-array-continuity} is a
hypothesis, not a consequence of convergence of overlaps.  Above it
follows from bounded slope and Gaussian exponential moments; the later
cavity applications verify the same bound before inserting capped rows.
Potentials with the negative quadratic tail of the perceptron are handled
by the cap-removal estimates in Section~\ref{sec:row-concavity}.

\subsection{The spherical functional}

Let $\sigma_N$ be normalized surface measure on
$\mathbb S_N=\{x\in\R^N:\|x\|^2=N\}$.  Take independent Gaussian fields
$z_i(\tau)$, $1\le i\le N$, each with off-diagonal profile $p$ and diagonal
$p_d\ge p_*$.  Define
\[
 S_N(p;p_d)=\frac1N\E\log\sum_\tau v_\tau
       \int_{\mathbb S_N}e^{\sum_{i=1}^N x_i z_i(\tau)}\,\dd\sigma_N(x),
 \qquad \mathcal S_N(p)=S_N(p;p_d)-\frac{p_d}{2}.
\]
As above, a residual diagonal is integrated separately for each replica.
It contributes $\exp\{N(p_d-p_*)/2\}$ inside each spin integral, because
$\|x\|^2=N$.  Thus $\mathcal S_N$ does not depend on the diagonal convention.
Gaussian interpolation, followed by the cascade change of measure for the
spin-integrated marks, gives
\begin{equation}
 |\mathcal S_N(p)-\mathcal S_N(\widetilde p)|
 \le\frac12\|p-\widetilde p\|_1.
 \label{eq:hier-spin-lipschitz}
\end{equation}
Indeed the diagonal term is removed by $p_d/2$ and the off-diagonal
derivative is
$-\frac12\E\langle R_{12}(\widetilde p-p)(U_{12})\rangle$,
where $R_{12}=N^{-1}x^1\cdot x^2$ and $|R_{12}|\le1$.  The tilted marginal
law of $U_{12}$ remains uniform by Lemma \ref{lem:hier-cascade-tilt}.
The lemma applies here because the entire vector of Gaussian increments
can be taken as the mark at each vertex; no coordinate factorization is
needed for this interpolation.
For general bounded profiles, $\mathcal S_N$ is defined by its unique
$L^1$-continuous extension from step profiles, whose existence follows from
\eqref{eq:hier-spin-lipschitz}.  Lemma \ref{lem:hier-array-continuity},
applied at fixed $N$, identifies this extension with the Gaussian functional
of the continuous rank array.

\begin{proposition}[Explicit spherical entropy]
\label{prop:hier-spherical-limit}
The limit $\mathcal S(p)=\lim_{N\to\infty}\mathcal S_N(p)$ exists and is
uniform over all nonnegative nondecreasing profiles bounded by any fixed
constant.  For a finite-step profile without residual diagonal, put
\[
 \Delta p_j=p_j-p_{j-1},\qquad
 I=\sum_{j=1}^k m_{j-1}\Delta p_j,
 \qquad
 \lambda_j=b-\sum_{i=j+1}^k m_{i-1}\Delta p_i\quad(0\le j\le k).
\]
Thus $\lambda_k=b$, $\lambda_0=b-I$, and
$\lambda_{j-1}=\lambda_j-m_{j-1}\Delta p_j$.  Define, for $b>I$,
\begin{equation}
 X(b)=-\frac12\log b+\frac{p_0}{2\lambda_0}
       +\sum_{j=1}^k\frac1{2m_{j-1}}
                         \log\frac{\lambda_j}{\lambda_{j-1}}.
 \label{eq:hier-gaussian-spin}
\end{equation}
There is exactly one $b>I$ satisfying
\begin{equation}
 -2X'(b)=\frac1b+\frac{p_0}{\lambda_0^2}
       +\sum_{j=1}^k\frac{p_j-p_{j-1}}{\lambda_j\lambda_{j-1}}=1.
 \label{eq:src2}
\end{equation}
At this value of $b$,
\begin{equation}
 \mathcal S(p)=X(b)+\frac b2-\frac12-\frac{p_k}{2}.
 \label{eq:hier-spherical-value}
\end{equation}
If $p\le B$, its saddle point satisfies
$1\le b<2(B+1)$.
\end{proposition}

\begin{proof}
The calculation has four steps.  We evaluate a Gaussian spin integral,
choose its quadratic multiplier so that one disorder-averaged spin has
variance one, and compare the resulting typical shell with the sphere.
The common Lipschitz bound then makes the limit uniform over bounded
profiles.

First replace the spherical integral by
\[
 Z_N^{\rm G}(b)=\sum_\tau v_\tau
 \int_{\R^N}\exp\{x\cdot z(\tau)-b\|x\|^2/2\}
                     \frac{\dd x}{(2\pi)^{N/2}}.
\]
The terminal logarithm factors into a sum of
$-\frac12\log b+z_i^2/(2b)$.  The elementary Gaussian integral
\[
 \frac1m\log\E\exp\!\left\{
 m\left(c+\frac{(z+\sqrt d\,G)^2}{2\lambda}\right)\right\}
 =c+\frac1{2m}\log\frac\lambda{\lambda-md}
       +\frac{z^2}{2(\lambda-md)}
\]
is valid when $\lambda>md$.  Iterating it and averaging the root Gaussian
gives \eqref{eq:hier-gaussian-spin}, with
$\E\log Z_N^{\rm G}(b)=NX(b)$.  The condition $b>I$ ensures every required
Gaussian exponential moment is finite.

Differentiate this explicit expression to get \eqref{eq:src2}.  Its
left-hand side is strictly decreasing and continuous on $(I,\infty)$ and
tends to zero at infinity.  At $b\downarrow I$ it diverges: if $p_0>0$, use
$p_0/\lambda_0^2$; otherwise use the summand corresponding to the first
positive increment; if every $p_j=0$, use $1/b$.  Hence the solution is
unique.  The equation also gives $b\ge1$.  If $p_k\le B$, then $I\le B$;
at $b=2(B+1)$ its left-hand side is at most
\[
 \frac1{2(B+1)}+\frac{B}{(B+2)^2}\le\frac58<1.
\]
Thus the solution is below $2(B+1)$.

We now justify replacing the Gaussian spin integral by the sphere.
For \emph{one} spin vector sampled from $Z_N^{\rm G}(b)$ and then averaged
over the disorder, the coordinates are iid.  Indeed the terminal function
is a sum of independent coordinate functions; Lemma
\ref{lem:hier-cascade-tilt} makes every path transition factor by coordinate,
and the independent root Gaussian marks are not tilted.  Conditional on
the terminal fields, the spins are independent Gaussians of mean $z_i/b$
and variance $1/b$.  Integrating the factorized tilted path law therefore
preserves independence and identical distribution.  Their common law can
be computed directly.  Write $f_j(z)=c_j+z^2/(2\lambda_j)$ for the quadratic
recursion.  Completing the square in its tilted transition gives
\begin{equation}
 Z_j\mid Z_{j-1}=z\ \sim
 \cN\!\left(\frac{\lambda_j}{\lambda_{j-1}}z,
        \frac{\Delta p_j\lambda_j}{\lambda_{j-1}}\right).
 \label{eq:hier-spin-transition}
\end{equation}
Consequently $Y_j=Z_j/\lambda_j$ is a centered Gaussian martingale with
independent increments of variances
$\Delta p_j/(\lambda_j\lambda_{j-1})$, starting with variance
$p_0/\lambda_0^2$.  The terminal spin is $Y_k$ plus an independent
$\cN(0,1/b)$ variable, so its variance is
\[
 \frac{p_0}{\lambda_0^2}
 +\sum_{j=1}^k\frac{\Delta p_j}{\lambda_j\lambda_{j-1}}+\frac1b=1
\]
by \eqref{eq:src2}.  Thus the averaged single-vector law is exactly
$\cN(0,I_N)$.  This assertion concerns a single disorder-averaged sample;
it does not assert independence at fixed disorder or between replicas.

Let $G_N^{\rm G}$ denote the $x$-marginal of the normalized Gibbs measure
associated with the partition function $Z_N^{\rm G}(b)$.
Fix $0<\delta<1$ and write
$A_{N,\delta}=\{|\|x\|^2/N-1|\le\delta\}$.  The law of large numbers gives
$\E G_N^{\rm G}(A_{N,\delta}^c)\to0$, and Markov's inequality gives
$G_N^{\rm G}(A_{N,\delta})\to1$ in disorder probability.  To pass to
logarithms, let $Z_{N,\delta}^{\rm G}$ be the integral restricted to this
shell.  Lemma \ref{lem:hier-cascade-tilt} expresses $\log Z_N^{\rm G}$ as
an iid coordinate sum at the root plus a difference of two log total
cascade masses.  More explicitly the root shift is
$\sum_{i=1}^N[c_0+Z_{0i}^2/(2\lambda_0)]$, where
\[
 c_0=-\frac12\log b+
          \sum_{j=1}^k\frac1{2m_{j-1}}\log\frac{\lambda_j}{\lambda_{j-1}}.
\]
At each vertex the edge multiplier has $m$th moment exactly one by
\eqref{eq:hier-mark-recursion}, also when the marks comprise $N$
coordinates.  Hence the transformed weight processes have exactly their
original intensities.  Conditionally on root marks, $T'$ has the same
$N$-independent marginal law as $T$; all scale changes have been included
in the root shift.  Its quadratic Gaussian summands have fixed second
moments, and Lemma \ref{lem:hier-cascade-tilt} bounds the second moments of
the two log masses.  Therefore
$\E|\log Z_N^{\rm G}|^2=O(N^2)$, with a constant allowed to depend on the
fixed step profile.  The restricted logarithm has a
deterministic lower bound $-C_{b,\delta}N$ for large $N$: at any radius the
sphere average of $e^{x\cdot z}$ is at least one by Jensen's inequality,
and the Gaussian volume of the shell is at least $e^{-C_{b,\delta}N}$.
Consequently
\[
 0\le\frac1N\log\frac{Z_N^{\rm G}}{Z_{N,\delta}^{\rm G}}
 \le\frac1N\log Z_N^{\rm G}+C_{b,\delta}
\]
is uniformly integrable.  The left side tends to zero in probability, so
\begin{equation}
 \E\log Z_{N,\delta}^{\rm G}=NX(b)+o(N).
 \label{eq:hier-shell-log}
\end{equation}

For completeness the radial comparison includes the normalization of
surface measure.  With $x=r\widehat x$, $\widehat x\in\mathbb S_N$, polar
coordinates give
\[
 \frac{\dd x}{(2\pi)^{N/2}}
 =c_N r^{N-1}\,\dd r\,\dd\sigma_N(\widehat x),
 \qquad c_N=\frac{2(N/2)^{N/2}}{\Gamma(N/2)},
 \qquad \frac1N\log c_N\longrightarrow\frac12.
\]
For every fixed field the sphere integral of $e^{r\widehat x\cdot z}$ is
nondecreasing in $r\ge0$, by symmetry and monotonicity of the hyperbolic
cosine.  On $\sqrt{1-\delta}\le r\le\sqrt{1+\delta}$, the logarithm of the
remaining radial factor, divided by $N$, is
$1/2-b/2+O_b(\delta)+o(1)$.  Its integral over this interval has the same
estimate.  Sandwich the spin integral between the values at the two shell
endpoints.  Scaling a field by $r$ scales its covariance by $r^2$, and
\eqref{eq:hier-spin-lipschitz}, including the diagonal correction, gives
\[
 |S_N(r^2p;r^2p_k)-S_N(p;p_k)|\le B|r^2-1|.
\]
Together with \eqref{eq:hier-shell-log}, the radial sandwich shows
$S_N(p;p_k)=X(b)+b/2-1/2+O_{b,B}(\delta)+o(1)$.  First let $N\to\infty$,
then $\delta\downarrow0$, to obtain \eqref{eq:hier-spherical-value}.

Finally, nonnegative nondecreasing profiles bounded by $B$ form a compact
set in $L^1$, and finite-step profiles are dense in it.  The common
Lipschitz estimate \eqref{eq:hier-spin-lipschitz} transfers convergence from
a finite net of step profiles to uniform convergence on this compact set.
This proves the proposition for general bounded profiles.
\end{proof}

\subsection{Concavity and the entropy conjugate}

For a trial overlap $q$, the pressure calculation in
Section~\ref{sec:gaussian-pressure} will couple the row functional to
the spherical functional through $p_q=\alpha D_q$.  Its two bounds first
lead to the expression
\[
 \alpha\Phi_u(q)+\mathcal S(p_q)+\frac12\int_0^1q(s)p_q(s)\,\dd s.
\]
To identify the last two terms, we compute the spherical response.
At a step profile $p$, the differential of $\mathcal S$ is represented
by $-r/2$, where $r$ is a nondecreasing profile strictly below one.
For general $p$ we construct a supergradient of this form.  The inverse
relation between $p$ and the selected $r$ will both prove concavity and
identify when the last two terms equal the entropy conjugate $E(q)$.

For any nonnegative nondecreasing $v$ with $v_*<1$, define its overlap
distribution and integrated distribution by
\[
 m_v(t)=\Prob(v(U)\le t),\qquad
 \chi_v(t)=\int_t^1m_v(a)\,\dd a,
 \qquad U\sim\operatorname{Unif}[0,1].
\]
The continuous form of the inverse map will be
\begin{equation}\label{eq:press-continuous-inverse}
 \mathcal I(v)(s)=\int_0^{v(s)}\frac{\dd t}{\chi_v(t)^2}.
\end{equation}
Its denominator is positive on the integration range, since
$\chi_v(t)\ge1-v_*$ for $t\le v_*$.  The next proposition proves the
inverse relation for the selected supergradient.

\begin{proposition}[Concavity and supergradients]
\label{prop:hier-spin-concavity}
The functional $\mathcal S$ is concave on bounded nonnegative nondecreasing
profiles.  At a finite-step profile its differential is
\begin{equation}
 \dd\mathcal S(p)=-\frac12\sum_{j=0}^k w_jr_j\,\dd p_j,
 \quad
 r_j=\frac{p_0}{\lambda_0^2}
       +\sum_{i=1}^j\frac{p_i-p_{i-1}}{\lambda_i\lambda_{i-1}},
 \quad r_k=1-\frac1b.
 \label{eq:hier-spin-gradient}
\end{equation}
For every general bounded profile $p$ there exists a nonnegative
nondecreasing profile $r$, with $r_*<1$, such that
\begin{equation}
 \mathcal S(\widetilde p)
 \le\mathcal S(p)-\frac12\int_0^1r(s)(\widetilde p-p)(s)\,\dd s
 \label{eq:hier-spin-tangent}
\end{equation}
for every bounded nonnegative nondecreasing $\widetilde p$.
The profile $r$ can be chosen so that $p=\mathcal I(r)$ almost everywhere.
If $p\le B$, $r$ may be chosen with
$r_*\le1-[2(B+1)]^{-1}$.
\end{proposition}

\begin{proof}
All the derivatives can first be taken with $p_0>0$ and strictly positive
increments, and then extended to boundary points by continuity.
Differentiate \eqref{eq:hier-gaussian-spin} and
\eqref{eq:hier-spherical-value}.  The coefficient of $\dd b$ is
$X'(b)+1/2=0$ by \eqref{eq:src2}.  Substituting
$\dd\lambda_j=\dd b-\sum_{i>j}m_{i-1}(\dd p_i-\dd p_{i-1})$
and collecting coefficients of $\dd p_j$ gives
\eqref{eq:hier-spin-gradient}.  The last equality there is precisely
\eqref{eq:src2}.

To check the Hessian sign without a sign convention ambiguity, invert this
gradient map.  Write $\chi_j=1/\lambda_j$.  The differences of
\eqref{eq:hier-spin-gradient} and the recursion for $\lambda$ give
\[
 r_j-r_{j-1}=\Delta p_j\chi_j\chi_{j-1},\qquad
 \chi_{j-1}-\chi_j=m_{j-1}(r_j-r_{j-1}),\qquad
 \chi_k=1-r_k.
\]
Telescoping yields the inverse formulas
\begin{equation}
 \begin{split}
 \chi_j&=1-m_jr_j-\sum_{i>j}w_ir_i,\\
 p_0&=r_0/\chi_0^2,\qquad
 p_j-p_{j-1}
 =\frac{\chi_j^{-1}-\chi_{j-1}^{-1}}{m_{j-1}}
 =\frac{r_j-r_{j-1}}{\chi_j\chi_{j-1}}.
 \end{split}
 \label{eq:hier-inverse-gradient}
\end{equation}
Conversely every $0\le r_0\le\cdots\le r_k<1$ gives positive $\chi_j$ and a
nonnegative nondecreasing $p$ by these formulas; putting $b=1/\chi_k$
recovers \eqref{eq:src2}.  Thus they are genuine inverse relations.

Here is the full quadratic-form calculation.  Put
\[
 T_j=\sum_{i\ge j}w_i\,\dd r_i,\quad T_{k+1}=0,
 \qquad L_j=m_j\,\dd r_j+T_{j+1}=-\dd\chi_j;
 \qquad L_0=T_0.
\]
Differentiating \eqref{eq:hier-inverse-gradient} gives
\[
 \dd p_0=\frac{\dd r_0}{\chi_0^2}
                  +\frac{2r_0T_0}{\chi_0^3},\qquad
 \dd(\Delta p_j)=\frac1{m_{j-1}}
 \left(\frac{L_j}{\chi_j^2}-\frac{L_{j-1}}{\chi_{j-1}^2}\right).
\]
Insert these into
$\sum_jw_j\dd r_j\dd p_j=T_0\dd p_0+\sum_{j\ge1}T_j\dd(\Delta p_j)$.
Summation by parts gives exactly
\begin{equation}
 \sum_{j=0}^k w_j\dd r_j\dd p_j
 =\frac{m_0(\dd r_0)^2-T_1^2/m_0}{\chi_0^2}
  +\frac{2r_0T_0^2}{\chi_0^3}
  +\sum_{j=1}^k\frac{w_jL_j^2}{m_{j-1}m_j\chi_j^2}.
 \label{eq:hier-positive-form}
\end{equation}
Since $\chi_j\le\chi_0$, replace $\chi_j^{-2}$ in the last, nonnegative
sum by $\chi_0^{-2}$.  The identity
\[
 \frac{w_jL_j^2}{m_{j-1}m_j}
 =w_j(\dd r_j)^2+\frac{T_j^2}{m_{j-1}}-\frac{T_{j+1}^2}{m_j}
\]
then telescopes and proves
\begin{equation}
 \sum_jw_j\dd r_j\dd p_j
 \ge\frac1{\chi_0^2}\sum_jw_j(\dd r_j)^2\ge0.
 \label{eq:hier-hessian-sign}
\end{equation}
For $k=0$ the same calculation has $m_0=1$ and empty sums.
Along a line in $p$, \eqref{eq:hier-spin-gradient} and
\eqref{eq:hier-hessian-sign} give $\dd^2\mathcal S\le0$.
Continuity proves concavity on the boundary and, by step approximation,
on the full domain.

For general $p\le B$, choose step profiles $p_n\le B$ converging to $p$.
Their gradients have nonnegative nondecreasing representatives $r_n$ and,
by Proposition \ref{prop:hier-spherical-limit},
\[
 (r_n)_*=1-1/b_n\le1-[2(B+1)]^{-1}.
\]
Pass to an $L^1$ convergent subsequence.  The finite-step tangent inequality
holds against every step competitor after refining to a common partition.
Letting $n\to\infty$, and then approximating an arbitrary bounded
competitor, gives \eqref{eq:hier-spin-tangent}.  The displayed uniform bound
passes to the limit.

This same construction gives the asserted inverse relation.  For a step
profile $v$, $\chi_v$ is constant below its lowest value and affine of
slope $-m_{j-1}$ between successive values.  Integrating $\chi_v^{-2}$
therefore gives $v_0/\chi_0^2$ and increments
$(v_j-v_{j-1})/(\chi_j\chi_{j-1})$.  Thus the finite inverse formulas
are exactly $p_n=\mathcal I(r_n)$.  To pass to the limit, Fubini gives
\begin{equation}\label{eq:hier-chi-quantile}
 \chi_v(t)=1-t-\int_0^1(v(s)-t)_+\,\dd s,
 \qquad
 \sup_t|\chi_v(t)-\chi_{\widetilde v}(t)|
       \le\|v-\widetilde v\|_1.
\end{equation}
For $v,\widetilde v\le Q<1$, their integrated distributions are at least
$1-Q$ on $[0,Q]$.  Comparing the integrands and the upper integration
limits in \eqref{eq:press-continuous-inverse} gives
\[
 \|\mathcal I(v)-\mathcal I(\widetilde v)\|_1
       \le C_Q\|v-\widetilde v\|_1.
\]
Consequently $p=\mathcal I(r)$.  This identifies the supergradient
obtained from the step approximations; no assertion about uniqueness of
all supergradients on the monotone cone is needed.
\end{proof}

\begin{proposition}[Entropy conjugacy]
\label{prop:hier-entropy-conjugacy}
Let $q$ be nonnegative and nondecreasing with $q_*<1$, and use $\chi_q$
defined above.  Then the following supremum is attained:
\begin{equation}
 \begin{split}
 E(q)&:=\sup_{\substack{p\ge0\text{ bounded}\\p\text{ nondecreasing}}}
       \left\{\mathcal S(p)+\frac12\int_0^1q(s)p(s)\,\dd s\right\}\\
 &=\frac12\left\{\log(1-q_*)+
                       \int_0^{q_*}\frac{\dd t}{\chi_q(t)}\right\},
 \qquad \chi_q(t)=\int_t^1m_q(v)\,\dd v.
 \end{split}
 \label{eq:src3}
\end{equation}
An optimizer may be chosen with
$p_*\le q_*/(1-q_*)^2$.  Equivalently,
\begin{equation}
 E(q)=\frac12\int_0^1
       \left\{\frac1{\chi_q(t)}-\frac1{1-t}\right\}\,\dd t,
 \label{eq:hier-entropy-integral}
\end{equation}
where the integrand is identically zero for $t>q_*$.  For profiles
$q,\widetilde q\le Q<1$,
\begin{equation}
 |E(q)-E(\widetilde q)|
 \le\frac{Q}{2(1-Q)^2}\|q-\widetilde q\|_1.
 \label{eq:hier-entropy-continuity}
\end{equation}
\end{proposition}

\begin{proof}
For a step profile $q$, set $r=q$ in
\eqref{eq:hier-inverse-gradient} and obtain $p$.
The tangent inequality in Proposition \ref{prop:hier-spin-concavity} gives,
for every competitor $\widetilde p$,
\[
 \mathcal S(\widetilde p)+\tfrac12\int q\widetilde p
 \le\mathcal S(p)+\tfrac12\int qp;
\]
thus $p$ attains the supremum.  In these inverse formulas
$\chi_j=\chi_q(q_j)$, $b=1/(1-q_k)$, and
\[
 p_0=\frac{q_0}{\chi_0^2},\qquad
 \Delta p_j=\frac{q_j-q_{j-1}}{\chi_j\chi_{j-1}}.
\]
The algebra in the value can be made explicit.  Abel summation and the
inverse formulas give
\begin{align*}
 p_k-\sum_jw_jq_jp_j
 &=p_0\chi_0+
   \sum_{j=1}^k\Delta p_j(\chi_{j-1}+m_{j-1}q_{j-1})\\
 &=\frac{q_0}{\chi_0}+
   \sum_{j=1}^k\left(\frac{q_j}{\chi_j}
                         -\frac{q_{j-1}}{\chi_{j-1}}\right)
 =\frac{q_k}{\chi_k}=b-1.
\end{align*}
Substitution into \eqref{eq:hier-gaussian-spin} and
\eqref{eq:hier-spherical-value} therefore yields
\[
 2E(q)=\log(1-q_k)+\frac{q_0}{\chi_0}
       +\sum_{j=1}^k\frac1{m_{j-1}}
                                   \log\frac{\chi_{j-1}}{\chi_j}.
\]
On $[0,q_0]$, $\chi_q=\chi_0$; on each interval
$[q_{j-1},q_j]$ it is affine of slope $-m_{j-1}$.  Integrating its reciprocal
on these intervals gives exactly \eqref{eq:src3}, also when an interval has
zero length.

For passage to general profiles, fix $Q<1$ with $q\le Q$ and approximate
by step profiles $q_n\le Q$ in $L^1$.  Their inverse profiles satisfy
\[
 (p_n)_*=\frac{(q_n)_0}{(\chi_n)_0^2}
   +\sum_j\frac{(q_n)_j-(q_n)_{j-1}}
                   {(\chi_n)_j(\chi_n)_{j-1}}
 \le\frac{Q}{(1-Q)^2},
\]
because every $\chi_n\ge1-Q$.  Compactness gives a subsequence $p_n\to p$
in $L^1$.  The finite-step tangent inequalities, first against step
competitors and then against arbitrary bounded competitors, pass to the
limit by the Lipschitz continuity of $\mathcal S$ and the uniform bounds.
They prove attainment for $q$ and show that its value is the limit of the
finite-step values.  This argument does not interchange an uncontrolled
supremum with a limit.

To identify the limiting expression, use \eqref{eq:hier-chi-quantile}.
For $q\le Q$,
$\chi_q(t)\ge1-Q$ on $[0,Q]$, and $\chi_q(t)=1-t$ on $[Q,1]$.
Consequently the right side of \eqref{eq:hier-entropy-integral} is continuous
under these approximations, with the explicit bound
\eqref{eq:hier-entropy-continuity}.  It agrees with
\eqref{eq:src3} because
$\int_0^{q_*}(1-t)^{-1}\,\dd t=-\log(1-q_*)$.
Approximating with $Q=q_*$ gives the stated optimizer bound; when $q_*=0$,
$p=0$ suffices.  This completes the proof.
\end{proof}


\section{Concavity and stability of the row functional}
\label{sec:row-concavity}

The pressure argument needs three properties of the row functional at a
fixed diagonal covariance $Q$.  Ordinary concavity in the off-diagonal
profile $q$ supplies the interpolation sign.  Strict concavity makes the
derivative profile $D_q$ strongly monotone; this will force a minimizing
overlap to equal the spherical response constructed in the preceding
section.  Finally, concavity must persist under both signs of a small
terminal perturbation, so that pressure bounds identify empirical row laws.

We first derive a Hessian identity for smooth profiles and terminals.
For a concave terminal it is a sum of nonpositive terms.  The hinge has
enough curvature on one fixed interval to make the bound uniformly
negative.  We then control the change in that Hessian under a compact
perturbation and remove the terminal and profile regularizations.

For the hinge terminal $u=-\beta V$, two perturbation regimes are used
below. Corollary~\ref{cor:conc-hinge-stability} permits arbitrary
$\psi\in C_c^\infty(\R)$ in $u+d\psi$, with $|d|$ small at each fixed
$\beta$. Corollary~\ref{cor:conc-convex-loss} treats
$-\beta[V+w\psi(1+z)]$ with $\psi\in C_c^\infty(( -\infty,0))$;
convexity of the loss gives an admissible $|w|$ independent of $\beta$.

\subsection{The parabolic representation and its regularity}

Suppose initially that $q\in C^\infty([0,1])$, $q(0)=0$, $q(1)=Q$, and
$\dot q>0$.  The Gaussian recursion defining the row functional is equivalent
to
\begin{equation}
 f_s=-\frac{\dot q(s)}2\bigl(f_{zz}+s f_z^2\bigr),
 \qquad f(1,z)=u(z),\qquad \Phi_u(q;Q)=f(0,0).
 \label{eq:src4}
\end{equation}
Here and below a dot denotes differentiation in rank time $s$.  In variance
time $t=q(s)$, write $F(t,z)=f(q^{-1}(t),z)$ and $m(t)=q^{-1}(t)$.  The equation
becomes
\begin{equation}
 F_t=-\frac12\bigl(F_{zz}+m(t)F_z^2\bigr),\qquad F(Q,z)=u(z).
 \label{eq:conc-variance-pde}
\end{equation}
When $m$ is constant, equal to $a\in[0,1]$, on an interval of length $T$, its
backward solution operator is
\[
 \mathcal T_{a,T}v(z)=
 \begin{cases}
 a^{-1}\log\E\exp\{a v(z+\sqrt T G)\},&a>0,\\
 \E v(z+\sqrt T G),&a=0,
 \end{cases}
 \qquad G\sim\cN(0,1).
\]
Consequently step functions $m$ reproduce exactly the finite Gaussian
recursion, including its initial heat interval and terminal convolution.

We record the analytic class in which subsequent differentiations are made.
Assume that $u$ is smooth, all its derivatives of order at least two are
bounded, and, for fixed constants $A,C\geq0$, $B\in\R$, and $\kappa\geq0$,
\begin{equation}
 -A(1+z^2)\leq u(z)\leq B,\qquad
 -C\leq u''(z)\leq\kappa,\qquad \kappa Q<1.
 \label{eq:conc-terminal-class}
\end{equation}
Constants in the following bounds depend only on these constants and $Q$,
unless a dependence on a higher derivative of $u$ is explicitly indicated.

\begin{lemma}[Parabolic bounds and differentiability]
\label{lem:conc-regularity}
Equation \eqref{eq:conc-variance-pde} has a unique solution of at most quadratic
growth with a spatial derivative of at most linear growth.  For continuous
$m:[0,Q]\to[0,1]$, the solution is classical and satisfies
\begin{equation}
 -C\leq F_{zz}(t,z)\leq\frac{\kappa}{1-\kappa(Q-t)},
 \qquad |F_z(t,z)|\leq C_1(1+|z|).
 \label{eq:conc-parabolic-bounds}
\end{equation}
For a fixed smooth terminal, each spatial derivative of order at least two
is bounded.  The corresponding bounds are uniform over $m\in[0,1]$.
For a smooth profile $q$ with $\dot q>0$ and a smooth direction $\eta$ with
$\eta(0)=\eta(1)=0$, the solution of \eqref{eq:src4} is twice differentiable
in the parameter $\varepsilon$ at $q+\varepsilon\eta$.
\end{lemma}

\begin{proof}
\emph{Construction and curvature bounds.}
For a step coefficient $m$, use the displayed Gaussian operators.  Their
integrals exist since $u$ is bounded above.  Jensen's inequality and iteration
give
\[
 -A(1+z^2+Q-t)\leq F(t,z)\leq B.
\]
The differentiated equation, with
$\mathcal D=\partial_t+\frac12\partial_{zz}+mF_z\partial_z$, is
\[
 \mathcal D F_z=0,\qquad
 \mathcal D F_{zz}=-m F_{zz}^2.
\]
Backward parabolic comparison, or forward comparison after the change of
time $\tau=Q-t$, bounds the second derivative below by $-C$ and above by
the solution of $K'=K^2$, $K(0)=\kappa$.  These comparisons can first be made
for bounded-slope terminals and on bounded spatial intervals, with quadratic
barriers, and then passed to expanding intervals.  Here is an explicit
bounded-slope approximation preserving the constants needed for that passage.
Keep $u$ on $[-R,R]$.  If $p=u'(R)\leq0$, continue affinely to the right.
If $p>0$, continue with derivative $p-C(z-R)$ until that derivative reaches
zero, and then continue constantly.  The maximum of this continuation is
$u(R)+p^2/(2C)\leq B$: the inequality follows by applying
$u(R+h)\geq u(R)+ph-Ch^2/2$ at $h=p/C$ to the original terminal.
On the left, continue affinely if $u'(-R)\geq0$; otherwise continue with
second derivative $-C$ until the derivative reaches zero, and then
constantly.  The same argument bounds its maximum by $B$.
If $C=0$, the original terminal is convex and bounded above, hence constant,
so no continuation is needed.  The resulting $C^{1,1}$ function has
curvature in $[-C,\kappa]$ distributionally.  Its value and derivative at
zero agree with those of $u$, so the lower curvature bound supplies a common
negative quadratic envelope on the whole line.  Convolution with a compact
nonnegative mollifier preserves both curvature bounds and the upper bound;
let its width tend to zero as $R\to\infty$.
The preceding quadratic barriers and Gaussian integrability now permit
passage to the original terminal.  The two-sided bound on $F$ near
$z=0$, together with the bound on $F_{zz}$, bounds $F_z(t,0)$; integration of
$F_{zz}$ then gives the linear-growth bound in
\eqref{eq:conc-parabolic-bounds}.

\emph{Higher spatial derivatives.}
These estimates do not require bounds on derivatives of $m$.  For example,
\[
 \mathcal D F_{zzz}=-3mF_{zz}F_{zzz}.
\]
The diffusion representation bounds this derivative by
$\|u'''\|_\infty\exp(3KQ)$, where $K$ bounds $|F_{zz}|$.
Each further spatial differentiation gives a linear equation whose
zeroth-order coefficient is bounded and whose source contains only already
bounded derivatives.  Induction gives the assertion.

\emph{General coefficients.}
Approximate a continuous $m$ uniformly by step functions.  If $F$ and $F'$
are the two step solutions, their difference solves a linear backward
equation with diffusion coefficient $1/2$, drift
$m(F_z+F'_z)/2$, and source
$-(m-m')(F'_z)^2/2$.  Its drift is globally Lipschitz and has linear growth.
The diffusion representation, with the terminal difference zero, gives on
each compact spatial interval
\[
 |F-F'|\leq C_R\|m-m'\|_\infty.
\]
The spatial derivative bounds and the equation give local compactness also
for the derivatives.  Thus the limit solves
\eqref{eq:conc-variance-pde}; the same difference equation proves uniqueness.
This constructs the claimed classical solution and identifies it with the
Gaussian recursion.  The argument also works with step approximations of
bounded measurable $m$, using convergence in measure and dominated
convergence in the diffusion representation whenever needed below.

\emph{Dependence on the covariance profile.}
Work in rank time and keep $\dot q$ positive in a neighborhood of
the parameter under consideration.  Subtracting two equations gives a
linear equation with drift
$\dot q_{\varepsilon}s(f_{\varepsilon,z}+f_z)/2$ and source
$-(\dot q_{\varepsilon}-\dot q)(f_{zz}+s f_z^2)/2$.
Its diffusion representation and its spatially differentiated equations
bound the difference and each fixed number of its spatial derivatives by
$C\|\dot q_{\varepsilon}-\dot q\|_\infty(1+z^2)$.
To see that differentiation does not increase this polynomial degree, keep
the linear-growth drift in the transport operator: every spatial derivative
of that drift of order at least one is bounded.  The commutator terms are
therefore bounded coefficients times derivatives already estimated, while
the source and all its fixed-order spatial derivatives have at most
quadratic growth.  Induction and the diffusion moment bounds preserve the
same quadratic envelope.
The constant here may depend on the fixed profile and on finitely many
terminal derivatives.  Applying the same bounds to difference quotients,
and then to the difference between those quotients and their linearized
equations, proves first and second parameter differentiability.  This
also justifies differentiation under the diffusion expectations below.
Indeed the diffusions have moments of every fixed order, and stopping at
$|Z|=R$ followed by $R\to\infty$ removes all spatial localizations.
\end{proof}

\subsection{An exact Hessian identity}

Return to rank time.  Put $a=f_z$, $b=f_{zz}$, and let $Z$ solve
\begin{equation}
 \dd Z_s=\sqrt{\dot q(s)}\,\dd W_s+
       \dot q(s)s a(s,Z_s)\,\dd s,\qquad Z_0=0.
 \label{eq:conc-diffusion}
\end{equation}
Its generator, including the time derivative, is
$\mathcal D=\partial_s+\dot q\partial_{zz}/2+\dot q s a\partial_z$.
All unqualified expectations in the rest of this subsection evaluate the
displayed functions along $Z_s$.  Differentiation of
\eqref{eq:src4} gives
\[
 \mathcal D a=0,\qquad \mathcal D b=-\dot q s b^2,
 \qquad U=b+s a^2,\qquad \mathcal D U=a^2.
\]
Thus $a(s,Z_s)$ is precisely the derivative martingale of the tilted
Gaussian recursion.

Let $v=\partial_\varepsilon f_{q+\varepsilon\eta}|_{\varepsilon=0}$ and
$w=\partial_\varepsilon^2 f_{q+\varepsilon\eta}|_{\varepsilon=0}$.
Their terminal values vanish, and
\begin{align}
 \mathcal Dv&=-\dot\eta U/2,\nonumber\\
 \mathcal Dw&=-\dot\eta(v_{zz}+2s a v_z)-\dot q s v_z^2.
 \label{eq:conc-linearized}
\end{align}
Since $\eta$ vanishes at both endpoints, integration by parts in
$\E U(s,Z_s)$ gives
\[
 v(0,0)=\frac12\int_0^1\dot\eta(s)\E U(s,Z_s)\,\dd s
       =-\frac12\int_0^1\eta(s)\E a(s,Z_s)^2\,\dd s.
\]
This is the differential formula \eqref{eq:src1} in the smooth setting.

For the second derivative, the term $\dot\eta$ in the equation for $v$
would obscure the sign.  Adding $\eta U/2$ cancels that term and leaves
$\mathcal D(v+\eta U/2)=\eta a^2/2$.  The spatial derivative of this
modified variation is the useful unknown in the Hessian identity.

\begin{lemma}[Hessian identity]
\label{lem:conc-hessian}
Set $\bar v=v+\eta U/2$ and $l=\bar v_z$.  Then
\begin{equation}
 \begin{split}
 w(0,0)&=\E\int_0^1
 \left(\dot q s l^2-2\eta a l-\frac12\eta^2b^2\right)\dd s,\\
 \mathcal Dl&=\eta a b-\dot q s b l,\qquad l(1,\cdot)=0.
 \end{split}
 \label{eq:src5}
\end{equation}
\end{lemma}

\begin{proof}
For $Mg=g_{zz}+2sag_z$, direct differentiation gives
\[
 \mathcal D(Mg)=M(\mathcal Dg)+(2a-\dot q s U_z)g_z.
\]
Integrate the first term on the right side of the second equation in
\eqref{eq:conc-linearized} by parts, using
$\frac{\dd}{\dd s}\E Mg(s,Z_s)=\E\mathcal D(Mg)(s,Z_s)$.
Substitute $v_z=l-\eta U_z/2$ and complete the square.  Apart from
$\dot q s l^2-2\eta a l$, the remaining integrand is
\[
 -\frac14\dot q s\eta^2 U_z^2+\eta^2aU_z
       +\frac12\eta\dot\eta\,MU.
\]
The second identity
\[
 \mathcal D(MU)=2b^2+4aU_z-\dot q s U_z^2
\]
and one further integration by parts turn this integrand into
$-\eta^2b^2/2$.  There are no boundary terms because $\eta(0)=\eta(1)=0$.
Finally $\mathcal D\bar v=\eta a^2/2$ and spatial differentiation, taking
account of the derivative $\dot q s b$ of the drift, give the equation
for $l$.
\end{proof}

\begin{proposition}[Concavity]
\label{prop:conc-concavity}
For every concave terminal in the class
\eqref{eq:conc-terminal-class}, $q\mapsto\Phi_u(q;Q)$ is concave on
smooth strictly increasing profiles.  For every concave globally Lipschitz
terminal, which need not be smooth or bounded above, and for
the quadratic hinge $u=-\beta V$, it is concave on the full convex set
of nonnegative nondecreasing profiles bounded by $Q$.
\end{proposition}

\begin{proof}
First suppose $u''\leq-\rho<0$.  The stochastic flow in
\eqref{eq:conc-diffusion}, started from $z$ at time $s$, has spatial
Jacobian
$J_{s,r}=\exp\{\int_s^r\dot q(t)t b(t,Z_t)\,\dd t\}$.
Differentiating the martingale representation
$a(s,z)=\E_{s,z}u'(Z_1)$ gives
\begin{equation}
 b(s,z)=\E_{s,z}\left[u''(Z_1)
       \exp\left\{\int_s^1\dot q(t)t b(t,Z_t)\,\dd t\right\}\right].
 \label{eq:conc-curvature-flow}
\end{equation}
Thus $k=-b$ is bounded and bounded away from zero.  The product and quotient
rules for $\mathcal D$ give
\[
 \mathcal D(l^2/k)=-2\eta a l+\dot q s l^2
       +\frac{\dot q}{k}(l_z-lk_z/k)^2.
\]
Since $l(1,\cdot)=0$, substitution into \eqref{eq:src5} yields the exact
identity
\begin{equation}
 \begin{split}
 w(0,0)={}&-\frac{l(0,0)^2}{k(0,0)}
 -\E\int_0^1\frac{\dot q}{k}(l_z-lk_z/k)^2\,\dd s\\
 &-\frac12\int_0^1\eta(s)^2\E b(s,Z_s)^2\,\dd s\leq0.
 \end{split}
 \label{eq:conc-negative-squares}
\end{equation}
For a merely concave $u$, apply this argument to $u-\lambda z^2$ and
let $\lambda\downarrow0$.  The parabolic comparison and moment bounds in
Lemma~\ref{lem:conc-regularity} give convergence on every smooth profile
and its segments.  Concavity passes to that limit.  Finally approximate
two arbitrary profiles and their segment together by smooth increasing
profiles with the prescribed endpoint values.  The $L^1$ continuity in
\eqref{eq:src1} proves this extension for Lipschitz terminals already in
the analytic class.

To cover a general concave globally $L$-Lipschitz terminal, convolve it with
a compact nonnegative mollifier of width $\delta$.  The resulting
$u_\delta$ is smooth and concave, has bounded derivatives of every order
at least two, and satisfies $\|u_\delta-u\|_\infty\leq L\delta$.
For each $\lambda>0$, $u_\delta-\lambda z^2$ belongs to
\eqref{eq:conc-terminal-class}, so concavity holds for smooth profiles.
Let $\lambda\downarrow0$ on each of the three profiles of a fixed chord.
This limit needs no bound on $\sup u_\delta$ uniform in $\lambda$.
Indeed, for $0<\lambda\leq1$, Jensen gives a common negative quadratic
lower bound for the solutions, while
$u_\delta(z)-\lambda z^2\leq u_\delta(0)+L|z|$ and
$\mathcal T_{a,T}v\leq\mathcal T_{1,T}v$ for $a\in[0,1]$ give the
upper bound
\[
 F_{\delta,\lambda}(t,z)
 \leq u_\delta(0)+L|z|+
       \log\E\exp\{L\sqrt Q\,|G|\}.
\]
The curvature bounds, hence the linear-growth drift and its diffusion
moments, are uniform for $\lambda\leq1$ at fixed $\delta$.
The linear equation comparing two such solutions therefore bounds their
difference by $C_\delta|\lambda-\lambda'|(1+z^2)$.
It identifies the limit with the Gaussian recursion for $u_\delta$;
alternatively the same conclusion follows by monotone convergence of its
terminal weights, with the displayed upper and lower integrable bounds.
Thus concavity holds for $u_\delta$ on smooth profiles.
Profile continuity from \eqref{eq:src1} extends it to all profiles, and
the $1$-Lipschitz dependence of $\Phi$ on the terminal uniform norm lets
$\delta\downarrow0$.  The quadratic hinge and its other required limits
are treated below.
\end{proof}

\subsection{Strict concavity uniform over profiles}

Fix $\beta>0$.  Let $u_L$ be the function $-\beta V$ on $[-L,\infty)$,
extended by its tangent line at $-L$ on $(-\infty,-L)$, where $L>4$.
Convolve with a nonnegative smooth mollifier of total mass one and support
in $[-\delta,\delta]$, $0<\delta\leq1/4$.  Denote the resulting terminal by
$u_{L,\delta}$.  For the present argument one may also subtract
$\lambda z^2$, $0<\lambda\leq1$.
All these terminals have uniform quadratic-growth bounds, and
\begin{equation}
 -\beta-2\leq u''\leq0,\qquad
 \int_\R|u'''(z)|\,\dd z\leq2\beta,
 \qquad u''\leq-\beta\quad\hbox{on }I=[-3,-2].
 \label{eq:conc-base-bounds}
\end{equation}
The added quadratic is used only to justify division by $-b$; it will
always be removed before an argument that requires a monotone terminal.

\begin{lemma}[Diffusion density bounds]
\label{lem:conc-density}
Let $X$ solve $\dd X_t=\dd B_t+A(t,X_t)\dd t$, with
$|\partial_zA|\leq K$ and $|A(t,0)|\leq K$.
Over an interval of length $T>0$, its transition density is bounded above
by $e^{KT}/\sqrt{2\pi T}$.  Moreover, when $X_0=0$,
$Q\in[1/2,2]$, and $K$ is fixed, the probability
$\Prob(X_Q\in I)$ is bounded below by a positive constant depending
only on $K$ and the nondegenerate bounded interval $I$.
\end{lemma}

\begin{proof}
Condition on the Brownian bridge and write the Brownian path over the
interval as $\mathsf b_t+(t/T)\xi$, where $\xi\sim\cN(0,T)$ is independent
of the bridge.  The endpoint $X_T=\mathcal F(\xi)$ is an increasing
function of $\xi$.  Differentiation of the integral equation gives
\[
 \mathcal F'(\xi)=\frac1T\int_0^T
       \exp\left\{\int_r^T\partial_zA(v,X_v)\,\dd v\right\}\dd r
       \in[e^{-KT},e^{KT}].
\]
The conditional change-of-variables formula proves the upper density
bound, which remains true after averaging over the bridge.
For the lower bound, restrict the bridge to an event on which its supremum
is at most a fixed constant.  Such an event has uniformly positive
probability for $Q\in[1/2,2]$, by Brownian scaling.  Gronwall's inequality
then bounds $|\mathcal F(0)|$ uniformly.  The inverse image of $I$ is an
interval of length at least $e^{-2K}|I|$ contained in a fixed bounded
interval.  Gaussian densities with variance $Q\in[1/2,2]$ have a positive
common lower bound there.  Multiplication by the bridge-event probability
proves the claim.
\end{proof}

\begin{proposition}[Uniform strictness and stability]
\label{prop:conc-stability}
There is $c_\beta>0$ such that, for every $Q\in[1/2,2]$, every regularized
base terminal just described, every smooth strictly increasing profile,
and every smooth endpoint-zero direction,
\begin{equation}
 \frac{\dd^2}{\dd\varepsilon^2}
 \Phi_u(q+\varepsilon\eta;Q)\big|_{\varepsilon=0}
 \leq-c_\beta\|\eta\|_2^2.
 \label{eq:conc-strict-hessian}
\end{equation}
For each $\psi\in C_c^\infty(\R)$ there is $d_0=d_0(\beta,\psi)>0$,
independent of $L$, $\delta$, $\lambda$, $Q$, and the profile, such that
\eqref{eq:conc-strict-hessian} holds with $c_\beta/2$ for the terminal
$u+d\psi$ whenever $|d|\leq d_0$.
\end{proposition}

\begin{proof}
In variance time the drift is $m(t)F_z(t,z)$.  Its Lipschitz and growth
bounds are uniform by Lemma~\ref{lem:conc-regularity}.  Thus
Lemma~\ref{lem:conc-density} gives $\Prob(Z_1\in I)\geq p_\beta>0$,
where $Z_1$ denotes the rank-time endpoint.  From
\eqref{eq:conc-curvature-flow}, the bound on $|b|$, and
\eqref{eq:conc-base-bounds}, for every $s$,
\[
 \E[-b(s,Z_s)]\geq\beta e^{-KQ}\Prob(Z_1\in I)
       \geq\beta e^{-2K}p_\beta.
\]
Jensen's inequality supplies a positive uniform lower bound on
$\E b(s,Z_s)^2$.  Equation~\eqref{eq:conc-negative-squares} proves
\eqref{eq:conc-strict-hessian}.  The lower bound does not depend on
$\lambda$, so it survives $\lambda\downarrow0$.

To prove stability, we compare slopes, then curvatures, then the modified
variation $l$ under a common Brownian coupling.  These three comparisons
control every term in the Hessian identity.  The curvature comparison uses
the total variation of $u''$, so its constants remain uniform as the
mollifier is removed.

\emph{Slope comparison.}
A tilde will denote the terminal $u+d\psi$.
Choose $|d|$ initially so that
$|d|\|\psi''\|_\infty<1/4$; then its positive curvature bound is strictly
less than $1/(2Q)$ for $Q\in[1/2,2]$, and all solution and drift bounds
remain uniform.  In this proof only, use variance-time notation
$a=F_z$, $b=F_{zz}$ and its tilde analogue.  At the same spatial point,
\[
 \widetilde{\mathcal D}(\widetilde a-a)
       =-m b(\widetilde a-a),\qquad
 (\widetilde a-a)(Q,\cdot)=d\psi'.
\]
The diffusion representation implies
\begin{equation}
 \sup_{t,z}|\widetilde a-a|\leq C|d|.
 \label{eq:conc-slope-stability}
\end{equation}
Couple both diffusions using the same Brownian motion.  Starting from a
common position at time $t$, Gronwall gives
$|\widetilde X_r-X_r|\leq C|d|(r-t)$ for $t\leq r\leq Q$.

\emph{Curvature comparison.}
We use estimates that depend on the total variation of
$u''$, and not on $\|u'''\|_\infty$.  The equation
$\mathcal D b_z=-3mbb_z$ and its diffusion representation give
\begin{equation}
 \int_\R|b_z(t,z)|\,\dd z\leq C,
 \qquad \sup_z|b_z(t,z)|\leq\frac{C}{\sqrt{Q-t}}\quad(t<Q).
 \label{eq:conc-curvature-bv}
\end{equation}
For the first bound integrate the representation over starting positions;
the pathwise spatial flow has Jacobian between $e^{-KQ}$ and $e^{KQ}$.
For the second use the upper density bound of
Lemma~\ref{lem:conc-density} and
$\int|u'''|\leq2\beta$.

Compare the flow representations \eqref{eq:conc-curvature-flow} for
$b$ and $\widetilde b$, started from the same point at time $t$.
The terminal perturbation contributes at most $C|d|$.
For the displacement in the base terminal, use the elementary inequality
\[
 |u''(y)-u''(x)|\leq
       \int_{x-h}^{x+h}|u'''(v)|\,\dd v\qquad(|y-x|\leq h).
\]
With $h=C|d|(Q-t)$ and the transition-density bound, its expectation is
at most $C|d|\sqrt{Q-t}$.
The difference of the two exponential Jacobians is at most a constant
times the integral of the difference of their curvatures along the paths.
Putting $A(t)=\sup_z|\widetilde b(t,z)-b(t,z)|$ and using
\eqref{eq:conc-curvature-bv} therefore gives
\[
 A(t)\leq C|d|+C\int_t^Q
       \left(A(r)+\frac{|d|(r-t)}{\sqrt{Q-r}}\right)\dd r.
\]
The singularity is integrable, and backward Gronwall yields
\begin{equation}
 \sup_{t,z}|\widetilde b(t,z)-b(t,z)|\leq C|d|.
 \label{eq:conc-curvature-stability}
\end{equation}

For paths coupled from the root, their displacement at variance time $t$
is at most $C|d|t$.  The same elementary inequality, now applied to $b(t,\cdot)$,
and its total variation and the density bound at time $t$, give
\[
 \E|b(t,\widetilde X_t)-b(t,X_t)|\leq C|d|\sqrt t.
\]
At $t=0$ the positional difference is zero.  Combining this estimate with
\eqref{eq:conc-curvature-stability} and the uniform boundedness of both
curvatures gives, for every fixed $p\geq1$,
\begin{equation}
 \sup_t\|\widetilde b(t,\widetilde X_t)-b(t,X_t)\|_{L^p}
       \leq C_p|d|^{1/p}.
 \label{eq:conc-coupled-curvature}
\end{equation}
The coupled slopes differ pointwise by $C|d|$ and have uniformly bounded
moments of every fixed order.

\emph{Comparison of the modified variations.}
Return to rank time and write $a_s=a(s,Z_s)$, $b_s=b(s,Z_s)$ and
$l_s=l(s,Z_s)$.  On the common Brownian filtration $\mathcal F_s$, the
second equation of \eqref{eq:src5} gives
\begin{equation}
 l_s=-\E\left[
   \left.\int_s^1\eta(r)a_rb_r
       \exp\left\{\int_s^r\dot q(v)v b_v\,\dd v\right\}\dd r
       \right|\mathcal F_s\right].
 \label{eq:conc-l-representation}
\end{equation}
This same filtration is used for the tilde representation; no comparison
of unrelated conditional laws is involved.  Both exponential factors are
uniformly bounded since $\int_0^1\dot q\,\dd s=Q\leq2$.
Conditional expectation is a contraction in $L^p$, so
\[
 \sup_s\|l_s\|_{L^p}+\sup_s\|\widetilde l_s\|_{L^p}
       \leq C_p\|\eta\|_1.
\]
H\"older's inequality, \eqref{eq:conc-coupled-curvature} with $p=4$,
the slope estimates, and Minkowski's inequality for the exponential
integrals give
\begin{equation}
 \sup_s\|\widetilde l_s-l_s\|_{L^2}
       \leq C|d|^{1/4}\|\eta\|_1.
 \label{eq:conc-l-stability}
\end{equation}
For example, the difference of the exponential factors has $L^4$ norm
at most
$C\int_s^r\dot q(v)\|\widetilde b_v-b_v\|_{L^4}\dd v
\leq C|d|^{1/4}$; multiplying by $a_rb_r$, whose $L^4$ norm is bounded,
proves the relevant $L^2$ bound.

\emph{Hessian comparison.}
Compare the three terms of \eqref{eq:src5}.  The $\dot q s l^2$
term changes by at most $C|d|^{1/4}Q\|\eta\|_1^2$; the $2\eta a l$ term
by at most $C|d|^{1/4}\|\eta\|_1^2$; and the $\eta^2b^2/2$ term by at
most $C|d|\|\eta\|_2^2$.  Since $\|\eta\|_1\leq\|\eta\|_2$,
\begin{equation}
 |\widetilde w(0,0)-w(0,0)|
       \leq C|d|^{1/4}\|\eta\|_2^2.
 \label{eq:conc-hessian-stability}
\end{equation}
Choosing $d_0$ so that $Cd_0^{1/4}\leq c_\beta/2$ proves the proposition.
\end{proof}

\subsection{Removing regularizations and general profiles}

We now pass the estimates to the terminals and profiles used in the
pressure formula.  The resulting statement is the row-functional input
for both minimizer identification and compact row perturbations.

\begin{corollary}[The uncapped row functional]
\label{cor:conc-hinge-stability}
Fix $\beta>0$ and $\psi\in C_c^\infty(\R)$.  There are constants
$c_\beta>0$ and $d_0=d_0(\beta,\psi)>0$ such that, for
$u_d=-\beta V+d\psi$, $|d|\le d_0$, $Q\in[1/2,2]$, and all
$q,r\in\mathcal Q_Q$,
\begin{equation}
 \Phi_{u_d}((1-t)q+tr;Q)
 \ge(1-t)\Phi_{u_d}(q;Q)+t\Phi_{u_d}(r;Q)
       +\frac{c_\beta}{4}t(1-t)\|r-q\|_2^2,
 \qquad 0\le t\le1.
 \label{eq:conc-chord}
\end{equation}
When $d=0$, the last coefficient can be replaced by $c_\beta/2$.
The derivative formula \eqref{eq:src1}, including its one-sided endpoint
interpretation, holds with the terminal $u_d$ and fixed diagonal $Q$.
Its nonnegative nondecreasing derivative profile satisfies
$\sup_{q,Q,d}\|D^{u_d}_q\|_\infty\le C_{\beta,\psi}$ in these ranges
and depends continuously on $q$ in $L^1$ at fixed $Q,d$.
The functional values and derivative profiles are limits of those for
$u_{L,\delta}+d\psi$, uniformly over $q,Q,d$, with derivative convergence
in $L^1(\dd s)$.  Here $\delta\downarrow0$ at each fixed cap before
$L\to\infty$.
\end{corollary}

\begin{proof}
The profile inequalities above are uniform in $\dot q$: their constants
use only $Q=\int\dot q$.  To approximate two arbitrary profiles $q,r$,
smooth their monotone extensions, add a small increasing linear function,
and insert short ramps at $0$ and $1$ to impose the values $0,Q$.
These approximations can be chosen to remain in $[0,Q]$ and converge
in $L^1$, hence also in $L^2$; their convex combinations approximate
the whole segment.  For a fixed capped and mollified terminal, first let
$\lambda\downarrow0$ on these smooth segments.  Then
\eqref{eq:src1} passes the strong concavity inequalities to all profiles.
Integrating the Hessian bound gives \eqref{eq:conc-chord} for each capped
and mollified terminal, with coefficient $c_\beta/2$ in the unperturbed
case and $c_\beta/4$ in the perturbed case.

We justify the remaining terminal limits uniformly in the profile.
If $\mu$ is any probability measure on $\R$, $w$ is nonnegative and
increasing with $0<\int w\,\dd\mu<\infty$, and $T$ is nonnegative and
decreasing, then
\begin{equation}
 \frac{\int Tw\,\dd\mu}{\int w\,\dd\mu}\leq\int T\,\dd\mu.
 \label{eq:conc-monotone-tilt}
\end{equation}
First take $T$ bounded.  The covariance of $T$ and $w$ equals half the integral of
$(T(x)-T(y))(w(x)-w(y))$ against $\mu\otimes\mu$, and is nonpositive.
Apply this to $T\wedge K$ and let $K\to\infty$ to obtain the general
case by monotone convergence.
To apply this inequality, condition first on the cavity measure and the
new Gaussian field, and let $\mu$ be the pushforward of the cavity measure
under that field.  Insertion of an increasing terminal $u$ replaces
$\mu$ by its tilt with weight $w=e^u$.  Only after applying the inequality
do we average over the independent field and the cavity randomness.
Before insertion, one sampled field value then has law $\cN(0,Q)$,
because its diagonal covariance is $Q$.  For the RPC row functional the
same construction includes any independent residual Gaussian in the base
space.  Adding $d\psi$ increases the right side of
\eqref{eq:conc-monotone-tilt} by at most the factor
$\exp(2|d|\|\psi\|_\infty)$.

The tangent cap satisfies the exact identity
\[
 0\leq u_L(z)+\beta V(z)=\frac\beta2(z+L)_-^2.
\]
Every terminal between $u_L$ and $-\beta V$ is increasing.  Interpolating
between them and using \eqref{eq:conc-monotone-tilt} proves
\begin{equation}
 \sup_q\big|\Phi_{u_L+d\psi}(q;Q)
              -\Phi_{-\beta V+d\psi}(q;Q)\big|
 \leq e^{2|d|\|\psi\|_\infty}
       \E\frac\beta2(\sqrt Q G+L)_-^2\longrightarrow0.
 \label{eq:conc-cap-limit}
\end{equation}
The convergence is also uniform for $Q\in[1/2,2]$.
For each fixed $L$, mollification converges uniformly in terminal value,
since $u_L$ is Lipschitz; the row functional is $1$-Lipschitz for the
uniform norm on its terminal.  Thus the limits are taken in the order
$\lambda\downarrow0$, $\delta\downarrow0$, and $L\to\infty$.
Equation~\eqref{eq:conc-chord}, including its perturbed form, survives them.

The derivative functional also converges uniformly in the profile.  Here
are the necessary uniform-integrability details.  The derivatives of
$u_L+d\psi$ agree with that of $u_d$ outside $(-\infty,-L)$ and their
absolute values there are bounded by a constant times $1+|z|$.
Equation~\eqref{eq:conc-monotone-tilt}, applied to higher negative-tail
moments, bounds their tilted moments of every fixed order uniformly.
Let $\nu_L$ denote a row measure tilted by $u_L+d\psi$ and let $\nu$
be the corresponding uncapped measure.  Both measures are defined on the
same marked base space: the cascade leaf, its full Gaussian field, and an
independent residual Gaussian integrated inside the terminal weight, when
present.  Include the common refined rank genealogy, inserting zero-variance
levels as in Section~\ref{sec:hierarchy}, or its equivalent auxiliary rank
variables.  In a two-replica expectation the continuous rank is then a
common measurable observable even when both replicas select the same
leaf of the original finite cascade.  Thus comparison of the full product
measures also controls every bounded rank test.
Their unnormalized density ratio
is $e^{-\beta V-u_L}\leq1$ and equals one on $[-L,\infty)$.
If $\varepsilon_L=\nu_L(( -\infty,-L))$, normalization and the triangle
inequality give
\[
 \|\nu-\nu_L\|_{\mathrm{TV},1}\leq4\varepsilon_L,
\]
where $\|\cdot\|_{\mathrm{TV},1}$ is total variation without the factor
$1/2$.  The averaged right side tends to zero uniformly by
\eqref{eq:conc-monotone-tilt}.  The same is true for two replicas.
Test the product of the two terminal derivatives against any bounded
function of rank, first truncating the derivatives and then using their
uniform higher moments.  The two-replica formula for $D_q$ therefore gives
uniform $L^1(\dd s)$ convergence of the capped derivative profiles to
$D^{u_d}_q$.  Mollification is handled similarly at each
fixed cap.  In particular $D^{u_d}_q$ is continuous in $q$ in $L^1$, by the
corresponding capped assertion and this uniform approximation.
Moreover its $L^\infty$ norm is uniformly bounded: the derivative martingale
and the terminal second-moment bound give
$D^{u_d}_q(s)\leq\E_{\mathrm{tilted}}(u_d')^2\leq C_{\beta,\psi}$.
Thus the integrated derivative formula \eqref{eq:src1} and its one-sided
derivative interpretation extend to $u_d$.
\end{proof}

For $u=-\beta V$, write $p_q=\alpha D_q$, $\alpha>0$.
Adding the two tangent inequalities furnished by
the unperturbed version of \eqref{eq:conc-chord}, and using the gradient
$-D_q/2$, proves
\begin{equation}
 \int_0^1(r-q)(p_r-p_q)\,\dd s
       \geq c_{\alpha,\beta}\|r-q\|_2^2,
 \qquad c_{\alpha,\beta}=2\alpha c_\beta>0.
 \label{eq:src6}
\end{equation}
All these assertions hold at the fixed diagonal $Q=1$ used below; the
uniform neighborhood $Q\in[1/2,2]$ is available when varying a diagonal
in a cavity calculation.

\begin{corollary}[Convex perturbations of the loss]
\label{cor:conc-convex-loss}
Let
\[
 W(z)=cV(z)+\chi(z),\qquad c>0,\qquad
 \chi\in C_c^\infty((-\infty,-1)),
\]
and suppose that $W\geq0$, $W'\leq0$, and $W''\geq0$ distributionally.
For every fixed $\beta>0$, the terminal $u=-\beta W$ has a concave row
functional on all profiles with fixed diagonal $Q\in[1/2,2]$.
The derivative formula \eqref{eq:src1}, continuity of $D_q$ in $L^1$, and
the uniform tangent-cap limits hold for this terminal.  In particular,
the Gaussian interpolation upper bound derived for smooth concave caps
applies to $-\beta W$ after removal of the cap.
\end{corollary}

\begin{proof}
The terminal is concave and increasing.  Its derivative is Lipschitz,
has at most linear growth, and vanishes on $[-1,\infty)$.
Choose $L$ so large that $-L<\inf\supp\chi$, with the condition omitted
when $\chi=0$, and extend $u$ tangentially below $-L$.
The resulting terminal $u_L$ is concave, increasing, and globally
Lipschitz.  Proposition~\ref{prop:conc-concavity} applies after
mollification and its removal.  Because the cutoff lies below the entire
support of $\chi$, its error is exactly
\[
 u_L(z)-u(z)=\frac{\beta c}{2}(z+L)_-^2.
\]
Every intermediate terminal is increasing, so
\eqref{eq:conc-monotone-tilt} bounds the difference of the row functionals
by $\frac{\beta c}{2}\E(\sqrt QG+L)_-^2$, uniformly in the profile.
The same comparison applies to a row insertion in the microscopic
Gaussian model.  The derivatives agree outside the cutoff region and
are bounded by $C_{\beta,W}(1+|z|)$ within it.  Higher Gaussian
negative-tail moments and the marked-space total-variation comparison
above prove uniform convergence of $D_q$ and its differential formula.
This proves all assertions, including passage of the interpolation bound
to the uncapped loss.
\end{proof}

If
$\psi\in C_c^\infty((-\infty,0))$ has support in $[-R,-a]$, $a>0$,
then $W_w(z)=V(z)+w\psi(1+z)$ satisfies its assumptions for both signs
of all sufficiently small \emph{fixed} $w$.  For example it suffices to
require
\[
 |w|\|\psi''\|_\infty\leq\tfrac12,\qquad
 |w|\|\psi'\|_\infty\leq\tfrac a2,\qquad
 |w|\|\psi\|_\infty\leq\tfrac{a^2}{4}.
\]
On the perturbation support these inequalities give
$W_w''\geq1/2$, $W_w'\leq-a/2$, and $W_w\geq a^2/4$; outside it
the loss equals $V$.  Hence the concavity assertion applies to
$-\beta W_w$ at every $\beta$, without requiring $|\beta w|\leq d_0$.
The scaled loss $(1+w)V$, $w>-1$, is covered by $\chi=0$, $c=1+w$.


\section{The Gaussian pressure}\label{sec:gaussian-pressure}

For each row terminal $u$ used below, define the finite-system Gaussian
pressure by
\[
 P_N(u)=\frac1N\E\log\int_{\mathbb S_N}
       \exp\left\{\sum_{\mu=1}^{M}u(g_\mu(x))\right\}\sigma_N(\dd x),
 \qquad M=\lfloor\alpha N\rfloor.
\]
Thus $P_N(-\beta V)$ is the pressure of the model at inverse temperature
$\beta$.

We prove, for $\alpha,\beta>0$, the Gaussian pressure formula
\[
 \lim_{N\to\infty}P_N(-\beta V)
 =\min_{q_*<1}\{\alpha\Phi_{-\beta V}(q)+E(q)\}.
\]
We first prove matching bounds for smooth caps of the row potential.
The upper bound inserts one independent row into a limiting overlap array
and uses concavity of its row functional. The lower bound adds a fixed
number of spin coordinates. Their Gaussian fields have an empirical
covariance built from row derivatives; deleting and reinserting two rows
identifies that covariance and its fluctuations. These two arguments use
the same cascade representation, established below from a vanishing
perturbation and the Ghirlanda--Guerra identities. After removing the cap,
the spin inverse relation turns the common row--spin expression into the
displayed entropy formula.

Throughout this section a profile has diagonal $1$, unless another diagonal is
explicitly specified. We write $\langle\cdot\rangle$ for the Gibbs expectation
of the system currently under consideration, and include all its randomness
in $\E$. The perturbation coefficients introduced below are held fixed inside
$\E$; their product-uniform average is denoted by $\E_c$.

We first work with a smooth concave tangent cap $u$ having bounded first three
derivatives and at most linear growth. The concentration and cavity arguments
need only these regularity bounds. The interpolation upper bound additionally
uses profile concavity from Section~\ref{sec:row-concavity}; it therefore also
applies to the small terminal perturbations constructed there.

\subsection{Perturbations and overlap arrays}

The perturbation will identify every subsequential overlap law needed in
the bounds, without requiring convergence of the unperturbed overlap law.
An overlap array is viewed as an element of the compact product space of its
entries. Convergence of its law means convergence of every finite subarray.
For spins $x^\ell\in\sqrt N\mathbb S^{N-1}$ put
$R_{\ell\ell'}=N^{-1}x^\ell\cdot x^{\ell'}$. When a finite trial cascade is
present, let $Q_{\ell\ell'}$ be its covariance overlap, including
$Q_{\ell\ell}=1$.

\Needspace{8\baselineskip}
\begin{lemma}[A vanishing perturbation]\label{lem:press-gg}
Fix $\rho\in(1/4,1/2)$ and put $s_N=N^\rho$. There is an independent centered
Gaussian perturbation $H_N^{\rm pert}$ with the following properties.
\begin{enumerate}[label=(\roman*)]
\item Its addition changes the expected pressure by at most $Cs_N^2/N$.
\item For the sphere-only system, every subsequential overlap-array limit
obtained at almost every perturbation parameter satisfies the
Ghirlanda--Guerra identities. With a finite trial cascade, the corresponding
identities hold for every continuous function of the pair $(R,Q)$:
\begin{equation}\label{eq:press-joint-gg}
 \E\langle f\psi(R_{1,n+1},Q_{1,n+1})\rangle
 =\frac1n\E\langle f\rangle\E\langle\psi(R_{12},Q_{12})\rangle
  +\frac1n\sum_{\ell=2}^n
   \E\langle f\psi(R_{1\ell},Q_{1\ell})\rangle.
\end{equation}
Here $n\ge2$ and $f$ is a bounded measurable function of the first $n$ joint
overlaps. For the sphere-only system, omit $Q$.
\item The extraction can be made simultaneously for any prescribed countable
collection of systems with finitely many rows removed, and for almost every
time in a bounded interpolation parameter interval.
\end{enumerate}
In (ii)--(iii) the precise extraction convention is this: from any dimension
subsequence one can first choose a further subsequence on which all the
identity errors vanish for almost every coefficient vector, and then, for
each such vector, take any further convergent array subsequence. No assertion
is made that averaging limiting arrays over coefficient vectors preserves
these identities.
\end{lemma}

\begin{proof}
Enumerate pairs $(j,\lambda)$, where $j\ge1$ and $\lambda$ belongs to a
countable dense subset of $[0,\infty)$ containing $0$. Let $h_{j,\lambda}$
be independent unit-variance Gaussian fields with covariance
\[
 K_{j,\lambda}((x,\tau),(x',\tau'))
 =\left(\frac{R(x,x')+\lambda Q(\tau,\tau')}{1+\lambda}\right)^j.
\]
These are positive semidefinite kernels: $R+\lambda Q$ is a Gram kernel and
integer powers preserve positive semidefiniteness by the Schur product
identity. The binomial expansion separates powers of $R$ and $Q$. Each power
of $Q$ has nonnegative increments along the cascade, so the fields can be
realized using independent child fields conditional on their ancestor fields.
Choose positive weights $w_{j,\lambda}$ with
$\sum w_{j,\lambda}^2j<\infty$, and set
\[
 H_N^{\rm pert}=s_N\sum_{j,\lambda}
       w_{j,\lambda}c_{j,\lambda}h_{j,\lambda},
 \qquad c_{j,\lambda}\in[1,2].
\]
In the sphere-only construction use just covariance kernels $R^j$.
The Gaussian covariance metric is continuous on each sphere, and the fields
are taken in their separable versions. Gaussian interpolation bounds the
pressure change by $Cs_N^2/N$.

Here is the concentration estimate needed for the perturbation argument:
\begin{equation}\label{eq:press-log-concentration}
 \E|\log Z-\E\log Z|\le C(\sqrt N+s_N).
\end{equation}
Without a cascade, Gaussian concentration gives this directly. Indeed the
gradient of $\sum_{\mu\le M}u(g_\mu(x))$ with respect to the row entries has
squared norm at most $M\|u'\|_\infty^2$, uniformly on the sphere. Any linear
Gaussian spin terms have squared reproducing-kernel norm $O(N)$, and the
perturbation has squared norm $O(s_N^2)$.

For a fixed finite cascade apply Lemma~\ref{lem:hier-cascade-tilt} to the
spin-integrated leaf partition marks. Conditional on shared root Gaussian
fields, the log partition is a recursive root shift plus the difference of
two log cascade masses. The latter have uniformly bounded second moments,
with a constant depending on the fixed cascade. Recursive logarithmic
moments preserve a uniform Lipschitz constant: if $|F(g,z)-F(g',z)|\le L\|g-g'\|$
for every child mark $z$, the same bound holds after applying
$m^{-1}\log\E_z e^{mF}$. Thus the root shift has Lipschitz constant
$C(\sqrt N+s_N)$ in its shared Gaussian fields. All child moments needed
here are finite, by bounded slope and Gaussian moments; at fixed dimension
one can also bound them by the supremum of the corresponding Gaussian
sphere process. This proves \eqref{eq:press-log-concentration}, uniformly in
interpolation time and after deletion of a fixed number of rows.

For completeness we derive the identities rather than invoke a separate
perturbation theorem. Fix one field $h$ with coefficient $swc$, where
$s=s_N$ and $w>0$ is fixed, and put $F(c)=\log Z(c)$. Then
\[
 F'(c)=sw\langle h\rangle,\qquad
 F''(c)=(sw)^2\langle(h-\langle h\rangle)^2\rangle.
\]
Gaussian integration by parts gives
\[
 \E\langle h\rangle=swc\,\E\langle K_{11}-K_{12}\rangle.
\]
Consequently integration over $c\in[1,2]$ bounds the mean thermal variance of
$h$ by a constant. The resulting integrated thermal absolute deviation is
$O(1)=o(s)$.

For the disorder deviation, the convex difference-quotient inequalities for
$F$ and $\E F$, at distance $d\in(0,1/4)$, give
\[
 \int_1^2\E|\langle h\rangle-\E\langle h\rangle|\,\dd c
 \le C_w\left\{\frac{\sqrt N+s}{sd}+sd\right\}.
\]
To justify the second term, integrate the increments of the monotone
function $(\E F)'$ over the enlarged interval $[1-d,2+d]$; their total is at
most $C_w d s^2$, before division by $sw$. The concentration bound applies
on this enlarged interval as well. Taking $d=N^{1/4-\rho}$ makes the last
display $o(s)$. In Gaussian integration by parts against a bounded replica
test $f$, one may therefore replace $h(x^1,\tau^1)$ by its disorder mean at
an error $o(s)$ after coefficient averaging. Division by $swc$ gives a
vanishing integrated absolute error in \eqref{eq:press-joint-gg}, first with
$\psi=K$.

For each total degree $j$, the polynomials $(R+\lambda Q)^j$ at $j+1$
distinct values of $\lambda$ span the homogeneous polynomials of degree $j$
in $(R,Q)$. Polynomial approximation gives all continuous $\psi$ and all
continuous finite-array tests; a monotone-class argument gives measurable
$f$. Choose countable determining families. The sum of their integrated
errors, with summable weights, tends to zero. A subsequence with summable
weighted errors has pointwise vanishing errors for almost every coefficient
vector. Including interpolation time and a countable list of deleted-row
systems in this extraction proves (iii). Compactness of the array space
then gives the asserted limiting identities.
\end{proof}

The only general ultrametricity theorem used here is the following
array formulation of Panchenko's theorem \cite[Theorem~1]{Panchenko2013}.
A weakly exchangeable positive semidefinite random array with fixed diagonal,
whose off-diagonal entries satisfy the Ghirlanda--Guerra identities for all
bounded measurable pair tests, is almost surely ultrametric:
\[
 R_{23}\ge\min\{R_{12},R_{13}\}.
\]
The equivalent random-measure formulation concerns independent samples from
a random probability measure on a separable Hilbert ball. Precisely, the
Gram-array representation \cite[Proposition~1]{Panchenko2010DS} says that a weakly
exchangeable positive semidefinite array with diagonal $1$ has a realization
\[
 R_{\ell\ell'}=v^\ell\cdot v^{\ell'}
       +\ind_{\{\ell=\ell'\}}(1-\|v^\ell\|^2),
\]
where, conditional on a random probability measure on the unit Hilbert ball,
the vectors $v^\ell$ are independent samples from that measure. To account for the latent diagonals, let $q^\dagger$ be the supremum of the
one-overlap support. The directing measure is almost surely supported on
$\|v\|^2=q^\dagger$; this is also
\cite[Theorem~2(a)]{Panchenko2010Sphere}. Here is the short argument needed
for this fact. If an overlap set $B$ has probability $c>0$, GG implies that
the probability of avoiding $B$ in all $n$ overlaps with replica $1$ equals
$\prod_{j=1}^n(1-c/j)$, which tends to zero. Conditional independent sampling
then shows that almost every directing point has positive mass of partners
with overlap in $B$. Apply this to $B=(t,q^\dagger]$, for every rational
$t<q^\dagger$. Since two directing points almost surely have inner product
at most $q^\dagger$, continuity and separability imply
$\|v\|^2\le q^\dagger$ throughout the support: a sufficiently small ball
around a counterexample would give larger inner products with positive
probability. Cauchy--Schwarz now excludes $\|v\|^2<q^\dagger$, since such a
point could not have partners with overlaps arbitrarily close to
$q^\dagger$. Thus the latent diagonals are deterministic, and the
random-measure GG identities follow from the array identities. Panchenko's
ultrametricity theorem applies. Adding a constant Gram kernel and rescaling
allows this reasoning for any bounded fixed diagonal before positivity is
known.

\begin{lemma}[Synchronization and identification]\label{lem:press-sync}
Suppose the joint array $(R,Q)$ is a subsequential limit from
Lemma~\ref{lem:press-gg}. Then $R\ge0$ off the diagonal. There are
nonnegative nondecreasing profiles $r,q$, on a common canonical rank array
$(U_{\ell\ell'})$, such that
\[
 (R_{\ell\ell'},Q_{\ell\ell'})
 \stackrel{\mathrm{law}}{=}
 (r(U_{\ell\ell'}),q(U_{\ell\ell'})),\qquad\ell\ne\ell'.
\]
If the finite trial cascade has prescribed profile $q$, this is that same
profile. In the sphere-only case the array is the RPC array determined by its
one-overlap marginal.
\end{lemma}

\begin{proof}
Apply ultrametricity separately to $R$, $Q$, and $R+Q$, using
\eqref{eq:press-joint-gg}. First, if $\Prob(R_{12}\le a)=m>0$ for some $a<0$,
ultrametricity makes $R_{ij}>a$ an equivalence relation on distinct replicas.
On the event that $n$ sampled replicas occupy distinct classes, GG gives
probability $(1-m)/n$ for the next replica to join any specified class.
These joining events are disjoint, so the probability of avoiding them all
is $m$. Induction gives positive probability of arbitrarily many pairwise
overlaps at most $a$. Positive semidefiniteness forbids this, since on that
event the sum of the entries of the $n\times n$ Gram matrix is at most
$n+n(n-1)a<0$ for sufficiently large $n$.

Let $\mu$ be the law of $(R_{12},Q_{12})$. If two points in its support were
strictly incomparable coordinatewise, choose disjoint neighborhoods $A,B$
with $R_A<R_B$ and $Q_A>Q_B$. GG with two old replicas gives
\[
 \Prob\big((R_{12},Q_{12})\in A, (R_{13},Q_{13})\in B\big)
 =\tfrac12\mu(A)\mu(B)>0.
\]
On this event scalar ultrametricity forces $R_{23}=R_{12}$ and
$Q_{23}=Q_{13}$. Thus $R_{23}+Q_{23}$ is strictly smaller than both
$R_{12}+Q_{12}$ and $R_{13}+Q_{13}$, contradicting ultrametricity of the sum.
The support of $\mu$ is therefore a chain. Each coordinate is a nondecreasing
function of the sum on that support, and the pair has comonotone quantiles.

It remains to identify the whole array, since a pair marginal alone would
otherwise be insufficient. At any threshold with lower mass $m$, GG says,
conditional on the old array, that the new replica joins an existing
ultrametric class containing $l$ old replicas with probability $(l-m)/n$.
For finitely many thresholds, the classes form a nested partition. These
joining probabilities, and their differences for a parent and its children,
determine the distribution of the new branch. Induction in the number of
replicas determines the law of the discretized array. By
Lemma~\ref{lem:hier-rank-array} it is exactly the finite RPC sampling law.
Refining thresholds at continuity points identifies the full scalar array.
Apply this to the sum and then to its two coordinate functions to obtain the
common rank representation.

Finally, at finite $N$ the cascade genealogy retains its prescribed law after
the spin-integrated tilt: all Gaussian fields in the model, including the
perturbation, have conditional child independence. This is precisely
Lemma~\ref{lem:hier-cascade-tilt}. Hence the limiting second marginal is the
prescribed trial profile.
\end{proof}

\begin{lemma}[Deleting and reinserting finitely many rows]
\label{lem:press-row-insertion}
Consider any convergent deleted-row overlap-array subsequence satisfying the
preceding lemmas. Reinsert finitely many independent Gaussian rows, with
potentials of bounded slope and with covariances given by continuous functions
of the available overlap kernels. Suppose those covariances are represented
by nondecreasing profiles on the common rank array. Then:
\begin{enumerate}[label=(\roman*)]
\item the entire common-rank genealogy retains its law after insertion;
\item conditional on this genealogy, separate inserted row-mark families are
independent and have their individually tilted RPC laws;
\item for a row of covariance profile $C$, the conditional mean of the product
of two terminal derivatives is $D_C(U_{12})$, with diagonal mean $D_{C,d}$.
\end{enumerate}
The conclusions include convergence of tilted polynomial-growth row
observables whenever their requisite moments are uniformly integrable.
In particular any convergent rowful array and any further limit with a fixed
finite set of its rows removed have the same profile.
\end{lemma}

\begin{proof}
Conditionally on the deleted-row Gibbs measure, the new row values sampled
on finitely many replicas are centered Gaussian vectors with the displayed
covariance arrays. They therefore converge jointly with the base arrays.
For every fixed $a<\infty$, their bounded-slope potentials obey
$\E\langle e^{a|u(g)|}\rangle\le C_a$ before insertion. The replica
truncation argument of Lemma~\ref{lem:hier-array-continuity} consequently
passes both the finite-row tilt and its observables to the limit.
For a finite RPC, claims (i)--(iii) are the marking transformation and
additive-mark factorization of Lemma~\ref{lem:hier-cascade-tilt}, together
with the derivative-martingale definition of $D_C$. In particular, the
marking transformation makes the transformed weight tree independent of
the separate row-mark families. Sampling replicas from those weights
therefore preserves their conditional factorization given the genealogy.
This is the depth-two calculation in the proof of
Lemma~\ref{lem:hier-cascade-tilt}, applied to the sum of the inserted-row
potentials; it uses no independence between a stable total and its
normalized jumps.

For the passage to general profiles, retain the canonical rank array itself
in the replica convergence, even on plateaus of the covariance profile.
The following conditional-kernel version of the truncation argument makes
the factorization passage precise. For a bounded continuous observable $F$
of the row marks at the first $n$ replicas, its conditional expectation
given their ranks is approximated in $L^1$ of the fixed rank law by the
finite-replica polynomial approximants to the truncated tilt numerator and
denominator. The truncation error is uniform over bounded profiles, by the
exponential moment bounds. At finite profile, invariance of the tilted rank marginal lets us
interpret these approximants as conditional kernels relative to the same
canonical rank law. A polynomial approximant uses a fixed finite number
of additional replicas. Integrating its Gaussian marks gives a bounded
function $K_j$ of all the ranks of those replicas. Under $C_j\to C$ in
$L^1$, the covariance entries converge almost surely along a subsequence;
Gaussian covariance continuity and dominated convergence give
$K_j\to K$ in $L^1$ of the fixed rank law. Averaging the additional ranks
conditional on the first $n$ ranks preserves this convergence, since
\[
 \big\|\E[K_j-K\mid(U_{ab})_{a,b\le n}]\big\|_1
       \le\|K_j-K\|_1.
\]
The conditional law of those additional ranks is independent of the
covariance profile. Thus the exact bounded conditional kernels converge in $L^1$.
Products of such kernels also converge in $L^1$. The finite-profile
factorization identities against any bounded cylinder test of the ranks
therefore pass to the limit. A monotone-class argument extends them to the
entire countable genealogy. This proves (i)--(iii) without presuming that
conditional independence is closed under weak convergence.

In particular, for two inserted rows $A,B$ with profile $C_j$ and for any
bounded cylinder test $H$ of the full rank array, the identity needed below is
\[
 \E\langle H a_A^1a_A^2a_B^1a_B^2\rangle
       =\E[H D_{C_j}(U_{12})^2].
\]
Its limit also follows directly from $D_{C_j}\to D_C$ in $L^1$ and bounded
slope. Single-replica analogues have constant right-hand factors.
Polynomial observables follow by truncation and the same exponential moment
bounds. The last assertion follows by first taking a joint subsequence of
the rowful and deleted-row laws, then representing the former as the tilt of
the latter. Claim (i) identifies their complete overlap-array laws.
\end{proof}

\subsection{The interpolation upper bound}

The comparison follows Guerra's interpolation scheme~\cite[\S4]{Guerra2003}:
we connect the interacting system to a hierarchical trial system and
control the sign of the overlap remainder. Here that sign follows from
the row-profile concavity proved in Section~\ref{sec:row-concavity}.
The synchronized cavity limits used for this nonlinear spherical model
are proved here; the SK bound itself is not an input for the present model.

Fix a finite trial profile $q$ reaching diagonal $1$, and put
$p=p_q=\alpha D_q$, with diagonal $p_d=\alpha D_{q,d}$. Since $D_q$ is constant on plateaus of
$q$, the off-diagonal field $p$ is a measurable function of the trial overlap
$Q$ itself.
Take independent fields $z_\mu(\tau)$, $z_i^p(\tau)$, and $Y(\tau)$ of
covariance profiles $q$, $p$, and $qp$, respectively. On the product of sphere
and cascade interpolate with exponent
\begin{equation}\label{eq:press-interpolating-exponent}
 \sum_{\mu\le M}u\big(\sqrt t\,g_\mu(x)+\sqrt{1-t}\,z_\mu(\tau)\big)
 +\sqrt{1-t}\sum_{i=1}^N x_i z_i^p(\tau)
 +\sqrt{tN}\,Y(\tau)+H_N^{\rm pert}.
\end{equation}
Write $\varphi_N(t)$ for its normalized expected log partition, including
coefficient averaging. For an individual row put
$a_\mu=u'(\sqrt t\,g_\mu+\sqrt{1-t}\,z_\mu)$.
Gaussian covariance differentiation cancels all diagonal terms and gives
\begin{equation}\label{eq:press-interpolation-derivative}
 \varphi_N'(t)
 =-\frac{M}{2N}\E_c\E\langle(R_{12}-Q_{12})a_\mu^1a_\mu^2\rangle
   +\frac12\E_c\E\langle(R_{12}-Q_{12})p(U_{12}^{\tau})\rangle.
\end{equation}
Here row exchangeability permits any fixed $\mu$, and $U^\tau$ denotes the
trial cascade rank. The last term is the combined derivative of the linear
spin and compensation fields: their covariance derivative is
$N(Q-R)p$, and their diagonal derivatives cancel.

To control a subsequential limit at almost every $t$ and coefficient vector,
remove row $\mu$. Lemmas~\ref{lem:press-gg} and \ref{lem:press-sync} give
comonotone cavity profiles $(r,q)$. The independent removed row has covariance
profile $C=tr+(1-t)q$ and diagonal $1$. Lemma~\ref{lem:press-row-insertion}
identifies the derivative limit as
\[
 -\frac\alpha2\int_0^1(r-q)D_C\,\dd s
 +\frac12\int_0^1(r-q)p\,\dd s.
\]
Concavity of $\Phi_u$ and \eqref{eq:src1} imply, for $t>0$,
\[
 \int_0^1(r-q)(D_C-D_q)\,\dd s\ge0.
\]
The limiting derivative is therefore nonpositive. For fixed bounded-slope
$u$, \eqref{eq:press-interpolation-derivative} is uniformly bounded. Starting
from any subsequence realizing the upper limit of the integrated derivative,
use the almost-everywhere extraction in Lemma~\ref{lem:press-gg}, and at each
such parameter-time pair use compactness to examine every further array
limit. Every resulting derivative limit is nonpositive. Reverse Fatou gives
$\limsup_N\int_0^1\varphi_N'(t)\,\dd t\le0$.

After removing the negligible perturbation, the endpoint $t=0$ separates into
$M\Phi_u(q)+NS_N(p)$ by additive-mark factorization. At $t=1$ the physical
sphere partition separates from the cascade compensation, whose expected
logarithm is
\[
 \frac N2\int_0^1s\,\dd(qp)(s)
 =\frac N2\left(p_d-\int_0^1qp\,\dd s\right).
\]
Thus, with $\mathcal S=\lim_N(S_N-p_d/2)$,
\begin{equation}\label{eq:src7}
 \limsup_{N\to\infty}P_N(u)
 \le\mathcal P_u(q)
 :=\alpha\Phi_u(q)+\mathcal S(p_q)
       +\frac12\int_0^1q p_q\,\dd s.
\end{equation}
Profile continuity extends this to every admissible $q$. The argument uses
profile concavity rather than concavity of the terminal itself, so it also
proves \eqref{eq:src7} for the small terminal perturbations of
Section~\ref{sec:row-concavity}.

\subsection{Identification of the empirical cavity fields}

Adding a spin coordinate couples its fresh Gaussian column to the vector
of row derivatives. Its conditional covariance between two replicas is
therefore a row average, rather than their original spin overlap.
The quadratic expansion also produces a scalar shift: the variance of
the new columns gives a term $u''$, while normalization of the old
projections gives $-g_\mu u'$. We identify both quantities before
comparing dimensions. In a sphere-only system write
$a_\mu(x)=u'(g_\mu(x))$ and
\[
 p^N_{\ell\ell'}=\frac1N\sum_{\mu\le M_N}
                    a_\mu(x^\ell)a_\mu(x^{\ell'}),\qquad
 b_N(x)=\frac1N\sum_{\mu\le M_N}
          \big[u''(g_\mu(x))-g_\mu(x)a_\mu(x)\big].
\]

The two-row deletion argument below identifies the first quantity as a
function of the limiting spin overlap and proves that the second has a
deterministic limit under the joint disorder and replica law.

\Needspace{4\baselineskip}
\begin{lemma}[Empirical cavity coefficients]\label{lem:press-empirical}
Take a subsequence on which the base overlap array converges to profile $q$
and the GG errors for the base, one-row-deleted, and two-row-deleted systems
vanish. In every joint array limit,
\begin{equation}\label{eq:src8}
 \begin{aligned}
 p_{\ell\ell'}&=\alpha D_q(U_{\ell\ell'}),&&\ell\ne\ell',\\
 p_d&=\alpha D_{q,d},\qquad
 b_N(x^\ell)\longrightarrow b=-p_d+\int_0^1q p_q\,\dd s.
 \end{aligned}
\end{equation}
In particular,
$\E\langle|b_N(x^\ell)-b|^2\rangle\to0$ along this subsequence.
\end{lemma}

\begin{proof}
For a continuous function $F$ of one overlap, remove a specified row and
reinsert it using Lemma~\ref{lem:press-row-insertion}. Row exchangeability gives
\[
 \lim_N\E\langle p^N_{12}F(R_{12})\rangle
 =\alpha\int_0^1D_q(s)F(q(s))\,\dd s.
\]
On expanding $(p^N_{12})^2$, the terms with equal row indices have total size
$O(N^{-1})$, since $u'$ is bounded. For the distinct indices, delete both
rows and use conditional factorization of their tilted mark families. Hence
\[
 \lim_N\E\langle(p^N_{12})^2\rangle
 =\alpha^2\int_0^1D_q(s)^2\,\dd s.
\]
The function $D_q$ is constant almost everywhere on each plateau of $q$.
Thus there is a bounded Borel function $d_q$ such that
$D_q(s)=d_q(q(s))$ almost everywhere. Approximate $d_q$ in $L^2$ of the overlap
marginal by continuous functions. The preceding two displays imply that any
joint limit $p_{12}$ satisfies
$\E(p_{12}-\alpha d_q(R_{12}))^2=0$. Exchangeability gives all off-diagonal
entries simultaneously.

The same one-row and two-row calculation for the single-replica mark
$a_\mu(x)^2$ identifies the diagonal limit. With only one sampled replica
there is no variable pair genealogy, so its conditional mean is the scalar
$D_{q,d}$. Apply the calculation also to
$u''(g_\mu)-g_\mu u'(g_\mu)$. This observable has at most linear growth.
Its moments after inserting one or two rows are uniformly bounded by
reciprocal Jensen, Gaussian exponential moments, and bounded slope; this
justifies truncation in the preceding argument and gives a deterministic
$L^2$ limit for $b_N$.

To compute that limit, Gaussian integration by parts in a single row gives
\[
 \E\langle g_\mu u'(g_\mu)\rangle
 =\E\langle u''(g_\mu)+u'(g_\mu)^2\rangle
  -\E\langle R_{12}u'(g_\mu^1)u'(g_\mu^2)\rangle.
\]
Multiplication by $M_N/N$ and the already identified row limits give exactly
$b=-p_d+\int q p_q$. All deletions and limits take place at a fixed admissible
coefficient vector; there is no mixture-of-profiles step.
\end{proof}

\subsection{The dimension-increment lower bound}

We compare systems of dimensions $N$ and $N+m$, using the block-cavity
organization of Aizenman--Sims--Starr
\cite[Theorem~1(ii), equations~(27)--(33)]{AizenmanSimsStarr2003} and
its spherical implementation by Chen~\cite[Theorem~2.1]{Chen2013SphericalASS}.
For this row Hamiltonian, the preceding lemma supplies the empirical
covariance identification; the expansion below supplies the corresponding
insertion and normalization terms. The spherical-shell comparison, nonlinear
and fresh-row estimates, and telescoping argument without superadditivity
are supplied locally below. Throughout the cavity expansion, $m$ and the shell width $\delta$ are fixed while $N\to\infty$.
The preceding lemma will identify the limiting insertion of these $m$
coordinates. Only after that limit do we let $m\to\infty$ to recover the
spherical functional, and finally let $\delta\downarrow0$.

\paragraph{Sphere coordinates and the cavity expansion.}
Let $A_N$ be the expected log partition including the sphere-only perturbation
and the coefficient average. Fix a positive integer $m$, put
$k=M_{N+m}-M_N=\alpha m+O(1)$, and fix $\delta\in(0,1)$.
Use coordinates
\[
 X=\left(\sqrt{\frac{N+m-\|\eps\|^2}{N}}\,x,\eps\right),
 \qquad x\in\sqrt N\mathbb S^{N-1},\quad\eps\in\R^m,
\]
and restrict $\eps$ to
$\mathcal A_{m,\delta}=\{m(1-\delta)\le\|\eps\|^2\le m(1+\delta)\}$.
Write $G_m\sim\cN(0,I_m)$ for a standard Gaussian vector.
At fixed $m,\delta$ the marginal surface density of $\eps$ converges relatively
uniformly on this compact shell to standard Gaussian density. The angular
law of $x$ is independent uniform surface measure.

The $N+m$ perturbation can be replaced here by the same $N$-spin perturbation
used in $A_N$. For any two shell configurations their full overlaps differ
from the old overlaps by $O_m(N^{-1})$. Since
$|r^j-r'^j|\le j|r-r'|$ on $[-1,1]$, covariance interpolation and
$\sum w_j^2j<\infty$ bound the change of expected log partition by
$O_m(N^{2\rho-1})=o(1)$. This includes the change from $s_{N+m}$ to $s_N$.

Let $B_{\mu i}$ be the new independent Gaussian columns. For each old row,
its new argument is
\[
 \sqrt{1-\frac{\|\eps\|^2}{N+m}}\,g_\mu(x)
      +\frac{B_\mu\cdot\eps}{\sqrt{N+m}}.
\]
Taylor expansion of its potential around $g_\mu(x)$ yields, after summing rows,
\begin{equation}\label{eq:press-cavity-expansion}
 \sum_{i=1}^m\eps_i z_i(x)+\frac12\|\eps\|^2b_N(x)+T(x,\eps)+o(1),
 \qquad z_i(x)=\frac1{\sqrt N}\sum_{\mu\le M_N}a_\mu(x)B_{\mu i},
\end{equation}
where
\[
 T(x,\eps)=\frac1{2N}\sum_{\mu\le M_N}u''(g_\mu(x))
       \big[(B_\mu\cdot\eps)^2-\|\eps\|^2\big].
\]
Writing $\operatorname{Rem}_N(x,\eps)$ for the $o(1)$ term in
\eqref{eq:press-cavity-expansion}, we have, for fixed $m,\delta$,
\[
 \E\sup_{\substack{x\in\mathbb S_N\\\eps\in\mathcal A_{m,\delta}}}
       |\operatorname{Rem}_N(x,\eps)|\longrightarrow0.
\]
To check this directly, use
$\sum_\mu g_\mu(x)^2\le\|A\|_{\rm op}^2$ and
$\E\|A\|_{\rm op}^j=O(N^{j/2})$, together with
$\E\sum_\mu\|B_\mu\|^j=O_m(N)$ for $j\le3$. For the cubic remainder, write the argument increment as
$d_\mu=c_N g_\mu+B_\mu\cdot\eps/\sqrt{N+m}$, where
$c_N=-\|\eps\|^2/(2N)+O_m(N^{-2})$. Then
\[
 \sum_\mu|d_\mu|^3
 \le C_m\left(N^{-3}\|A\|_{\rm op}^3
                   +N^{-3/2}\sum_\mu\|B_\mu\|^3\right).
\]
This uses $|a+b|^3\le4(|a|^3+|b|^3)$, so no mixed fourth moment is needed.
The quadratic mixed term is at most
\[
 C_m N^{-3/2}\|A\|_{\rm op}
                    \left(\sum_\mu\|B_\mu\|^2\right)^{1/2},
\]
whose mean is $O_m(N^{-1/2})$. The shrinkage-square term is bounded by
$C_mN^{-2}\|A\|_{\rm op}^2=O_{L^1,m}(N^{-1})$, and denominator replacements
use only the first two displayed column moments. These estimates prove the
claimed uniform remainder. The operator-norm
estimate follows from finite bilinear sphere nets and Gaussian tail bounds.
The $k$ fresh rows can similarly be replaced by independent old-dimensional
fields $g_a^{\rm new}(x)$: their new-column contribution is $O_m(N^{-1/2})$
and their shrinkage error is bounded by $O_m(N^{-1})$ times a row norm.

\paragraph{The centered quadratic term.}
For the lower bound we may discard $T$ at expected cost $o(1)$, although its
supremum need not be small. Here is the required estimate. Condition the
Gaussian shell measure to be a probability, and let $\langle\cdot\rangle_0$
include it and the old Gibbs measure. Put
\[
 W(x,\eps)=\eps\cdot z(x)+\tfrac12\|\eps\|^2b_N(x)
                    +\sum_{a=1}^k u(g_a^{\rm new}(x)).
\]
Jensen's inequality gives
\[
 \log\langle e^{W+T}\rangle_0-\log\langle e^W\rangle_0
 \ge-\langle|T|\rangle_W,
 \qquad
 \E\langle|T|\rangle_W
 \le\E\langle |T_1|e^{W_1-W_2}\rangle_0.
\]
The second inequality uses
$\langle e^W\rangle_0^{-1}\le\langle e^{-W}\rangle_0$.
Conditional on the old disorder and $(x,\eps)$, the centered Gaussian squares
in $T$ are independent across rows; hence
$\E\langle T_1^2\rangle_0=O_m(N^{-1})$.
Furthermore $\E\langle e^{a(W_1-W_2)}\rangle_0$ is bounded for every fixed
$a,m,\delta$. Indeed the Gaussian fields $z_i$ have covariance bounded by
$(M_N/N)\|u'\|_\infty^2$, fresh row potentials have linear growth, and
\[
 |b_N(x)|\le C\left(1+\frac{\|A\|_{\rm op}}{\sqrt N}\right).
\]
Cauchy--Schwarz proves the desired $o(1)$ bound. Here is the lower bound needed later for coefficient averaging. For fixed
coefficients put $F_N(c)=\E\log Z_{N+m}(c)-\E\log Z_N(c)$. All the preceding
sphere-density, perturbation, Taylor, and fresh-row errors are $o(1)$ uniformly
in $c$, with $m,\delta$ fixed. Before tilting, $\E\langle z_i\rangle_0=0$
and $\E\langle T\rangle_0=0$, while the fresh-row potential means are finite
constants and $\E\langle b_N\rangle_0\ge-C$. Jensen therefore gives
\[
 F_N(c)\ge \log\Prob(G_m\in\mathcal A_{m,\delta})-C_{m,\delta}-o(1),
\]
uniformly in $c$. This is a bound on the actual dimension increment,
not only on its limiting RPC expression.

\paragraph{The limiting insertion.}
At fixed $m$, the row count $k$ has at most two possible values; pass first
to a dimension subsequence on which it is constant. Now fix an admissible
coefficient vector and extract a base-array limit.
Conditionally on the old disorder, the $z_i$ are independent centered Gaussian
fields with covariance array $p^N$. Lemma~\ref{lem:press-empirical} identifies
that array and $b_N$, while the fresh rows have covariance array $R$.
To handle the unbounded retained observable $b_N$, set
$K_N=C(1+\|A\|_{\rm op}/\sqrt N)$. Clipping $b_N$ to $[-L,L]$ changes $W$
uniformly in $(x,\eps)$ by at most
$C_mK_N\ind_{\{K_N>L\}}$. Thus the expected log integral changes by at most
$C_m\E[K_N\ind_{\{K_N>L\}}]$, which tends to zero uniformly in $N$ and
$c$. Apply Lemma~\ref{lem:hier-array-continuity} with this bounded clipping,
then let $L\to\infty$. Together with the exponential bounds for the Gaussian
terms, this passes the logarithmic increment to its RPC limit. The independent fresh-row contributions separate, giving
$k\Phi_u(q)$. If $p$ has a diagonal residual, its independent Gaussian
convolution contributes precisely half that variance times $\|\eps\|^2$.
Thus the remaining integral is the radial Gaussian version of the spin
functional with covariance $p=p_q$ and scalar shift $b$.

On the shell, write $\eps=r\omega$, where
$\omega\in\sqrt m\mathbb S^{m-1}$ and $r^2\in[1-\delta,1+\delta]$.
Let $Z_{\rm sph}(r)$ be the cascade sum of the spherical exponential
integrals at radius $r$, including any residual-diagonal factor, and let
$Z_{\rm shell}$ be the corresponding Gaussian shell integral including
the scalar shift $b\|\eps\|^2/2$. For each fixed field realization,
$Z_{\rm sph}(r)$ increases with $r\ge0$ by symmetry of $\omega$ and
nonnegativity of the residual variance. Thus, with $r_- =\sqrt{1-\delta}$,
\[
 Z_{\rm shell}\ge
 \Prob(G_m\in\mathcal A_{m,\delta})
 \exp\{mb/2-|b|m\delta/2\}\,Z_{\rm sph}(r_-).
\]
The expected logarithm of the last factor is
$mS_m((1-\delta)p;(1-\delta)p_d)$. Covariance interpolation changes it
to $mS_m(p;p_d)$ at cost at most $C\delta m$, uniformly over bounded
profiles $p$. Adding the fresh-row contribution $k\Phi_u(q)$ and using
$|k-\alpha m|\le1$ and the uniform bound on $|\Phi_u(q)|$ gives
\begin{equation}\label{eq:press-increment-lower}
 \begin{split}
 \liminf_{N\to\infty}(A_{N+m}-A_N)
 &\ge \inf_q\{\alpha m\Phi_u(q)+mS_m(p_q)+mb(q)/2\}\\
 &\quad -C(1+\delta m)+\log\Prob(G_m\in\mathcal A_{m,\delta}),
 \end{split}
\end{equation}
with $C$ independent of $m,\delta$. The uniformity follows from
boundedness of $p_q$, $b(q)$, and $\Phi_u(q)$.
More explicitly, take a subsequence realizing the left lower limit; extract
GG errors almost everywhere in the coefficients for the base and both row
deletions; then take pointwise array subsequences. Every resulting functional
is bounded below by the displayed infimum. The uniform finite lower bound
on increments permits Fatou in the coefficients. This establishes
\eqref{eq:press-increment-lower} without requiring a common profile for
different coefficient vectors or different dimensions.

\paragraph{From increments to pressure.}
The insertion bound still contains the finite-dimensional spin functional
$S_m$. Telescoping the dimension increments and then sending $m$ to infinity
will replace it by $\mathcal S$, giving the trial functional in the upper
bound. To carry this out,
for any real sequence $(A_N)$ with the present linear bounds, telescoping
separately in the $m$ residue classes gives
\[
 \liminf_{N\to\infty}\frac{A_N}{N}
 \ge\liminf_{N\to\infty}\frac{A_{N+m}-A_N}{m}.
\]
Indeed, for any number below the right side, every sufficiently late increment
is at least $m$ times that number, and summing proves the inequality.
Substitute $b=-p_d+\int qp_q$, use the uniform convergence
$S_m-p_d/2\to\mathcal S$ of Proposition~\ref{prop:hier-spherical-limit}, and
divide the error in \eqref{eq:press-increment-lower} by $m$.
For fixed $\delta$,
$\Prob(G_m\in\mathcal A_{m,\delta})\to1$ as $m\to\infty$. Sending first
$m\to\infty$ and then $\delta\downarrow0$ proves
\[
 \liminf_{N\to\infty}P_N(u)\ge\inf_q\mathcal P_u(q).
\]
Together with \eqref{eq:src7}, this proves the pressure formula for smooth
concave caps.

\subsection{Removing the cap and identifying the minimizing profile}

The uniform tail comparison of Section~\ref{sec:row-concavity} removes the
cap. For clarity, the mechanism is row insertion: conditional on a cavity
spin, a fresh Gaussian projection is standard normal; insertion with an
increasing potential decreases expectations of decreasing nonnegative tail
functions. The error between a tangent cap below $-L$ and
$u=-\beta(-1-z)_+^2/2$ is bounded by a Gaussian negative-tail quadratic
expectation tending to zero with $L$. A bounded terminal perturbation only
multiplies this bound by $e^{2|d|\|\psi\|_\infty}$.
The same estimate applies to the RPC row functional and to its tilted
derivative moments. Consequently $\Phi_u$, $D_q$, and $\mathcal P_u$ have
uniform cap limits on the compact space of monotone profiles, with its $L^1$
topology. We obtain
\[
 P(u):=\lim_{N\to\infty}P_N(u)=\min_q\mathcal P_u(q).
\]
Continuity and profile compactness give the stated minimum.

\begin{proposition}[Entropy form of the pressure]\label{prop:press-formula}
For Gaussian rows, $u=-\beta V$, and $\alpha>0$,
\begin{equation}\label{eq:src9}
 P(u)=\min_{q:\,q_*<1}\{\alpha\Phi_u(q)+E(q)\},
 \qquad q_*:=\esssup q.
\end{equation}
Every minimizer of $\mathcal P_u$ has $q_*<1$ and satisfies the spin inverse
relations of Section~\ref{sec:hierarchy} with $p=p_q$.
\end{proposition}

\begin{proof}
Let $q$ minimize $\mathcal P_u$. By
Proposition~\ref{prop:hier-spin-concavity}, $\mathcal S$ has at $p_q$ a
supergradient $-r/2$, where $r$ is a nonnegative nondecreasing profile bounded
strictly below $1$ and $p_q=\mathcal I(r)$ for the inverse map
\eqref{eq:press-continuous-inverse}. Put $q_t=q+t(r-q)$ and $p_t=p_{q_t}$.
The supporting inequality for the concave row functional and the
supergradient inequality for $\mathcal S$ give
\[
 \begin{split}
 \alpha\big[\Phi_u(q_t)-\Phi_u(q)\big]
   &\le-\frac t2\int(r-q)p_q,\\
 \mathcal S(p_t)-\mathcal S(p_q)
   &\le-\frac12\int r(p_t-p_q).
 \end{split}
\]
Adding the exact difference of the coupling terms yields
\[
 \mathcal P_u(q_t)-\mathcal P_u(q)
 \le-\frac{1-t}{2}\int(r-q)(p_t-p_q).
\]
By \eqref{eq:src6}, for $0<t<1$ the integral on the right is at least
$c_{\alpha,\beta}t\|r-q\|_2^2$. Minimality therefore forces $r=q$.
In particular $q_*<1$ and $p_q=\mathcal I(q)$, proving the asserted spin
inverse relation for the original profile $p_q$. The equality of this spin
supergradient with $-q/2$
means that $p_q$ attains the entropy conjugate:
\[
 E(q)=\mathcal S(p_q)+\frac12\int q p_q.
\]
For any competing profile $q'$ with $q'_*<1$, the definition of $E$ gives
$\mathcal P_u(q')\le\alpha\Phi_u(q')+E(q')$. Equality at the minimizing $q$
proves \eqref{eq:src9}.
\end{proof}

One useful endpoint estimate follows directly from minimality. Let $a_0$
be the root first derivative before the initial Gaussian variance is added,
and recall $\chi_q(0)=\int_0^1(1-q(s))\,\dd s$. Then
\begin{equation}\label{eq:src10}
 \alpha a_0^2\chi_q(0)\le\frac1{\chi_q(0)}-1.
\end{equation}
Indeed, for $a\ge0$ small, set $q^{(a)}=a+(1-a)q$. Formula~\eqref{eq:src3}
gives the exact entropy change
\[
 E(q^{(a)})-E(q)=\tfrac12\log(1-a)
                +\frac{a}{2(1-a)\chi_q(0)}.
\]
The row derivative from \eqref{eq:src1} is
$-\alpha\int(1-q)D_q/2\le-\alpha a_0^2\chi_q(0)/2$, because the derivative
martingale gives $D_q\ge a_0^2$. The nonnegative right derivative of the
minimized objective is thus at most half the right side of
\eqref{eq:src10} minus half its left side, proving the claim.

\subsection{Touching supports for row observables}

The pressure formula concerns the base quadratic loss. To identify row
observables, we also need an upper bound for perturbed terminals that
agrees with the base pressure at zero perturbation. No pressure formula
for the perturbed model is required.

Fix $\psi\in C_c^\infty(\R)$. For $|d|$ sufficiently small,
Corollary~\ref{cor:conc-hinge-stability} and its regularized proof give
profile concavity for $u+d\psi$ and all its approximations. The interpolation upper bound applies at any base
minimizer $q$. With $p=\alpha D_{q,u+d\psi}$, the entropy conjugate bounds
$\mathcal S(p)+\int qp/2\le E(q)$. After removing the cap this gives the
upper support inequality
\begin{equation}\label{eq:src11}
 \limsup_{N\to\infty}P_N(u+d\psi)
 \le\alpha\Phi_{u+d\psi}(q)+E(q).
\end{equation}
The permissible size of $d$ may depend on $\beta$ and $\psi$, as required
for the terminal perturbation argument.


\Needspace{4\baselineskip}
\begin{corollary}[Upper bound for convex loss perturbations]
\label{cor:press-convex-loss-upper}
Let $\psi\in C_c^\infty(( -\infty,0))$. There is $w_0>0$, independent
of $\beta$, such that for both signs of $w$ with $|w|<w_0$ the loss
$L_w(z)=V(z)+w\psi(1+z)$ is nonnegative, convex, and nonincreasing. For every
fixed $\beta>0$ and every fixed profile $q$ with $q_*<1$,
\begin{equation}\label{eq:press-convex-loss-upper}
 \limsup_{N\to\infty}P_N(-\beta L_w)
 \le\alpha\Phi_{-\beta L_w}(q)+E(q).
\end{equation}
More generally this upper bound holds for any nonpositive, concave,
nondecreasing terminal of at most quadratic growth which agrees with a
negative quadratic hinge outside a compact set.
\end{corollary}

\begin{proof}
The support of $\psi(1+z)$ is a compact subset of $z<-1$, where $V''=1$,
$V>0$, and $V'<0$. Taking $|w|$ sufficiently small preserves all three
asserted inequalities, including at the boundary of that support.
The threshold for $w$ is a property of the loss and does not depend on
$\beta$. For the fixed concave terminal $-\beta L_w$, take tangent caps
and concave mollifications. Ordinary concave-terminal profile concavity
from Section~\ref{sec:row-concavity} gives the interpolation upper bound
for each such cap. The entropy conjugate bounds its spin terms by $E(q)$,
so the right side is simply $\alpha\Phi(q)+E(q)$ at this fixed trial.
Mollification errors vanish uniformly in terminal value for a fixed cap.
Beyond the compact perturbation support the cap error is the original
negative Gaussian quadratic-tail error, controlled by monotone row
insertion both microscopically and in the cascade. Sending the cap to
infinity proves the displayed bound. Only pressure and row-functional
values are passed to the limit; no derivative-profile uniformity in $w$
and no complete pressure formula for the perturbed loss are required.
\end{proof}


\section{Universality at fixed temperature}\label{sec:universality}

Write $\nu_\eps=\frac12\cN(0,1-\eps)+\frac12\cN(0,1+\eps)$, and let
$P_{N,\eps}(w)$ denote the pressure when the entries of $A$ have law
$\nu_\eps$.  Thus $P_{N,0}=P_N$. Surface measure on each sphere is normalized to
have mass one.  We use unit spins $y=x/\sqrt N$ when comparing
matrix entries.  Fix $0<B_0<\infty$ and $u(z)=-\frac\beta2(-1-z)_+^2$.

\begin{theorem}[Pressure and upper supports]\label{thm:univ-pressure}
There is $\eps_0(B_0)>0$, independent of $\beta>0$, with the following
property.  If $0<\alpha\le B_0$ and $0\le\eps<\eps_0(B_0)$, then
\[
 \lim_{N\to\infty}P_{N,\eps}(u)=P(u).
\]
If $q$ is any Gaussian minimizer in \eqref{eq:src9} and
$\psi\in C_c^\infty(\R)$, then, for all sufficiently small $|d|$,
\[
 \limsup_{N\to\infty}P_{N,\eps}(u+d\psi)
 \le \alpha\Phi_{u+d\psi}(q;1)+E(q).
\]
The allowed size of $d$ can depend on $\beta,\psi$ and $B_0$.
\end{theorem}

The principal issue is that a unit spin may place a positive fraction of
its squared norm on a few coordinates. Moment matching controls the
remaining coordinates, but gives no vanishing error for these large ones.
We therefore bound separately the slices on which their values are fixed.
A slice with squared norm $t$ in the large coordinates loses spherical
entropy. Its rows acquire a random variance change of size at most
$\eps t$. The decisive estimate compares both effects on the same scale,
$t/\chi_q(0)$, with constants independent of temperature. This permits
one choice of $\eps_0(B_0)$ before any later zero-temperature limit.
After proving the pressure bounds, we use their derivatives in $d$ to
identify the empirical row law and the normalized-force observables.

\subsection{Concentration, truncation, and replacement}

We record the quantitative elementary estimates needed below.  Constants
in this subsection may depend on a fixed smooth capped potential, on
$\beta$, and on $B_0$, but are uniform in $N$ and $0\le\eps\le1/2$.
Every entry is subgaussian with variance proxy at most $3/2$.  Applying the
exponential Markov inequality to a fixed bilinear form and taking
$1/4$-nets of the two unit spheres gives
\begin{equation}\label{eq:univ-opnorm}
 \Prob\{\|A\|_{\mathrm{op}}>K\sqrt N\}
 \le 2\exp(C(M+N)-cK^2N).
\end{equation}
Indeed the nets have at most $9^{M+N}$ pairs, each bilinear form has a
subgaussian tail, and their maximum controls the operator norm within an
absolute factor.  Integrating this estimate also gives every fixed moment
bound for $\|A\|_{\mathrm{op}}/\sqrt N$.

A mixture entry is a Lipschitz function of a standard Gaussian.  To see
this, let $T=F_\eps^{-1}\circ\Phi$, where $F_\eps$ is the mixture
CDF and $\Phi$ the standard Gaussian CDF, and put
$\sigma_+=\sqrt{1+\eps}$.  For $z\ge0$,
$0\le T(z)\le\sigma_+z$, and the mixture density $f_\eps$ satisfies
\[
 f_\eps(T(z))\ge \frac1{2\sigma_+}
       \phi(T(z)/\sigma_+)\ge\frac{\phi(z)}{2\sigma_+}.
\]
Consequently $T'\le2\sigma_+$; symmetry gives the same bound on the
negative half-line.  Product Gaussian concentration therefore applies to
a Lipschitz function of the entries with only this absolute change in its
Lipschitz constant.

For a potential $w$ with bounded first three derivatives and any
nonrandom nonzero finite measure $\lambda$ supported on unit spins, set
\[
 L(A)=\log\int\exp\Big\{\sum_{\mu=1}^M w(A_\mu y)\Big\}\,\lambda(\dd y).
\]
The total mass of $\lambda$ need not be one.  Differentiation gives
$\|\nabla_A L\|_F\le\|w'\|_\infty\sqrt M$.  If
$|y_i|\le b_N$ on the integration domain, differentiating three times in
$A_{\mu i}$ gives
\begin{align*}
 \partial_{\mu i}^3L={}&\langle w'''(g_\mu)y_i^3\rangle
 +3\Cov\big(w''(g_\mu)y_i^2,w'(g_\mu)y_i\big)\\
 &+\left\langle\big(w'(g_\mu)y_i-
             \langle w'(g_\mu)y_i\rangle\big)^3\right\rangle .
\end{align*}
Its absolute value is at most $C_wb_N^3$.  Taylor's theorem with integral
remainder, replacing entries one at a time, thus shows that two entry
laws with matching first two moments and uniformly bounded third absolute
moments satisfy
\begin{equation}\label{eq:univ-replacement}
 |\E L(A)-\E L(\widetilde A)|\le C_wMN b_N^3.
\end{equation}
This is the smooth-function Lindeberg replacement argument
\cite[Theorem~1.1]{Chatterjee2006Lindeberg}. The coordinate bound $b_N$
is what makes its error small; the localization argument below handles
spins for which that bound fails. The estimate is unchanged by adding
arbitrary, fixed row biases, and also holds conditionally on such biases.

Here and below the hinge can be capped and smoothed as in
Section~\ref{sec:row-concavity}.  More explicitly, writing $r=(-1-z)_+$,
replace $-\beta r^2/2$ for $r>L$ by its tangent
$-\beta Lr+\beta L^2/2$.  The difference from $u$ is
$\frac\beta2(r-L)_+^2$.  Mollification with a symmetric kernel, followed
by a constant of size
$O(\beta\delta^2)$, gives a smooth upper approximation with bounded first
three derivatives and arbitrarily small smoothing error.  The pressure
error of the cap tends to zero uniformly in $N$, including on any fixed
restricted spin domain.  For completeness, insert the rows one at a time.
For a fixed row and any cavity measure, weighting by a decreasing function
of $r$ decreases the mean of any increasing function of $r$.  This is the
identity
\[
 2\Cov(a(R),b(R))
 =\E[(a(R)-a(R'))(b(R)-b(R'))]\le0
\]
for increasing $a$, decreasing $b$, and independent copies $R,R'$.
Interpolate between the cap and the hinge and apply this observation to
the cap error.  Before row insertion, the fresh projection conditional on
the spin is Gaussian with a random variance in $[1-\eps,1+\eps]$.
The expected error is therefore bounded by
$\frac\beta2\sup_{v\in[1/2,3/2]}\E[(-1-\sqrt v Z-L)_+^2]$ per row,
which tends to zero.  Addition of the bounded test $d\psi$ changes this
bound by at most $e^{2|d|\|\psi\|_\infty}$.  For the row functional, interpolate the terminal potentials in the same
way.  The derivative is the tilted expectation of the cap error.  Every
interpolating base terminal has nonnegative first derivative, so its
tilted diffusion has nonnegative drift; driven by the same Brownian
motion, its endpoint is at least the Brownian endpoint.  Its residual
therefore has no larger upper tail than the fresh Gaussian residual.
The same Gaussian tail bound applies, uniformly over profiles and
diagonals in $[1/2,3/2]$.  The bounded-terminal-perturbation comparison
proved below supplies the factor $e^{2|d|\|\psi\|_\infty}$ for row
functionals as well.

Choose once and for all
\begin{equation}\label{eq:univ-a-range}
  \frac13<a<\frac12,\qquad
  \mathcal D_N=\{y\in S^{N-1}:\max_i|y_i|\le N^{-a}\}.
\end{equation}
For Gaussian disorder, rotation invariance implies that the
\emph{disorder-averaged} Gibbs spin is uniform on the sphere.  Spherical
coordinate tails and a union bound give
$\E\langle\ind_{\mathcal D_N^c}\rangle
\le CN\exp(-cN^{1-2a})$.  In particular the probability that the Gibbs
mass of $\mathcal D_N$ is below $1/2$ is exponentially small on this scale.
On the complementary event the loss in log partition is at most $\log2$.
On the exceptional event, Jensen's inequality on the restricted surface
measure bounds the restricted log partition below by
$\log\sigma_N(\mathcal D_N)-C_\beta(M+\|A\|_{\mathrm{op}}^2)$.
The full log partition has the same polynomial moment bound.  By
Cauchy--Schwarz and \eqref{eq:univ-opnorm}, the expected restriction loss
is $o(N)$.  This argument applies to the hinge and to its caps.
For a fixed smooth cap, \eqref{eq:univ-replacement} costs
$O(N^{2-3a})=o(N)$ on $\mathcal D_N$.  Removing the cap afterwards proves
\begin{equation}\label{eq:univ-lower}
 \liminf_{N\to\infty}P_{N,\eps}(u)\ge P(u).
\end{equation}

\subsection{Reduction of the upper bound to a localized slice}

Fix a smooth upper cap $w$ of $u+d\psi$.  For a unit spin define its
large-coordinate set
$I(y)=\{i:|y_i|>N^{-a}\}$.  Its cardinality is at most $N^{2a}$.
The number of possible sets is at most
$\exp(CN^{2a}\log N)=\exp(o(N))$.  Fix such a set $I$ of size $k$,
write $c=y_I$, $t=\|c\|^2$, and set $n=N-k$.  Conditional on $c$, the
remaining coordinates have the form $\sqrt{1-t}\,z$ with
$z\in S^{n-1}$.  For $k>0$ the marginal density of $c$ is
\[
 a_{N,k}(1-t)^{(n-2)/2}\ind_{\{t<1\}},\qquad
 a_{N,k}=\frac{\Gamma(N/2)}{\pi^{k/2}\Gamma(n/2)}.
\]
Stirling's formula gives $|\log a_{N,k}|=O(k\log N)=o(N)$.
For $k=0$ there is one slice, $c=0,t=0$.

Fix $\eta>0$.  The total contribution of $t>1-\eta$, after summing over
$I$, has normalized log at most
$\frac12\log\eta+B_0\sup w+o(1)$, since ball volumes and density
prefactors contribute $o(N)$.  The lower bound
$P(u)\ge- C\beta B_0$ follows from Jensen's inequality, and the desired
upper-support value is at least $P(u)-\alpha|d|\|\psi\|_\infty$.
Hence $\eta$ can be chosen, at the fixed temperature and test perturbation
under consideration, so that this part lies strictly below that value.  Dependence of $\eta$ on
$\beta$ will have no bearing on $\eps_0$.

For $t\le1-\eta$, the original restriction on the remaining coordinates
implies $|z_i|\le\eta^{-1/2}N^{-a}$.  Enlarge their integration domain
to this fixed restriction, independently of $c$.  Cover the $c$-ball by
a Euclidean $N^{-2}$-net, retaining centers with
$\|c\|^2\le1-\eta/2$.  Across all $I$ the number of centers is still
$\exp(O(N^{2a}\log N))=\exp(o(N))$.
On $\|A\|_{\mathrm{op}}\le K\sqrt N$, changing a center by $h$
changes the row arguments in Euclidean norm by at most
$C_\eta K\sqrt N\|h\|$.  Cauchy--Schwarz over the rows shows that the
log integral changes by at most $C_{w,\eta,K}N\|h\|$.  The logarithm
of the marginal density changes by at most $C_\eta N\|h\|$ as well.
Both errors are $o(N)$ on these nets.

For each center, the restricted log integral is
$C_w\sqrt N$-Lipschitz in the matrix.  Gaussian concentration through
$T$ and a union bound make this uniform over all centers.  Explicitly, put
$\delta_N=N^{-(1-2a)/4}$.  The probability that any centered log integral
exceeds $N\delta_N$ in absolute value is at most
\[
 2\exp\{C N^{2a}\log N-cN^{(1+2a)/2}\}\longrightarrow0.
\]
Conditional on the localized columns, replace the remaining
entries by Gaussians using \eqref{eq:univ-replacement}; their coefficients
are bounded by $C_\eta N^{-a}$, so the error is $o(N)$ uniformly over the
centers.  Average this conditional inequality over the localized columns
\emph{before} maximizing over centers.  The maximum is therefore over
deterministic, fully averaged expected log partitions; no maximum of
conditional expectations is taken.  Then enlarge the remaining spin
domain to the full sphere.
Consequently it suffices to bound, for every deterministic sequence of
centers, the unrestricted hybrid expected log partition plus
$\frac N2\log(1-t)$.  There is no compactness assertion uniform in an
unspecified growing collection here: if the required bound failed for the
maximum, one could choose a maximizing center for each $N$ and apply the
following subsequence argument to that deterministic sequence.
The passage from probability to expected pressure is legitimate because
$N^{-1}\log Z$ is bounded above for these potentials, and its negative
part has uniformly bounded moments by Jensen and
\eqref{eq:univ-opnorm}.

For each localized row, condition on independent fair signs $s_i$ in the
variance-mixture representation.  Its bias is Gaussian of variance
$t+\eps D$, where
\begin{equation}\label{eq:univ-mark}
 D=\sum_{i\in I}c_i^2s_i,\qquad |D|\le t.
\end{equation}
These marks and biases are independent from row to row and from the
Gaussian remaining columns.  Along a subsequence, $t$ converges and the
laws of $D$ converge weakly to a probability measure supported on
$[-t,t]$.  We use $t,D$ for these limits.  The next calculation is valid
for any such law; in particular it does not need a lower bound on the
number of large coordinates.

\subsection{The marked-row interpolation}

Here we give the extension of the Gaussian upper interpolation from
Section~\ref{sec:gaussian-pressure}, including the terms caused by the
localized coordinates.  For two remaining spins with overlap $R$, put
$R_f=t+(1-t)R$.  Conditional on the row mark $D$, the actual row field
has covariance $R_f+\eps D$ and diagonal
$v=1+\eps D$.  Its common Gaussian variance $t+\eps D$ is nonnegative.
For a finite trial profile $q$ reaching diagonal $1$, use a row trial
field of profile $vq$ and diagonal $v$, and put
\[
 p=\alpha\E_D D_{vq},\qquad
 p_d=\alpha\E_D D_{vq,d},
\]
where the row derivative profile is calculated with terminal $w$.
The spin field must have covariance $NR_fp$. To realize it, write
$x=\sqrt n\,z$ for the remaining spin and take independent cascade fields
$z_0^p,z_1^p,\ldots,z_n^p$, each with profile $p$ and diagonal $p_d$.
Use
\[
 \sqrt{Nt}\,z_0^p(\tau)
 +\sqrt{\frac{N(1-t)}n}\sum_{i=1}^n x_i z_i^p(\tau).
\]
The first term is common to all remaining spins and supplies covariance
$Ntp$; the second supplies $N(1-t)Rp$. Add a compensation field with
covariance $Nqp$.  At interpolation time $s\in[0,1]$, using
independent fields conditional on the marks, each row input is $\sqrt{s}$
times the actual field plus $\sqrt{1-s}$ times the trial field; the spin
and compensation fields enter the exponent with factors $\sqrt{1-s}$
and $\sqrt{s}$, respectively, while $H_N^{\rm pert}$ is unchanged.
At finite $N$ one may use the current mark
law and $M/N$ in this definition; continuity of the finite cascade
recursion gives the displayed limits.  Since $n/N\to1$, the difference
between the $n$-spin and $N$-spin normalizations vanishes.

Here are the fluctuation and deletion details needed to apply the cavity
argument of Section~\ref{sec:gaussian-pressure}.  All limits in this
paragraph are taken at a fixed finite trial grid and a fixed cap.  Keep
$s$ fixed, and let $Q$ be the covariance of two cascade leaves.
Changing one mark from $D$ to $D'$ changes that row's
covariance by
\[
 \Delta C=\eps(D'-D)[s+(1-s)Q],\qquad
 \Delta C_{\mathrm{diag}}=\eps(D'-D).
\]
The profile $p$ is defined from the deterministic mark law, rather than
the realized marks, and is held fixed in this comparison.  Gaussian
covariance differentiation of the expected log partition, conditional on
all marks, gives the change-rate bound
\[
 \frac{\eps|D'-D|}{2}
       \bigl(\|w''\|_\infty+2\|w'\|_\infty^2\bigr).
\]
Indeed its two terms are the diagonal expectation of $w''+(w')^2$ and
the off-diagonal expectation of $[s+(1-s)Q]w'_1w'_2$; the latter
coefficient lies in $[0,1]$.  The bound is $O(1)$ uniformly in the other
marks, $N$, and $s$.  Conditional on the marks, the independent row fields
give a Gaussian Lipschitz constant at most
$\|w'\|_\infty\sqrt{3M/2}$; spin and compensation fields add
$O(\sqrt N)$ and the perturbation adds $O(s_N)$.  The finite-cascade
log-mass estimate from Section~\ref{sec:gaussian-pressure} handles the
weight randomness.  Bounded differences for the marks therefore imply
an unconditional expected fluctuation of $O(\sqrt N+s_N)$, also after
any fixed number of row deletions.

Apply the GG argument after averaging the marks as well as the Gaussian
disorder.  There is no assertion of GG for every quenched mark vector.
Deleting row $1$ removes its row field and mark but leaves the
deterministic $p$, the spin and compensation fields, and the perturbation
fixed.  Thus at each $N$ the joint law of the removed mark and deleted-row
replica array is exactly
$\rho_N\otimes\mathop{\rm Law}(\text{deleted-row array})$.
At fixed almost-everywhere perturbation parameters and interpolation
time, take deterministic subsequences of these distributions.  Every
limit retains this product structure.  The independent mark may
therefore be adjoined to the limiting cavity array and integrated after
Gaussian row reinsertion.  This argument does not select a subsequence
depending on the realized removed mark.  Bounded derivatives at the
fixed cap justify the extra integration by dominated convergence.
The positivity and comonotone-profile conclusions of the Gaussian
cavity argument now apply.  Finite-grid limits are taken before profile
approximation, so no concentration bound uniform in cascade depth is
required.

To make the cancellation explicit, let $Q$ denote the overlap of two
cascade leaves, and let $a_\mu^\ell$ be $w'$ evaluated at the interpolating
field of row $\mu$ in replica $\ell$.  Gaussian integration by parts gives
the following exact derivative of the normalized expected log partition
at fixed perturbation coefficients (with $\E$ averaging the row marks,
Gaussian fields, and cascade weights):
\[
 -\frac{M}{2N}\E\left\langle
  [(R_f-Q)+\eps D_1(1-Q)]a_1^1a_1^2\right\rangle
 +\frac12\E\langle(R_f-Q)p(Q)\rangle.
\]
Here $p(Q)$ means the value of the spin covariance at the same cascade
rank; repeated covariance levels can be refined without changing the
fields.  The diagonal differences are zero in both terms.

Let $r$ be the limiting remaining-spin cavity profile.  It is
nonnegative and comonotone with the prescribed cascade profile $q$.
Thus $r_f=t+(1-t)r$ satisfies $r_f+\eps D\ge(1-\eps)t\ge0$.
The interpolating conditional profile is
$C_D(s)=s(r_f+\eps D)+(1-s)vq$, with fixed diagonal $v$.
Row concavity, with the small perturbation allowed by
Section~\ref{sec:row-concavity}, bounds the row derivative by
\[
 -\frac\alpha2\int_0^1
    \E_D\big[((r_f-q)+\eps D(1-q))D_{vq}\big]\,\dd a.
\]
The spin and compensation derivatives cancel exactly the term
$-\frac12\int(r_f-q)p$.  The remaining upper bound is
$-\frac{\alpha\eps}{2}\int(1-q)\E_D[DD_{vq}]$.
As in the Gaussian proof, this pointwise subsequential derivative bound
integrates by Fatou after the perturbation parameters are averaged;
the fixed-cap derivatives are bounded, and the perturbation contributes
$o(N)$.

At the trial endpoint the row contribution is
$\alpha\E_D\Phi_w(vq;v)$.  The common spin field contributes
$t(p_d-\int p)/2$.  The remaining-spin field contributes, at finite $N$,
$(n/N)S_n((N/n)(1-t)p)$; by uniform convergence on bounded profiles,
this converges to $\mathcal S((1-t)p)+(1-t)p_d/2$.  The second term
restores the diagonal contribution removed in the definition of
$\mathcal S$.
Compensation subtracts
$(p_d-\int qp)/2$.  Using $|D|\le t$, $D_{vq}\ge0$ and the definition
of $p$, the sum of these terms, the residual derivative bound, and the
slice density is at most
\begin{equation}\label{eq:univ-slice-bound}
 \alpha\E_D\Phi_w(vq;v)+\frac12\log(1-t)
 +\mathcal S((1-t)p)+\frac12\int_0^1(q_e-t)p,\qquad
 q_e=\eps t+(1-\eps t)q.
\end{equation}
In this cancellation the restored diagonal term $(1-t)p_d/2$ combines
with the common-field contribution and the compensation, leaving no
diagonal term.  This also checks the normalization of the common field.

We now fix $q$ to be a Gaussian base minimizer, so that $q_*<1$.
This includes every minimizer of \eqref{eq:src9}: the conjugate bound
gives $\mathcal P_u(q)\le\alpha\Phi_u(q)+E(q)=P(u)$, whereas the
pressure formula gives $\mathcal P_u(q)\ge P(u)$.  Thus such a $q$ is
also a minimizer of $\mathcal P_u$ and inherits its root estimates.
At a fixed cap first pass from finite trial profiles to this $q$ by
profile continuity.  Since $p\ge0$, clipping
$(q_e-t)/(1-t)$ below at zero only increases its pairing with $p$.
The spin conjugate \eqref{eq:src3} therefore bounds the non-row terms by
\begin{equation}\label{eq:src12}
 \begin{split}
 &\frac12\log(1-t)+
 E\left(\max\left\{0,\frac{q_e-t}{1-t}\right\}\right)\\
 &\qquad=E(q_e)-\frac12\int_0^t\frac{\dd z}{\chi_{q_e}(z)}\\
 &\qquad\le E(q)+\frac12\log(1-\eps t)
 -\frac{(1-\eps)t}{2(1-\eps t)\chi_q(0)}.
 \end{split}
\end{equation}
Here is a direct check, including clipping.  If
$\widehat q=(q_e-t)_+/(1-t)$, then for $z\ge0$
\[
 m_{\widehat q}(z)=m_{q_e}(t+(1-t)z),\qquad
 \chi_{\widehat q}(z)=\frac{\chi_{q_e}(t+(1-t)z)}{1-t}.
\]
Substitution into \eqref{eq:src3} proves the equality.  When $t$ exceeds
the essential supremum of $q_e$, use
$\chi_{q_e}(z)=1-z$ above that supremum; the same equality reduces to
$E(\widehat q)=E(0)=0$.
For $b=\eps t$, the affine change $q_e=b+(1-b)q$ gives
\[
 E(q_e)=E(q)+\tfrac12\log(1-b)
       +\frac{b}{2(1-b)\chi_q(0)},\qquad
 \chi_{q_e}(0)=(1-b)\chi_q(0).
\]
Since $\chi_{q_e}$ is nonincreasing,
$\int_0^t\dd z/\chi_{q_e}(z)\ge t/\chi_{q_e}(0)$, proving the inequality.
Crucially, the resulting estimate contains no $p$.  We can now remove the
cap, using its uniform row-functional error over $v\in[1/2,3/2]$.

\subsection{A variance estimate uniform in temperature}

The entropy estimate \eqref{eq:src12} gives a loss of at least
$t/(4\chi_q(0))$ when $\eps\le1/2$. We now bound the possible row gain
by $C(B_0)\eps t/\chi_q(0)$ with no temperature dependence in $C(B_0)$.
The needed residual estimate is a property of the single-row diffusion.

\begin{lemma}[Residual moments under a decrease of variance]
\label{lem:univ-residual-variance}
Let $q$ be a nonnegative nondecreasing profile with diagonal $1$, and
let $u(z)=-\beta(-1-z)_+^2/2$, $\beta>0$.
Put $a_0=f_z(0,0)$ for its row solution before the root Gaussian variance.
Write $\E_{u,v}$ for expectation under the tilted endpoint law for profile
$vq$ and diagonal $v$. For $1/2\le v\le1$,
\begin{equation}\label{eq:src13}
 \sqrt v\,\beta\E_{u,v}(-1-g)_+\le a_0,
 \qquad
 \E_{u,v}(-1-g)_+^2
 \le \sqrt{\pi/2}\,\E_{u,v}(-1-g)_+.
\end{equation}
\end{lemma}

\begin{proof}
For the first inequality put the variance scale into the terminal, so the
Gaussian variance profile is $q$ on the fixed time interval $[0,1]$.
Let $A_v$ and $A_1$ be the derivative solutions for the terminals
$u(\sqrt v z)$ and $u(z)$, respectively.  Their difference $H=A_v-A_1$
satisfies
\[
 H_t+\tfrac12H_{zz}+m_q A_vH_z+m_q(A_1)_zH=0,
 \qquad H(1,z)=\sqrt v\,\beta(-1-\sqrt v z)_+
                       -\beta(-1-z)_+\le0.
\]
The zeroth-order coefficient lies in $[-\beta,0]$, and the drift is
Lipschitz with at most linear growth.  The maximum principle gives
$H\le0$.  To handle the hinge, convolve the two terminal derivative
functions with the same nonnegative mollifier, preserving their order,
and pass to the limit.  When $q(0)>0$, $m_q=0$ before that root variance,
so the same comparison includes the initial heat interval.  The derivative
martingale at $(0,0)$ identifies $A_v(0,0)$ with the expectation on the
left of \eqref{eq:src13}.

For the second inequality use variance time, writing the tilted diffusion
as
\[
 \dd Z_s=\dd B_s+m(s)a(s,Z_s)\,\dd s,\qquad Z_0=0,
 \qquad 0\le s\le v.
\]
The base derivative $a=f_z$ is nonnegative, nonincreasing, and convex in
position.  Nonnegativity and monotonicity follow from the terminal and
parabolic comparison.  Convexity follows by differentiating once more:
$c=f_{zzz}$ has nonnegative terminal value (a positive measure for the
hinge) and solves the linear equation
$(\partial_s+\frac12\partial_{zz}+ma\partial_z)c=-3mf_{zz}c$.
Smooth concave approximations with convex first derivative and no far cap
justify the claim; their diffusion limits give the hinge.
Thus the drift is nonnegative, nonincreasing, and convex.

Condition on the Brownian bridge $\omega_s=B_s-(s/v)B_v$ and write
$b=B_v\sim\cN(0,v)$, independent of the bridge.  The endpoint map
$T_\omega(b)=Z_v$ is increasing, convex, and $1$-Lipschitz.  Indeed its
first variation $J_s$ solves
$J_s'=1/v+d_z(s,Z_s)J_s$, $J_0=0$, for $d=ma$; hence
$0\le J_s\le s/v$ and, more precisely,
\[
 T_\omega'(b)=\frac1v\int_0^v
       \exp\left\{\int_r^v d_z(s,Z_s)\,\dd s\right\}\dd r
 \ge\frac{1-e^{-\beta v}}{\beta v}>0.
\]
Its second variation solves
$K_s'=d_zK_s+d_{zz}J_s^2$, $K_0=0$, and is nonnegative.
These deterministic integral-equation variations remain valid for every
continuous bridge path by smooth approximation.  Also $T_\omega(b)\ge b$
by nonnegative drift.  The strictly positive lower slope bound makes
$T_\omega$ onto, so it has a unique crossing $b_0\le-1$ of $-1$.  Set
$R(s)=-1-T_\omega(b_0-s)$, $s\ge0$.  Then $R$ is increasing and
concave, $R(0)=0$, and $0\le R(s)\le s$.
Consequently $R(s)/s$ is nonincreasing.  Under the Gaussian left-tail
measure with density proportional to
$\exp(-(b_0-s)^2/(2v))\ind_{s\ge0}$, decreasing-likelihood-ratio
comparison gives
\[
 \frac{\E[R(s)^2]}{\E[R(s)]}
 \le \frac{\E[sR(s)]}{\E[R(s)]}
 \le \frac{\E[s^2]}{\E[s]}
 \le\sqrt{\pi/2}.
\]
The middle inequality follows from the same covariance identity used
above, now relative to the probability density proportional to
$s\exp(-(b_0-s)^2/(2v))$.  Relative to the density proportional to
$s e^{-s^2/2}$, the logarithmic derivative of its likelihood ratio is
$b_0/v-(1/v-1)s\le0$.  The last ratio is therefore at most
$\int_0^\infty s^2e^{-s^2/2}\dd s/\int_0^\infty se^{-s^2/2}\dd s
=\sqrt{\pi/2}$.  Multiply by the conditional first moment and integrate
over the bridge to obtain the second assertion of \eqref{eq:src13}.
The lower slope bound is used only for the crossing and approximation;
its dependence on $\beta$ does not enter the moment constant.  Mollified
drifts retain all three shape properties and the common Lipschitz bound
$\beta$.  Approximate a general monotone $m$ by steps in $L^1$; the linear
PDE difference equation and uniform diffusion moments give local
convergence of solutions, and the flow maps converge on compact endpoint
sets.  Finally $(-1-Z_v)_+\le(-1-B_v)_+$ supplies uniform integrability
for all residual moments.  These observations justify the hinge and
coefficient limits in the bridge argument.
\end{proof}

\paragraph{Bounded terminal tests.}
A bounded terminal perturbation changes a tilted endpoint
expectation of a nonnegative observable by at most the factor
$e^{2|d|\|\psi\|_\infty}$.  On the same realized cascade and Gaussian
field, the normalized Gibbs measure with terminal $u+d\psi$ has density
\[
 \frac{e^{d\psi(g)}}{\langle e^{d\psi(g)}\rangle_u}
 \in[e^{-2|d|\|\psi\|_\infty},e^{2|d|\|\psi\|_\infty}]
\]
relative to the measure with terminal $u$.  Average this pointwise bound
over the cascade and field.  The endpoint expectation here is precisely
the PDE tilted expectation: for a bounded continuous terminal direction
$h$, differentiating $\Phi_{u+\lambda h}$ at zero gives
$\E\langle h(g)\rangle_u$ directly, while differentiating the recursion
gives the following linear equation, now denoting the fixed diagonal by
$Q$:
\[
 J_t+\tfrac12J_{zz}+m f_zJ_z=0,\qquad J(Q,z)=h(z).
\]
Its diffusion representation identifies the same expectation with the
tilted endpoint law.  The initial $m=0$ heat interval includes common
root variance.  For residual diagonal variance, augment the base measure
at each leaf by the residual Gaussian variable and integrate it inside
that leaf weight.  Each replica draws a fresh residual variable, even
when leaf labels coincide, so this convention adds no off-diagonal
variance.  The terminal $m=1$ interval is exactly this convolution.
Profile approximation preserves the equality for bounded tests and the
likelihood-ratio bound; monotone truncation extends the latter to
nonnegative observables.  In particular there is no factor depending on
cascade depth.

\paragraph{Comparison with the entropy loss.}
Put $w=u+d\psi$ after removing the cap.  Differentiating the variance
scale in the terminal gives
\[
 \partial_v\Phi_w(vq;v)=\frac1{2v}\E_{w,v}[g w'(g)].
\]
This identity follows first by differentiating the finite cascade
recursion and then by the uniform quadratic moment bounds.  If
$r=(-1-g)_+$, then $g u'(g)=-\beta(r^2+r)$.  For $v\ge1$ this part of
the derivative is nonpositive.  For $1/2\le v\le1$, the preceding
perturbation bound and \eqref{eq:src13} show that its absolute value is
at most $C_0a_0$, provided $|d|\|\psi\|_\infty\le1$, where
$C_0=e^2\sqrt2(1+\sqrt{\pi/2})$.
The test part is bounded by $|d|\sup_z|z\psi'(z)|/(2v)$.
Impose also $|d|\sup_z|z\psi'(z)|\le1$.  Since $|v-1|\le\eps t$,
\[
 \alpha\big(\Phi_w(vq;v)-\Phi_w(q;1)\big)
 \le \eps t(C_0\alpha a_0+\alpha)
 \le \frac{(C_0\sqrt{B_0}+B_0)\eps t}{\chi_q(0)}.
\]
The final inequality uses \eqref{eq:src10}, which implies
$\alpha a_0\le\sqrt\alpha/\chi_q(0)$, and $\chi_q(0)\le1$.
The constant $C_0$ is absolute: the dependence on $\beta$ has disappeared.
For $\eps\le1/2$ the negative last term in \eqref{eq:src12} has magnitude
at least $t/(4\chi_q(0))$, and its logarithmic term is nonpositive.
Choose, for example, any
\[
 0<\eps_0(B_0)<
 \min\left\{\frac12,\frac1{8(C_0\sqrt{B_0}+B_0)}\right\}.
\]
The entropy loss then dominates the row improvement in every slice.
To specify the quantifiers, choose this $\eps_0(B_0)$ before choosing
$\beta$, a minimizer, a test perturbation, or a cap.  For each fixed
$\beta$ and $\psi$, Proposition~\ref{prop:conc-stability} and
Corollary~\ref{cor:conc-hinge-stability} allow both
signs of $d$ with
\[
 |d|\le\min\left\{d_{\mathrm{stab}}(\beta,\psi),
       \frac1{1+\|\psi\|_\infty},
       \frac1{1+\sup_z|z\psi'(z)|}\right\}.
\]
The stability radius is uniform in the trial profile, the diagonal in
$[1/2,2]$, and the cap and mollifier.  The thermodynamic limit is taken
at each fixed cap and finite trial grid; profile approximation follows,
then removal of mollification at fixed cap, and finally removal of the
cap.  The entropy factor $1/\chi_q(0)$ may depend on $\beta$ but appears
on both sides of the comparison and cancels.  No convergence rate
uniform in $\beta$ is required.  The same fixed mixture law therefore
remains valid throughout the later ordered limits.
Equations \eqref{eq:univ-slice-bound}--\eqref{eq:src13} prove the upper
support in Theorem~\ref{thm:univ-pressure}.  Combined with
\eqref{eq:univ-lower} at $d=0$, they prove its pressure assertion.

\subsection{The inner thermodynamic observables}

Let $\mu_{\beta,\alpha}$ be the Gaussian tilted endpoint law associated
with a minimizer $q$.  The following argument proves, in particular, that
this law does not depend on the choice of minimizer.  Define the random
empirical projection measure
\[
 L_N=\frac1M\sum_{\mu=1}^M\delta_{A_\mu y}.
\]
Convergence in probability below refers to the joint law of disorder and
one Gibbs sample.

\begin{proposition}[Empirical row law and force normalization]
\label{prop:univ-observables}
For either the Gaussian law or any mixture in
Theorem~\ref{thm:univ-pressure}, at fixed $\beta>0$ and
$0<\alpha\le B_0$, $L_N$ converges weakly in probability to
$\mu_{\beta,\alpha}$.  This measure is atomless, assigns positive mass
to $(-\infty,-1)$, and has finite moments of its negative part.
The gap and normalized-force observables in the problem have inner
$N\to\infty$ limits at every fixed argument, and those limits agree
for all these entry laws.
\end{proposition}

\begin{proof}
First, the normalized random log partition concentrates for the uncapped
potentials $u+d\psi$.  On the convex operator-norm ball
$\|A\|_{\mathrm{op}}\le K\sqrt N$, differentiation gives
\[
 \|\nabla_A\log Z\|_F
 \le \beta(\|A\|_{\mathrm{op}}+\sqrt M)
        +|d|\|\psi'\|_\infty\sqrt M=O(\sqrt N).
\]
Extend its restriction from this ball with the same Lipschitz constant
and use Gaussian concentration, through $T$ for mixture entries.
The complement has probability exponentially small in $N$ for a large
fixed $K$, by \eqref{eq:univ-opnorm}.  Jensen's lower bound
$\log Z\ge-C_\beta(M+\|A\|_{\mathrm{op}}^2)-M|d|\|\psi\|_\infty$
and the corresponding upper bound show that the extension changes its
expectation by $o(N)$.  Hence $N^{-1}\log Z-P_{N,\eps}(u+d\psi)$
converges to zero in probability.

For a fixed test $\psi$, put
$S(d)=\alpha\Phi_{u+d\psi}(q;1)+E(q)$. Theorem~\ref{thm:univ-pressure}
gives $S(0)=P(u)$ and
$S'(0)=\alpha\int\psi\,\dd\mu_{\beta,\alpha}$.
If $b>S'(0)$, choose a sufficiently small $d>0$ and $c>0$ such that
$S(d)-S(0)<db-2c$. The expected upper support and base pressure limit,
together with the preceding concentration, imply with disorder probability
tending to one that
\[
 \log Z(u+d\psi)-\log Z(u)\le N(db-c).
\]
This uses only a limsup bound at $d$; a thermodynamic pressure limit for
the perturbed terminal has not been assumed.  Exponential Markov under the Gibbs
measure gives
\[
 \left\langle\ind_{\{N^{-1}\sum_\mu\psi(g_\mu)\ge b\}}\right\rangle
 \le\exp\{\log Z(u+d\psi)-\log Z(u)-Ndb\}.
\]
On that event the right side is at most $e^{-Nc}$.  The same argument with $d<0$ gives the lower
bound.  Therefore
$\int\psi\,\dd L_N\to\int\psi\,\dd\mu_{\beta,\alpha}$ in joint
probability for every $\psi\in C_c^\infty$.  A countable determining
class, together with tightness obtained by smooth cutoffs whose
$\mu_{\beta,\alpha}$-masses tend to one, proves weak convergence.
The same limits for every base minimizer prove their equality of row
marginals.

Here are direct bounds ensuring that threshold tests and normalization
are permitted.  Remove one row and let $\langle\cdot\rangle_0$ be its
cavity measure.  Its reweighting denominator is
$D=\langle e^{u(g)}\rangle_0\le1$, and Jensen gives
$D\ge\exp\langle u(g)\rangle_0$.  Since the row is independent of this
measure, $\E[-\langle u(g)\rangle_0]\le C_\beta$.  Thus
$\Prob(D<e^{-K})\le C_\beta/K$.  For an interval $J$, conditional on a
unit cavity spin the fresh projection has density bounded by an absolute
constant, and $e^u\le1$.  Consequently its reweighted marginal satisfies
\[
 \E\langle\ind_{\{g\in J\}}\rangle
 \le C e^K|J|+C_\beta/K.
\]
This estimate is uniform in $N$ and passes to the deterministic limiting
law by continuous tests.  Sending $|J|\downarrow0$ and then $K\to\infty$
proves that the law has no atoms.
Conversely, for a fixed nonempty interval $J\Subset(-\infty,-1)$,
$D\le1$ gives
\[
 \E\langle\ind_{\{g\in J\}}\rangle
 \ge \inf_{g\in J}e^{u(g)}
     \inf_{v\in[1/2,3/2]}\Prob(\sqrt v Z\in J)>0.
\]
Hence its limiting negative-gap probability is positive.
Finally the single-row covariance inequality used above gives, for every
fixed $p>0$,
\[
 \E\langle(-1-g)_+^p\rangle
 \le\sup_{v\in[1/2,3/2]}\E(-1-\sqrt v Z)_+^p<\infty.
\]
Using an exponent larger than $p$ gives uniform integrability, including
for empirical averages.

Set
\[
 p_-=\mu_{\beta,\alpha}(( -\infty,-1)),\qquad
 m_-=\int(-1-g)_+\,\mu_{\beta,\alpha}(\dd g),\qquad
 \bar r=m_-/p_-.
\]
Both $p_-$ and $m_-$ are strictly positive.  Weak convergence,
atomlessness, and uniform integrability imply in joint probability
\[
 \frac{|C(y)|}{M}\longrightarrow p_-,\qquad
 \frac1M\sum_{\mu\in C(y)}(-1-g_\mu)\longrightarrow m_-.
\]
For the second assertion, the second-moment bound gives $\E\int r\ind_{\{r>K\}}\,\dd L_N\le C/K$, where
$r=(-1-g)_+$.  Markov's inequality makes this truncation error small in
joint probability.  The bounded continuous truncated residual converges
by weak convergence, and then $K\to\infty$ gives the assertion.
Thus the configuration-dependent conditional mean residual converges to
$\bar r$.  For every $s>0$, sandwich its random threshold between
$s(\bar r-\delta)$ and $s(\bar r+\delta)$, apply weak convergence to
these fixed thresholds, and send $\delta\downarrow0$.  Atomlessness at
$g=-1-s\bar r$ justifies the last step.  Since the observables lie in
$[0,1]$, convergence in probability also gives convergence of their
expectations.  Explicitly,
\begin{align*}
 \lim_{N\to\infty}G_{N,\beta,\alpha}(\delta,u_+)
 &=\mu_{\beta,\alpha}((\delta-1,u_+-1]),\qquad 0<\delta<u_+,\\
 \lim_{N\to\infty}F_{N,\beta,\alpha}(s)
 &=\frac{\mu_{\beta,\alpha}
       (\{-1-s\bar r\le g<-1\})}{p_-}.
\end{align*}
The convention when $C(y)$ is empty is immaterial, since its probability
tends to zero.  These formulas prove the proposition with exactly the
force normalization specified in the problem.
\end{proof}


\section{Ground energy and the sharp satisfiability threshold}
\label{sec:threshold}

This section separates an energetic issue from a geometric one.  The
pressure formula determines when the minimum energy is subextensive.  A
uniform repair lemma then turns subextensive energy into exact
satisfiability after deleting a fixed fraction of the constraints.
There are two geometric obstacles.  Solving for the violated rows can
create violations in other rows, so the repair enlarges the set of rows
whose projections it controls.  The repaired vector may then have norm
less than one, so a second displacement restores its norm without
spoiling those controlled projections.  The negative margin permits
radial contraction when the repaired vector instead has norm at least one.

Write $A\in\R^{M\times N}$ for the pattern matrix, and use unit-spin
coordinates $y=x/\sqrt N$.  Set
\[
 r_A(y)=(-\mathbf 1-Ay)_+,\qquad
 H_A(y)=\frac12\|r_A(y)\|_2^2,\qquad
 R(A)=\min_{\|y\|_2=1}\|r_A(y)\|_2.
\]
All vector norms without a subscript are Euclidean, and $\|A\|$ denotes
the operator norm between Euclidean spaces.  Let
$P_{N,\alpha}^{\nu}(\beta)=N^{-1}\E\log\int e^{-\beta H_A(y)}\,
\dd\sigma(y)$, where $M=\lfloor\alpha N\rfloor$, $\sigma$ is normalized
surface measure on the unit sphere, and $\nu$ is the entry law.  We use
$\nu_0=\mathcal N(0,1)$ and the variance mixtures $\nu_\eps$ from the
introduction.  The pressure and universality results of
Sections~\ref{sec:gaussian-pressure} and~\ref{sec:universality} are the
only input from the variational analysis: the Gaussian pressure has a
limit at every fixed $(\alpha,\beta)$; for every finite $B_0$, the same
limit holds for $0\leq\eps<\eps_0(B_0)$ and
$0<\alpha\leq B_0$.

We will repeatedly use a uniform matrix-norm estimate.  For
$\eps\leq1/2$, an entry satisfies
$\E e^{tA_{\mu i}}\leq e^{3t^2/4}$.  Independence gives the same bound
for $v^{\mathsf T}Ay$ when $\|v\|=\|y\|=1$.  Nets of radius $1/4$ on
the two spheres have at most $9^{M+N}$ pairs, and approximating both
arguments gives $\|A\|\leq2\max|v^{\mathsf T}Ay|$ on these nets.  Thus
\begin{equation}
 \Prob(\|A\|>2t)\leq2\,9^{M+N}e^{-t^2/3},\qquad t>0.
 \label{thresh:eq:operator-tail}
\end{equation}
In particular $\E\|A\|^2\leq C(M+N)$, and sufficiently large constant
multiples of $\sqrt M+\sqrt N$ bound $\|A\|$ with exponential error.
The case $M=0$ is immediate.

\subsection{The ground-energy limit}

\begin{lemma}[Uniform approximation by the pressure]
\label{thresh:pressure-energy}
For every finite $B_0$ there is a constant $C=C(B_0)$ such that, for
$N\geq2$, $0\leq M\leq B_0N$, $0\leq\eps\leq1/2$, $\beta>0$, and
$0<\delta<1/2$,
\begin{equation}
 0\leq -\frac{P_{N,\alpha}^{\nu_\eps}(\beta)}{\beta}
       -\frac1N\E\min_{\|y\|=1}H_A(y)
 \leq C\delta+\frac1\beta\log\frac C\delta.
 \label{thresh:eq:pressure-energy}
\end{equation}
Consequently the Gaussian limit
\[
 e(\alpha)=\lim_{N\to\infty}\frac1N\E\min_{\|y\|=1}H_A(y)
          =\lim_{\beta\to\infty}
                    -\frac{P_{\alpha}^{\nu_0}(\beta)}{\beta}
\]
exists.  On every range where the fixed-temperature pressures agree,
the mixture ground-energy limit also exists and equals $e(\alpha)$.
Moreover, $e$ is nondecreasing and $1$-Lipschitz, and $e(\alpha)=0$ for
$0\leq\alpha\leq1$.
\end{lemma}

\begin{proof}
The positive-part map is $1$-Lipschitz, so for unit $y,y'$,
\[
 \|r_A(y)-r_A(y')\|\leq\|A\|\,\|y-y'\|,
 \qquad \|r_A(y)\|\leq\sqrt M+\|A\|.
\]
It follows that $H_A$ has Lipschitz constant at most
$L_A=\|A\|(\sqrt M+\|A\|)$ on the sphere.  A spherical cap of Euclidean
radius $\delta$ about a minimizer has surface probability at least
$(c\delta)^N$, for an absolute $c>0$.
For completeness, rotate its center to the last coordinate vector and
parametrize the upper hemisphere by
$z\mapsto(z,\sqrt{1-\|z\|^2})$.  The ball
$\|z\|\leq\delta/2$ maps inside the cap, and its surface Jacobian is at
least one.  The ratio of the $(N-1)$-ball volume to the
$(N-1)$-sphere area is bounded below by $c_1^N$, as follows directly
from their gamma-function formulas; this proves the stated bound.
Integrating over this cap and also using $H_A\geq\min H_A$ gives the
deterministic estimate
\[
 0\leq-\frac1{\beta N}\log\int e^{-\beta H_A}\,\dd\sigma
             -\frac{\min H_A}{N}
 \leq\frac{L_A\delta}{N}+\frac1\beta\log\frac C\delta.
\]
Equation~\eqref{thresh:eq:operator-tail} gives
$\E\|A\|^2\leq C(B_0)N$; hence $\E L_A/N\leq C(B_0)$.
This proves~\eqref{thresh:eq:pressure-energy}.

At fixed $N$, the quantity $-P_{N,\alpha}^{\nu_0}(\beta)/\beta$ is
nonincreasing in $\beta$ and nonnegative.  Indeed,
$\beta\mapsto\log\int e^{-\beta H_A}\dd\sigma$ is convex and vanishes at
zero.  Its limiting pressure inherits these properties, so its
zero-temperature limit exists.  First let $N\to\infty$ in the
inequality, then let $\beta\to\infty$ with, for example,
$\delta=\beta^{-1/2}$.  The liminf and limsup of the expected ground
energy coincide with this limit.  The identical argument applies to
any law with the same fixed-temperature pressures.

Couple matrices of different row counts by appending independent rows.
Appending rows cannot decrease the minimum energy.  Conversely,
evaluate the new rows at a measurable approximate minimizer of the old
matrix and then let its approximation error tend to zero.  For every
fixed unit $y$, a fresh projection $g=A_\mu y$ has mean zero and variance
one, and
\[
 \E V(g)\leq\tfrac12\E(1+g)^2=1.
\]
Thus the expected minimum energy increases by at most one per appended
row.  After dividing by $N$ and passing to the limit this proves the
monotonicity and Lipschitz assertions.  Approximate minimizers can be
selected measurably by taking the first adequate point in a fixed
countable dense subset of the sphere, since $H_A$ is continuous.
If $M<N$, the row kernel contains a unit vector, at which every gap is
one.  Hence $e(\alpha)=0$ for $\alpha<1$; continuity gives $e(1)=0$.
\end{proof}

\begin{lemma}[Positive energy at a finite density]
\label{thresh:large-density}
There is an absolute finite $B>1$ such that the Gaussian ground energy
satisfies $e(B)>0$.  The same finite-dimensional lower bound used to
prove this assertion holds uniformly for $0\leq\eps\leq1/2$.
\end{lemma}

\begin{proof}
Conditioned on the variance marks of a mixture matrix, a fixed unit-spin
projection is Gaussian with variance in $[1/2,3/2]$.  Thus
\[
 \Prob(A_\mu y<-2)\geq p,
 \qquad p=\Prob(Z<-2\sqrt2)>0,
 \quad Z\sim\mathcal N(0,1).
\]
At a deterministic unit $y$, row independence and the elementary
Chernoff bound give
\[
 \Prob\bigl(\#\{\mu:A_\mu y<-2\}<pM/2\bigr)
 \leq e^{-c_pM}.
\]
Outside this event, $\|r_A(y)\|\geq\sqrt{pM/2}$.
There is an absolute $K_0$ such that
$\|A\|\leq K_0(\sqrt N+\sqrt M)$ with probability tending to one,
uniformly in these entry laws.  Choose a fixed
$\delta>0$ so small that $2K_0\delta\leq\frac12\sqrt{p/2}$, and take a
$\delta$-net of the unit sphere of cardinality at most
$(1+2/\delta)^N$.  For $M\geq N$, the Lipschitz bound for $r_A$ transfers
the residual lower bound on the net to
\[
 R(A)\geq\tfrac12\sqrt{pM/2}.
\]
The probability that a net point fails is at most
$\exp\{N\log(1+2/\delta)-c_pM\}$.  Choosing $B$ sufficiently large and
setting $M=\lfloor BN\rfloor$ makes this probability tend to zero.
Therefore $\min H_A/N\geq pM/(16N)$ with probability tending to one, and
$e(B)\geq pB/16>0$.
\end{proof}

Fix henceforth $B_0>B$ and reduce the pressure-comparison constant, if
necessary, so that $\eps_0(B_0)\leq1/2$.  Define
\begin{equation}
 \alpha_J=\sup\{\alpha\geq0:e(\alpha)=0\}.
 \label{thresh:eq:threshold-definition}
\end{equation}
The preceding lemmas show that $1\leq\alpha_J<B<B_0$.  By continuity
and monotonicity, $e(\alpha)=0$ for $\alpha\leq\alpha_J$ and
$e(\alpha)>0$ for $\alpha>\alpha_J$.  It remains to show that vanishing
energy density yields exact satisfiability at each strictly smaller
density.

\subsection{Uniform geometry of the rows}

The repair needs a uniformly bounded right inverse for every set of
slightly fewer than $N$ rows.  It also needs a bound on how many rows
can be close to contact at an arbitrary spin.  The next lemma supplies
both facts before the spin is selected from the random matrix.  This
uniformity is what permits us to apply it at a ground minimizer.


\begin{lemma}[Uniform row estimates]
\label{thresh:uniform-geometry}
Fix $B_0<\infty$ and $0<\rho<1/8$.  There are constants
$K<\infty$, $c>0$, $u_0>0$, $a>0$, and $C<\infty$, depending only on
$(B_0,\rho)$, such that, for all sufficiently large $N$, all
$M\leq B_0N$, and every $0\leq\eps\leq1/2$, the following hold
simultaneously with probability at least $1-Ce^{-aN}$:
\begin{enumerate}
\item $\|A\|\leq K\sqrt N$;
\item for every row set $I$ with $1\leq|I|\leq(1-\rho)N$,
\begin{equation}
 \|A_I^{\mathsf T}v\|\geq c\sqrt N\,\|v\|
       \quad\text{for all }v\in\R^{|I|};
 \label{thresh:eq:surjection}
\end{equation}
\item every unit $y$ satisfies
\begin{equation}
 \#\{\mu:|1+A_\mu y|\leq u_0\}\leq(1+\rho/2)N.
 \label{thresh:eq:gap-count}
\end{equation}
\end{enumerate}
\end{lemma}

\begin{proof}
Choose $K=K(B_0)$ large enough in
\eqref{thresh:eq:operator-tail} that the operator-norm failure
probability is exponentially small, and work on this event for the two
net transfers below.

For a fixed $I$, $k=|I|$, and fixed $v\in\R^k$ with $\|v\|=1$, the
coordinates of $A_I^{\mathsf T}v$, conditional on all variance marks,
are independent centered Gaussians with variances at least $1/2$.
Their joint density is at most $\pi^{-N/2}$.  The Euclidean ball-volume
formula therefore gives
\[
 \Prob\bigl(\|A_I^{\mathsf T}v\|\leq2c\sqrt N\bigr)
       \leq(C_1c)^N.
\]
Use a $(c/K)$-net of the unit sphere in $\R^k$, of cardinality at most
$(3K/c)^k$, taking $c<K$.  If~\eqref{thresh:eq:surjection} fails while
$\|A\|\leq K\sqrt N$, some net point has image norm at most
$2c\sqrt N$.  Summing over dimensions and row sets bounds the failure
probability by
\[
 N\,2^{B_0N}(3K/c)^{(1-\rho)N}(C_1c)^N.
\]
Its logarithm divided by $N$ is bounded by a constant depending on
$(B_0,K)$ plus $\rho\log c+o(1)$.  Choose $c$ sufficiently small to
make this bound negative.

For the gap-count estimate take a net of radius
$\eta=u_0\sqrt\rho/(4K)$.  If $y$ and its net representative $y_0$
satisfy $\|y-y_0\|\leq\eta$, then
\[
 \#\{\mu:|A_\mu(y-y_0)|>u_0\}
 \leq\frac{\|A(y-y_0)\|^2}{u_0^2}
 \leq\rho N/16.
\]
If~\eqref{thresh:eq:gap-count} fails, more than
$(1+7\rho/16)N$ rows at $y_0$ therefore satisfy
$|1+A_\mu y_0|\leq2u_0$; in particular at least
$\lceil(1+\rho/4)N\rceil$ do.  For fixed $y_0$, each row projection has
density at most $1/\sqrt\pi$, conditionally on its variance marks and
also unconditionally.  Hence its probability of this interval is at
most $C_2u_0$.  Row independence, a union bound over row subsets, and
then a union bound over the spin net give, for $C_2u_0<1$,
\[
 \Prob(\text{gap-count failure},\ \|A\|\leq K\sqrt N)
 \leq
 (1+8K/(u_0\sqrt\rho))^N\,2^{B_0N}
       (C_2u_0)^{(1+\rho/4)N}.
\]
Its logarithm divided by $N$ is at most a constant depending on
$(B_0,K,\rho)$ plus $(\rho/4)\log u_0$.  Choose $u_0$ sufficiently small.
These choices establish all three assertions with exponential error.
\end{proof}

The following elementary fact supplies a direction with uniformly
small row projections.  It does not require any randomness.

\begin{lemma}[A direction inside finitely many slabs]
\label{thresh:slab-direction}
Let $W$ be a $D$-dimensional Euclidean space, $D\geq1$, and
$L:W\to\R^m$ a linear map.  There exists a unit $v\in W$ such that
\begin{equation}
 \|Lv\|_\infty\leq\frac{\|L\|}{\sqrt D}.
 \label{thresh:eq:slab-direction}
\end{equation}
\end{lemma}

\begin{proof}
If $L$ has a nonzero kernel, take a unit kernel vector.  Otherwise the
set $Q=\{w\in W:\|Lw\|_\infty\leq1\}$ is compact and has nonempty
interior.  Let $w_*$ maximize $\|w\|$ over $Q$.  The active linear
functionals, those satisfying $|(Lw_*)_j|=1$, span $W^*$.  Indeed, if
they had a common nonzero annihilator $h$, both $w_*+th$ and $w_*-th$
would remain in $Q$ for sufficiently small $t>0$: active inequalities
would be unchanged and the finitely many inactive inequalities have
positive slack.  But at least one of these two points has squared norm
at least $\|w_*\|^2+t^2\|h\|^2$, a contradiction.
There are consequently at least $D$ distinct active rows, and
\[
 D\leq\|Lw_*\|^2\leq\|L\|^2\|w_*\|^2.
\]
The vector $v=w_*/\|w_*\|$ satisfies the required bound.
\end{proof}

\subsection{Repair after deleting constraints}

\begin{lemma}[Deterministic repair]
\label{thresh:repair}
Fix $K,c,u_0>0$ and $0<\rho<1/8$.  There exists
$d_*=d_*(K,c,u_0,\rho)>0$ such that the following holds for all
sufficiently large $N$.  Suppose $A$ satisfies the first two assertions
of Lemma~\ref{thresh:uniform-geometry}, and some unit $y$ satisfies
\[
 \|r_A(y)\|\leq d\sqrt N,
 \qquad
 S=\{\mu:1+A_\mu y\leq u_0\},
 \qquad |S|\leq(1-2\rho)N,
 \quad 0<d\leq d_*.
\]
Then there is a unit $\widehat y$ with $1+A_\mu\widehat y\geq0$ for
every row $\mu$.
\end{lemma}

\begin{proof}
We first correct the deficient gaps in the ambient vector space, keeping
all uncontrolled rows a fixed distance above zero.  We then return the
corrected vector to the unit sphere.

\emph{Correct the gaps.}
Choose $d_*>0$ sufficiently small that, for $d\leq d_*$,
\begin{equation}
 \frac dc\leq\frac12,\qquad
 L_d:=\frac{16K^2d^2}{c^2u_0^2}\leq\frac\rho4,\qquad
 K\sqrt{\frac{4d}{c\rho}}\leq\frac{u_0}{2}.
 \label{thresh:eq:repair-smallness}
\end{equation}
Write $h=\mathbf 1+Ay$ and $r=(-h)_+=r_A(y)$; these vectors
remain fixed throughout the construction.  For each imposed set
$T\supseteq S$, consider the equations $A_Sz=r_S$ and
$A_{T\setminus S}z=0$.  When they are solvable, let $z_T$ be their
minimum-norm solution.  Whenever $|T|\leq(1-\rho)N$, the
surjection bound gives existence and
\begin{equation}
 \|z_T\|\leq\frac{\|r_S\|}{c\sqrt N}\leq\frac dc.
 \label{thresh:eq:correction-norm}
\end{equation}
The empty-set case is understood as $z_\varnothing=0$.

Start at $T=S$.  If
\[
 J=\{j\notin T:|A_jz_T|>u_0/4\}
\]
is nonempty, add all of $J$ to $T$, impose $A_Jz=0$, and solve again.
Here is a count that verifies that every solve is permitted.
If $z'$ is the new solution, then $z'-z_T\in\ker A_T$ while
$z_T$ belongs to the row span of $A_T$.  Thus
\[
 \|z'\|^2-\|z_T\|^2=\|z'-z_T\|^2
 \geq\frac{\|A_Jz_T\|^2}{K^2N}
 >\frac{|J|u_0^2}{16K^2N}.
\]
Summing over completed additions and using
\eqref{thresh:eq:correction-norm} shows that they contain at most
$L_dN$ rows altogether.  The next, not yet imposed, set also has at
most $L_dN$ rows, since
\[
 |J|u_0^2/16<\|A_Jz_T\|^2
 \leq K^2N\|z_T\|^2\leq K^2N(d/c)^2.
\]
Consequently, before every proposed solve,
\[
 |T\cup J|\leq(1-2\rho)N+2L_dN
             \leq(1-3\rho/2)N<(1-\rho)N.
\]
This proves the induction without assuming existence of the proposed
solution.  Each nonterminal step adds at least one row, so the
procedure terminates after finitely many steps.  At termination write
$T$ for the imposed set and $z=z_T$.  In particular
$|T|\leq(1-\rho)N$ and $\|z\|\leq d/c$.

Put $y'=y+z$.  On $S$ its gaps are
$1+A_Sy'=h_S+(-h_S)_+=(h_S)_+\geq0$.
On $T\setminus S$ the correction has zero image and the original gaps
exceed $u_0$.  Every remaining row has original gap greater than $u_0$
and $|A_jz|\leq u_0/4$, so its new gap is at least $3u_0/4$.
Thus $y'$ is feasible in the ambient space.

\emph{Restore the norm.}
If $n=\|y'\|\geq1$, set $\widehat y=y'/n$.  Each gap becomes
\[
 1+A_\mu\widehat y=(1-n^{-1})+n^{-1}(1+A_\mu y')\geq0.
\]
If $n<1$, take
$W=\ker A_T\cap(y')^\perp$, whose dimension satisfies
$D\geq N-|T|-1\geq\rho N-1\geq\rho N/2$ for large $N$.
Apply Lemma~\ref{thresh:slab-direction} to the remaining rows restricted
to $W$.  There is a unit $v\in W$ such that
\[
 \|A_{T^c}v\|_\infty\leq K\sqrt{N/D}
                      \leq K\sqrt{2/\rho}.
\]
Here the infinity norm of an empty vector is zero.  Set
$\widehat y=y'+\sqrt{1-n^2}\,v$.
Orthogonality gives $\|\widehat y\|=1$, and imposed gaps are unchanged.
Since $n\geq1-\|z\|\geq1-d/c$, we have
$1-n^2\leq2d/c$.  The absolute change of every remaining gap is at most
$K\sqrt{4d/(c\rho)}\leq u_0/2$, leaving it at least $u_0/4$.
This completes the repair.
\end{proof}

\begin{proposition}[Subextensive parent energy implies satisfiability]
\label{thresh:deletion}
Let $0<\alpha<\alpha'\leq B_0$, and let the entries have any one of the
laws $\nu_\eps$, $0\leq\eps\leq1/2$.  If
\[
 \frac1N\E\min H_{\lfloor\alpha'N\rfloor\times N}\longrightarrow0,
\]
then a matrix with $\lfloor\alpha N\rfloor$ rows is satisfiable with
probability tending to one.
\end{proposition}

\begin{proof}
Deleting rows serves one precise purpose: it reduces the number of
small or negative gaps from at most slightly more than $N$ to strictly
fewer than $N$.  The remaining rows can then be repaired using the
surplus dimension in Lemma~\ref{thresh:repair}.

The constants are chosen in the following order.  First fix
$\eta\in(0,1-\alpha/\alpha')$, so that for all sufficiently large $N$
the retained row count $m=\lfloor\alpha N\rfloor$ is at most
$(1-\eta)M$, where $M=\lfloor\alpha'N\rfloor$.  Next choose
$0<\rho<\min(1/8,\eta/8)$ and obtain $K,c,u_0$ from
Lemma~\ref{thresh:uniform-geometry} for the parent matrix.  Finally fix
$d>0$ satisfying~\eqref{thresh:eq:repair-smallness} and
$d^2/u_0^2\leq\rho/2$.

Choose a measurable approximate minimizer $y$ of the parent matrix,
with energy at most $\min H_A+N^{-1}$.  Markov's inequality and the
hypothesis show that $\|r_A(y)\|\leq d\sqrt N$ with probability tending
to one.  On the uniform geometric event, the set of problematic parent
rows
$S_{\mathrm{par}}=\{\mu:1+A_\mu y\leq u_0\}$ has size at most
$(1+\rho)N$.  Indeed, the rows with $|1+A_\mu y|\leq u_0$ number at most
$(1+\rho/2)N$, while those with $1+A_\mu y<-u_0$ number at most
$\|r_A(y)\|^2/u_0^2\leq\rho N/2$.

Independently of the parent matrix, retain a uniformly chosen set
$I$ of exactly $m$ rows.  Conditional on the parent and $y$, the random
variable $X=|I\cap S_{\mathrm{par}}|$ is hypergeometric, with
\[
 \E[X\mid A,y]\leq(1-\eta)(1+\rho)N\leq(1-7\rho)N.
\]
Writing $X$ as a sum of inclusion indicators gives, with
$k=|S_{\mathrm{par}}|$ and $r=m/M$,
\[
 \Var(X\mid A,y)=kr(1-r)\frac{M-k}{M-1}\leq\frac{k}{4}
                  \leq\frac{(1+\rho)N}{4}.
\]
Here the covariance of two distinct inclusion indicators is
$-r(1-r)/(M-1)$; the formula is also immediate when $k=0$.
Chebyshev's inequality therefore gives
\[
 \Prob\{X>(1-2\rho)N\mid A,y\}
 \leq\frac{1+\rho}{100\rho^2N}.
\]
Thus $X\leq(1-2\rho)N$ with probability tending to one, uniformly over
the parent on the stated event.

The retained matrix inherits the parent's operator-norm and
surjection bounds, and its residual at $y$ is no larger than the
parent residual.  Lemma~\ref{thresh:repair} makes it exactly feasible
on the unit sphere.  Finally, an independently chosen subset of rows
of an iid matrix, listed in increasing order, has precisely the law
of an iid $m\times N$ matrix.  This proves the assertion for the desired
row count.
\end{proof}

\subsection{The common threshold and subsequent ordered limits}

\begin{theorem}[Sharp common satisfiability threshold]
\label{thresh:sharp}
For Gaussian entries and, with the choice of $B_0$ above, for every
fixed $0\leq\eps<\eps_0(B_0)$, the number $\alpha_J$ in
\eqref{thresh:eq:threshold-definition} satisfies
$1\leq\alpha_J<\infty$ and
\[
 \lim_{N\to\infty}\Prob(\mathrm{SAT})=
 \begin{cases}
 1,&0<\alpha<\alpha_J,\\
 0,&\alpha>\alpha_J.
 \end{cases}
\]
In particular the threshold itself agrees for all these laws.
No assertion at $\alpha=\alpha_J$ is needed here.
\end{theorem}

\begin{proof}
If $\alpha<\alpha_J$, choose a fixed
$\alpha'\in(\alpha,\alpha_J)$.  The expected parent ground energy is
$o(N)$ for both laws by Lemma~\ref{thresh:pressure-energy} and pressure
universality.  Proposition~\ref{thresh:deletion} proves satisfiability
with probability tending to one.

For $\alpha_J<\alpha\leq B_0$, the same ground-energy limit is strictly
positive.  The map $A\mapsto R(A)$ is $1$-Lipschitz in Frobenius norm:
for every fixed unit $y$,
\[
 \bigl|\|r_A(y)\|-\|r_{A'}(y)\|\bigr|
 \leq\|(A-A')y\|\leq\|A-A'\|_{\mathrm F},
\]
and taking minima preserves this bound.  Gaussian concentration gives
$\Var R(A)\leq C$ for Gaussian entries.  For mixture entries use the
coordinatewise Gaussian quantile representation from
Section~\ref{sec:universality}, whose Lipschitz constant is at most
$2\sqrt{1+\eps}$; the same variance bound follows with an absolute
constant for $\eps\leq1/2$.  Since
\[
 \frac{\E R(A)^2}{N}=2\frac{\E\min H_A}{N}\longrightarrow2e(\alpha),
\]
we obtain
$\E R(A)/\sqrt N\to\sqrt{2e(\alpha)}>0$ by subtracting the bounded
variance.  Chebyshev's inequality now implies
\[
 \Prob(\mathrm{SAT})=\Prob(R(A)=0)
 \leq\frac{\Var R(A)}{(\E R(A))^2}\longrightarrow0.
\]
For larger densities, retain the first
$\lfloor\alpha_*N\rfloor$ rows at any fixed
$\alpha_*\in(\alpha_J,B_0)$ and use monotonicity of satisfiability.
\end{proof}

\begin{corollary}[Transfer of the remaining ordered limits]
\label{thresh:ordered-limits}
Fix $0\leq\eps<\eps_0(B_0)$.  On a full interval above the common
threshold, the thermodynamic gap and normalized-force observables at
every fixed positive temperature and fixed argument agree with their
Gaussian counterparts.  Consequently each further ordered limit in
the problem exists for the mixture law if and only if it exists for
the Gaussian law, and their values agree whenever it exists.
\end{corollary}

\begin{proof}
The fixed-temperature row-marginal result in
Section~\ref{sec:universality} gives equality of the inner
$N\to\infty$ limits for every fixed $\beta$, argument, and
$\alpha\in(\alpha_J,B_0]$.  Limits of these identical functions, taken
successively in the prescribed order, have identical existence and
values.  This assertion makes no exchange of limits; existence of the
Gaussian zero-temperature and jamming limits is established in the
following sections.
\end{proof}


\section{Zero temperature and mechanical bounds}
\label{sec:zero-temperature}

The pressure formula supplies variational minimizers at each positive
temperature.  We now extract their zero-temperature limits and obtain the
mechanical estimates that will control the subsequent approach to the
threshold.  There are two endpoints to distinguish: the diffusion immediately
before the concentrated terminal layer, and the physical row
gap after that layer.  We first derive their relation.  A supporting
functional then identifies the negative part of the physical law with
ground-minimizer residuals.  Finally the contact count, spectral estimate,
and small-gap bounds give control uniform as the density approaches the
threshold.  Section~\ref{sec:critical} will use that control to preserve
mass and moments in the critical limit.

Throughout this section the rows are Gaussian and
\(\alpha_J<\alpha\le B_0\) is fixed.  Write \(e=e(\alpha)>0\).
Constants with a subscript \(e\) may depend on this fixed density;
constants called uniform are uniform on a fixed bounded interval of
UNSAT densities adjacent to \(\alpha_J\).  All assertions about a random
matrix are in probability as \(N\to\infty\), at fixed values of every
cutoff occurring in the assertion.  No uniqueness of the limiting
variational parameters is assumed in this section.

\subsection{Compactness of the variational parameters}

Let \(q_\beta\) be any minimizer in \eqref{eq:src9}.  Denote its
variance-time row solution by $f_\beta$; thus, in the notation of
\eqref{eq:conc-variance-pde},
\[
 (f_\beta)_t=-\tfrac12\bigl((f_\beta)_{zz}
                  +m_{q_\beta}(t)(f_\beta)_z^2\bigr),\qquad
 f_\beta(1,z)=-\beta V(z),\qquad
 f_\beta(0,0)=\Phi_{-\beta V}(q_\beta).
\]
Put
\[
 g_\beta=f_\beta/\beta,\qquad
 \gamma_\beta(t)=\beta m_{q_\beta}(t),\qquad
 \bar\chi_\beta(t)=\beta\chi_{q_\beta}(t),\qquad
 q_{\beta,*}=\esssup q_\beta.
\]
Thus \(g_{\beta,z}\ge0\), \(-1\le g_{\beta,zz}\le0\), and
\(\gamma_\beta=\beta\) on \((q_{\beta,*},1)\).
The next lemma shows that this last interval shrinks while its mass
$\beta(1-q_{\beta,*})$ stays bounded away from zero.  The limiting
scaled measure therefore retains an endpoint mass, denoted by $\Delta$
below; mass from just below $q_{\beta,*}$ can also contribute to it.

\begin{lemma}[UNSAT compactness bounds]\label{zero:compactness}
For all sufficiently large \(\beta\),
\begin{equation}
 0<c\le\bar\chi_\beta(0)\le C_e,
 \qquad 0<c_e\le\beta(1-q_{\beta,*})\le C_e.
 \label{eq:src14}
\end{equation}
Consequently every sequence \(\beta\to\infty\) has a subsequence for which
\[
 \gamma_\beta(t)\,\dd t\ \Longrightarrow\
 \gamma(t)\,\dd t+\Delta\delta_1,
 \qquad
 \bar\chi(t)=\Delta+\int_t^1\gamma(v)\,\dd v,
 \qquad L=\bar\chi(0),
\]
where \(\gamma\ge0\) is nondecreasing and integrable and \(\Delta\ge c_e\).
The convergence of \(\gamma_\beta\) is in \(L^1([0,T])\) for every
\(T<1\), and \(\bar\chi_\beta(t)\to\bar\chi(t)\) for \(t<1\).
\end{lemma}

Here $\gamma_\beta$ and $\gamma(t)$ are variational profiles; the
unindexed scalar $\gamma$ in Theorem~\ref{thm:main} denotes the gap
exponent.

\begin{proof}
The entropy \(E(q)\) is nonnegative.  Since
\(P(-\beta V)/\beta\to-e\), at a minimizer
\(\alpha g_\beta(0,0)\le-e/2\) for large \(\beta\).
Jensen's inequality in the row recursion gives
\(g_\beta(0,z)\ge-\E V(z+G)\), where \(G\sim\cN(0,1)\).
Choose \(R_e>0\) so that \(\E V(R_e+G)<e/(4\alpha)\).
Concavity then implies
\(g_{\beta,z}(0,0)\ge e/(4\alpha R_e)\).
After dividing \eqref{eq:src10} by the appropriate powers of \(\beta\),
\[
 \alpha g_{\beta,z}(0,0)^2\bar\chi_\beta(0)^2\le1,
\]
which proves the upper bound for \(\bar\chi_\beta(0)\).
Conversely \(E(q_\beta)\le C\beta\), since the pressure is nonpositive
and \(\Phi_{-\beta V}\ge-\beta\E V(G)\).  The integral version of
\eqref{eq:src3}, restricted to \([0,1/2]\), gives
\[
 E(q_\beta)\ge\frac1{4\chi_{q_\beta}(0)}-\frac{\log2}{2},
\]
and hence the lower bound.

We shall also use the following consequences of the spin inverse
relations:
\begin{equation}\label{zero:spin-endpoint}
 \frac1{1-q_{\beta,*}}
 =\frac1{\chi_{q_\beta}(0)}+\int_0^1 s\,\dd p_{q_\beta}(s),
 \qquad
 q_{\beta,*}\le\chi_{q_\beta}(0)^2p_{q_\beta}(1-).
\end{equation}
These follow directly from the continuous inverse relation
\eqref{eq:press-continuous-inverse}.  For brevity write $q=q_\beta$,
$p=p_q$, and $\chi=\chi_q$.  Stieltjes integration by parts and
Fubini's theorem give
\[
 p(1-)=\int_0^{q_*}\frac{\dd t}{\chi(t)^2},\qquad
 \int_0^1s\,\dd p(s)
 =p(1-)-\int_0^1p(s)\,\dd s
 =\int_0^{q_*}\frac{m_q(t)}{\chi(t)^2}\,\dd t.
\]
Since $\chi'=-m_q$ and $\chi(q_*)=1-q_*$, the last integral is
$(1-q_*)^{-1}-\chi(0)^{-1}$.  Also $\chi(t)\le\chi(0)$ in the
first integral, so $p(1-)\ge q_*/\chi(0)^2$.  This proves both
relations, including general profiles and their top limits.

Let \(Z_t\) solve
\(\dd Z_t=m_{q_\beta}(t)f_{\beta,z}(t,Z_t)\,\dd t+\dd B_t\),
\(Z_0=0\), and put
\(\mathscr P_\beta(t)=\alpha\E f_{\beta,z}(t,Z_t)^2\).
The derivative martingale and It\^o's isometry imply
\(\mathscr P_\beta'(t)=\alpha\E f_{\beta,zz}(t,Z_t)^2\le\alpha\beta^2\),
and \(p_{q_\beta}(s)=\mathscr P_\beta(q_\beta(s))\).
The Stieltjes chain rule, including jumps by integration over the
intervening variance interval, yields
\[
 \int s\,\dd p_{q_\beta}(s)
 \le\alpha\beta^2\int s\,\dd q_\beta(s)
 \le\alpha\beta^2\chi_{q_\beta}(0).
\]
Thus the first relation in \eqref{zero:spin-endpoint} gives
\[
 \frac1{\beta(1-q_{\beta,*})}
 \le\frac1{\bar\chi_\beta(0)}+\alpha\bar\chi_\beta(0).
\]
The reverse endpoint bound follows from
\(1-q_{\beta,*}\le\chi_{q_\beta}(0)\).

Finally, bounded mass and monotonicity give subsequential weak
compactness of \(\gamma_\beta\,\dd t\).  On every compact interval
before 1, monotonicity bounds the densities uniformly and gives
\(L^1\) compactness.  Any singular mass is therefore at 1.
The terminal interval has mass \(\beta(1-q_{\beta,*})\ge c_e\),
so the endpoint atom has mass at least \(c_e\).
\end{proof}

\begin{proposition}[Zero-temperature variational principle]
\label{zero:variational}
For every subsequential limit in Lemma~\ref{zero:compactness}, let \(g\)
be the solution of
\[
 g_t=-\tfrac12\bigl(g_{zz}+\gamma(t)g_z^2\bigr),
 \qquad g(1,z)=-\frac{V(z)}{1+\Delta}.
\]
It satisfies
\begin{equation}
 \mathcal E_\alpha(\gamma,\Delta)
 :=\alpha g(0,0)+\frac12\int_0^1\frac{\dd t}{\bar\chi(t)}
 =-e,
 \label{eq:src15}
\end{equation}
and minimizes this functional over nonnegative nondecreasing
integrable \(\gamma\) and \(\Delta>0\).
Here solutions have quadratic growth, nonnegative first derivative,
and second derivative in \([-1,0]\); an integrable time coefficient
is interpreted by approximation in \(L^1\).
\end{proposition}

\begin{proof}
We give the endpoint comparison explicitly.  For $t<1$, set
\[
 U_\beta(t,z)=\sup_{w\in\R}
 \left\{-V(w)-\frac{(z-w)^2}{2\bar\chi_\beta(t)}\right\}
 =-\frac{V(z)}{1+\bar\chi_\beta(t)}.
\]
At $t=1$ use the continuous extension $U_\beta(1,z)=-V(z)$.
This solves \(U_t=-\gamma_\beta U_z^2/2\) and has curvature in
\([-1,0]\).  Subtraction from the viscous equation for \(g_\beta\)
produces a linear equation with drift
\(\gamma_\beta(g_{\beta,z}+U_{\beta,z})/2\) and a source of absolute
value at most \(1/2\).  Its diffusion representation, or comparison
after smoothing the hinge, gives
\begin{equation}\label{zero:terminal-comparison}
 \sup_z|g_\beta(t,z)-U_\beta(t,z)|\le(1-t)/2,
 \qquad
 \sup_z|g_{\beta,z}(t,z)-U_{\beta,z}(t,z)|\le C\sqrt{1-t}.
\end{equation}
For the gradient bound use the value bound and the common Lipschitz
bound on the two gradients, comparing difference quotients with
step \(\sqrt{1-t}\).  The identical comparison holds for \(g\), with
\(\bar\chi\) in place of \(\bar\chi_\beta\).

On \([0,T]\), subtraction of two viscous equations again gives a
linear equation.  Its source is bounded by the coefficient difference
times \(C(1+z^2)\), and its drift has a spatial Lipschitz constant
bounded by an integrable function with uniformly bounded integral.
The associated diffusions have uniform moments of every fixed order,
by Gronwall's inequality.  The diffusion representation therefore
proves stability under \(L^1\) coefficient convergence and local
uniform terminal convergence with a common quadratic growth bound.
Applying this at \(T<1\), using \eqref{zero:terminal-comparison}, and
then sending \(T\uparrow1\), proves local uniform convergence of
values and gradients on compact time intervals before 1.  The same
argument constructs, and proves uniqueness of, the stated solution
for an integrable \(\gamma\).

For the entropy, \eqref{eq:src3} gives
\[
 \frac{E(q_\beta)}\beta
 =\frac12\int_0^{q_{\beta,*}}\frac{\dd t}{\bar\chi_\beta(t)}
   +\frac{\log(1-q_{\beta,*})}{2\beta}.
\]
On this integration interval the denominator is at least
\(\beta(1-q_{\beta,*})\ge c_e\).  Dominated convergence applies,
and the logarithmic term is \(O_e(\log\beta/\beta)\).  This proves
\eqref{eq:src15}.

For an arbitrary trial pair, use the finite-temperature CDF
\(m_\beta(t)=\min\{\gamma(t)/\beta,1\}\) until
\(1-\Delta/\beta\), and set \(m_\beta=1\) thereafter.
This is nondecreasing and its scaled measure converges to
\(\gamma\,\dd t+\Delta\delta_1\); integrability of \(\gamma\)
ensures that its mass on the shrinking terminal interval vanishes.
The comparisons just proved show convergence of its row term and
entropy.  The finite-temperature minimum is at most this trial value,
so \(-e\le\mathcal E_\alpha(\gamma,\Delta)\).
\end{proof}

\subsection{Endpoint laws and the negative marginal}

\begin{proposition}[Subsequential endpoint law]\label{zero:endpoint}
Let \(Z_0=0\) and
\(\dd Z_t=\gamma(t)g_z(t,Z_t)\,\dd t+\dd B_t\), and put \(H=1+Z_1\).
Along the subsequence defining \((\gamma,\Delta)\), the tilted row gap
converges in distribution to
\begin{equation}
 h^*=\begin{cases}
 H,&H\ge0,\\[2pt]
 H/(1+\Delta),&H<0.
 \end{cases}
 \label{eq:src16}
\end{equation}
Every fixed moment of the negative part converges.  Both \(H\) and
\(h^*\) are atomless and assign positive probability to every nonempty
open interval.  This statement is subsequential; uniqueness is addressed
separately below and in Section~\ref{sec:critical}.
\end{proposition}

\begin{proof}
On each constant-CDF interval the row tilt is the Gaussian Doob
transform.  Composing these transforms, then approximating a general
CDF, identifies its physical gap with $1+Z_1^\beta$, where
\[
 Z_t^\beta=B_t+\int_0^t
       \gamma_\beta(v)g_{\beta,z}(v,Z_v^\beta)\,\dd v.
\]
Couple $Z^\beta$ and $Z$ with the same Brownian motion.  Before a fixed time \(T<1\), the coefficient and
gradient convergence above and Gronwall's inequality imply convergence
in expected supremum norm.  On \([T,1]\), replace the gradient by
\(U_{\beta,z}\).  The error has expected size at most
\(C_e\sqrt{1-T}\), by \eqref{zero:terminal-comparison} and the uniform
bound on \(\int\gamma_\beta\); the Brownian increment has the same
bound.  The resulting deterministic flow fixes positive gaps and
multiplies a negative gap between times \(T\) and \(t\) by
\((1+\bar\chi_\beta(t))/(1+\bar\chi_\beta(T))\).
First send \(\beta\to\infty\), and then \(T\uparrow1\).
The limiting integrable drift on \([T,1]\) has vanishing total
mass, whereas the total mass concentrating near time $1$ in
$\gamma_\beta(t)\,\dd t$ is $\Delta$.  Since
$\bar\chi_\beta(1)=0$ and $\bar\chi_\beta(T)\to\bar\chi(T)$,
the displayed flow factor tends to $(1+\Delta)^{-1}$ as
$\beta\to\infty$ and then $T\uparrow1$.  This proves
\eqref{eq:src16}.
Since the drift is nonnegative, the negative part of the gap is
bounded by \((1+B_1)_-\); this gives uniform integrability of every
moment of the negative part.

For atomlessness, condition on the Brownian bridge
\(\widetilde B_t=B_t-tB_1\).  The map from \(b=B_1\) to \(Z_1\)
is increasing and has slope between \(\exp(-\int\gamma)\) and 1.
Indeed its variation satisfies
\(J'(t)=1+\gamma(t)g_{zz}(t,Z_t)J(t)\), \(J(0)=0\).
Approximation proves the same bounds when derivatives exist only
almost everywhere.  The conditional endpoint law is thus the
pushforward of a Gaussian under a globally bi-Lipschitz map.
That map is onto \(\R\), so the law has a bounded density and
assigns positive probability to every nonempty open interval.  The piecewise linear map in
\eqref{eq:src16} preserves these properties.
\end{proof}

The variable $H$ is the pre-layer gap; $h^*$ is the gap sampled by the
microscopic observable.  In particular their signs agree, but a negative
gap is reduced by the factor $1+\Delta$.  Define
\[
 p=\Prob(h^*<0)=\Prob(H<0),\qquad
 \mu=\E(h^*)_-,\qquad z_c=\alpha p.
\]
The parameter $z_c$ will be identified with the number of active
constraints per spin coordinate.  That identification requires a
small-residual estimate: convergence against tests supported away from
zero alone would miss constraints whose negative gaps tend to zero.

\begin{proposition}[Identification of negative gaps]
\label{zero:negative-law}
For every subsequential endpoint law in Proposition~\ref{zero:endpoint},
the restriction of the law of $h^*$ to $(-\infty,0)$ is the same.  For
$\psi\in C_c^\infty((-\infty,0))$, uniformly over Gaussian ground
minimizers $y_N$,
\[
 \frac1N\sum_{i=1}^M\psi(1+A_i y_N)
       \longrightarrow\alpha\E\psi(h^*)
 \quad\text{in probability}.
\]
Moreover,
\begin{equation}\label{zero:energy-moment}
 \alpha\E(h^*)_-^2=2e.
\end{equation}
\end{proposition}

\begin{proof}
We perturb the loss in a compact negative-gap interval.  For either
sign of the perturbation, the pressure trial gives a lower support for
the ground energy.  Evaluating the perturbed energy at an unperturbed
minimizer gives the opposite inequality, and the two one-sided slopes
identify its empirical test average.  The perturbation interval must
remain fixed as $\beta\to\infty$.

Let \(\psi\) be smooth with
support in a compact negative-gap interval \([-R,-a]\), \(a>0\), and
put \(W_w(z)=V(z)+w\psi(1+z)\).  Choose \(w_0>0\) so that
\[
 w_0\|\psi''\|_\infty\le\tfrac12,\qquad
 w_0\|\psi'\|_\infty\le a/2,\qquad
 w_0\|\psi\|_\infty\le a^2/4.
\]
For \(|w|\le w_0\), \(W_w\) is convex, nonnegative, and
nonincreasing, and agrees with \(V\) outside that compact interval.
In particular \(-\beta W_w\) is concave for every \(\beta>0\).
This is a fixed interval of loss perturbations, independent of
\(\beta\).

Here only a trial upper bound is needed.  The Gaussian interpolation
with smooth concave tangent caps, followed by the spin conjugate
inequality (Corollary~\ref{cor:press-convex-loss-upper}), gives for
every \(q_*<1\)
\begin{equation}\label{zero:fixed-loss-pressure-trial}
 \limsup_{N\to\infty}P_N(-\beta W_w)
 \le\alpha\Phi_{-\beta W_w}(q)+E(q).
\end{equation}
This uses Corollary~\ref{cor:conc-convex-loss}, namely ordinary
profile concavity for a concave terminal, not the
small-terminal-perturbation radius in \eqref{eq:src11}.  The tangent
caps have the required concavity, and removal of the cap is valid at
each fixed \(\beta,w\), because \(W_w\) has the same quadratic tail
as \(V\).  Neither strict profile concavity nor identification of a
perturbed minimizer is required for
\eqref{zero:fixed-loss-pressure-trial}.

For a base limiting pair \((\gamma,\Delta)\), let \(g_w\) solve the
PDE with coefficient \(\gamma\) and terminal value
\[
 g_w(1,z)=\sup_{v\in\R}
 \left\{-W_w(v)-\frac{(z-v)^2}{2\Delta}\right\},
 \qquad
 S(w)=-\alpha g_w(0,0)-\frac12\int_0^1\frac{\dd t}{\bar\chi(t)}.
\]
Write \(\widehat e_N(w)=N^{-1}\min_{\|y\|=1}\sum_iW_w(A_i y)\)
and \(e_N(w)=\E\widehat e_N(w)\).
The surface-cap pressure--energy estimate applies to \(W_w\), whose
curvature and growth bounds are uniform for \(|w|\le w_0\).
Use the recovery profiles of Proposition~\ref{zero:variational} in
\eqref{zero:fixed-loss-pressure-trial}, take \(N\to\infty\), and
then \(\beta\to\infty\).  For completeness, the endpoint comparison
for this nonquadratic loss uses
\[
 U_{\beta,w}(t,z)=\sup_v\left\{-W_w(v)
               -\frac{(z-v)^2}{2\bar\chi_\beta(t)}\right\},\qquad t<1,
\]
with continuous terminal value $U_{\beta,w}(1,z)=-W_w(z)$.
Both \(U_{\beta,w}\) and the scaled viscous solution
\(g_{\beta,w}\) have curvature in \([-C,0]\), uniformly for
\(|w|\le w_0\), and \(U_{\beta,w,t}=-\gamma_\beta U_{\beta,w,z}^2/2\).
The difference equation has forcing bounded by \(C/2\), so
\[
 \sup_z|g_{\beta,w}(t,z)-U_{\beta,w}(t,z)|\le C(1-t),\qquad
 \sup_z|g_{\beta,w,z}(t,z)-U_{\beta,w,z}(t,z)|\le C\sqrt{1-t}.
\]
Recovery has \(\int\gamma_\beta=L+o(1)\); hence its drift Lipschitz
constants have bounded integrals and its diffusions have uniform
moments.  Compact-time \(L^1\) coefficient stability, followed by
\(t\uparrow1\), gives exactly the terminal supremum defining
\(g_w\), and the entropy converges as before.  We obtain
\begin{equation}\label{zero:fixed-loss-support}
 \liminf_{N\to\infty}e_N(w)\ge S(w),\qquad S(0)=e.
\end{equation}
No existence of a perturbed thermodynamic ground-energy limit is
being asserted or used.

The function \(S\) is differentiable at zero.  The quadratic
penalization in the endpoint supremum is strictly concave in \(v\),
so the optimizer is unique; its envelope derivative is minus
\(\psi\) at the proximal gap.  Linearization of the PDE propagates
this bounded terminal derivative along the tilted diffusion.
The coefficient stability follows from the difference equation
used above.  Hence
\begin{equation}\label{zero:fixed-loss-derivative}
 S'(0)=\alpha\E\psi(h^*).
\end{equation}

For every base ground minimizer \(y_N\), let
\(T_N=N^{-1}\sum_i\psi(1+A_i y_N)\).  Direct evaluation gives
\[
 \widehat e_N(w)\le\widehat e_N(0)+wT_N.
\]
For each fixed \(w\), the centered ground energy
\(\widehat e_N(w)-e_N(w)\) tends to zero in probability.
Indeed on \(\|A\|\le K\sqrt N\) its matrix Lipschitz constant is
\(C_w/\sqrt N\), uniformly over spins; extend Lipschitz and use
Gaussian concentration, and then control the complement by the
quadratic growth bound.  Thus \eqref{zero:fixed-loss-support} and
the known convergence \(\widehat e_N(0)\to e\) imply, for fixed
\(w>0\) and \(v<0\), uniformly over all base minimizers,
\[
 \frac{S(w)-e}{w}-o_{\Prob}(1)
 \le T_N\le
 \frac{S(v)-e}{v}+o_{\Prob}(1).
\]
Sending \(w\downarrow0\) and \(v\uparrow0\) proves
\(T_N\to S'(0)\) in probability.  Apply this argument to any two
base limiting pairs: the same empirical tests have the two asserted
limits, so their derivatives agree.  Smooth compact negative tests
determine the restriction of a measure to \((-\infty,0)\), proving
the asserted uniqueness and, at the same time, microscopic
negative-test convergence.

For the energy identity use \(W_w=(1+w)V\), \(|w|<1/2\).
The same trial-support argument applies, while now
\(\widehat e_N(w)=(1+w)\widehat e_N(0)\) identically.
The derivative squeeze therefore gives
\(e=S'(0)=\alpha\E V(h^*-1)\), proving
\eqref{zero:energy-moment}.  Quadratic growth in this differentiation
is justified by the uniform diffusion moments.
\end{proof}

The first moment of the negative part also converges at microscopic
ground minimizers: outside residual size \(R\), its empirical
contribution is at most \(2H_{\min}/(MR)\).  Together with the compact
negative-test convergence just proved, this controls the far tail.
For any subsequential thermal limit at fixed \(\alpha\), weak
convergence to \eqref{eq:src16}, atomlessness at zero, and first-moment
convergence identify the force CDF as
\[
 F_{\alpha,\infty}(s)
 =\frac1p\Prob\left(0<(h^*)_-\le s\frac\mu p\right),\qquad s>0.
\]
The conditional mean residual is $\mu/p$; $p,\mu>0$ by
\eqref{zero:energy-moment}.  Proposition~\ref{zero:negative-law} makes
this expression independent of the temperature subsequence.  Identification with ground-minimizer
force statistics additionally uses the active-mass cutoff proved in
Lemma~\ref{zero:negative-cutoffs}.  Positive gap tests will be identified
uniquely only in Section~\ref{sec:critical}.

\subsection{Support stability}

\begin{lemma}[Endpoint stability]\label{zero:stability}
For every subsequential zero-temperature optimizer, the contact
density parameter \(z_c=\alpha\Prob(H<0)\) satisfies
\begin{equation}
 z_c\le(1+1/\Delta)^2.
 \label{eq:src17}
\end{equation}
\end{lemma}

\begin{proof}
Let \(Z^\beta\) denote the finite-temperature tilted diffusion.
At finite temperature vary the overlap distribution, with CDF \(m\),
by moving mass to a point below 1.  Linearization of the PDE and
\eqref{eq:src3} gives, for variations supported below a cutoff \(T<1\),
\[
 \delta(\alpha\Phi+E)=\frac12\int_0^T\delta m(t)J(t)\,\dd t,
 \qquad
 J(t)=\mathscr P_\beta(t)-\int_0^t\frac{\dd v}{\chi_{q_\beta}(v)^2}.
\]
Let \(\nu\) be the minimizing overlap distribution, choose
\(q_{\beta,*}<T<1\), and fix \(s\in[0,T)\).
The feasible direction
\(\nu_\eps=(1-\eps)\nu+\eps\delta_s\) has
\(m_\eps(t)-m(t)=\eps(\ind_{\{s\le t\}}-m(t))\) for \(t\le T\),
and vanishes for \(t>T\).  With \(K(s)=\int_s^T J(t)\,\dd t\),
Fubini's theorem makes its directional derivative
\[
 \frac12\left(K(s)-\int K(r)\,\nu(\dd r)\right)\ge0.
\]
Therefore \(K\) is at least its \(\nu\)-average everywhere on
\([0,T)\), and equals that average on the support, by continuity.
Thus every support point in \((0,T)\) is a local minimum of \(K\).
Since \(K'=-J\) and \(K''=-J'\), at each such point
\[
 \alpha\E f_{\beta,zz}(s,Z^\beta_s)^2\chi_{q_\beta}(s)^2\le1.
\]
The derivatives \(J'\) and \(K''\) are continuous before time 1,
including across jumps of \(m\).  Indeed the mild heat equation and
the local spatial Lipschitz bound on \(f_{\beta,z}^2\) give joint
continuity of \(f_{\beta,zz}\) away from the terminal time; the
second derivative of the heat kernel applied to a Lipschitz function
has an integrable time singularity.  The diffusion moment formula
then gives continuity of \(\mathscr P_\beta'\), and explicitly
\[
 J'(s)=\alpha\E f_{\beta,zz}(s,Z_s^\beta)^2
                   -\chi_{q_\beta}(s)^{-2}.
\]  No continuity of
\(f_{\beta,t}\) at a jump of \(m\) is asserted.

Select continuity times \(t_n\uparrow1\) with
\((1-t_n)\gamma(t_n)\to0\); this follows from integrability and
monotonicity.  A diagonal choice \(t_\beta\uparrow1\) ensures
\(t_\beta<q_{\beta,*}\),
\(\bar\chi_\beta(t_\beta)\to\Delta\), and
\((1-t_\beta)\gamma_\beta(t_\beta)\to0\).
Let \(s_\beta\) be the next overlap-support point at or above
\(t_\beta\).  The CDF is constant between these points, so
\[
 \bar\chi_\beta(t_\beta)-\bar\chi_\beta(s_\beta)
 =\gamma_\beta(t_\beta)(s_\beta-t_\beta)
 \le(1-t_\beta)\gamma_\beta(t_\beta)\longrightarrow0.
\]
Consequently \(\bar\chi_\beta(s_\beta)\to\Delta\).
Moreover, for fixed \(T<1\),
\[
 \int_T^{s_\beta}\gamma_\beta(t)\,\dd t
 \longrightarrow\int_T^1\gamma(t)\,\dd t.
\]
The drift over this interval is negligible after \(T\uparrow1\),
so \(Z^\beta_{s_\beta}\) converges to the pre-atom endpoint \(Z_1\).

By \eqref{zero:terminal-comparison},
\(g_{\beta,z}(s_\beta,z)\to(-1-z)_+/(1+\Delta)\) locally uniformly.
These gradients are convex in \(z\): differentiating the PDE three
times gives a linear equation preserving the nonnegative third
spatial derivative of a smoothed hinge terminal, and passage to the
limit preserves convexity.  Secant bounds for convex functions
therefore give locally uniform convergence of their derivatives
on each compact interval avoiding \(z=-1\): the limit derivative
is continuous there, and a finite secant mesh gives a uniform bound.
Since these derivatives are bounded and \(Z_1\) is atomless, the
finite-temperature stability inequality passes to
\[
 \alpha\frac{\Delta^2}{(1+\Delta)^2}\Prob(Z_1<-1)\le1,
\]
as asserted.
\end{proof}

\subsection{Mechanical contacts and uniform cutoff estimates}

We now relate these variational quantities to finite-dimensional
minimizers.  Three separate facts are needed.  A positive-energy
minimizer has at least $N$ strictly violated rows and almost surely no
row at exactly zero gap.  A uniform small-residual bound then prevents
an extensive number of those violated rows from disappearing at zero
in the large-$N$ limit.  A different bound, proved on near-ground
configurations later in the section, controls small positive gaps in
the prescribed thermal order.  Write \(\Pi_y=I-yy^T\), \(r_i(y)=(-1-A_i y)_+\), and
\(\nabla_S H(y)=\Pi_y\nabla H(y)\) on the unit sphere.

\begin{lemma}[Mechanical contact count]\label{zero:mechanical}
Almost surely, every positive-energy Gaussian ground minimizer has
no zero gap outside its active set \(C=\{i:r_i>0\}\), and
\(k=|C|\ge N\).  At such a minimizer,
\begin{equation}\label{zero:mechanical-hessian}
 A_C^T(A_Cy+\mathbf1)=\lambda y,
 \qquad \lambda=\sum_{i\in C}r_i(1+r_i)>0,
 \qquad
 v^T(A_C^TA_C-\lambda I)v\ge0\quad(v\perp y).
\end{equation}
\end{lemma}

\begin{proof}
For a fixed \(C\), Gaussian genericity gives simple nonzero singular
values of \(A_C\) and a nonzero component of \(A_C^T\mathbf1\) in
each corresponding right singular direction.  These properties hold
almost surely, since their exceptional sets are proper algebraic
sets.  Stationarity implies the first identity in
\eqref{zero:mechanical-hessian}; its scalar product with \(y\) gives
\(\lambda>0\).  The genericity just noted precludes \(\lambda\)
from being an eigenvalue of \(A_C^TA_C\).  Thus the stationary vector
is obtained by inversion.  Its unit-norm equation is rational in
\(\lambda\), is not identically satisfied (it tends to zero at
infinity), and has finitely many roots.  Conditional on \(A_C\),
any row outside \(C\) almost surely avoids gap zero at all these
finitely many vectors.  A union over the finitely many sets \(C\)
proves the first assertion, even if such a row was initially allowed
to have zero gap.  The energy is now smooth near the minimizer, so
its constrained Hessian is nonnegative on the tangent space.
Finally stationarity with \(\lambda>0\) puts \(y\) in the row span
of \(A_C\).  A nontrivial kernel of \(A_C\) would be tangent and
would have Hessian value \(-\lambda\), a contradiction.  Hence
\(\operatorname{rank} A_C=N\), and \(k\ge N\).
\end{proof}

We record the common counting estimate behind the cutoff bounds.
Work on \(\|A\|\le K\sqrt N\), whose probability tends to one for
fixed sufficiently large \(K=K(B_0)\).  Both the residual map and the
tangent gradient obey
\begin{equation}\label{zero:lipschitz}
 \|r(y)-r(y')\|_2\le K\sqrt N\|y-y'\|,
 \qquad
 \|\nabla_S H(y)-\nabla_S H(y')\|\le C N\|y-y'\|.
\end{equation}
At a deterministic unit vector \(y_0\), condition on all scalar
projections \(A_i y_0\).  The tangent row components remain independent
standard Gaussians, and therefore
\[
 \nabla_S H(y_0)\mid(A_i y_0)_i
 \sim\cN\bigl(0,\|r(y_0)\|_2^2 I_{N-1}\bigr).
\]
For \(\|r(y_0)\|_2\ge t\sqrt N/2\), its small-ball bound is
\begin{equation}\label{zero:gradient-small-ball}
 \Prob\bigl(\|\nabla_S H(y_0)\|\le CaN\mid(A_i y_0)_i\bigr)
 \le(Ca/t)^{N-1}.
\end{equation}
A deterministic net of mesh proportional to \(a\) has at most
\((C/a)^N\) points.  Thus a projection event of probability at most
\(C^N t^k w^j\), involving specified row sets, has union probability
at most
\begin{equation}\label{zero:net-master}
 a^{-1}C^N t^{k-N+1}w^j
\end{equation}
provided the original projection event transfers to that event
at a net point, and the configurations satisfy
\(\|\nabla_S H\|\le CaN\) and have residual size in
\([t\sqrt N,2t\sqrt N]\).  Choose the mesh at most $t/(2K)$;
then \eqref{zero:lipschitz} ensures the required residual lower
bound $t\sqrt N/2$ at the net point.  The factor $t^{k-N+1}$ records the useful cancellation: forcing $k$
active projections into a ball of radius $O(t\sqrt N)$ contributes
$t^k$, while making the tangent gradient small costs $t^{-(N-1)}$.
With at least $N$ active rows, no negative power growing with $N$
remains as $t\downarrow0$.
Here all choices of row subsets, at most
\(3^M\), have been absorbed into \(C^N\).  The mesh \(a>0\) is
chosen after \(t,w\) and every other cutoff, and is then held fixed
as \(N\to\infty\).  In particular the factor \(a^{-1}\) does not
affect exponential decay.  The operator-norm event is used only to
transfer a configuration to the net, never in the conditional
Gaussian probability calculation.

\begin{lemma}[Negative cutoffs]\label{zero:negative-cutoffs}
There are constants $c,C>0$, independent of the density, such that
$\alpha\mu\ge c\sqrt e$, where $\mu=\E(h^*)_-$.  At each fixed
UNSAT density, uniformly over Gaussian ground minimizers,
\[
 \frac{k}{N}\longrightarrow z_c\quad\text{in probability}.
\]
There is \(s_0>0\), independent of the density, such that their
limiting mean-one force law obeys
\(F_{\alpha,\infty}(s)\le C/|\log s|\) for all \(0<s<s_0\),
uniformly as \(\alpha\downarrow\alpha_J\).
Furthermore \(z_c\to1\) in this limit, and every weak subsequential
jamming force law has total mass one and no atom at zero.
\end{lemma}

\begin{proof}
Set $t=\sqrt e$.  Since $H_{\min}/N\to e$, with probability
tending to one,
\[
 \frac{3e}{4}\le\frac{H_{\min}}N\le\frac{5e}{4},\qquad
 t\sqrt N\le\|r\|_2\le2t\sqrt N.
\]
Thus a single residual annulus suffices at every fixed density.
Moreover $t\le\sqrt{B_0}$, since $e(\alpha)\le\alpha\le B_0$.
All exponential bases in the counting estimates can consequently be
chosen independently of the density; the net mesh may still depend
on its fixed value of $t$.
If \(\|r\|_1\le c_1tN\), transfer to a net of sufficiently small
mesh so that the same active set \(C\) has
\(\sum_{i\in C}|1+A_i y_0|\le2c_1tN\).
The joint projection density is bounded by an absolute constant to
the power \(k\), and the \(k\)-dimensional cross-polytope volume is
\(2^k(2c_1tN)^k/k!\).  Since \(k\ge N\), the projection probability
is at most \((Cc_1t)^k\).  Combining with
\eqref{zero:gradient-small-ball} and the net gives
\(a^{-1}C^N c_1^k t^{k-N+1}\), which is exponentially small for a
uniformly small \(c_1\), because \(t\le\sqrt{B_0}\) and \(N\le k\le B_0N\).
Thus \(\|r\|_1\ge c_1tN\); the previously established first-moment
convergence proves the bound for \(\mu\).

Suppose at least \(\eta N\) active residuals are at most \(st\),
where \(0<\eta<1/2\).  Choose
\(a\ll st\sqrt\eta/K\).  By the squared projection displacement
bound, at least \(j=\lfloor\eta N/2\rfloor\) corresponding net gaps
have absolute value at most \(2st\).  The remaining coordinates of
\(C\) lie in a Euclidean ball of radius \(Ct\sqrt N\).
Its volume, together with the \(j\) intervals, bounds the projection
probability by \(C^N t^k s^j\); here \(k-j\ge N/2\) bounds the
ball's dimension factor.  Equation~\eqref{zero:net-master} excludes
this event for \(s\le\exp(-C'/\eta)\), with \(C'\) uniform.
At fixed \(e>0\), this estimate and convergence of compact negative
tests show that no extensive active mass is lost at zero, proving
\(k/N\to\alpha\Prob(h^*<0)=z_c\).

The mean residual among active rows is at most
\(\|r\|_2/\sqrt k\le2t\).
A normalized force at most \(s\) therefore has residual at most
\(Cst\).  The preceding estimate, divided by \(k\ge N\), gives
\(F_{\alpha,\infty}(s)\le C/|\log s|\).
The mean-one laws are tight by Markov's inequality; in fact their
second moments are uniformly bounded, since
\(k\|r\|_2^2/\|r\|_1^2\le C\).  Their weak limits therefore
retain mean one.  The displayed logarithmic bound prevents any
subsequential limit from putting mass at zero.

For the contact density, use the uniform all-spin near-gap estimate
in Lemma~\ref{thresh:uniform-geometry}.  For each
fixed \(\rho>0\), it supplies \(u_0>0\) such that every ground
minimizer satisfies
\[
 k\le(1+\rho)N+2H_{\min}/u_0^2.
\]
Take \(N\to\infty\), then \(\alpha\downarrow\alpha_J\), and finally
\(\rho\downarrow0\).  This gives \(\limsup z_c\le1\), while
Lemma~\ref{zero:mechanical} gives \(z_c\ge1\).
\end{proof}

\Needspace{8\baselineskip}
\subsection{A uniform spectral bound}

The precise Mar\v{c}enko--Pastur input needed here is the following
Gaussian special case~\cite[Theorem~1 and Example~1,
pp.~459--461]{MarcenkoPastur1967}: if \(k_N/N\to z\ge1\)
and \(A_N\) is a \(k_N\)-by-\(N\) matrix of independent standard
Gaussians, the empirical law of the eigenvalues of \(A_N^TA_N/N\)
converges weakly in probability to the probability measure with
density
\[
 \frac{\sqrt{(b_z-x)(x-a_z)}}{2\pi x}\,
 \ind_{[a_z,b_z]}(x),
 \qquad a_z=(\sqrt z-1)^2,\quad b_z=(\sqrt z+1)^2.
\]
In particular every interval immediately to the right of \(a_z\)
has positive limiting mass.  No extreme-eigenvalue theorem is needed.

\begin{lemma}[Spectral and susceptibility bounds]\label{zero:spectral}
For every \(\eps>0\), with probability tending to one, simultaneously
for all row subsets \(C\) with \(N\le|C|\le B_0N\), at least two
singular values of \(A_C/\sqrt N\) are at most
\(\sqrt{|C|/N}-1+\eps\).  Consequently
\begin{equation}\label{zero:susceptibility-bounds}
 2e+\alpha\mu\le(\sqrt{z_c}-1)^2\le\Delta^{-2},
 \qquad
 \Delta\le C e^{-1/4},
 \qquad
 \frac1{\sqrt{2e}}\le L\le\frac C{\sqrt e}.
\end{equation}
In particular \(\Delta/L\to0\) as \(\alpha\downarrow\alpha_J\).
\end{lemma}

\begin{proof}
Fix \(\eps>0\).  For each aspect ratio in \([1,B_0]\), choose a
nonnegative Lipschitz function of the singular value, supported
strictly below \(\sqrt z-1+\eps\), with positive integral under the
limiting singular-value law.  By continuity of the spectral edges and
density, each fixed test stays below the cutoff and retains a positive
limiting integral throughout the closure of a sufficiently small ratio
neighborhood.  A finite collection of these neighborhoods covers the
compact ratio interval, giving a common positive lower bound for the
relevant limiting integrals.  Convergence of their expectations is uniform
within the chosen neighborhoods: otherwise a violating sequence of
ratios would have a convergent subsequence and contradict the stated
spectral limit.  The Frobenius singular-value inequality gives
\[
 A\longmapsto\frac1N\sum_{i=1}^N
 \varphi\bigl(s_i(A/\sqrt N)\bigr)
 \quad\hbox{Lipschitz constant at most }\Lip(\varphi)/N.
\]
Gaussian concentration therefore bounds a fixed positive downward
deviation by \(2\exp(-c_\eps N^2)\).  Union over at most \(2^{B_0N}\)
subsets still has probability tending to zero.  A positive normalized
test average implies a positive multiple of \(N\) singular values
below the cut, hence at least two for large \(N\).

The span of two corresponding right singular vectors intersects
\(y^\perp\).  The Hessian
inequality \eqref{zero:mechanical-hessian} therefore yields
\[
 \frac\lambda N\le(\sqrt{k/N}-1+\eps)^2.
\]
Since \(\lambda/N=2H_{\min}/N+\|r\|_1/N\), first take
\(N\to\infty\), then \(\eps\downarrow0\), and use the contact and
moment convergence.  The upper bound \(\Delta^{-2}\) follows from
\eqref{eq:src17} and \(z_c\ge1\).  The estimate for \(\mu\) in
Lemma~\ref{zero:negative-cutoffs} then gives \(\Delta\le Ce^{-1/4}\).

The second inequality in \eqref{zero:spin-endpoint}, together with
\(p_{q_\beta}(1-)\le\alpha\E f_{\beta,z}(1,Z^\beta_1)^2\), passes to
\(1\le L^2\alpha\E(h^*)_-^2=2eL^2\).
Finally the derivative martingale identifies
\(g_{\beta,z}(0,0)\to\mu\), so \eqref{eq:src10} implies
\(\alpha\mu^2L^2\le1\).  The lower bound for \(\mu\) proves the
upper bound for \(L\).  Dividing the two estimates proves the final
assertion.  This conclusion concerns the thermal endpoint atom;
it does not preclude additional layers concentrating at the endpoint
in the subsequent density limit.
\end{proof}

\Needspace{4\baselineskip}
\begin{corollary}[Cutoffs after susceptibility normalization]
\label{zero:normalized-force-cutoff}
For every subsequential endpoint law, put
\(R=L(h^*)_-\) and \(p=\Prob(h^*<0)\).
There are constants \(c,C,u_0>0\), independent of the density, such
that throughout the fixed UNSAT interval and for \(0<u<u_0\),
\[
 c\le \E R\le C,\qquad
 \Prob(0<R\le u)\le\frac C{|\log u|}
 \qquad \alpha p\longrightarrow1
 \quad(\alpha\downarrow\alpha_J).
\]
\end{corollary}

\begin{proof}
The bounds \(L\ge(2e)^{-1/2}\), \(\alpha\mu\ge c\sqrt e\), and
\(\alpha\mu^2L^2\le1\) give the two bounds for \(\E R=L\mu\).
Moreover \(p=z_c/\alpha\) is bounded away from zero because
\(z_c\ge1\) and \(\alpha\le B_0\).  Thus the conditional mean
\(\E[R\mid R>0]=L\mu/p\) is bounded above and below.
Apply Lemma~\ref{zero:negative-cutoffs} to the corresponding
mean-one force to obtain the displayed small-value bound.
The final assertion is \(z_c\to1\).
\end{proof}

\subsection{Positive gaps in the prescribed temperature order}

\begin{proposition}[Positive-gap cutoffs]\label{zero:positive-cutoffs}
Constants \(c,C,u_0>0\), independent of the density, exist such
that, for \(0<\delta<u<1\),
\(G_{N,\beta,\alpha}(\delta,u)\ge c(u-\delta)\).
For every fixed UNSAT density in the interval under consideration
and every \(0<u<u_0\),
\begin{equation}\label{zero:direct-positive-cutoff}
 \limsup_{\beta\to\infty}\limsup_{N\to\infty}
 \E\left\langle\frac1M\sum_{i=1}^M
       \ind_{\{0<h_i\le u\}}\right\rangle
 \le\frac C{|\log u|}
\end{equation}
with the same constants \(C,u_0\) at every such density.
In particular every subsequential endpoint law in
\eqref{eq:src16} satisfies the direct bound
\(\Prob(0<h^*\le u)\le C/|\log u|\).
The same estimate bounds the ordered zero-temperature and density
limsups of \(G_{N,\beta,\alpha}(\delta,u)\), uniformly in
\(0<\delta<u\).
The upper bound requires no identification of positive gap tests at
exact ground minimizers.
\end{proposition}

\begin{proof}
Remove a tested row.  Its insertion denominator is at most one,
whereas its penalty vanishes on the tested positive interval.
Conditional on the cavity spin, the removed projection is standard
Gaussian.  Its density is bounded below on \([-1,0]\), which proves
the lower bound.  The same proof works for the small mixtures,
conditional on their variances, with a uniform lower density bound.

For the upper bound, exact ground minimizers alone are insufficient:
the observable first takes $N\to\infty$ at positive temperature.
We therefore prove the counting bound throughout a thin energy sublevel
and only then concentrate the Gibbs measure on it.  The needed
replacement for the exact contact count $k\ge N$ is a lower bound
$k\ge(1-2\eta')N$ on that whole sublevel.  We derive it before counting
the positive gaps.

Fix the UNSAT density and a tolerance \(\rho>0\).  On the good operator-norm
event the surface-cap argument implies that, for all sufficiently
large \(\beta\), Gibbs mass outside
\[
 \mathcal L_\rho=\{y:H(y)\le H_{\min}+\rho N\}
\]
is exponentially small in \(N\).  Along any great circle the energy
derivative has Lipschitz constant \(CN\).  Descending in the negative
tangent-gradient direction therefore gives
\begin{equation}\label{zero:near-ground-gradient}
 \|\nabla_S H(y)\|\le C\sqrt\rho\,N
 \qquad(y\in\mathcal L_\rho).
\end{equation}
Indeed a derivative of magnitude \(d\) permits an energy decrease
at least \(c d^2/N\), contradicting near-minimality if
\(d>C\sqrt\rho N\).

Set $t=\sqrt e$ and require $\rho\le e/4$.
On the event $H_{\min}/N\in[3e/4,5e/4]$, every point of
$\mathcal L_\rho$ then satisfies
$t\sqrt N\le\|r(y)\|_2\le2t\sqrt N$.  Choose from the outset
\[
 0<\eta'\le\min\left\{\frac14,
                   \frac1{2(1+|\log t|)}\right\}.
\]
At a sufficiently small fixed cutoff \(v=v(e,\eta')>0\), and then sufficiently small
\(\rho>0\), no point of \(\mathcal L_\rho\) has \(\eta'N\) absolute
gaps at most \(v\).  To see this, use the net and conditional-gradient
argument with projection factor \((Cv)^{\lfloor\eta'N/2\rfloor}\)
and residual lower bound \(t\sqrt N\).  Its union probability is
at most \(a^{-1}C^N t^{-(N-1)}v^{\lfloor\eta'N/2\rfloor}\).
Choose \(v\) small, then \(a\ll\min\{t,v\sqrt{\eta'}\}\), then
\(\sqrt\rho\ll a\).  This is exponentially small by
\eqref{zero:near-ground-gradient}.

Every point of \(\mathcal L_\rho\) must consequently have
\(k\ge(1-2\eta')N\) active rows, after decreasing \(\rho\) if
necessary.  Otherwise impose zero image on all active and all
absolute-near-zero rows and intersect their kernel with \(y^\perp\).
This subspace has dimension at least \(\eta'N-1\).
Lemma~\ref{thresh:slab-direction} gives a unit
vector \(v_0\) there with \(\|Av_0\|_\infty\le C/\sqrt{\eta'}\).
For small fixed \(b>0\), depending on \(v\) and \(\eta'\), set
\(y_b=\sqrt{1-b^2}y+bv_0\).  With $c_b=\sqrt{1-b^2}$,
\[
 h_i(y_b)=c_b h_i(y)+(1-c_b)+bA_i v_0.
\]
Choose $b$ so small that $c_b\ge1/2$ and
$bC/\sqrt{\eta'}<v/4$.  Positive gaps above $v$ then remain
positive.  All rows with $|h_i(y)|\le v$ have $A_i v_0=0$, so
their nonnegative gaps remain nonnegative.  The same vanishing
projection on every active row gives
$r_i(y_b)\le c_b r_i(y)$.  Therefore
\(H(y_b)\le(1-b^2)H(y)\), an improvement of at least
\(b^2 eN/2\) with high probability.  This contradicts membership
in \(\mathcal L_\rho\) when \(\rho<b^2e/2\).

Now suppose that \(\eta N\) positive gaps are at most \(u\).
Transfer to a sufficiently fine net, retaining at least
\(j=\lfloor\eta N/2\rfloor\) of these rows with absolute gap at
most \(2u\); these row indices are disjoint from the active set of
the original point.  The active coordinates at the net point lie
in a Euclidean ball of radius \(Ct\sqrt N\), with the same
$t=\sqrt e$ fixed above.  Since \(k\ge N/2\), the joint
projection probability is at most \(C^N t^k u^j\).
The net estimate \eqref{zero:net-master} still applies, by choosing
\(\rho\) after the mesh.  Since $t=\sqrt e\le\sqrt{B_0}$,
$k\ge(1-2\eta')N$, and $2\eta'|\log t|\le1$,
\[
 t^{k-N}\le
 \max\left\{\max(1,\sqrt{B_0})^{B_0N},
               \exp\bigl(2\eta'N|\log t|\bigr)\right\}
 \le C^N,
\]
with a uniform $C$, even as $e\downarrow0$.
The union probability is consequently bounded by
\(a^{-1}C^N u^{\lfloor\eta N/2\rfloor}\), and tends to zero for
\(u\le\exp(-C'/\eta)\), where \(C'\) is uniform in the density.

The sequence of choices is: density and \(\eta\); then \(\eta'\),
\(v\), and the corresponding descent step; then the uniform
exponential cutoff \(u\); then the net mesh and \(\rho\); finally
\(\beta\) large enough for Gibbs concentration, and \(N\to\infty\).
Thus no interchange of the temperature and size limits is used.
The row count bound \(\eta N\), divided by \(M\), is at most a
uniform multiple of \(\eta\), since the densities stay bounded away
from zero.  Taking \(\eta=C/|\log u|\) proves
\eqref{zero:direct-positive-cutoff}, since the counting event included
all positive gaps in \((0,u]\).  At fixed density, the finite-temperature
row marginal and Proposition~\ref{zero:endpoint} converge along the
chosen subsequence at both endpoints 0 and \(u\), which are atomless.
Thus the bound passes directly to every subsequential law of
\(h^*\), without removing a positive lower cutoff or identifying
positive tests at exact ground minimizers.
\end{proof}

\subsection{A variational characterization of the threshold}

\begin{corollary}\label{zero:threshold-characterization}
With \(g\) and \(\bar\chi\) as in Proposition~\ref{zero:variational},
\begin{equation}\label{zero:threshold-ratio}
 \alpha_J=\inf_{\substack{\gamma\ge0\ \mathrm{nondecreasing},\ \gamma\in L^1\\
                         \Delta>0}}
 \frac{\displaystyle\int_0^1\bar\chi(t)^{-1}\,\dd t}
      {-2g(0,0)}.
\end{equation}
Equivalently one may impose \(\bar\chi(0)=1\) and replace the
terminal condition by \(g(1,z)=-V(z)/\Delta\).
\end{corollary}

\begin{proof}
The denominator is positive.  The diffusion representation writes
\(g(0,0)\) as a negative terminal expectation minus
\(\frac12\E\int\gamma g_z^2\); atomlessness with a positive density
on negative intervals makes the terminal expectation strictly
negative.  Recovery trials in Proposition~\ref{zero:variational}
remain valid when \(e=0\).  In that regime all trial functionals
are nonnegative, so \(\alpha\) is at most every ratio in
\eqref{zero:threshold-ratio}.  For \(e>0\), an attained negative
functional gives a ratio strictly below \(\alpha\).
Approaching the threshold from the two sides proves the formula.

To verify the normalization assertion, scale a fixed pair by
\(c>0\).  If \(g_c\) is the solution for \((c\gamma,c\Delta)\),
then \(c g_c\) solves the equation with coefficient \(\gamma\) and
terminal value \(-V/(\Delta+1/c)\).  The ratio for the scaled pair
is therefore the ratio for this rescaled equation.  It decreases
as \(c\uparrow\infty\) and converges to the ratio with terminal
\(-V/\Delta\), by the PDE comparison and stability already proved.
Every shape can be normalized to \(\bar\chi(0)=1\), and every
normalized limiting ratio is approached by its finite scalings.
The two infima are equal.
\end{proof}

The next section uses the susceptibility bounds of
Lemma~\ref{zero:spectral} and the two direct cutoff estimates to
construct the normalized critical law.  Joint convexity will first
supply the fixed-density uniqueness still missing for positive gaps.


\Needspace{10\baselineskip}
\section{A convex critical problem and the limiting microscopic laws}
\label{sec:critical}

At a fixed UNSAT density, Section~\ref{sec:zero-temperature} describes
subsequential zero-temperature laws by a susceptibility and its PDE.
We first prove that this description is unique. We then approach the
threshold, retaining the notation $\alpha_c=\alpha_J$, and construct a
unique normalized pair consisting of a susceptibility $\chi$ and a
nonnegative terminal force variable $R$. The limiting state $H_1$ will
be the gap variable: $H_1\ge0$, and $RH_1=0$ separates positive gaps
from contacts.

There are two different convergence questions. Convexity on one Wiener
space identifies every weak subsequential limit, but does not preserve
contact probabilities or force normalization. To obtain those facts we
prove that the critical susceptibility vanishes at the endpoint, upgrade
the terminal convergence to $L^2$, and apply the uniform small-gap and
small-force estimates from the preceding section. The resulting laws,
with the limits taken in the prescribed order, are stated in
Theorem~\ref{crit:ordered-laws}.

The stochastic variational viewpoint is related to
Bou\'e--Dupuis~\cite{BoueDupuis1998} and
Auffinger--Chen~\cite{AuffingerChen2015}.  More specifically, the use of
terminal duality to obtain a variational formula over Brownian martingales,
with the optimizer identified by the derivative of the Parisi PDE along
its diffusion, has a close analogue in Mourrat~\cite[Lemma~2.2 and
Theorem~2 in the preprint]{MourratUninverting2025}; see also Chen, Issa and
Mourrat~\cite[Proposition~8.1 and Theorem~1.5]{ChenIssaMourrat2025}.
Proposition~\ref{crit:duality} establishes the representation needed here
for a concave quadratic terminal function and a nonnegative terminal
variable in $L^2$.  Rewriting its quadratic terms using $s=1/\chi$ gives
the joint convexity in Lemma~\ref{crit:convexity}.  The possibility that
$\chi$ vanishes at the endpoint requires the additional compactness and
equality arguments below; the critical uniqueness statement is not a
direct consequence of the cited results.

\subsection{The probability space and the admissible class}

Fix canonical Wiener space, with its completed natural filtration
$(\mathcal F_t)_{0\le t\le1}$ and coordinate Brownian motion $B$.
Realize each fixed-density zero-temperature optimizer's forward diffusion
as the strong solution driven by this same $B$. Thus the terminal variables
for different densities will belong to one Hilbert space, rather than
merely having laws that can be compared. Expectations in this section
are on that space. For every
$R\in L^2(\mathcal F_1)$, martingale representation gives a unique predictable
$b\in L^2(\Omega\times[0,1])$ such that
\[
 r_t:=\E[R\mid\mathcal F_t]
      =c+\int_0^t b_v\,\dd B_v,\qquad c:=\E R.
\]
The maps $R\mapsto c$ and $R\mapsto b$ are bounded linear maps, and
$\E R^2=c^2+\int_0^1\E b_t^2\,\dd t$.

The susceptibility class extends the positive-energy class in
\eqref{eq:src15} to allow a zero endpoint value. An admissible
susceptibility is
\[
 \chi(t)=\Delta+\int_t^1\gamma(v)\,\dd v,
 \qquad \Delta\ge0,\qquad \chi(t)>0\quad(t<1),
\]
where $\gamma\ge0$ is nondecreasing and belongs to $L^1(0,1)$.
Thus $\chi$ is concave, nonincreasing, and continuous on $[0,1]$; the endpoint
value is $\chi(1)=\Delta$.  Set $s=1/\chi$ on $[0,1)$ and require
$\int_0^1s<\infty$.  The value of $s$ at the single endpoint is immaterial.
We explicitly require $s\in L^2(0,1)$ whenever that stronger condition
is needed.  A pair $(R,s)$ is admissible when its susceptibility is
admissible and $R\ge0$ almost surely belongs to $L^2(\mathcal F_1)$.

For $\lambda\ge0$ define the joint functional
\begin{equation}\label{crit:def-functional}
 \begin{split}
 J_{\alpha,\lambda}(R,s)
  ={}&\alpha\E\left[R(1+B_1)+\frac\lambda2R^2\right]
   +\frac\alpha2\left\{\frac{c^2}{s(0)}
         +\int_0^1\frac{\E b_t^2}{s(t)}\,\dd t\right\}
   +\frac12\int_0^1s(t)\,\dd t .
 \end{split}
\end{equation}
All terms are finite.  In particular $\chi\le\chi(0)$ bounds the quadratic
martingale term.  The elementary identity
\begin{equation}\label{crit:quadratic-identity}
 c^2\chi(0)+\int_0^1\chi(t)\E b_t^2\,\dd t
   =\Delta\E R^2+\int_0^1\gamma(t)\E r_t^2\,\dd t
\end{equation}
follows from $\E r_t^2=c^2+\int_0^t\E b_v^2\,\dd v$ and Tonelli's theorem.

\subsection{Duality and positive-energy uniqueness}

For fixed $\chi$, quadratic duality replaces the PDE value by a minimum
over terminal force variables. Passing from $\chi$ to its reciprocal
$s$ then makes the susceptibility and terminal variable jointly convex.
This change of variables is what allows two physical optimizers to be
compared on the Wiener space just fixed.

\begin{proposition}[Joint martingale formulation]\label{crit:duality}
For an original susceptibility $\chi$ with $\Delta>0$, the functional in
\eqref{eq:src15} equals the minimum over $R\ge0$ in $L^2(\mathcal F_1)$ of
\begin{equation}\label{eq:src18}
 J_{\alpha,1}(R,s)
 =\alpha\E\left[R(1+B_1)+\frac12R^2\right]
   +\frac\alpha2\left[\frac{c^2}{s(0)}
          +\int_0^1\frac{\E b_t^2}{s(t)}\,\dd t\right]
   +\frac12\int_0^1s(t)\,\dd t .
\end{equation}
This minimizer is unique for fixed $s$.
\end{proposition}

\begin{proof}
Use gap coordinates $h=1+z$ and write $f(t,h)=g(t,h-1)$ and $a=f_h$.
The zero-temperature equation and forward diffusion are
\[
 f_t+\tfrac12f_{hh}+\tfrac12\gamma f_h^2=0,
 \quad f(1,h)=-\frac{h_-^2}{2(1+\Delta)},
 \qquad H_t=1+B_t+\int_0^t\gamma(v)a(v,H_v)\,\dd v,
\]
where $h_-=\max\{-h,0\}$.  The coefficient is time integrable; $a$ has
linear growth and an integrable Lipschitz constant after multiplication by
$\gamma$.  The diffusion therefore has finite second moments.  It\^o's
formula, first on $[0,T]$ and then with $T\uparrow1$, gives
\[
 f(0,1)=\E\left[-\frac{(H_1)_-^2}{2(1+\Delta)}
                    -\frac12\int_0^1\gamma(t)a(t,H_t)^2\,\dd t\right].
\]
For every nonnegative $R$, the pointwise quadratic dual inequality is
\[
 -\frac{h_-^2}{2(1+\Delta)}
       \le Rh+\frac{1+\Delta}{2}R^2.
\]
Insert $h=H_1$, condition $R$ against $\mathcal F_t$ in the drift term, and
use $r_ta-a^2/2\le r_t^2/2$.  This proves
\[
 f(0,1)\le
 \E\left[R(1+B_1)+\frac{1+\Delta}{2}R^2\right]
                 +\frac12\int_0^1\gamma(t)\E r_t^2\,\dd t.
\]
For $R=a(1,H_1)=(H_1)_-/(1+\Delta)$ the differentiated equation makes
$a(t,H_t)$ a square-integrable martingale with terminal value $R$.
Consequently $r_t=a(t,H_t)$, and both inequalities are equalities.
Equation~\eqref{crit:quadratic-identity} proves the claimed formula.
Strict convexity of $\E R^2/2$ proves uniqueness for fixed $s$.
All uses of It\^o's formula can be justified with smooth terminal data and
bounded smooth coefficients; the curvature and growth bounds above dominate
the terms when those approximations are removed.
\end{proof}

\begin{lemma}[Convexity and its equality cases]\label{crit:convexity}
The set of admissible $s$ is convex, as is its normalized part $s(0)=1$.
The functional $J_{\alpha,\lambda}$ is jointly convex in $(R,s)$.
If two pairs attain equality in this convexity inequality at a nontrivial
convex combination, then
\begin{equation}\label{crit:perspective-equality}
 \frac{c_1}{s_1(0)}=\frac{c_2}{s_2(0)},\qquad
 \frac{b_{1,t}}{s_1(t)}=\frac{b_{2,t}}{s_2(t)}
       \quad\text{for }\dd t\otimes\Prob\text{-almost every }(t,\omega).
\end{equation}
If $\lambda>0$, equality also forces $R_1=R_2$ in $L^2$.
\end{lemma}

\begin{proof}
For $0<q<1$, the reciprocal of $qs_1+(1-q)s_2$ is the weighted harmonic
mean
$\mathsf H(x,y)=(q/x+(1-q)/y)^{-1}$ evaluated at $(\chi_1,\chi_2)$.
This function is increasing in each coordinate and concave on the positive
quadrant: differentiation gives a negative semidefinite Hessian.  Hence
$\mathsf H(\chi_1(t),\chi_2(t))$ is concave and nonincreasing in $t$.
It extends continuously to the endpoint, and its negative derivative is
nonnegative, nondecreasing, and integrable.  Integrability of the new $s$
is immediate from its definition as a convex combination.  The same
argument preserves $L^2$ integrability and the normalization.

For a real Hilbert space $\mathcal H$, the perspective
$(x,u)\mapsto\|x\|_{\mathcal H}^2/u$, $u>0$, is convex.  Its convexity
defect vanishes exactly when $x_1/u_1=x_2/u_2$.  Apply this once with
$\mathcal H=\R$ to $c$, and for almost every $t$ with
$\mathcal H=L^2(\Omega)$ to $b_t$.  Integration preserves nonnegative
defects, so equality in the sum implies the asserted almost-everywhere
equalities.  The remaining terms are linear, except for
$\alpha\lambda\E R^2/2$, whose strict convexity gives the last assertion.
\end{proof}

\begin{corollary}[Uniqueness in the UNSAT phase]\label{crit:unsat-unique}
At each fixed $\alpha>\alpha_c$ the minimizing physical susceptibility in
\eqref{eq:src15} is unique.  Thus the entire zero-temperature endpoint law
in \eqref{eq:src16}, including its restriction to positive gaps, is unique.
\end{corollary}

\begin{proof}
Two minimizing pairs have the same $R$ and therefore the same $c$ and $b$.
For a PDE optimizer,
$b_t=f_{hh}(t,H_t)<0$ almost surely for $t<1$.  To see the strict inequality,
the curvature is nonpositive, and its differentiated diffusion
representation propagates the strictly negative terminal curvature on
$(-\infty,0)$.  On every interval ending before time~$1$ the drift is
Lipschitz, its flow derivative is positive, and the diffusion has positive
probability of reaching that interval.  Equivalently the strong maximum
principle applied after smoothing gives strict negativity; the
representation retains strictness on removing the smoothing.
Equality in \eqref{crit:perspective-equality} therefore gives
$s_1=s_2$ almost everywhere, hence everywhere before the endpoint by
continuity.  Their continuous susceptibilities have the same endpoint
value as well.  The PDE and its forward diffusion are then identical.
The subsequential description in \eqref{eq:src16} consequently has only
one possible limit.
\end{proof}

\subsection{Normalization and compactness at the threshold}

The initial physical susceptibility $\bar\chi(0)$ diverges as the
density decreases to the threshold. Normalize that value before taking
the limit.
Use bars for the fixed-density quantities in \eqref{eq:src15}:
$\bar\gamma=-\bar\chi'$, $\bar\Delta=\bar\chi(1)$,
$\bar f(t,h)=g(t,h-1)$, and $\bar a=\bar f_h$.
The dual optimizer is $\bar R=\bar a(1,H_1)=(h^*)_-$ by
\eqref{eq:src16}. Put $L=\bar\chi(0)$ and define
\[
 \begin{aligned}
 \lambda&=L^{-1},&\chi&=L^{-1}\bar\chi,&s&=L/\bar\chi,\\
 \gamma&=L^{-1}\bar\gamma,&\Delta&=L^{-1}\bar\Delta,&
 R&=L\bar R,\\
 f&=L\bar f,&a&=L\bar a.&&
 \end{aligned}
\]
Hence $\chi(0)=s(0)=1$, the terminal condition is
$f(1,h)=-h_-^2/[2(\lambda+\Delta)]$, and the physical residual is
$\lambda R$. The state process $H$ is unchanged on each Brownian path:
its drift satisfies $\gamma a=\bar\gamma\bar a$.
Below $\chi,\gamma,\Delta,f,a,R$ denote these normalized quantities.

Multiplying the original functional by $L$ gives
$J_{\alpha,\lambda}$. Thus the normalized physical pair minimizes this
functional among normalized pairs at the fixed value $\lambda=L^{-1}$.
Write $\bar s=1/\bar\chi$. Moreover
\begin{equation}\label{eq:src19}
 J_{\alpha,\lambda}(R,s)=-\frac{e(\alpha)}\lambda,
 \qquad \int_0^1s(t)^2\,\dd t
            =\alpha\E R^2=\frac{2e(\alpha)}{\lambda^2}.
\end{equation}
Indeed, variation of the original $\bar\Delta>0$ in both directions, with
$\bar R$ fixed in \eqref{eq:src18}, yields
$\int\bar s^2=\alpha\E\bar R^2$.  The second-moment identity for the
physical residual in Section~\ref{sec:zero-temperature} says
$\alpha\E\bar R^2=2e(\alpha)$.  These identities give
\eqref{eq:src19} after normalization.
Here optimality is initially over $\Delta>0$.  It extends to the
admissible normalized class with $\Delta=0$ by replacing $\chi$ with
$(\chi+\eta)/(1+\eta)$ and keeping $R$ fixed: the quadratic terms converge
by domination and the entropy by monotone convergence followed by the
factor $1+\eta$.

The spectral and residual estimates in
Lemma~\ref{zero:spectral} and
Corollary~\ref{zero:normalized-force-cutoff} give, uniformly as
$\alpha\downarrow\alpha_c$:
\begin{equation}\label{crit:imported-bounds}
 \lambda\asymp\sqrt{e(\alpha)}\longrightarrow0,
 \qquad \int_0^1s^2+\E R^2\le C,
 \qquad 0<c_0\le\E R\le C_0.
\end{equation}
The contact-count limit and the two logarithmic cutoff bounds will enter
only when we identify the endpoint masses. For compactness we need only
\eqref{crit:imported-bounds}. In particular we do not assume a critical
power law or an endpoint boundary condition for $\chi$.

\begin{proposition}[Existence of a normalized critical pair]
\label{crit:compactness}
The functional $J_{\alpha_c,0}$ is nonnegative on the whole admissible
class, without a normalization constraint. Every sequence of normalized physical
pairs with $\alpha_n\downarrow\alpha_c$ has a subsequence such that
\[
 \chi_n\to\chi\text{ uniformly on every }[0,T],\ T<1,
 \qquad R_n\rightharpoonup R\text{ in }L^2(\mathcal F_1).
\]
Its limit is normalized, satisfies $s\in L^2(0,1)$ and $c\ge c_0$, and
attains $J_{\alpha_c,0}(R,s)=0$.
\end{proposition}

We will prove uniqueness for the limits supplied by this proposition.
That conclusion concerns physical limits; it will not require a
classification of all admissible zero-value pairs.

\begin{proof}
\emph{Nonnegativity of the limiting functional.}
First suppose $\Delta>0$.  For $A>0$, direct substitution gives the exact
homogeneity identity
\begin{equation}\label{crit:homogeneity}
 A J_{\alpha,1}(R/A,s/A)
      =J_{\alpha,0}(R,s)+\frac{\alpha}{2A}\E R^2.
\end{equation}
The reciprocal of $s/A$ is $A\chi$, so this is an admissible
zero-temperature trial.  At $\alpha_c$, zero ground energy and
Proposition~\ref{crit:duality} make its value nonnegative.  Letting
$A\to\infty$ proves the assertion.  If $\Delta=0$, replace $\chi$ by
$\chi+\eta$.  The quadratic term converges by dominated convergence,
while $1/(\chi+\eta)\uparrow1/\chi$; let $\eta\downarrow0$.

\emph{Compactness on the common space.}
Normalization and concavity give
$1\ge\chi_n(t)\ge1-t$.  On each $[0,T]$ their slopes are bounded by a
constant depending only on $T$: monotonicity of $\gamma_n$ gives
$\gamma_n(t)\le(1-t)^{-1}$.  Arzel\`a--Ascoli and a diagonal selection
give the stated local uniform convergence.  The limit is concave and
nonincreasing; its continuous endpoint value is
$\Delta=\lim_{t\uparrow1}\chi(t)$, which is not initially identified with
$\lim_n\chi_n(1)$.  Its derivative gives an admissible integrable
$\gamma=-\chi'$.  At differentiability points of $\chi$, the concavity
inequalities imply $\gamma_n\to\gamma$; the same local bound then gives
$L^1(0,T)$ convergence for every $T<1$.  Fatou's lemma gives $s\in L^2$.

The bound on $R_n$ gives a weakly convergent subsequence on the fixed
Wiener space.  The cone of nonnegative $L^2$ variables is weakly closed.
Martingale representation, being a bounded linear isometry after the
constant component is included, gives weak convergence of $b_n$ and
convergence of $c_n$.  In particular $c\ge c_0$.

\emph{Passage to the minimum.}
For each $T<1$, uniform convergence of the positive weights $\chi_n$ on
$[0,T]$ and weak convergence of $b_n$ imply
\[
 \int_0^T\chi(t)\E b_t^2\,\dd t
   \le\liminf_n\int_0^T\chi_n(t)\E b_{n,t}^2\,\dd t.
\]
Let $T\uparrow1$ for the corresponding global lower semicontinuity.
In contrast the entropy term actually converges: local convergence of
$s_n$ combines with
$\int_T^1s_n\le C\sqrt{1-T}$ from \eqref{crit:imported-bounds}.
The linear term converges weakly, and
$\lambda_n\E R_n^2\to0$.  Equations~\eqref{eq:src19} and
\eqref{crit:imported-bounds} therefore give
\[
 J_{\alpha_c,0}(R,s)
       \le\liminf_nJ_{\alpha_n,\lambda_n}(R_n,s_n)=0.
\]
Nonnegativity proves equality.
\end{proof}

\subsection{Interior PDE representation}

A weak terminal limit need not by itself retain the PDE representation.
We now recover that representation before comparing two critical limits.
The strict curvature and the time regularity proved here will make the
common-noise uniqueness argument possible.

For the normalized physical pairs use gap coordinates and denote the
scaled PDE by $f_n$.  Its equation and terminal data are
\[
 (f_n)_t+\tfrac12(f_n)_{hh}
               +\tfrac12\gamma_n(t)(f_n)_h^2=0,
 \qquad f_n(1,h)=-\frac{h_-^2}{2(\lambda_n+\chi_n(1))}.
\]
Write $a_n=(f_n)_h$.  The comparison estimates below control the terminal error uniformly as
the density changes. They will also specify the boundary condition of
the critical equation if $\chi(1)=0$:
\begin{equation}\label{eq:src20}
 \begin{gathered}
 -s_n(t)\le-\frac1{\lambda_n+\chi_n(t)}
                      \le (a_n)_h(t,h)\le0,
 \qquad a_n\ge0,\qquad (a_n)_{hh}\ge0,\\
 \left|f_n(t,h)+\frac{h_-^2}{2(\lambda_n+\chi_n(t))}\right|
                          \le\frac12\int_t^1s_n(v)\,\dd v,\\
 \left|a_n(t,h)-\frac{h_-}{\lambda_n+\chi_n(t)}\right|
                          \le C\sqrt{s_n(t)\int_t^1s_n(v)\,\dd v}.
 \end{gathered}
\end{equation}

\begin{lemma}[Interior limits and strict curvature]\label{crit:pde}
For each subsequence in Proposition~\ref{crit:compactness}, $f_n$ and $a_n$
converge locally uniformly on $[0,1)\times\R$ to functions $f$ and $a=f_h$
determined by $\chi$.  They satisfy \eqref{eq:src20} with
$\lambda_n+\chi_n$ replaced by $\chi$.  They are smooth in the spatial
variable in the interior and locally Lipschitz in time there, together
with each spatial derivative used below.  For $t<1$,
\[
 f_{hh}(t,h)<0,\qquad f_{hhh}(t,h)>0,
 \qquad \lim_{h\to+\infty}a(t,h)=0.
\]
On the common Wiener space the limiting martingale and diffusion obey
\begin{equation}\label{eq:src21}
 H_t=1+B_t+\int_0^t\gamma(v)r_v\,\dd v,
 \qquad r_t=a(t,H_t),\qquad b_t=f_{hh}(t,H_t),\qquad t<1.
\end{equation}
Both $H_t$ and $r_t$ extend continuously and in $L^2$ to $t=1$, with
$r_1=R$.
\end{lemma}

\begin{proof}
\emph{Comparison estimates.}
The curvature bounds are parabolic comparison bounds.  Specifically,
if $b=(f_n)_{hh}$, then
$b_t+\tfrac12b_{hh}+\gamma_n a_n b_h+\gamma_n b^2=0$.
The spatially constant function
$-(\lambda_n+\chi_n(t))^{-1}$ solves this equation, and zero is the upper
comparison solution.  Differentiating once more gives the linear
equation
$d_t+\tfrac12d_{hh}+\gamma_n a_n d_h+3\gamma_n b d=0$ for
$d=(f_n)_{hhh}$, which preserves nonnegativity backwards from the convex
terminal derivative.  These comparisons follow first for smooth
terminal data and smooth coefficients; the bounds survive approximation.
The first-order solution is
$q_n(t,h)=-h_-^2/(2(\lambda_n+\chi_n(t)))$.
Its missing diffusion term is bounded in absolute value by $s_n(t)/2$.
Comparison with $q_n\pm\tfrac12\int_t^1s_n$ proves the value bound.
The derivative bound follows by taking secants of length
$\sqrt{(\int_t^1s_n)/s_n(t)}$ and using the curvature bound; the
zero-length limiting case is immediate.

\emph{Convergence with a terminal cutoff.}
At a
fixed cutoff $T<1$ the value estimates compare $f_n(T,\cdot)$ and
$f_m(T,\cdot)$ with quadratic functions whose coefficients converge.
The error left after this quadratic comparison is bounded uniformly
in space by $C\sqrt{1-T}$, uniformly in $n,m$.  On $[0,T]$, the difference $w=f_n-f_m$ solves a linear equation
with drift $\gamma_n(a_n+a_m)/2$ and source
$(\gamma_n-\gamma_m)a_m^2/2$.  The drift has uniformly bounded
Lipschitz and linear-growth constants on that interval, and its diffusion
has uniformly bounded second moments on bounded starting sets.  Its
Feynman--Kac representation propagates a constant terminal error without
amplification.  The quadratic terminal error tends to zero at fixed
$T$, as does the source error because $\gamma_n\to\gamma$ in $L^1(0,T)$.
For each fixed $T$, first let $n,m\to\infty$: only the constant
terminal error remains, bounded on every earlier compact set by
$C\sqrt{1-T}$. Then let $T\uparrow1$. This proves the Cauchy property
and shows that the limit depends only on $\chi$.  Curvature bounds and secants give convergence of the first
derivatives.

\emph{Regularity for a measurable time coefficient.}
On a compact time interval ending strictly before $1$, $\gamma$ is
bounded. We need spatial smoothness and local time-Lipschitz bounds,
including at jumps of $\gamma$, for the inverse-curvature map used below.
We give the regularity argument in a form that does
not require continuity of this coefficient in time.  Reverse time and
multiply $a$ by smooth space and time cutoffs supported in a larger
interior rectangle, with the time cutoff zero at its initial edge.
The resulting inhomogeneous heat equation has bounded source: the
nonlinear term is $\gamma aa_h$, and the cutoff terms involve only
$a$ and $a_h$, already bounded on that rectangle.  For the heat
semigroup $P_u$, differentiation of the Gaussian kernel gives
\[
 \|\partial_hP_uF\|_\infty\le C u^{-1/2}\|F\|_\infty,
 \qquad
 \|\partial_{hh}P_uF\|_\infty
      \le C_\eta u^{-1+\eta/2}[F]_{C^\eta_h},
       \quad0<\eta<1.
\]
The second estimate uses the zero integral of the differentiated
kernel to subtract $F(h)$.  The first estimate and its spatial
increment version in Duhamel's formula give a uniform spatial
$C^{1,\eta}$ bound for every $\eta<1$: splitting the time integral
at the square of the spatial increment bounds the gradient increment
by a constant times $|z|(1+|\log|z||)$.  The source is therefore
spatially $C^\eta$.  The second displayed estimate is integrable at
$u=0$ and yields $C^{2,\eta'}$ on a smaller rectangle for
$0<\eta'<\eta$.  More generally, if $a$ is already bounded in
$C_h^{m,\eta}$, $m\ge1$, then the localized source is bounded in
$C_h^{m-1,\eta}$: its derivatives of order $m-1$ use derivatives
of $a$ of order at most $m$.  Apply the second heat-kernel estimate
to $\partial_h^{m-1}F$ to obtain $C_h^{m+1,\eta'}$ on a smaller
rectangle.  This induction uses no derivative of $a$ that it is
trying to establish.  It gives every spatial derivative on nested
interior rectangles.  Multiplication by the bounded function $\gamma(t)$
preserves all spatial H\"older bounds, so this bootstrap is valid
for measurable time coefficients.  Smooth coefficient approximation
justifies the intermediate differentiations with these same bounds.
Finally the PDE bounds the almost-everywhere time derivative of every
spatial derivative on smaller rectangles.  This gives the claimed
local time-Lipschitz property, including across jumps of $\gamma$.
The uniform spatial H\"older and time-Lipschitz bounds make these
derivatives jointly continuous.  At the start of this argument the
source uses the bounded weak derivative $a_h$ supplied by the
curvature estimate; the distributional PDE passes from $a_n$ by
local uniform convergence and $L^1$ convergence of $\gamma_n$.
Thus the initial Duhamel representation does not presume smoothness.

\emph{Strict curvature and identification of the martingale.}
The value bound on $h>0$ and $a\ge0$ imply $a(t,h)\to0$ as
$h\to+\infty$.  On the negative half-line, \eqref{eq:src20} makes $a$
unbounded.  Thus $a$ is not affine, its nonnegative second derivative is
nontrivial at every later interior time, and its nonpositive first
derivative is nontrivial there.  The strong maximum principle, on an
interval between two strictly interior times, gives
$f_{hhh}>0$ and $f_{hh}<0$ at the earlier time.

The coupled forward diffusions converge in $L^2$ uniformly on each
compact time interval, by the local gradient convergence,
$L^1$ convergence of $\gamma_n$, and the common linear-growth bounds.
For precision, those deterministic growth bounds and Gronwall's
inequality give uniformly bounded moments of every fixed finite
order for $\sup_{t\le T}|H_{n,t}|$ and
$\sup_{t\le T}|a_n(t,H_{n,t})|$, for each $T<1$.
The fourth-moment bound permits spatial localization to be removed
in $L^2$ after the coupled convergence is proved on bounded sets.
Consequently $a_n(t,H_{n,t})\to a(t,H_t)$ strongly in $L^2$ for every
$t<1$.  But $a_n(t,H_{n,t})=\E[R_n\mid\mathcal F_t]$ converges weakly
to $\E[R\mid\mathcal F_t]$.  This identifies $r_t$; It\^o's formula
identifies $b_t=f_{hh}(t,H_t)$ and proves \eqref{eq:src21}.
Finally Doob's inequality gives
$\E\sup_{t\le1}|r_t|^2\le4\E R^2$.
Since $\gamma\in L^1$, the drift integral in \eqref{eq:src21} converges
both almost surely and in $L^2$ at $1$.  The martingale itself has its
continuous $L^2$ extension with terminal value $R$.
\end{proof}

\subsection{Uniqueness at zero energy}

\begin{proposition}[Uniqueness of the normalized critical limit]
\label{crit:unique-critical}
Any two subsequential limits supplied by
Proposition~\ref{crit:compactness} coincide on the common Wiener space.
In particular the normalized susceptibility and the law of $(H_1,R)$
are independent of the density sequence.
\end{proposition}

\begin{proof}
\emph{Equality identifies the normalized curvatures along the states.}
Both pairs minimize the globally nonnegative critical functional.
Their midpoint is admissible and also has value zero.
Lemma~\ref{crit:convexity} and $s_i(0)=1$ yield $c_1=c_2=:c>0$ and
$\chi_1b_1=\chi_2b_2$ almost everywhere.  Put
$k_i(t,h)=-\chi_i(t)(f_i)_{hh}(t,h)$.  The proof uses this matched
normalized curvature to express one state process as a deterministic
transform of the other.  Their common Brownian noise forces the transform
to have spatial derivative one; the common martingale mean then identifies
the susceptibilities.  Interior continuity upgrades the curvature
identity to
\[
 k_1(t,H_{1,t})=k_2(t,H_{2,t}),\qquad 0<t<1,
\]
on a single event of probability one.  Each $k_i(t,\cdot)$ is strictly
decreasing with strictly negative derivative by Lemma~\ref{crit:pde}.
Each $H_{i,t}$ has full support in $\R$: on every interior compact
time interval its additive-noise diffusion has locally bounded
Lipschitz drift, and localized Girsanov equivalence, or the Brownian
support property followed by continuity of the solution map, reaches
every open interval.  Thus the support of $k_i(t,H_{i,t})$ is the closure
of the range of $k_i(t,\cdot)$.  Equality of their laws identifies the
two ranges, which are open intervals.  We may therefore define
\[
 \Psi(t,h)=k_2(t,\cdot)^{-1}(k_1(t,h)),\qquad
 H_{2,t}=\Psi(t,H_{1,t}).
\]
These inverse functions have uniform local bounds.  At a fixed
$(t_0,h_0)$ put $y_0=\Psi(t_0,h_0)$.  Strict decrease gives
\[
 k_2(t_0,y_0-1)>k_1(t_0,h_0)>k_2(t_0,y_0+1).
\]
Joint continuity preserves these inequalities on a neighborhood of
$(t_0,h_0)$, trapping $\Psi$ between $y_0-1$ and $y_0+1$.
Compactness and uniqueness of the inverse give continuity there.
The negative derivative of $k_2$ is bounded away from zero on a
slightly smaller compact rectangle.  The implicit derivative
$\Psi_h(t,h)=(k_1)_h(t,h)/(k_2)_h(t,\Psi(t,h))$
is therefore jointly continuous, all needed
spatial derivatives are bounded locally, and the mean-value theorem
with the time-Lipschitz bounds on $k_1,k_2$ proves a uniform local
time-Lipschitz bound on $\Psi$.  Finite covers give the same assertions
on every compact subset of $(0,1)\times\R$.

\emph{The common noise forces a spatial translation.}
Taking quadratic covariations with the driving Brownian motion on
an arbitrary interval $[u,t]\subset(0,1)$ gives
\[
 t-u=[H_2,B]_t-[H_2,B]_u
               =\int_u^t\Psi_h(v,H_{1,v})\,\dd v.
\]
This chain rule needs no continuous time derivative.  After
localization, the temporal increments in a partition contribute at
most $C\sum\Delta t\,|\Delta B|$, which tends to zero.  The spatial
Taylor remainders contribute at most
$C\max|\Delta B|\sum|\Delta H_1|^2$, which tends to zero in
probability, since quadratic-variation sums of a continuous
semimartingale are tight.  The first-order sums give the displayed
integral.  Rational interval endpoints and continuity give all these
identities on one event of probability one.
It follows that $\Psi_h(t,H_{1,t})=1$ for almost every time, almost surely.
Full support gives $\Psi_h(t,h)=1$ for every $h$ at almost every $t$;
continuity extends it to every interior $t$.  Hence
$\Psi(t,h)=h+d(t)$.

\emph{The martingale mean identifies the susceptibility.}
The equality of $k$'s now holds for all $h$. Integrate the matched
curvatures from $h$ to $+\infty$, using gradient decay there, to obtain
\[
 \chi_1(t)a_1(t,h)
          =\chi_2(t)a_2(t,h+d(t)).
\]
Along the states this gives $\chi_1r_{1,t}=\chi_2r_{2,t}$.
Taking expectations yields $c\chi_1=c\chi_2$, so $\chi_1=\chi_2$.
The PDE is determined by $\chi$, and the two forward diffusions have
the same initial condition and Brownian motion.  They agree, and their
martingales, including their $L^2$ terminal values, agree as well.
\end{proof}

\subsection{Variational conditions and the vanishing endpoint mass}

From now on, ``the critical pair'' means the unique limit of the
normalized physical optimizers just constructed. To reach the endpoint
we must show that $\Delta=\chi(1)$ vanishes. The convergence so far is
local in time: even a vanishing sequence of endpoint values
$\chi_n(1)$ would not exclude mass of $\gamma_n$ concentrating near~1.
We instead use optimality of the critical pair in the full admissible
class.

The first-order conditions are expressed by the stationarity defect
$J(t)$ and its integral $K(t)$, defined by
\[
 P(t)=\alpha_c\E r_t^2,\qquad
 I(t)=\int_0^t s(v)^2\,\dd v,\qquad
 J(t)=P(t)-I(t),\qquad K(t)=\int_t^1J(v)\,\dd v.
\]
The quantity $K(t)/2$ is the directional derivative for the admissible
positive variation $\gamma\mapsto\gamma+\eps\ind_{[t,1)}$, as proved below.
Here $J(t)$ is a function of time, distinct from the joint functional
$J_{\alpha,\lambda}$, and $P,I,J$ extend continuously to $1$.  In the interior,
$P'=\alpha_c\E b_t^2$ is continuous.  Choose the right-continuous
nondecreasing representative of $\gamma$.  Let $\nu=\dd\gamma$ denote the
nonnegative Stieltjes measure on $[0,1)$ with
$\nu([0,t])=\gamma(t)$; this convention includes a possible initial
value of $\gamma$ as an atom at zero.

\Needspace{4\baselineskip}
\begin{lemma}[Complete first-order conditions]\label{crit:kkt}
The critical pair satisfies
\begin{gather}
 H_1+\Delta R\ge0,\qquad R(H_1+\Delta R)=0
                         \quad\text{almost surely},
                  \label{crit:complementarity}\\
 K(t)\ge0\quad(0\le t\le1),\qquad
 \int_0^1\gamma(t)J(t)\,\dd t=\int_{[0,1)}K(t)\,\nu(\dd t)=0.
                  \label{crit:kkt-profile}
\end{gather}
The profile $\gamma$ vanishes on an initial interval.  At every interior
point of $\supp\nu$ one has $J=0$ and $J'\le0$; at every non-isolated
such point, $J'=0$.  Finally $J(1)=0$.
\end{lemma}

\begin{proof}
The critical pair minimizes without normalization, since its value is
the global lower bound zero.  By \eqref{crit:quadratic-identity}, the
Fr\'echet derivative in $R$ is the $L^2$ variable
$\alpha_c(H_1+\Delta R)$.  Explicitly the quadratic operator is
\[
 \mathcal A Q=\Delta Q+
          \int_0^1\gamma(t)\E[Q\mid\mathcal F_t]\,\dd t,
 \qquad \|\mathcal A\|_{L^2\to L^2}\le\chi(0).
\]
Conditional expectations are self-adjoint contractions, so
$\mathcal A$ is bounded and self-adjoint; the functional's gradient is
$\alpha_c(1+B_1+\mathcal AR)$, as asserted.
In particular $\|H_1\|_2\le1+\|B_1\|_2+\|\gamma\|_1\|R\|_2$
also when $\Delta=0$.  Minimization on the nonnegative cone means
that this variable is nonnegative and has zero pairing with $R$;
these are exactly \eqref{crit:complementarity}.

Adding $\eps\ind_{[t,1)}$ to $\gamma$ is an admissible positive
variation.  Its susceptibility increment at $v$ is
$\eps(1-\max\{t,v\})$, bounded by $\eps\chi(v)/\chi(0)$ by
concavity.  Differentiating the entropy term is therefore justified
by domination: if $q(v)=1-\max\{t,v\}$, its difference quotient
is bounded by $s(v)/\chi(0)$ for the positive variations used here.
For scaling $\gamma$, the susceptibility increment is
$q=\chi-\Delta\le\chi$, and the difference quotient is bounded
by $2s$ for $|\eps|<1/2$.  Both bounds require only $s\in L^1$.
Tonelli's theorem in the quadratic term gives
the directional derivative $K(t)/2$.  This proves $K\ge0$.
Scaling $\gamma$ by $1+\eps$ permits either sign for small $\eps$,
and gives $\int\gamma J=0$.  Fubini gives
$\int\gamma J=\int K\,\dd\nu$: its absolute version is bounded by
$\|J\|_\infty\int\gamma<\infty$.  Equivalently the Stieltjes
boundary term at $1$ vanishes because
$K(t)=O(1-t)$ and $(1-t)\gamma(t)\to0$.
Thus $K$ vanishes on $\supp\nu$.

Since $J(0)=\alpha_c c^2>0$ and $K'=-J$, nonnegativity of $K$ forces
$K(0)>0$.  Hence $\nu$ vanishes on an initial interval, including
its possible atom at zero.  An interior support point is a local
minimum of the nonnegative $C^2$ function $K$, so $J=0$ and $J'\le0$.
At a non-isolated support point, differentiability of $J$ and nearby
zeros give $J'=0$.

If $\Delta>0$, two-sided variation of $\Delta$ gives
$0=(P(1)-I(1))/2$.  If $\Delta=0$ and the support did not accumulate
at $1$, then $\gamma$ would be constant on a final interval and
$\chi(t)$ would be a positive constant times $1-t$.  This contradicts
$s\in L^2$.  Thus in this case interior support points approach $1$,
and continuity of $J$ again gives $J(1)=0$.
\end{proof}

\begin{proposition}[No critical endpoint susceptibility]
\label{crit:no-endpoint}
The critical susceptibility satisfies $\Delta=0$.  Its increase measure
$\nu$ accumulates at $1$, and with $k=-\chi f_{hh}$ one has, at every
interior support point,
\begin{equation}\label{eq:src22}
 \alpha_c\E k(t,H_t)=1+\alpha_c c\chi(t),
 \qquad \alpha_c\E k(t,H_t)^2\le1.
\end{equation}
The second inequality is an equality at non-isolated interior support
points.
\end{proposition}

\begin{proof}
Introduce $W(t)=s(t)+\alpha_c\E b_t$.  The curvature equation along
the forward diffusion gives
$\frac{\dd}{\dd t}\E b_t=-\gamma(t)\E b_t^2$ almost everywhere.
Since $s'=\gamma s^2$, we have $W'=-\gamma J'$.
The Stieltjes product rule and $J\,\dd\gamma=0$ now give
\begin{equation}\label{crit:W-identity}
 W(t)=W(0)-\gamma(t)J(t)\quad\text{almost everywhere},
 \qquad \int_0^1W(t)\,\dd t=W(0).
\end{equation}
The initial constant uses the initial zero interval of $\gamma$.

It\^o's formula for $H_tr_t$ and
\eqref{crit:complementarity} yield
\[
 -\Delta P(1)=\alpha_c c
                  +\int_0^1\bigl(\gamma(t)P(t)+\alpha_c\E b_t\bigr)\,\dd t.
\]
Every term is integrable: $P$ is bounded, $\gamma\in L^1$, and
$|\E b_t|\le s(t)\in L^1$.  The stochastic products are justified by
localization and the $L^2$ maximal bounds for $H$ and $r$.
Moreover \eqref{crit:kkt-profile} and integration by parts give
\[
 \int_0^1\gamma P=\int_0^1\gamma I
                   =\int_0^1s-\Delta I(1).
\]
Since $P(1)=I(1)$, these identities show
\begin{equation}\label{crit:W0}
 W(0)=-\alpha_c c<0.
\end{equation}

Suppose for contradiction that $\Delta>0$.  Then $s\le1/\Delta$,
and the value and derivative bounds in \eqref{eq:src20} give
$a(t,h)\to h_-/\Delta$ as $t\uparrow1$, uniformly on bounded sets.
Convex secant bounds give
$f_{hh}(t,h)\to-\ind_{\{h<0\}}/\Delta$ locally away from zero.
The endpoint $H_1$ has no atom at zero.  Here is a direct justification:
condition on the Brownian bridge and write $B_t=tZ+\mathcal B_t$, where
$Z$ has a Gaussian density independent of the bridge.  The map
$Z\mapsto H_1$ is increasing, and its derivative is
\[
 \int_0^1\exp\left\{\int_u^1\gamma(v)a_h(v,H_v)\,\dd v\right\}\dd u
      \ge\exp\left\{-\Delta^{-1}\int_0^1\gamma\right\}>0.
\]
The estimate follows first by smoothing and then by the Lipschitz
flow comparison.  Thus the conditional endpoint law, and hence the
unconditional one, is atom-free.  With $p=\Prob(H_1<0)$, bounded
convergence gives
\[
 W(1-)=\frac{1-\alpha_c p}{\Delta},\qquad
 J'(1-)=\frac{\alpha_c p-1}{\Delta^2}.
\]
Because $J(1)=0$ and $J'$ is bounded when $\Delta>0$,
$|J(t)|\le C(1-t)$.  Integrability and monotonicity of $\gamma$ imply
$\gamma(t)J(t)\to0$.  Hence \eqref{crit:W-identity}--\eqref{crit:W0}
give $W(1-)=-\alpha_c c<0$, so $\alpha_cp>1$ and $J'(1-)>0$.
It follows that $J(t)<0$, and therefore $K(t)<0$, for $t$ sufficiently
close to $1$.  This contradicts Lemma~\ref{crit:kkt}.

Thus $\Delta=0$, and the support accumulation follows from that lemma.
At a support point $J=0$, \eqref{crit:W-identity} and
\eqref{crit:W0} give the first identity in \eqref{eq:src22}.
Although \eqref{crit:W-identity} was initially obtained almost
everywhere, it applies here by continuity of $W$: approach the support
point through times where that identity holds, and use local
boundedness of $\gamma$ and $J(t)\to0$.
Multiplying $J'=\alpha_c\E b_t^2-s^2\le0$ by $\chi^2$ gives the
second.  The equality statement follows from $J'=0$ at non-isolated
support points.
\end{proof}

\subsection{Strong convergence through the endpoint}

We have now identified the only possible weak limit and proved
$\chi(1)=0$. The remaining analytic task is to strengthen convergence
of the terminal force. A trial obtained by rescaling the critical
minimizer bounds the physical normalization in terms of
$\alpha-\alpha_c$; comparison at that normalization then yields
convergence of the full $L^2$ norm. Positivity probabilities will be
treated separately using the cutoff bounds.

\Needspace{5\baselineskip}
\begin{proposition}[Strong critical convergence]\label{crit:strong}
As $\alpha_n\downarrow\alpha_c$, the normalized physical variables
satisfy
\[
 R_n\to R\text{ in }L^2,\qquad
 \gamma_n\to\gamma\text{ in }L^1(0,1),\qquad
 H_{n,1}\to H_1\text{ in }L^2.
\]
The physical endpoint gap is
$h_n^*=H_{n,1}+\chi_n(1)R_n$, its residual is
$(h_n^*)_-=\lambda_nR_n$, and $h_n^*\to H_1\ge0$ in $L^2$.
\end{proposition}

\begin{proof}
Use stars temporarily for the unique critical pair and put
$S(s)=\tfrac12\int_0^1s$, $d_n=\alpha_n-\alpha_c$.
The critical pair is a legitimate limiting normalized trial at any
$\lambda>0$: replace $\chi_*$ by
$(\chi_*+\eta)/(1+\eta)$, keep $R_*$ fixed, and send $\eta\downarrow0$.
Its entropy and quadratic terms converge by domination and monotone
convergence.  Since its critical value is zero, its value at
$(\alpha_n,\lambda)$ is
\[
 -\frac{d_n}{\alpha_c}S_*+
                 \frac{\alpha_n\lambda}{2}\E R_*^2.
\]
The original unnormalized trial has an additional factor $\lambda$.
Choosing
$\lambda=d_nS_*/(\alpha_c\alpha_n\E R_*^2)$ gives
\[
 e(\alpha_n)\ge
 \frac{d_n^2S_*^2}{2\alpha_c^2\alpha_n\E R_*^2}.
\]
Here $S_*>0$ and $\E R_*^2>0$, so
\eqref{crit:imported-bounds} implies $\lambda_n\ge c_1d_n$.

Compare now at the physical value $\lambda_n$.  Optimality and the
same trial give
\begin{equation}\label{crit:strong-inequality}
 \left(1+\frac{d_n}{\alpha_c}\right)
       J_{\alpha_c,0}(R_n,s_n)
 +\frac{\alpha_n\lambda_n}{2}
                   (\E R_n^2-\E R_*^2)
 \le\frac{d_n}{\alpha_c}(S_n-S_*).
\end{equation}
The first term is nonnegative, $S_n\to S_*$ by uniform integrability,
and $d_n/\lambda_n$ is bounded.  Therefore
$\limsup\E R_n^2\le\E R_*^2$.  Weak convergence and lower
semicontinuity give the reverse inequality.  The Hilbert-space identity
for $\|R_n-R_*\|_2^2$ proves strong convergence.
Uniqueness of all weak subsequential limits makes this conclusion
valid for the original sequence.

Local $L^1$ convergence of $\gamma_n$ becomes global because
$\int_T^1\gamma_n\le\chi_n(T)\to\chi(T)$, and
$\chi(T)\downarrow0$ by Proposition~\ref{crit:no-endpoint}.
The same estimate gives $\chi_n(1)\to0$.
Doob's inequality applied to $\E[R_n-R\mid\mathcal F_t]$ gives
$L^2$ convergence uniformly in time of the martingales.  In
\eqref{eq:src21}, split the drift difference into
$\gamma_n(r_n-r)+(\gamma_n-\gamma)r$.  Conditional-expectation
contraction and Minkowski's inequality give the explicit bound
\[
 \left\|\int_0^1(\gamma_nr_n-\gamma r)\,\dd t\right\|_2
 \le\|\gamma_n\|_1\|R_n-R\|_2
                      +\|\gamma_n-\gamma\|_1\|R\|_2\longrightarrow0.
\]
Hence $H_{n,1}\to H_1$ in $L^2$.

The terminal dual optimizer satisfies
$R_n=(H_{n,1})_-/(\lambda_n+\chi_n(1))$.  On its positive half-line
the physical gap equals $H_{n,1}$; on its negative half-line the
proximal formula \eqref{eq:src16} gives $h_n^*=-\lambda_nR_n$.
These two cases are exactly
$h_n^*=H_{n,1}+\chi_n(1)R_n$.  The asserted limits follow.
Finally \eqref{crit:complementarity} with $\Delta=0$ gives
$H_1\ge0$ and $RH_1=0$.
\end{proof}

\subsection{Contacts, atoms, and the prescribed order of limits}

Strong $L^2$ convergence preserves the first two moments, but a positive variable may still
converge to zero. To identify contacts we therefore use both uniform
cutoff estimates from Section~\ref{sec:zero-temperature}: the force
estimate preserves $\Prob(R>0)$, and the positive-gap estimate identifies
$\Prob(H_1=0)$. Only after these masses are known can we normalize the
limiting force law.

\begin{lemma}[Contact and noncontact mass]\label{crit:contacts}
At the critical pair,
\[
 \Prob(R>0)=\Prob(H_1=0)=\frac1{\alpha_c},
 \qquad \{R>0\}=\{H_1=0\}\quad\text{modulo null sets}.
\]
\end{lemma}

\begin{proof}
Let $p_n=\Prob(R_n>0)$. Corollary~\ref{zero:normalized-force-cutoff}
gives $\alpha_np_n\to1$ and the direct cutoff bound
\[
 \lim_{u\downarrow0}\sup_n\Prob(0<R_n\le u)=0
\]
after discarding a finite initial part of the density sequence.
Apply strong convergence at positive continuity cutoffs and then
send the cutoff to zero.  This gives
$\Prob(R>0)=\lim_np_n=1/\alpha_c$.

Similarly the direct bound in \eqref{zero:direct-positive-cutoff},
passed to the atomless endpoint law at each fixed UNSAT density, gives
$\lim_{u\downarrow0}\sup_n\Prob(0<h_n^*\le u)=0$.
Its constants and its upper cutoff are uniform in the density.
This is a bound on the whole interval $(0,u]$ at fixed density;
no interchange between a lower gap cutoff and the density limit is used.
At every positive continuity point $u$ of the law of $H_1$,
the $L^2$ convergence $h_n^*\to H_1$ gives convergence of the
probabilities of $(-\infty,u]$.  Subtracting the negative mass
$p_n$ and then sending $u\downarrow0$ proves
$\Prob(H_1=0)=1/\alpha_c$; fixed-density endpoint laws have no atom
at zero, as established in \eqref{eq:src16}.
Complementarity gives the inclusion $\{R>0\}\subset\{H_1=0\}$.
Equality of their probabilities proves the final assertion.
\end{proof}

The contact atom at zero is now identified. To pass the microscopic
observables at \emph{every} positive cutoff, we still need to exclude
atoms away from zero. The two variables require different estimates:
nondegeneracy in deterministic variance time for $R$, and a bounded
state drift away from zero for $H_1$.

\begin{lemma}[No positive atoms]\label{crit:no-positive-atoms}
Neither $R$ nor $H_1$ has an atom at any strictly positive point.
\end{lemma}

\begin{proof}
\emph{The force variable.}
Monotonicity of $s$ gives
\[
 s(t)\int_t^1s(v)\,\dd v\le\int_t^1s(v)^2\,\dd v\longrightarrow0.
\]
Thus \eqref{eq:src20} approximates $a(t,h)$ by $s(t)h_-$ with
an error $\eps_t\to0$ uniformly in $h$.
Fix $x>0$.  When $x/2\le a(t,h)\le2x$ and $t$ is sufficiently
close to $1$, one has $h<0$ and
$-h\ge(x/2-\eps_t)/s(t)$.  A forward secant of length
$x/(8s(t))$ stays negative.  Convexity of $a$ gives
\[
 a_h(t,h)\le-s(t)+16\eps_t s(t)/x\le-s(t)/2.
\]
Consequently $1/2\le k(t,h)\le1$ in that region, whereas
$0\le k\le1$ globally.

Change deterministic time to $v=I(t)=\int_0^ts^2$ and put
$T=I(1)<\infty$.  The martingale $r$ then has Brownian coefficient
$-k$: indeed, with $t(v)=I^{-1}(v)$,
\[
 \widetilde B_v=\int_0^{t(v)}s(q)\,\dd B_q,
 \qquad \dd r_{t(v)}
       =\frac{b_{t(v)}}{s(t(v))}\,\dd\widetilde B_v
       =-k(t(v),H_{t(v)})\,\dd\widetilde B_v.
\]
The process $\widetilde B$ is Brownian in the filtration
$\mathcal F_{t(v)}$, since its integrand and time change are
deterministic and its quadratic variation is $v$.
The coefficient has absolute value at most one and at least $1/2$ whenever
$|r-x|\le x/2$ in a sufficiently short final time interval.
For $\eta>0$ write $\ell=T-v+\eta^2$ and consider
\[
 \Phi(v,r)=\ell^{-\zeta/2}
             \exp\{-a_0(r-x)^2/\ell\},
 \qquad a_0=\tfrac14,\quad\zeta=\tfrac1{16}.
\]
Its drift divided by $\Phi$ is
\[
 \frac{\zeta/2-a_0k^2}{\ell}
      +\frac{(-a_0+2a_0^2k^2)(r-x)^2}{\ell^2}.
\]
On $|r-x|\le x/2$ both terms are nonpositive.  Off this interval
the second is at most $-x^2/(32\ell^2)$ and the first at most
$1/(32\ell)$, so their sum is nonpositive when $\ell\le x^2$.
Choose $\delta>0$ so that $v_0=T-\delta$ lies in that final region
and $2\delta\le x^2$.  For $\eta^2\le\delta$ one has
$\ell\le2\delta$.  At this fixed $\eta$, $\Phi$ and its stochastic
derivative are bounded, so It\^o's formula gives
\[
 \E\Phi(T,R)\le\E\Phi(v_0,r_{t(v_0)})
                       \le\delta^{-\zeta/2},
\]
uniformly in $\eta$.
Since $\Phi(T,R)\ge e^{-a_0}\eta^{-\zeta}$ on
$|R-x|\le\eta$,
\[
 \Prob(|R-x|\le\eta)\le C_x\eta^{\zeta}.
\]
This proves the first assertion.

\emph{The gap variable.}
We show that the state drift is bounded on spatial intervals away from
the contact point.  Positivity of the drift and monotonicity of $a$ imply, for
$t<t'<1$,
\[
 a(t,h)=\E a(t',H_{t'}^{t,h})
              \le\E a(t',h+B_{t'-t}).
\]
Fix $x>0$, write $\tau=1-t$, and set
$t_j=1-2^{-j}\tau$ and
$x_j=x/2+(x/2)2^{-j/4}$.  Split the last expectation according to
whether the Brownian increment is below
$-(x_j-x_{j+1})$.  The complement costs at most $a(t_{j+1},x_{j+1})$.
On the exceptional event, the global linear-growth bound in
\eqref{eq:src20} and the Gaussian tail estimate bound the cost by
\[
 C_x(1+s(t_{j+1}))
                  \exp\{-c_x2^{j/2}/\tau\}.
\]
Concavity and $\chi(1)=0$ give
$s(t_{j+1})\le2^{j+1}s(t)$.  The remaining term
$a(t_j,x_j)$ tends to zero by the uniform error $\eps_{t_j}$.
Summing the rapidly decreasing tail series proves
\[
 \sup_{h\ge x}a(t,h)
           \le C_xs(t)\exp\{-c_x/(1-t)\}
\]
for $t$ close to $1$.  Since
$\chi(t)=\int_t^1\gamma(v)\,\dd v\ge(1-t)\gamma(t)$,
the drift $\gamma(t)a(t,h)$ is bounded uniformly near time~$1$
on every compact interval contained in $(0,\infty)$; explicitly it
is at most $C_x(1-t)^{-1}\exp\{-c_x/(1-t)\}$ for $h\ge x$.

Start at a fixed late time, stop on leaving such an interval, and
set the drift to zero after exit.  The resulting diffusion has
bounded adapted drift up to time~$1$.  The exponential martingale
in Girsanov's theorem is integrable because that drift is bounded;
its law is equivalent to Brownian motion with its initial value.
Its terminal law therefore has no atoms.  The original diffusion
agrees with it on the event of no exit.  Continuity of $H$ implies
that every event $\{H_1=x\}$ with $x>0$ is covered by a countable
union of these no-exit events, using rational interval endpoints
and rational starting times.  Thus $\Prob(H_1=x)=0$.
\end{proof}

\begin{theorem}[The ordered limiting observables]\label{crit:ordered-laws}
The Gaussian observables in the problem have all their stated ordered
limits, at every $u>0$ and $s>0$, and are given by
\begin{equation}\label{crit:limiting-laws}
 G_J(u)=\Prob(0<H_1\le u),\qquad
 F_J(s)=\Prob\left(\frac{R}{\alpha_c c}\le s\,\middle|\,R>0\right).
\end{equation}
The force law has mean one and no atom at zero.  These limiting laws
are also the laws for all Gaussian scale mixtures in the
fixed-temperature comparison range.
More explicitly, for every $0<\delta<u$ and $s>0$, the individual
limits at each level of
\begin{gather*}
 \lim_{\alpha\downarrow\alpha_c}\lim_{\beta\to\infty}
       \lim_{N\to\infty}G_{N,\beta,\alpha}(\delta,u)
                  =\Prob(\delta<H_1\le u),\\
 \lim_{\alpha\downarrow\alpha_c}\lim_{\beta\to\infty}
       \lim_{N\to\infty}F_{N,\beta,\alpha}(s)
                  =F_J(s)
\end{gather*}
exist; the additional limit $\delta\downarrow0$ in the first line
gives $G_J(u)$.
\end{theorem}

\begin{proof}
\emph{Thermodynamic and zero-temperature limits at fixed density.}
First take $N\to\infty$ at fixed $(\beta,\alpha)$, using the inner
empirical-marginal result proved earlier.  At fixed UNSAT density,
Corollary~\ref{crit:unsat-unique} and \eqref{eq:src16} give the
zero-temperature endpoint law, with no atoms.  For completeness,
let $h_{\beta,\alpha}$ denote the inner limiting gap variable and put
$p_{\beta,\alpha}=\Prob(h_{\beta,\alpha}<0)$ and
$\mu_{\beta,\alpha}=\E(h_{\beta,\alpha})_-$.
Proposition~\ref{zero:endpoint} supplies uniform integrability of
the negative first moment, and the absence of an atom at zero gives
\[
 p_{\beta,\alpha}\longrightarrow p_\alpha>0,\qquad
 \mu_{\beta,\alpha}\longrightarrow\mu_\alpha>0.
\]
Positivity follows from
$\alpha\E(h_\alpha^*)_-^2=2e(\alpha)>0$.
The inner limiting force CDF at $s>0$ is
\[
 \frac1{p_{\beta,\alpha}}
 \Prob\left(-s\frac{\mu_{\beta,\alpha}}{p_{\beta,\alpha}}
                          \le h_{\beta,\alpha}<0\right).
\]
Both limiting interval endpoints are atomless under the law of
$h_\alpha^*$, so this expression converges at every $s>0$.
Thus the full $\beta\to\infty$ limit exists for all the required
gap and force arguments, and its conditioning mass is preserved
before the density is changed.  Uniqueness of every subsequential
endpoint law yields convergence along the full temperature limit.

\emph{The critical-density limit and the final gap cutoff.}
Now approach $\alpha_c$ from above. Proposition~\ref{crit:strong}
gives convergence of the physical endpoint gaps and the scaled
residuals, and Lemma~\ref{crit:no-positive-atoms} permits every positive
cutoff, rather than just a selected continuity sequence.

For gaps keep $0<\delta<u$ fixed throughout these three limits;
the answer is $\Prob(\delta<H_1\le u)$.  Only then let
$\delta\downarrow0$, which removes the contact atom and gives the
first formula in \eqref{crit:limiting-laws}.
For forces the conditional mean at density $\alpha_n$ is $c_n/p_n$,
which converges to $\alpha_cc$.  Lemma~\ref{crit:contacts} preserves
the conditioning mass, and its cutoff argument loses no positive
residual mass to zero.  Subtracting the noncontact mass and applying
the positive-atom assertion gives the second formula at every $s>0$.
Since $\E R=c$ and $\Prob(R>0)=1/\alpha_c$, its mean is one.
Its conditional probability of $\{0\}$ is zero; the uniform cutoff
bound also explicitly prevents an atom from appearing there in
the limit of the physical conditional laws.

Finally the Gaussian and mixture inner-limit observables agree at
each fixed temperature and density in a neighborhood above the common
threshold, by Section~\ref{sec:universality}.  Equality of these
functions passes through exactly the same ordered limits.
The limits $u\downarrow0$ and $s\downarrow0$, which determine their
exponents, are taken only in the subsequent sections.
\end{proof}


\section{Uniform control of the critical layer}\label{sec:scales}

Throughout this section, $(\chi,R)$ is the normalized critical optimizer,
$\alpha=\alpha_c$, $c=\E R>0$, and $Y=H_1$. The inputs from
Section~\ref{sec:critical} are the PDE and diffusion representation of
Lemma~\ref{crit:pde}, the first-order conditions of Lemma~\ref{crit:kkt},
and the endpoint conclusions of Proposition~\ref{crit:no-endpoint} and
Lemma~\ref{crit:contacts}. In particular,
\[
 \begin{gathered}
 \chi(1)=0,\qquad \int_0^1\chi^{-2}<\infty,\\
 \{R>0\}=\{Y=0\}\quad\text{modulo null sets},\qquad
 \Prob(R>0)=1/\alpha.
 \end{gathered}
\]
The increase support of $\gamma$ accumulates at $1$, and $\gamma=0$
on an initial interval. At its interior support points we have
$J=0$, $J'\le0$, and \eqref{eq:src22}, where $J=P-I$ is the
variational defect defined in Section~\ref{sec:critical}.
All constants below may depend on this fixed optimizer and on a fixed
interior starting time; they do not depend on the distance to the
endpoint. Constants with a subscript may also depend on that parameter.

Write $\tau=1-t$. Whenever $\tau$ is the time argument, $\chi(\tau)$
means the original susceptibility at physical time $1-\tau$; the same
convention applies to the other deterministic time-dependent quantities. Thus
$\chi_\tau=\gamma\ge0$ and $\chi(\tau)\to0$ as $\tau\downarrow0$;
the susceptibility is concave in this clock. The physical-time
normalization remains $\chi(t=0)=1$.
In this clock write
$\mathcal S=\{\tau\in(0,1):1-\tau\in\supp\dd\gamma\}$.
It has points arbitrarily close to zero. A support gap means a component
$(\tau_1,\tau_2)$ of its complement, on which $\chi$ is affine.
When extending estimates across gaps, we work below a fixed point of
$\mathcal S$; every intervening gap then has both endpoints in
$\mathcal S$. Choosing this point sufficiently small puts both endpoints
in the range of any support-point estimate under consideration.
Set
\[
 \begin{aligned}
 s&=\chi^{-1},& m(t,h)&=\chi(t)a(t,h),\\
 k&=-m_h,& l&=1-k,\\
 X(t,h)&=h+m(t,h),& w(t)&=\int_t^1s(v)^2\,\dd v,\\
 D&=\chi\sqrt{w/\tau}.&&
 \end{aligned}
\]
The same symbol $X_t$ denotes $X(t,H_t)$. In particular,
\[
 \dd r_t=-s(t)k(t,H_t)\,\dd B_t,\qquad
 \dd X_t=l(t,H_t)\,\dd B_t,
 \qquad r_1=R,\quad X_1=Y.
\]
The curvature bounds give $0\le k,l\le1$ and $k_h<0$, hence
$0<k<1$ at every finite spatial point. Since $|l|\le1$, $X_t$ is
a square-integrable martingale. Its $L^2$ endpoint limit is $Y$,
because $H_t\to Y$, $\chi\to0$, and $r_t\to R$ in $L^2$.
Thus $X_t=\E[Y\mid\mathcal F_t]\ge0$, just as
$r_t=\E[R\mid\mathcal F_t]\ge0$. Finally
$w(\tau)\ge\tau/\chi(\tau)^2$ gives $D\ge1$.
Here $m=\chi a$ is the rescaled force field; the earlier
finite-temperature CDF is $m_q$.

The identity $H_t=X_t-\chi r_t$ explains the susceptibility factor
in $D$. We first prove that $D$ stays bounded, making the spatial
scales $\chi\sqrt w$ and $\sqrt\tau$ comparable, and derive the
endpoint mass estimates \eqref{eq:src31}. We then use the
inverse-slope bounds \eqref{eq:src32} to exclude terminal gaps in
the increase support and obtain \eqref{eq:src34}. These are the
inputs for Section~\ref{sec:contraction}.

The equation for $m$ in remaining time is
\begin{equation}\label{eq:scale-m-pde}
 m_\tau=\tfrac12m_{hh}+d\,m(1+m_h),
 \qquad d=\gamma/\chi,\qquad 0\le d\le\tau^{-1}.
\end{equation}
The last inequality follows from $\tau\gamma(\tau)\le\chi(\tau)$.

\subsection{Value barriers and inverse flows}

Let $\phi$ and $\Phi$ be the standard normal density and distribution
function. The following estimates do not require any information about
$D$ beyond $D\ge1$.
\begin{lemma}[Pointwise barriers]\label{lem:scale-barriers}
For $z=h/\sqrt\tau$,
\begin{equation}\label{eq:src23}
 \sqrt\tau\,[\phi(z)-z\Phi(-z)]\le m(t,h)
 \le\sqrt\tau\,\frac{\phi(z)}{\Phi(z)}.
\end{equation}
Consequently, with fixed positive constants $C,c_0$,
\begin{equation}\label{eq:src24}
 k(t,h)\le C e^{-c_0z^2}\quad(z\ge1),\qquad
 l(t,h)\le C|z|^{-2}\quad(z\le-1),\qquad
 \int_{\R}kl\,\dd h\le C\sqrt\tau.
\end{equation}
\end{lemma}
\begin{proof}
First $m\ge h_-$. Indeed $X_h=l\ge0$, and the bounded difference
between $f(t,h)$ and $-h_-^2/(2\chi)$ in \eqref{eq:src20}
forces $X(t,h)\to0$ as $h\to-\infty$: a nonzero limit of the
monotone function $X$ would make the integral of
$f_h+h/\chi=X/\chi$ unbounded on that half-line. Thus $X\ge0$;
on the other half-line use $m\ge0$.

The heat evolution of the hinge is the lower function in
\eqref{eq:src23}. It is a subsolution of \eqref{eq:scale-m-pde}, since
$m(1+m_h)\ge0$. For the upper function $U$, direct differentiation of
$\rho(z)=\phi(z)/\Phi(z)$, using
$\rho'=-z\rho-\rho^2$, gives
$U_\tau=U_{hh}/2+\tau^{-1}U(1+U_h)$ and $-1\le U_h\le0$.
It is therefore a supersolution for $d\le\tau^{-1}$.

Here comparison must start before the singular endpoint. At
$\tau'>0$, \eqref{eq:src20} gives
\[
 0\le m(\tau',h)-h_-\le C\chi(\tau')\sqrt{w(\tau')}
 \quad\hbox{uniformly in }h.
\]
For the difference between two solutions, the zeroth-order coefficient
in the linear comparison equation is at most $d$. Thus the positive
error at $\tau'$ is amplified by at most
$\exp(\int_{\tau'}^\tau d)=\chi(\tau)/\chi(\tau')$.
The resulting error is at most $C\chi(\tau)\sqrt{w(\tau')}$, which
tends to zero as $\tau'\downarrow0$. Comparison may first be made with
smooth coefficients and on bounded spatial intervals; the uniform
linear growth and curvature bounds permit the limits.

Finally, for a convex decreasing function $m$, its negative derivative
at $h$ is bounded above by the backward secant over a distance
$\sqrt\tau$. Applying the upper value barrier gives the positive
Gaussian tail for $k$. Apply the forward secant to the convex increasing
function $X$ on the negative half-line, using
$X(t,h)\le C\tau/|h|$ there; this gives the stated bound for $l$.
The integral estimate follows by using $kl\le1$ on $|h|\le\sqrt\tau$
and integrating the two tails.
\end{proof}

We next turn these pointwise bounds into estimates under the law of
$H_t$. The inverse flow retains the Gaussian density at an early time;
its Jacobian records how much the critical drift compresses that law.

Choose $t_0>0$ in the interval on which $\gamma=0$. Then
$\chi(t_0)=1$ and $H_{t_0}$ has a strictly positive Gaussian density
$p_0$, independently of subsequent Brownian increments. Conditional
on those increments, the flow $F_t$ taking $H_{t_0}$ to $H_t$ is
increasing and convex: its first derivative solves an equation with
coefficient $-dk\in[-d,0]$, and its second derivative has nonnegative
forcing $d m_{hh}$. If $\mathcal I_t=F_t^{-1}$ and $\mathcal J_t=\mathcal I_t'$, then
\begin{equation}\label{eq:scale-inverse-density}
 1\le \mathcal J_t\le\chi(t)^{-1},\qquad \mathcal J_t\text{ is nonincreasing},\qquad
 p_t(h)=\E[p_0(\mathcal I_t(h))\mathcal J_t(h)]\le C/\chi(t).
\end{equation}
In particular the critical state has a density at every interior time.

We shall also use the lower bound
\begin{equation}\label{eq:scale-density-lower}
 \inf_{|h|\le K\chi(t)}p_t(h)\ge c_K>0.
\end{equation}
Here is a uniform moment justification. The inverse flow, read in
increasing remaining time $u$, solves the additive-noise equation
$\dd\widetilde H_u=-d(u)m(u,\widetilde H_u)\,\dd u+\dd\widetilde B_u$.
Its positive part is bounded by its initial positive part and twice the
Brownian supremum. For $N=-\widetilde H/\chi$,
\[
 \dd N=-s\,\dd\widetilde B+\frac{\gamma X(u,\widetilde H)}{\chi^2}\,\dd u.
\]
The value barrier gives $(-h)_+X(u,h)\le Cu$. Applying It\^o's formula
to $N_+^2$ therefore bounds its expected growth by $Cs^2\,\dd u$,
since $u\gamma/\chi\le1$. Its initial second moment is at most $K^2$.
The inverse position at $t_0$ consequently has a uniformly bounded
second moment. On a fixed interval its Gaussian density is bounded
below, and Markov's inequality puts a fixed positive fraction of the
inverse positions in that interval. Equation
\eqref{eq:scale-inverse-density}, with $\mathcal J_t\ge1$, proves
\eqref{eq:scale-density-lower}. All these inverse statements can also
be obtained pathwise from the additive-noise integral equation, so no
independence between the inverse position and its Jacobian is used.

At an interior point of the increase support of $\gamma$,
\eqref{eq:src22} gives $\E kl\ge c\chi$. Combining this with
\eqref{eq:src24} and $p_t\le C/\chi$ gives
\begin{equation}\label{eq:scale-chi-quarter}
 \chi(\tau)\le C\tau^{1/4}.
\end{equation}
This holds throughout a terminal interval. Indeed on a complementary
support interval $\chi$ is affine, whereas $\tau^{1/4}$ is concave,
so the endpoint inequalities imply the inequality in the interval.

\subsection{The ratio of the two natural scales}

Let
$\pi(t,h)=\Prob(R>0\mid H_t=h)$ and $q_g=1-\pi$.
The Markov property supplies these conditional probabilities as
measurable functions of the state. The curvature equation gives
\[
 \dd k(t,H_t)=k_h(t,H_t)\,\dd B_t-d(t)k(t,H_t)l(t,H_t)\,\dd t.
\]
Its endpoint limit is $\ind_{\{R>0\}}$. On $Y>0$ this follows from
\eqref{eq:src24}; on $R>0$, the uniform derivative error in
\eqref{eq:src20} is $o(1)$, so
$H_t=-\chi(t)R+o(\chi(t))$. A secant of length a fixed fraction of
$\chi(t)R$ on the negative side then gives $k(t,H_t)\to1$.
The two events exhaust the space. Conditional bounded-supermartingale
convergence thus gives
\begin{equation}\label{eq:scale-probability-comparison}
 k\ge\pi,\qquad
 q_g(t,h)\ge\Phi\!\left(\frac{h}{\chi\sqrt w}\right)
 \quad\hbox{for almost every }h.
\end{equation}
For the second inequality, run the inverse flow from position zero at
$t'>t$. The variable $N$ in the preceding argument has nonnegative
drift, and hence dominates, under the same noise, a centered Gaussian
of variance $w(t)-w(t')$. The monotonicity of the flow shows
\[
 \Prob(H_{t'}>0\mid H_t=h)
 \ge\Phi\!\left(\frac{h}{\chi(t)\sqrt{w(t)-w(t')}}\right).
\]
Now let $t'\uparrow1$. On the force event the sign is eventually
negative by the just established secant argument; on the gap event it
is eventually positive. Conditional dominated convergence proves the
claim.

A second useful consequence of complementarity is, for
$v=r/\sqrt w\ge2$,
\begin{equation}\label{eq:src26}
 X\le C\sqrt\tau\,v\Phi(-v),\qquad
 l\le C\frac{v^2}{D}\Phi(-v).
\end{equation}
Indeed the event $Y>0$ requires the martingale $r$ to hit zero. Its
remaining quadratic variation is at most $w$, so the reflection
principle after Brownian time change bounds this event's conditional
probability by $q\le2\Phi(-v)$. The martingale $X$ has remaining
quadratic variation at most $\tau$. Integration of its Gaussian tail
gives, for any event $A$ of probability $q$,
\[
 \E[|Y-X|\ind_A\mid H_t=h]
 \le C\sqrt\tau\,q\sqrt{\log(2/q)}.
\]
As $Y$ vanishes off $A=\{Y>0\}$, this bounds
$X(1-q)$. The normal tail estimates give the first claim, also with
adjusted constants for all $v\ge1$. To bound $l=X_h$, take its forward
secant over $\delta h=\chi\sqrt w/v$. The variable $r/\sqrt w$ at the
new position is at least $v-v^{-1}$, because $|r_h|\le1/\chi$.
The ratio of the corresponding normal tails is bounded by an absolute
constant. Dividing the value bound by $\delta h$ proves the second
claim. Continuity extends estimates derived at almost every state.

\begin{lemma}[Subpower scale ratio]\label{lem:scale-subpower}
As $\tau\downarrow0$,
\begin{equation}\label{eq:src25}
 \log D=o(|\log\tau|).
\end{equation}
\end{lemma}
\begin{proof}
At a support point, stability says
$\E(k^2-\pi)\le0$. Suppose instead that along support points
$D\ge\tau^{-\eps}$ for some $\eps>0$.
Since $h=X-\chi r\ge-\chi r$, \eqref{eq:scale-probability-comparison}
gives $q_g\ge\Phi(-v)$. Thus \eqref{eq:src26} implies
$k^2-\pi\ge q_g/2$ when
$2\le v\le a\sqrt D$, for a fixed sufficiently small $a>0$.
When $v>a\sqrt D$ its negative part is bounded by
$kl\le C\exp(-c_0D)$.
When $v<2$ and $h\le-K\sqrt\tau$, take $K$ sufficiently large in
\eqref{eq:src24}; again $k^2-\pi\ge q_g/2$.
In particular, on
\[
 I_-=\left[-\tfrac12\chi\sqrt w,-\tfrac14\chi\sqrt w\right]
\]
there is a fixed positive lower bound for $k^2-\pi$ once $D$ is large.
This follows from $q_g\ge\Phi(-1/2)$, the negative-tail slope bound,
and the upper value barrier, which gives $v\le1$ on this interval.
Writing $p_-:=\inf_{I_-}p_t$, the positive contribution is at least
$c_0\chi\sqrt w\,p_-$. Equation
\eqref{eq:scale-density-lower} gives $p_-\ge c_0$ for late times.

For clarity, the required density comparison has the following
quantitative form. If $h_2>h_1$, $\delta=(h_2-h_1)/\chi$, and $P$ is
fixed, then
\begin{equation}\label{eq:scale-truncated-density}
 p_t(h_2)\le
 \exp\!\{C_P[\delta\sqrt{|\log\tau|}+\delta^2+1]\}\,p_t(h_1)
 +\tau^P.
\end{equation}
To prove it, use the decreasing inverse Jacobian, the bound
$|\mathcal I_t(h_2)-\mathcal I_t(h_1)|\le\delta$, and the Gaussian density ratio on
$|\mathcal I_t(h_2)|\le C_P\sqrt{|\log\tau|}$. On the complement use
$\mathcal J_t\le\chi^{-1}\le C\tau^{-1/2}$; the last inequality follows from
$\tau/\chi^2\le w\le w(t_0)$. Increasing $C_P$ makes the discarded
Gaussian density at most $\tau^P$.

Apply this with $h_1\in I_-$ and
$-K\sqrt\tau\le h_2\le C_P\sqrt{\tau|\log\tau|}$.
Here $\delta\le C_P\sqrt{w(1+|\log\tau|)}$, so the exponential factor
is $\tau^{-o(1)}$. The negative contribution in this layer is at most
$C\sqrt\tau\,\tau^{-o(1)}p_-$ by
$(k^2-\pi)_-\le kl$ and \eqref{eq:src24}.
Above the layer the positive Gaussian tail of $k$ makes the contribution
an arbitrarily small power of $\tau$. Below it the only remaining
negative contribution is $C\exp(-c_0D)$. Both are negligible compared
with $\chi\sqrt w\,p_-=D\sqrt\tau\,p_-$, a contradiction.

It remains to cross a support gap $\tau_1<\tau<\tau_2$.
Affineness gives the exact identity
\begin{equation}\label{eq:scale-affine-gap}
 w(\tau)=w_1+\frac{\tau-\tau_1}{\chi_1\chi(\tau)}.
\end{equation}
If $\tau_2/\tau_1\le2$, concavity gives $\chi/\chi_1\le2$,
so the assertion transfers immediately. Otherwise
$\chi_2/\chi_1\le2D_2^2$, by applying
\eqref{eq:scale-affine-gap} at $\tau_2$. Hence
\[
 D(\tau)^2\le
 D_1^2\frac{\tau_1}{\tau}(2D_2^2)^2+2D_2^2.
\]
For every fixed $\eta>0$ the endpoint estimates bound $D_i$ by
$C_\eta\tau_i^{-\eta}$. The factor $\tau_1/\tau$ absorbs the
possible disparity between $\tau_1$ and $\tau$; since
$\tau_2\ge\tau$, for $0<\eta<1/2$ we obtain explicitly
\[
 D(\tau)^2\le
 C_\eta\frac{\tau_1^{1-2\eta}}{\tau}\tau_2^{-4\eta}
       +C_\eta\tau_2^{-2\eta}
 \le C_\eta(\tau^{-6\eta}+\tau^{-2\eta}).
\]
As $\eta$ is arbitrary, this proves the same subpower conclusion in
the gap.
\end{proof}

\begin{lemma}[Strict improvement of the endpoint power]\label{lem:scale-strict-power}
There exist $\eta>0$ and $C<\infty$ such that
$s(\tau)\le C\tau^{-1/2+\eta}$. In particular $w(\tau)\le C\tau^{2\eta}$.
\end{lemma}
\begin{proof}
In the inverse formula,
$\log \mathcal J_t=\int d(u)k(u,\widetilde H_u)\,\dd u$.
Starting at $h\in[-a\sqrt\tau,0]$, the inverse drift is nonpositive,
so at every reverse time its probability of a negative position is at
least $1/2$. The barriers and secants give a uniform lower bound
$k\ge\kappa>0$ for negative positions. Moreover the inverse moment
bound proved above gives $\E\log p_0(\mathcal I_t(h))\ge-C$ uniformly: its
initial normalized negative position has square at most
$a^2\tau/\chi^2\le a^2w$.
Jensen's inequality therefore yields
$p_t(h)\ge C^{-1}\chi^{-\kappa/2}$.
Here both $k$ and $l$ have fixed positive lower bounds on
$[-a\sqrt\tau,0]$ if $a>0$ is sufficiently small. To see this
uniformly in time, the heat barrier gives
$m(0)\ge\phi(0)\sqrt\tau$, while the upper barrier makes
$m(K\sqrt\tau)/\sqrt\tau$ arbitrarily small for fixed large $K$.
The secant from $0$ to $K\sqrt\tau$ bounds $k(0)$ below, and
$k(h)\ge k(0)$ for $h\le0$. Likewise
$X(h)\ge X(0)-|h|\ge(\phi(0)-a)\sqrt\tau$ on that interval,
whereas $X(-K\sqrt\tau)\le C\sqrt\tau/K$.
The backward secant from $-K\sqrt\tau$ to $h$ therefore gives
\[
 l(h)\ge\frac{X(h)-X(-K\sqrt\tau)}{h+K\sqrt\tau}
       \ge\frac{\phi(0)-a-C/K}{K}>0
\]
after fixing $a<\phi(0)/2$ and then a sufficiently large $K$.
Together with \eqref{eq:scale-chi-quarter} this gives
\[
 \E kl\ge C^{-1}\sqrt\tau\,\chi^{-\kappa/2}
 \ge C^{-1}\tau^{b_0},\qquad b_0=\tfrac12-\tfrac\kappa8<\tfrac12.
\]
The drift identity for $k$, including its endpoint limit, gives
$\E k-1/\alpha=\int_t^1d\,\E kl$.
Thus, if $\tau_2$ is a late support point and $0<\tau_1<\tau_2$,
\begin{equation}\label{eq:scale-power-growth}
 \chi(\tau_2)\ge C^{-1}\tau_1^{b_0}
 \log\frac{\chi(\tau_2)}{\chi(\tau_1)}.
\end{equation}

Put $b=\limsup_{\tau\downarrow0}\log\chi(\tau)/\log\tau$.
The bounds already proved imply $1/4\le b\le1/2$.
Fix $0<\eps<1$ and choose $\tau_n\downarrow0$ realizing this limsup.
For sufficiently small $\delta>0$, eventually
$\chi(u)\ge u^{b+\delta}$ at all small $u$, whereas along a
subsequence $\chi(\tau_n)\le\tau_n^{b-\delta}$.
At $u_n=\tau_n^{1-\eps}$, choosing
$(2-\eps)\delta<b\eps/4$ gives
\[
 \frac{\chi(u_n)}{\chi(\tau_n)}
 \ge \tau_n^{-b\eps+(2-\eps)\delta}
 \ge \tau_n^{-3b\eps/4}.
\]
If $u_n$ lies in a gap, move to its smaller endpoint $\tau_{2,n}$;
otherwise put $\tau_{2,n}=u_n$. The loss in susceptibility is
subpower in $1/u_n$. Indeed it is at most two on a gap of time ratio
at most two; on any other gap it is at most $2D(\tau_2)^2$, by
\eqref{eq:scale-affine-gap}, and $\tau_2\ge u_n$.
Lemma~\ref{lem:scale-subpower} therefore gives, for all large $n$,
\[
 \frac{\chi(\tau_{2,n})}{\chi(\tau_n)}
       \ge \tau_n^{-b\eps/2}.
\]
In particular $\tau_n<\tau_{2,n}\le u_n$. This is the ordering
needed to apply \eqref{eq:scale-power-growth} with lower endpoint
$\tau_n$. Concavity then gives
\[
 \chi(\tau_n)\ge\frac{\tau_n}{\tau_{2,n}}\chi(\tau_{2,n})
 \ge C^{-1}\tau_n^{b_0+\eps}|\log\tau_n|.
\]
Therefore $b\le b_0+\eps$. Letting $\eps\downarrow0$ gives
$b\le b_0<1/2$. Choose any exponent strictly between $b_0$ and
$1/2$ and use the defining limsup to obtain the asserted bound on $s$;
integration gives the bound on $w$.
\end{proof}

\subsection{A martingale bin comparison}

We give the probability estimate needed to replace subpower control
by a uniform bound. Let $(Z_\upsilon)_{0\le\upsilon\le T}$ be a nonnegative
continuous martingale on a finite deterministic horizon, with
$\dd Z_\upsilon=\sigma_\upsilon\,\dd W_\upsilon$ and
$|\sigma_\upsilon|\le1$. Write $u=T-\upsilon$
for its remaining time. Suppose that, for each fixed $a>0$, there is
$c_a>0$, independent of $\upsilon$ and the sample path, such that
\begin{equation}\label{eq:scale-nondegenerate}
 Z\ge a\sqrt u\quad\Longrightarrow\quad |\sigma|\ge c_a>0
 \quad(a>0).
\end{equation}
We also assume that its law has a positive continuous density on
$(0,\infty)$ at each strictly interior time $0<\upsilon<T$.
Continuity of its paths
and absence of positive atoms then make each moving-bin probability
continuous for $0<u<T$.
For $r$ take variance time $\int_0^t s^2$, remaining time $u=w$,
and $|\sigma|=k$. Condition \eqref{eq:scale-nondegenerate} follows
from $m/\sqrt\tau\ge aD\ge a$, the upper barrier, and a secant:
the corresponding $z$ is bounded above, and $k$ is bounded below on
that half-line. No upper bound on $D$ is needed here.

\begin{lemma}[Bin comparison]\label{lem:scale-bins}
Under these assumptions there exist $0<u_0<T$ and finite constants
$C,P$, depending on the process and the constants $c_a$, such that
$M_j(u):=\Prob(j\sqrt u\le Z_{T-u}\le(j+1)\sqrt u)$ satisfies
\begin{equation}\label{eq:src27}
 M_j(u)\le Cj^P M_1(u),\qquad j\ge2,\quad 0<u\le u_0.
\end{equation}
\end{lemma}
\begin{proof}
We first prove a conditional steering estimate; the martingale need
not itself have the Markov property. After any stopping time
$S$, apply Brownian time change to the shifted continuous martingale
$Z_{S+v}-Z_S$ in the filtration $\mathcal F_{S+v}$. Extending it with
independent Brownian increments after the terminal time if necessary,
this gives $Z_{S+v}=Z_S+W_{Q_v}$, where
$Q_v=\int_S^{S+v}\sigma^2\,\dd t$ and $W$ is Brownian with respect
to a filtration containing $\mathcal F_S$ at time zero. Its Brownian
law thus holds conditionally on $\mathcal F_S$. All estimates below
are uniform in that initial information. Conditional spatial and time
scaling permits us to set the remaining time at $S$ equal to one.

Consider first a single step with time ratio
$\rho\in[\sqrt2,2]$, elapsed time $L=1-\rho^{-1}$, and starting
position $x\in[1,2]$. Put
\[
 L_{\min}=1-1/\sqrt2,\qquad \kappa=c_{1/2},\qquad
 q_* =\kappa^2L_{\min}/2,\qquad y_\rho=\frac{3}{2\sqrt\rho}.
\]
The terminal first bin is $[\rho^{-1/2},2\rho^{-1/2}]$.
Take the fixed spatial corridor $(1/2,3)$. Up to its first exit,
\eqref{eq:scale-nondegenerate} gives $\kappa^2\le\dot Q\le1$,
since the remaining time is at most one. Let $g_{x,\rho}(q)$ move
linearly from $x$ to $y_\rho$ on $[0,q_*]$, and stay at $y_\rho$
on $[q_*,1/2]$. It lies in $[1,2]$. On the Brownian event
\[
 \sup_{0\le q\le1/2}|x+W_q-g_{x,\rho}(q)|<1/16,
\]
no first corridor exit can occur before elapsed time $L$: at such an
exit, $Q_v\le v\le L\le1/2$, whereas the displayed event places
the exit position strictly inside the corridor. This first-exit
argument establishes the lower clock bound without assuming it on
the event being proved. In particular $Q_L\ge\kappa^2L\ge2q_*$,
so $Z_{S+L}$ is within $1/16$ of $y_\rho$ and lies in the terminal
first bin. The paths $g_{x,\rho}-x$ form a compact family of
piecewise linear paths starting at zero. Brownian motion has positive
probability to stay in a fixed uniform neighborhood of each such
path; a finite cover by neighborhoods of radius $1/32$ gives one
positive lower bound for the displayed event, independent of $x$
and $\rho$. This elementary tube fact follows, for example, by
requiring finitely many Gaussian increments to lie in prescribed open
intervals and the intervening Brownian bridges to stay in prescribed
open strips.

For any total ratio $u'/u\ge\sqrt2$, partition it into factors in
$[\sqrt2,2]$ and iterate the conditional estimate. The number of
steps is at most $1+\log_2(u'/u)$, so the resulting probability is
at least $c(u/u')^{P_1/2}$ for a fixed $P_1$. On the first step the
starting interval may instead be $[1.5,J]$: use corridor
$(1/2,J+1)$ and the same piecewise linear path, obtaining a first-step
constant depending only on fixed $J$, with the same subsequent power.
There is also a uniform bound when starting at exactly $1.5$ and
$u'/u\in[1,\sqrt2]$. In normalized units the event
$\sup_{q\le L_{\min}}|W_q|<1/16$ suffices, since
$[1.5-1/16,1.5+1/16]$ lies inside every corresponding terminal first
bin. Here only $Q_L\le L\le L_{\min}$ is used; no lower clock bound
is necessary.

The steering estimates now give a polynomial lower bound for reaching the
first bin across the scale ratios used below.  We use that bound in an
induction over decreasing dyadic slabs of remaining time, comparing the
other bins to the first and controlling distant starting bins by Gaussian
tails.  On $[u_0/2,u_0]$, $M_1$ has a positive minimum, so the assertion follows
from $M_j\le1$. Suppose the assertion has been established on earlier
slabs and let $u$ lie in the next slab. For $j\ge2$, set
$u_a=\min(j^2u,u_0)\ge2u$.
If $u_a=u_0$, steering from its first bin gives
$M_1(u)\ge c(u/u_0)^{P_1/2}M_1(u_0)$. Since $j^2u\ge u_0$,
the trivial bound $M_j(u)\le1$ proves the desired inequality with
power $P_1$ and a fixed constant.

In the uncapped case $u_a=j^2u$, condition at $u_a$. Any path that
starts below $1.5\sqrt{u_a}$ and ends in bin $j$ at time $u$ must
hit the \emph{moving} boundary
$Z=1.5\sqrt{\text{remaining time}}$. Let $A$ be the event of starting below this boundary and hitting it
before the terminal observation, and let $S$ be the first hitting
time. On $A$ its remaining time $u_h$ belongs to $[u,u_a]$.
The conditional steering estimate just proved, restarted at $S$,
gives a probability at least $cj^{-P_1}$ of ending in the first bin,
since $u_h/u\le u_a/u=j^2$; its short-ratio case covers
$u_h/u\in[1,\sqrt2]$. Therefore
\[
 M_1(u)\ge
 \E\!\left[\ind_A\Prob(Z\text{ ends in the first bin}\mid\mathcal F_S)
       \right]\ge cj^{-P_1}\Prob(A).
\]
The target contribution from starting positions below the boundary
is at most $\Prob(A)$, because every such target path must hit it.
This proves the desired bound $Cj^{P_1}M_1(u)$ without conditioning
on later target success.
Starting positions between $1.5\sqrt{u_a}$ and $J\sqrt{u_a}$ obey
the same bound, with $C_J$, by direct steering.

For an ancestor in bin $i\ge J$ at $u_a$, ending in the target bin
requires a martingale displacement at least a constant times
$i\sqrt{u_a}$, once $J$ is a fixed sufficiently large integer:
the upper target endpoint divided by $\sqrt{u_a}$ is
$(j+1)/j\le3/2$. The exponential martingale bound therefore gives
probability at most $C e^{-c_0i^2}$.
The induction hypothesis at $u_a$ bounds the ancestor mass by
$C_*i^P M_1(u_a)$, and steering gives
$M_1(u_a)\le Cj^{P_1}M_1(u)$. Thus the remaining contribution is at
most
\[
 C_*Cj^{P_1}M_1(u)\sum_{i\ge J}i^P e^{-c_0i^2}.
\]
Choose $P\ge P_1$, then $J$ making the coefficient of $C_*$ at most
$1/2$, and finally $C_*$ large enough to absorb the nonrecursive terms
and the initial slab. These choices are independent of the slab and
close the induction.
\end{proof}

\begin{proposition}[Bounded scale ratio]\label{prop:scale-bounded-D}
There is a constant $C$ such that $1\le D\le C$ near the endpoint.
Consequently $w\asymp\tau/\chi^2$ and
$\int_t^1s\asymp\tau/\chi$.
\end{proposition}
\begin{proof}
Repeat the stability argument in Lemma~\ref{lem:scale-subpower}, now
along an arbitrary support sequence on which $D\to\infty$.
Use the force bin $1\le r/\sqrt w\le2$ for the positive contribution.
The first estimate in \eqref{eq:src26}, valid also for $v\ge1$,
places its state interval below $-K\sqrt\tau$ when $D$ is large.
There $q_g\ge\Phi(-2)$ and $l\to0$ uniformly, so its contribution
is at least $c_0M_1(w)$. Its interval has length at least
$\chi\sqrt w$, because $|r_h|\le1/\chi$.

Lemma~\ref{lem:scale-strict-power} makes the exponential factor in
\eqref{eq:scale-truncated-density} uniformly bounded when comparing
this bin to $[-K\sqrt\tau,C_P\sqrt{\tau|\log\tau|}]$.
Using the infimum density on the bin, \eqref{eq:src24} bounds the
negative layer contribution by $CM_1(w)/D$. Errors $\tau^P$ are
negligible compared with $M_1(w)$: the density lower bound and the
bin length give $M_1(w)\ge c_0\chi\sqrt w\ge c_0\sqrt\tau$.
For $v>a\sqrt D$, sum \eqref{eq:src26} over the force bins and use
Lemma~\ref{lem:scale-bins}. This contribution is at most
$CM_1(w)\sum_{j\ge a\sqrt D-1}j^{P+2}e^{-c_0j^2}=o(M_1(w))$.
The other regions have nonnegative contribution as before. Stability
is contradicted. Formula \eqref{eq:scale-affine-gap} transfers the
result across gaps, with the same two cases as in the subpower proof.
The estimate for $\int s$ follows from Cauchy--Schwarz for the upper
bound and monotonicity for the lower bound.
\end{proof}

\subsection{Density comparison and endpoint masses}

Bounded $D$ implies $\dd\log w/\dd\log\tau=D^{-2}\ge\delta_0>0$.
Combining this with $w\asymp\tau/\chi^2$ gives, for
$0<\tau\le u\le1-t_0$,
\begin{equation}\label{eq:src28}
 \frac{\chi(u)}{\chi(\tau)}\le C(u/\tau)^{(1-\delta_0)/2},
 \qquad
 \int_\tau^{1-t_0}\frac{\gamma(u)}{\sqrt u}\,\dd u
 \le C\frac{\chi(\tau)}{\sqrt\tau}.
\end{equation}
For the integral use $\gamma(u)\le\chi(u)/u$ and integrate the first
bound; the power in the resulting integral is $-1-\delta_0/2$.

We need the additional gradient estimate $|k_h|\le C/\sqrt\tau$.
Differentiating \eqref{eq:scale-m-pde} gives
$k_\tau=k_{hh}/2+d m k_h+dkl$. On $[\tau/2,\tau]$, the drift has
spatial Lipschitz constant at most $2/\tau$, and the source is bounded
by $2/\tau$. For completeness, the semigroup estimate used here is
\[
 |\partial_x\E_x g(Z_T)|\le
 \frac{e^{LT}}{\sqrt T}\|g\|_\infty
\]
for $\dd Z=b(v,Z)\,\dd v+\dd B$, $|b_x|\le L$.
If $J_v=\partial_xZ_v$, Brownian integration by parts represents the
derivative by
$\E[g(Z_T)T^{-1}\int_0^T J_v\,\dd B_v]$.
Indeed contraction of the noise derivative with the adapted process
$J_v/T$ produces the endpoint variation $J_T$; its divergence is its
It\^o integral. For smooth simple adapted integrands this is Gaussian
integration by parts, and the general statement follows by
$L^2$ approximation. Cauchy--Schwarz and $|J_v|\le e^{LT}$ give the
estimate. Apply it to the bounded datum $k(\tau/2,\cdot)$ and to
Duhamel's source integral; $\int_0^{\tau/2}v^{-1/2}\,\dd v$ gives
$|k_h|\le C/\sqrt\tau$. Smooth approximation and localization justify
the argument for the time-measurable coefficient and linear-growth
drift.

Couple inverse flows from $h$ and $0$. At reverse time $u$, their
separation is at most $|h|\chi(u)/\chi(\tau)$. The logarithmic
Jacobian formula, the preceding gradient estimate, and
\eqref{eq:src28} consequently give
\[
 |\log \mathcal J_t(h)-\log \mathcal J_t(0)|\le C|h|/\sqrt\tau,
 \qquad |\mathcal I_t(h)-\mathcal I_t(0)|\le |h|/\chi.
\]
These imply
\begin{equation}\label{eq:src29}
 p_t(h)\asymp_K p_t(0)\quad(|h|\le K\sqrt\tau),
 \qquad
 p_t(h)\le Cp_t(0)\exp\{C|z|+\eps(\tau)z^2\},
\quad \eps(\tau)\longrightarrow0.
\end{equation}
Here and subsequently $\asymp_K$ denotes two positive constants
depending only on $K$ and the fixed optimizer.
The first comparison is uniform in all sufficiently late times for each
fixed $K$; the second holds for every $z\in\R$ at each such time,
with a single function $\eps(\tau)$ independent of $z$.
To check the global assertion, write
$\delta=|h|/\chi\le C\sqrt w\,|z|$.
The Gaussian density satisfies, for $0<\eta\le1/2$,
\[
 p_0(\mathcal I_t(h))\le C p_0(\mathcal I_t(0))^{1-\eta}
                  \exp(C\delta^2/\eta).
\]
H\"older's inequality with measure $\mathcal J_t(0)\,\dd\Prob$ bounds its
expectation by
$p_t(0)^{1-\eta}\chi^{-\eta}$.
Take $\eta=\sqrt w$; Lemma~\ref{lem:scale-strict-power} makes
$\chi^{-\eta}$ bounded, while \eqref{eq:scale-density-lower} bounds
$p_t(0)$ below. This proves the global claim with
$\eps(\tau)=C\sqrt w$ after enlarging constants. Apply the same
argument with $h$ and $0$ interchanged for the compact lower bound;
the Jacobian estimate is symmetric, and powers $p_t(0)^{\pm\eta}$
are uniformly bounded because $c_0\le p_t(0)\le C/\chi$.

Put $M_{\rm lay}=p_t(0)\sqrt\tau$. The first force bin has mass
comparable to $M_{\rm lay}$. Indeed $D$ is bounded, the barriers put
its endpoints in a fixed scaled spatial interval, and $k$ is bounded
below there. Thus its spatial length is comparable to
$\chi\sqrt w\asymp\sqrt\tau$, and \eqref{eq:src29} applies.
The same bin lemma applies to $X$ with remaining time $\tau$ and
coefficient $l$. If $X\ge a\sqrt\tau$, its scaled position is bounded
below by the upper barrier for $X$ on the negative half-line; a secant
then bounds $l$ below. Its first bin also lies in a fixed scaled
spatial interval. Therefore both families of bins have mass at most
$Cj^P M_{\rm lay}$, and
\begin{equation}\label{eq:src30}
 \E kl\asymp M_{\rm lay},\qquad
 \E(rX)\asymp (\tau/\chi)M_{\rm lay}.
\end{equation}
For the lower bounds integrate over a fixed central interval: both
$m/\sqrt\tau$ and $X/\sqrt\tau$ are bounded below there by the heat
barrier, and so are $k$ and $l$ by secants.
For the upper bounds split into large force bins, large gap bins,
and their complement. In force bin $j$, \eqref{eq:src26} bounds
$kl$ and $\chi rX/\tau$ by a polynomial in $j$ times
$e^{-c_0j^2}$; bounded $D$ is used here. In gap bin $j$, the positive
barrier places $z$ above a fixed positive multiple of $j$ for large
$j$, and \eqref{eq:src23}--\eqref{eq:src24} give the same estimate.
The polynomial bin masses are summable against these tails. The
complement is in a fixed scaled spatial interval and has mass
$O(M_{\rm lay})$.

Recall $J=P-I$ from the critical variational conditions, where
$P(t)=\alpha\E r_t^2$ and $I(t)=\int_0^ts^2$.
The identities for $W$ and $J'$ give globally
\[
 \frac{\alpha\E kl}{\chi}=\alpha c+\gamma J-\chi J'
                         =\alpha c-(\chi J)'.
\]
It\^o's formula for $rX$, using its covariation $-skl\,\dd t$ and
$RY=0$, now gives
\begin{equation}\label{eq:scale-mixed-moment}
 \E(r_tX_t)=c\tau+\chi(t)J(t)/\alpha.
\end{equation}
The terminal boundary term is zero: $J$ is bounded by
$\alpha\E R^2+\int_0^1s^2$ and $\chi\to0$.
At support points $J=0$, so \eqref{eq:src30} implies
$M_{\rm lay}\asymp\chi$ there.

These comparisons hold throughout a terminal interval. First the
ratios $\tau_2/\tau_1$ across support gaps are bounded. Otherwise
\eqref{eq:scale-affine-gap} and bounded $D$ give
$\chi_2/\chi_1\le C$ along a sequence of gaps with unbounded time
ratio. Their affine slope satisfies $d\le C/\tau_2$.
The inverse flow from $0$ at $\tau_1$ to $\tau_2$ therefore has a
bounded Jacobian and a position whose ratio to $\sqrt{\tau_2}$ has
a uniform Gaussian tail. To see the latter directly, the barrier gives
$|m(u,h)|\le C(\sqrt{\tau_2}+|h|)$ for $u\le\tau_2$;
Gronwall's inequality bounds the inverse position by a constant times
$\sqrt{\tau_2}$ plus the Brownian supremum.
Use the inverse density formula on this interval and
\eqref{eq:src29} at $\tau_2$. The linear exponential integrates
against this Gaussian tail, as does the quadratic exponential once
$\eps(\tau_2)$ is sufficiently small. Thus
$p_{1-\tau_1}(0)\le Cp_{1-\tau_2}(0)$.
But the support comparisons make their ratio comparable to
$(\chi_1/\chi_2)\sqrt{\tau_2/\tau_1}$, which tends to infinity.
This contradiction bounds the time ratios.

For a point inside a gap, compare to its larger endpoint by the same
inverse formula. The bounded time ratio and $d\le1/\tau$ give
bounded Jacobians and uniformly tight scaled inverse positions.
Equation \eqref{eq:src29} gives the upper comparison. For the lower
comparison restrict to a fixed compact containing a positive
probability of inverse positions and use its compact lower bound and
$\mathcal J_t\ge1$. As $\chi$ and $\sqrt\tau$ also have bounded ratios,
\begin{equation}\label{eq:scale-layer-mass}
 p_t(0)\sqrt\tau\asymp\chi(t)
\end{equation}
throughout a terminal interval.

The two endpoint windows and the atoms they exclude are shown in
Figure~\ref{fig:scale-endpoint-windows}.
\begin{proposition}[Two-sided endpoint mass estimates]\label{prop:scale-endpoint-masses}
There exist $t_1<1$ and constants $0<c_1<C_1<\infty$ such that,
for every $t_1<t<1$,
\begin{equation}\label{eq:src31}
 \Prob(0<R\le\sqrt w)\asymp\chi(t),\qquad
 \Prob(0<Y\le\sqrt\tau)\asymp\chi(t).
\end{equation}
Both comparisons use the same fixed constants $c_1,C_1$, after
enlarging their range if necessary; the time $t$ need not be a support
point.
\end{proposition}
\begin{proof}
For a lower bound, start the relevant martingale in the scaled interval
$[0.4,0.6]$. The barriers, bounded $D$, and
\eqref{eq:src29}--\eqref{eq:scale-layer-mass} put at least $c_0\chi$
mass in this interval. Its remaining quadratic variation is at most
the square of the corresponding scale. Brownian time change shows
that its displacement stays within $0.2$ of that scale with a fixed
positive probability: it suffices that standard Brownian motion stay
in $[-0.2,0.2]$ up to time one. Its endpoint is then strictly positive
and below the desired cutoff.

For the force upper bound, paths starting in force bin $j\ge2$ must
make a displacement of at least $(j-1)\sqrt w$ to finish below
$\sqrt w$, and hence pay a Gaussian tail. Paths starting in gap bin
$j\ge2$ must make a displacement at least $j\sqrt\tau$ to have
$R>0$, since that event forces $Y=0$. Sum the polynomial bin estimates
against these Gaussian tails. The remaining region
$\{r\le2\sqrt w,\ X\le2\sqrt\tau\}$ is in a fixed scaled spatial
interval and has mass $O(\chi)$.
For the gap upper bound interchange the two martingales, using
$Y>0\Rightarrow R=0$. This proves both statements.
\end{proof}

\begin{figure}[htbp]
\centering
\setlength{\unitlength}{1pt}
\begin{minipage}{.47\textwidth}
\centering
\begin{picture}(175,86)
\put(20,27){\vector(1,0){146}}
\put(25,27){\circle*{5}}
\put(110,23){\line(0,1){8}}
\put(25,42){\line(1,0){85}}
\put(25,39){\line(0,1){6}}
\put(110,39){\line(0,1){6}}
\put(67,49){\makebox(0,0){mass $\asymp\chi$}}
\put(25,14){\makebox(0,0){$0$}}
\put(110,14){\makebox(0,0){$\sqrt\tau$}}
\put(166,14){\makebox(0,0){$Y$}}
\put(85,75){\makebox(0,0){Positive gaps}}
\put(85,0){\makebox(0,0){atom at zero: $1/\alpha$}}
\end{picture}
\end{minipage}\hfill
\begin{minipage}{.47\textwidth}
\centering
\begin{picture}(175,86)
\put(20,27){\vector(1,0){146}}
\put(25,27){\circle*{5}}
\put(110,23){\line(0,1){8}}
\put(25,42){\line(1,0){85}}
\put(25,39){\line(0,1){6}}
\put(110,39){\line(0,1){6}}
\put(67,49){\makebox(0,0){mass $\asymp\chi$}}
\put(25,14){\makebox(0,0){$0$}}
\put(110,14){\makebox(0,0){$\sqrt w$}}
\put(166,14){\makebox(0,0){$R$}}
\put(85,75){\makebox(0,0){Positive forces}}
\put(85,0){\makebox(0,0){atom at zero: $1-1/\alpha$}}
\end{picture}
\end{minipage}
\caption{The unconditional endpoint laws of $Y$ and $R$ and the
windows in \eqref{eq:src31}. Each bracket
denotes an interval open at zero, so its indicated mass excludes the
atom. These are probability windows, with no density assumption or
common horizontal scale. The force law $F_J$ conditions on $R>0$ and
rescales $R$ to conditional mean one. The martingales $X_t$ and $r_t$ have remaining
quadratic variations at most $\tau$ and $w$, respectively. In the
identity $H_t=X_t-\chi r_t$, the corresponding spatial scales are
$\sqrt\tau$ and $\chi\sqrt w$. Proposition~\ref{prop:scale-bounded-D}
makes them comparable.}
\label{fig:scale-endpoint-windows}
\end{figure}

\subsection{Quantitative inverse-slope barriers}

For the rest of this section put $V=-\sqrt\tau\,k_h>0$.
Since $k$ decreases strictly from $1$ to $0$, we may also regard $V$
as a function of $k\in(0,1)$ at a fixed time. Subscripts $z$ below
mean differentiation in $z=h/\sqrt\tau$, so $V=-k_z$.

\begin{lemma}[Inverse-slope bounds]\label{lem:scale-inverse-slope}
At every interior critical time,
\begin{equation}\label{eq:src32}
 \frac97 k(1-k)^{3/2}\le V\le\frac65[k(1-k)]^{3/4}.
\end{equation}
\end{lemma}
\begin{proof}
We first establish the claim for a regularization that makes all
endpoint comparisons classical. In the approximants of
\eqref{eq:src20}, replace $m$ by $(\lambda_n+\chi_n)a_n$.
Its terminal value is the hinge, and it solves
\eqref{eq:scale-m-pde} with
$d=\gamma_n/(\lambda_n+\chi_n)$, $0\le\tau d\le1$.
At fixed $n$ this coefficient is integrable. Approximate it in $L^1$
by bounded coefficients satisfying the same inequality. Next replace
the initial hinge by its heat smoothing of variance $e'>0$ and
replace the bounded coefficient $d(\tau)$ by
$d(\tau)\tau/(\tau+e')$. Put $T=\tau+e'$, and denote the replaced
coefficient again by $d$. Then $D_0=Td\in[0,1]$.
These operations converge on interior compact sets, first as
$e'\downarrow0$, then in the coefficient approximation, then as
$n\to\infty$: the linear difference equation and the uniform
curvature and value bounds give local uniform convergence, and the
semigroup gradient estimates give convergence of spatial derivatives.
It is therefore enough to prove the bounds uniformly for this
regularization, with $\sqrt T$ in the definition of $V$.

The slopes remain in $[-1,0]$ and are strictly monotone. Their endpoint
values are $-1$ and $0$. Write $v=-k_h$, and let $h=h(\tau,k)$ be
the inverse position. From the equation for $k$,
\[
 h_\tau=\tfrac12v_k-dm+\frac{d k(1-k)}v,\qquad h_k=-1/v.
\]
Differentiate the first identity in $k$, use $m_k=k/v$, and compare
with the time derivative of $h_k$. For $V=\sqrt T\,v$ the result is
\begin{equation}\label{eq:scale-inverse-pde}
 \partial_{\log T}V=\tfrac12V^2V_{kk}
  +[\tfrac12+D_0(1-3k)]V-D_0k(1-k)V_k.
\end{equation}
The endpoint values are $V(0)=V(1)=0$, uniformly on compact time
intervals in this regularization. To justify this at spatial infinity,
$v$ solves
$v_\tau=v_{hh}/2+d m v_h+d(1-3k)v$, with Gaussian initial datum,
Lipschitz drift of at most linear growth, and bounded potential.
Its reverse diffusion representation shows that $v$ vanishes uniformly
at spatial infinity on a compact time interval: starting positions
tending to infinity remain there in probability, by the bounded
integrated Lipschitz norm. The corresponding equations for $k$ and
$1-k$ show that their fixed intermediate levels stay in bounded
spatial intervals. These facts give the asserted endpoint behavior
in inverse coordinates.

Let $t_k=k(1-k)$ and $b_+=(6/5)t_k^{3/4}$. Direct differentiation gives
\[
 -b_+b_+''=\frac{36}{25}
       \left(\frac34\sqrt{t_k}+\frac{3}{16\sqrt{t_k}}\right),\qquad
 1-3k-t_k b_+' /b_+=\tfrac14-\tfrac32k.
\]
For $k\ge1/6$, $-b_+b_+''\ge27/25>1$ suffices to make the
right-hand side of \eqref{eq:scale-inverse-pde} nonpositive at $b_+$.
For $k<1/6$, the required inequality is
\[
 \frac{27}{100}\frac{1+4t_k}{\sqrt{t_k}}
 \ge\frac32\sqrt{1-4t_k}.
\]
After squaring, it follows from
$x(1-x)/(1+x)^2\le1/8$ for $0\le x\le1$, with $x=4t_k$.
Thus $b_+$ is a stationary supersolution for every $D_0\in[0,1]$.

For $b_-=Ck(1-k)^{3/2}$, the right-hand side divided by $b_-$ is
bounded below by
\[
 \frac{1-k}{2}
 \left[1-C^2\left(3k(1-k)-\frac34k^2\right)\right].
\]
The quadratic in parentheses has maximum $3/5$; with $C=9/7$,
its product with $C^2$ is $243/245<1$. This proves the subsolution
inequality.

At the heat-smoothed initial time, the inverse profile is
$V_0(k)=\phi(\Phi^{-1}(k))$ and satisfies $V_0''=-1/V_0$.
Both barriers have zero endpoint values. Their stationary inequalities
at $D_0=0$ give $b_-b_-''\ge-1$ and $b_+b_+''\le-1$.
If $b_--V_0$ had a positive interior maximum, then
$b_-''\le V_0''=-1/V_0<-1/b_-$, a contradiction.
The same argument at a positive maximum of $V_0-b_+$ proves
$V_0\le b_+$. Parabolic comparison in
\eqref{eq:scale-inverse-pde} now yields the bounds at all later times.
One may apply the maximum principle on $[\delta,1-\delta]$ and then
let $\delta\downarrow0$. The zeroth-order coefficient in the
linearized difference is bounded above: it involves
$(V+b)b''$, which is bounded above for the supersolution and, for the
subsolution, on the only relevant region $V<b_-$. Multiplying the
difference by a sufficiently decreasing exponential in time deals
with this bound. The uniform endpoint vanishing deals with the
spatial boundary. Smooth approximation handles measurable time
coefficients. Passing through the regularizations proves
\eqref{eq:src32}.
\end{proof}

We record the compactness and tail information used below. Uniformly
near the endpoint,
\begin{equation}\label{eq:scale-V-tails}
 |V(z)|+|V_z(z)|\le C e^{-c_0z^2}.
\end{equation}
For $V$ use \eqref{eq:src32}, the positive Gaussian tail in
\eqref{eq:src24}, and the negative tail in \eqref{eq:src26}; bounded
$D$ makes $r/\sqrt w$ comparable to $|z|$ on the negative half-line.
For $V_z$ apply the semigroup gradient argument given above to the
linear equation
$v_\tau=v_{hh}/2+dmv_h+d(1-3k)v$ for $v=-k_h$ on
$[\tau/2,\tau]$, using spatial unit $\sqrt\tau$ and time unit
$\tau$. In these units the drift has uniformly bounded Lipschitz
constant and at most linear growth, and the potential is bounded.
The bound just proved for $V$ holds at every time in this interval,
so both the initial datum and the Duhamel source (the potential times
$\sqrt\tau v$) have uniform Gaussian decay in the scaled position.
A reverse diffusion started at $z$ stays at distance at least
$c|z|-C$ from zero except on an event of probability
$Ce^{-c_0z^2}$, by the flow separation estimate and the Brownian
supremum bound. Thus the second moment of either decaying datum
evaluated along that diffusion is at most $Ce^{-c_0z^2}$.
Cauchy--Schwarz in the integration-by-parts gradient formula bounds
its gradient by $CT^{-1/2}e^{-c_0z^2}$, with a possibly smaller
$c_0$. The source cost is integrable because
$\int_0^{1/2}T^{-1/2}\,\dd T<\infty$.
This proves \eqref{eq:scale-V-tails}.
Repeating the spatial derivative estimates on smaller time intervals
gives local bounds and equicontinuity of $k,V,V_z$. In particular
sequences of scaled spatial profiles have subsequences converging
locally with these derivatives. Time differentiability of the
coefficient is not needed in these semigroup estimates.

The limiting density ratios have a useful order property. For
$h_1<h_2$ with bounded scaled positions, the inverse displacement is
$O(\sqrt w)$. On $|\mathcal I_t(h_2)|\le w^{-1/4}$, the Gaussian density ratio
$p_0(\mathcal I_t(h_2))/p_0(\mathcal I_t(h_1))$ is at most $1+o(1)$, uniformly on the
chosen scaled compact. The complement contributes at most
$C\chi^{-1}e^{-c_0/\sqrt w}=o(p_t(0))$, using the strict power
bound. The inverse Jacobian is decreasing. Hence
\[
 \frac{p_t(h_2)}{p_t(0)}\le
 \frac{p_t(h_1)}{p_t(0)}+o(1).
\]
The compact two-sided bounds in \eqref{eq:src29} and the elementary
selection theorem for monotone functions now give subsequential
limits
\begin{equation}\label{eq:scale-density-limit}
 \frac{p_t(z\sqrt\tau)}{p_t(0)}\longrightarrow P_\infty(z)
\end{equation}
at continuity points, where $P_\infty$ is positive, nonincreasing,
and bounded above and below on every compact interval. One can apply
the selection theorem to the decreasing upper envelopes on each
compact, whose difference from the original functions tends uniformly
to zero, and then use a diagonal subsequence.

\subsection{A strict curvature inequality}

\begin{lemma}[Weighted curvature inequality]\label{lem:scale-curvature}
There is $c_1>0$ such that, at all sufficiently late times, with
$A_0=\tau\gamma/\chi\in[0,1]$,
\begin{equation}\label{eq:src33}
 \E\!\left[V_z^2+6A_0(2k-1)V^2+6A_0^2k^2(1-k)^2\right]
 \ge c_1p_t(0)\sqrt\tau.
\end{equation}
\end{lemma}
\begin{proof}
We first prove an unweighted inequality for every limiting profile
satisfying \eqref{eq:src32}. Put
\[
 F_A(z)=V_z^2+6A(2k-1)V^2+6A^2k^2(1-k)^2,
 \qquad 0\le A\le1.
\]
Its cumulative integral from $-\infty$ is nonnegative until $k=1/2$,
since every term is then pointwise nonnegative. For an endpoint with
$k=k_0\le1/2$, change variables from $z$ to $k$ to obtain
\[
 \int_{-\infty}^{z(k_0)}F_A(z)\,\dd z=H+6AC_1+6A^2B_1,
\]
where all following integrals are over $[k_0,1]$ and
\[
 \begin{aligned}
 H&=\int V V_k^2,& C_1&=\int(2k-1)V,& B_1&=\int t_k^2/V,\\
 T_1&=\int t_k,& U_b&=\frac65\int t_k^{3/4},&
 B_b&=\frac79\int k\sqrt{1-k}.
 \end{aligned}
\]
Cauchy--Schwarz applied to $(V^{3/2})'$ on $[1/2,1]$ gives
\[
 H\ge\frac89 V(1/2)^3\ge\frac{81}{5488\sqrt2}>\frac1{100}.
\]
Put $b=T_1/U_b$. Integration by parts and weighted
Cauchy--Schwarz give
\begin{align*}
 C_1&=V(k_0)\{t_{k_0}-bV(k_0)/2\}
       +\int V(t_k/V-b)V_k\\
 &\ge V(k_0)\{t_{k_0}-bV(k_0)/2\}
       -\sqrt{HB_1\left(1-\frac{T_1^2}{U_bB_b}\right)}.
\end{align*}
Indeed the square integral in this application is
$B_1-2bT_1+b^2\int V\le B_1-T_1^2/U_b$;
\eqref{eq:src32} gives $\int V\le U_b$ and $B_1\le B_b$.

Here are explicit bounds for the ratio appearing in this estimate.
Under probability weight $t_k/T_1$, the mean of $t_k^{-1/4}$ on
$[k_0,1]$, $0\le k_0\le1/2$, is at most its full-interval mean.
This follows by symmetry and monotonic decrease on $[0,1/2]$:
the integral of $t_k$ times its deviation from the full mean over
$[0,k_0]$ is nonnegative. Cauchy--Schwarz bounds the full mean by
$\sqrt{3\pi/4}$, since $\int_0^1\sqrt{k(1-k)}=\pi/8$
and $\int_0^1k(1-k)=1/6$.
The mean of $(1-k)^{-1/2}$ increases with the lower cutoff; at
$k_0=1/2$ it is $7\sqrt2/5<2$.
If $k_0<1/50$, its value is less than $17/10$: integrate
$k\sqrt{1-k}$ explicitly and evaluate at $k_0=1/50$.
Consequently, using $\pi<22/7$,
\[
 \frac{U_bB_b}{T_1^2}<\frac{28}{15}\sqrt{\frac{33}{14}}
       <\frac{287}{100},
 \qquad
 \frac{U_bB_b}{T_1^2}<\frac{119}{75}\sqrt{\frac{33}{14}}
       <\frac{61}{25}\quad(k_0<1/50).
\]

The endpoint term is nonnegative unless $t_{k_0}<1/64$.
Indeed $b\le5/(6\sqrt2)$ and $V(k_0)\le(6/5)t_{k_0}^{3/4}$.
When $t_{k_0}<1/64$, necessarily $k_0<1/50$, and the endpoint term
is still greater than $-1/20000$. To verify the latter, its dependence
on $V(k_0)$ is concave, so its minimum on the permitted interval is
at an endpoint. With $x=t_{k_0}^{1/4}$, the nonzero endpoint is
bounded below by
$(6/5)x^6(x-1/(2\sqrt2))$. The absolute value of this polynomial's
minimum on $x\ge0$ is
$6^7/(5\cdot7^7\cdot1024\sqrt2)<1/20000$.

Complete the square in $A$ in $H+6AC_1+6A^2B_1$ and use $A\le1$
for the endpoint loss. The result is at least
\[
 H\left[1-\frac32\left(1-\frac{T_1^2}{U_bB_b}\right)\right]
 -\frac6{20000}\ind_{\{t_{k_0}<1/64\}}.
\]
The preceding rational bounds make this greater than
\[
 \min\left\{\frac{13}{57400},\frac7{6100}-\frac3{10000}\right\}
 >\frac1{5000}.
\]
Thus, if $C_A(z)=\int_{-\infty}^zF_A$, then $C_A\ge0$ everywhere
and $C_A\ge1/5000$ after the unique point $z_*$ at which $k=1/2$.

Let $P$ be any positive nonincreasing limiting density in
\eqref{eq:scale-density-limit}. Integration by parts against its
Stieltjes measure gives
\[
 \int_\R F_A P\,\dd z
 =C_A(+\infty)P(+\infty)+\int_\R C_A\,\dd(-P)
 \ge \frac1{5000}P(z_*)>0.
\]
The negative boundary product is zero by the Gaussian tail of $C_A$
and the exponential bound for $P$; one may first truncate both tails.
The inequality includes the case $P(+\infty)>0$, which contributes
the displayed boundary term.

Finally suppose \eqref{eq:src33} failed along a sequence to the
endpoint. Extract a limit of $A_0\in[0,1]$, of the spatial profiles,
and of the density ratios. Equations
\eqref{eq:src29} and \eqref{eq:scale-V-tails} give an integrable
Gaussian majorant for the normalized integrands, so dominated
convergence applies. The level $z_*$ stays in a fixed compact by the
value and slope barriers, and the limiting density is uniformly
bounded below there. The preceding strictly positive lower bound
contradicts failure. This proves \eqref{eq:src33}.
\end{proof}

\subsection{Full increase support near the endpoint}

\begin{proposition}[Absence of terminal support gaps]\label{prop:scale-full-support}
The increase support of $\gamma$ contains a terminal interval $(t_*,1)$.
On this interval $J=J'=0$, and $\gamma$ has a continuous version with
\begin{equation}\label{eq:src34}
 A_0=\frac{\E V^2}{2\E k^2(1-k)}.
\end{equation}
\end{proposition}
\begin{proof}
On an open gap, $\gamma$ is constant. Write $b=f_{hh}<0$ and
$H_b=\E b(t,H_t)^2$. The diffusion equations and It\^o's formula give
\[
 H_b'=\E b_h^2-2\gamma\E b^3,\qquad
 H_b''=\E b_{hh}^2-12\gamma\E(b b_h^2)+6\gamma^2\E b^4.
\]
For example $\dd b=b_h\,\dd B-\gamma b^2\,\dd t$ and
$\dd b_h=b_{hh}\,\dd B-3\gamma b b_h\,\dd t$; these yield both
identities by the product rule. Localization and the spatial tails
justify taking expectations.

Substitute $b=-sk$, $b_h=sV/\sqrt\tau$,
$b_{hh}=sV_z/\tau$, and $\gamma s\tau=A_0$. Then
\begin{align*}
 \frac{2H_bH_b''-3(H_b')^2}{2H_bs^2/\tau^2}
 ={}&\E V_z^2+12A_0\E(kV^2)+6A_0^2\E k^4\\
 &-\frac{3}{2\E k^2}(\E V^2+2A_0\E k^3)^2.
\end{align*}
Use $\E k^3/\E k^2\le1$ and the identity
\[
 \E k^4-\frac{(\E k^3)^2}{\E k^2}
 =\E[k^2(1-k)^2]-\frac{(\E[k^2(1-k)])^2}{\E k^2}.
\]
The last display is therefore at least the left-hand side of
\eqref{eq:src33} minus
\[
 \frac{6A_0^2(\E[k^2(1-k)])^2}{\E k^2}
 +\frac{3(\E V^2)^2}{2\E k^2}.
\]
Both errors are $O(\chi^2)$: \eqref{eq:src30},
\eqref{eq:scale-layer-mass}, and the Gaussian derivative tails bound
each numerator's unsquared expectation by $C\chi$.
Moreover $\E k^2\to\Prob(R>0)=1/\alpha>0$ by bounded convergence.
The main term in \eqref{eq:src33} is at least $c_0\chi$.
Hence $2H_bH_b''-3(H_b')^2>0$ on every sufficiently late open gap.
Equivalently $\psi=(\alpha H_b)^{-1/2}$ is strictly concave there.

Choose a support point $t_*$ beyond the time at which this determinant
inequality holds uniformly. Every support gap to the right of $t_*$
has both endpoints in the increase support: the left endpoint lies
at or after $t_*$, and accumulation at $1$ excludes a final gap.
For any such gap,
stability gives $\psi\ge\chi$ at those endpoints, whereas $\chi$ is
affine on the gap. Strict concavity consequently gives $\psi>\chi$
at every interior point. But then
$J'=\alpha H_b-s^2<0$ throughout the gap, contradicting $J=0$ at
both endpoints. Thus there are no complementary intervals to the
right of $t_*$. Closedness of the support proves that it contains
$(t_*,1)$.

On this interval the variational conditions give
$J=J'=0$, $\E k^2=1/\alpha$, and $\E kl=c\chi$.
Apply It\^o's formula to $k(t,H_t)^2$:
\[
 0=\frac{\dd}{\dd t}\E k^2
   =\frac1\tau\E V^2-\frac{2\gamma}{\chi}\E[k^2(1-k)]
 \quad\hbox{for almost every }t.
\]
The denominator is strictly positive. Interior spatial regularity,
the diffusion representation, and dominated convergence make the
quotient on the right of \eqref{eq:src34} continuous in $t$.
It supplies the continuous version of $A_0$, and hence of $\gamma$,
and proves the asserted identity everywhere on the terminal interval.
\end{proof}


\section{Selection of the critical similarity profile}
\label{sec:contraction}

The estimates of Section~\ref{sec:scales} leave a scalar quantity to
identify: the limiting logarithmic derivative of the susceptibility.
We prove a rigidity statement for every limit obtained by translating
the critical layer in logarithmic time.  Its quantitative part consists
of finitely many inequalities for one-dimensional functions.  The
analytic implications of those inequalities are proved here; their
finite verification and rounding-error bounds are given in
Section~\ref{sec:certificate}.
The stationary scaling equations have their antecedents in
\cite[Section IX]{CKPUZ2014} and \cite[Section V]{FPSUZ2017}.
Here their selection is deduced from the entire limits of the critical
problem constructed in Sections~\ref{sec:critical}--\ref{sec:scales}.

The coefficient reduction first confines $A$ to $[0,.41]$. With this
bound, stationary-reference comparisons give curvature and density
barriers for every entire critical limit. Using its critical integral
identities, signed tests then narrow the coefficient interval to
$[.285,.303]$, where a strict contraction around the root forces
$A\equiv a_*$. Compactness finally transfers this identification to the
original critical layer. The opening table in Section~\ref{sec:certificate}
maps the finite checks to these analytic steps.

\subsection{Entire limits and their identities}

We regard $\chi$, $\gamma$, and $m$ as functions of remaining time:
their value at $\tau$ in this section is their previous value at
physical time $t=1-\tau$.  The density $p_t$ retains its physical-time
index.  Write $v=\log\tau$, $u=-v$, and $q=1/2$.  Throughout this section the
letters $M,X,P$ refer to rescaled quantities:
\[
 \begin{aligned}
 M(v,z)&=\frac{m(\tau,z\sqrt\tau)}{\sqrt\tau},&
 X(v,z)&=z+M(v,z),\\
 P(u,z)&=\frac{\sqrt\tau}{\chi(\tau)}p_{1-\tau}(z\sqrt\tau),&
 A&=\frac{\tau\gamma}{\chi}.
 \end{aligned}
\]
Thus $P$ is a positive density divided by a mass scale and need not
integrate to one.  Set $k=-M_z$, $l=1-k$, and $V=-k_z=M_{zz}>0$.
Spatial derivatives in this section are derivatives in $z$, except
where inverse-slope coordinates are explicitly specified.  Direct
change of variables gives
\begin{equation}
 \begin{split}
 M_v&=\tfrac12M_{zz}+BM_z+(A-q)M,\\
 B&=qz+AM,\\
 P_u&=\tfrac12P_{zz}-(BP)_z+AP.
 \end{split}
 \label{eq:src35}
\end{equation}
The opposite time directions are essential: $M$ is parabolic in $v$,
whereas $P$ is parabolic in $u$.  The last term in the density equation
comes from $\chi_u=-A\chi$.

\begin{lemma}[Compactness in logarithmic time]
\label{lem:contr-entire}
Every sequence $u_n\to\infty$ has a subsequence along which the
translated rescaled profiles converge locally to a solution of
\eqref{eq:src35} for all $u\in\R$.  The limiting coefficient $A$ is
continuous and takes values in $[0,1]$.  The value barriers
\eqref{eq:src23}, the inverse-curvature bounds \eqref{eq:src32}, and the
Gaussian bounds for $k$ on the positive half-line, $l$ on the negative
half-line, and $V,V_z$ on both half-lines hold uniformly in entire time.
Moreover $P>0$, $P\le C\exp(C|z|)$, and on every bounded spatial
interval $P$ has a positive lower bound independent of entire time.
The limits satisfy
\begin{equation}
 \int_\R MXP\,\dd z=\int_\R klP\,\dd z=c,
 \qquad
 A=\frac{\displaystyle\int_\R V^2P\,\dd z}
          {\displaystyle2\int_\R k^2lP\,\dd z},
 \label{eq:src36}
\end{equation}
where $c=\E R>0$ is the critical normalization constant.
\end{lemma}

\begin{proof}
The density estimate \eqref{eq:src29}, together with
$p_{1-\tau}(0)\sqrt\tau/\chi\asymp1$, bounds $P$ above and below on
spatial compact sets.  The same estimate bounds it globally by
$C\exp(C|z|+o(1)z^2)$, which is integrable against the Gaussian tail
bounds already proved.  The value and curvature barriers bound the
coefficients of \eqref{eq:src35} locally, uniformly after sufficiently
late translation.  Interior estimates for the localized heat equation,
applied forward in $v$ to $M$ and forward in $u$ to $P$, give local
bounds and equicontinuity for their spatial derivatives.  Applying the
equations on a slightly smaller cylinder gives time equicontinuity.
For precision, one may first differentiate the localized heat-kernel
formula in space: the initial $C^1$ bound controls the drift term, and
iteration gives the derivatives needed below.  Bounded, initially
merely measurable, time coefficients suffice in this argument.

On the full-support endpoint interval, \eqref{eq:src34} expresses $A$
as a quotient of spatial integrals.  Its denominator, after division
by $\chi$, is bounded below uniformly.  Indeed the value barriers
place an interior interval of slope values in a bounded spatial
interval; inversion and \eqref{eq:src32} give positive spatial width,
and central positivity of $P$ gives a positive integral of $k^2lP$.
Gaussian tails and the density bound permit uniform truncation of both
integrals.  Consequently the same compactness gives compactness of
$A$ itself.  Pass first in the integral form of the equations; the
quotient then shows that the limiting $A$ is continuous.

The critical identities in physical units are
$\E[r_tX_t]=c\tau$ and $\E[kl]=c\chi$.  They become the first two
integrals in \eqref{eq:src36}.  Their passage to the limit, and passage
in the quotient, use the value barriers, the negative-tail estimate
\eqref{eq:src26}, and the Gaussian domination just described.  The
constants in all these arguments are uniform under translation, so
the conclusions hold on the whole time axis.
\end{proof}

We now use $Q$ for the logarithmic spatial slope of $P$ throughout this
section. There is a global, translation-uniform bound:
\begin{equation}
 Q=(\log P)_z\in[-C,0],\qquad
 Q_u=\tfrac12Q_{zz}+(Q-B)Q_z-(q-Ak)Q-AV.
 \label{eq:src37}
\end{equation}
Here nonincrease of $P$ follows from the inverse-flow comparison in
Section~\ref{sec:scales}.  To obtain the lower bound for $Q$, compare
two inverse positions corresponding to fixed rescaled positions
$z_1,z_2$.  Their displacement is $O(\sqrt w)$, so the ratio of the
initial Gaussian densities tends to one after truncating the inverse
positions as in that comparison; the complement is negligible.  The
inverse-Jacobian logarithmic ratio is bounded by $C|z_1-z_2|$, with
$C$ independent of the chosen compact set.  Hence in the limit
$|\log P(z_1)-\log P(z_2)|\le C|z_1-z_2|$ for all $z_1,z_2$.
Differentiating the density equation gives \eqref{eq:src37}.

All subsequent integrations by parts can be made first with compact
cutoffs.  To justify removal of the cutoffs also after differentiation,
use cylinders of spatial radius $(1+|z|)^{-1}$ and time length
$(1+|z|)^{-2}$.  Rescaling bounds the drift and its required spatial
derivatives there.  Interior derivative estimates thus introduce only
polynomial losses into the Gaussian bounds for $V$ and its derivatives
and into the exponential bounds for $P$.  These losses remain
integrable in every expression below.

\subsection{An absolute improvement of the coefficient bound}

Let $\mathcal L=\tfrac12\partial_z^2+B\partial_z$.  Differentiating the
$M$ equation gives
\[
 (\partial_u+\mathcal L)k=-Akl,
 \qquad
 (\partial_u+\mathcal L)V=-(q+A-3Ak)V.
\]
For a smooth integrable function $F=F(u,z)$, the density equation
therefore supplies the useful differentiation rule
\[
 \frac{\dd}{\dd u}\int FP
   =\int\bigl(F_u+\mathcal LF+AF\bigr)P.
\]
Apply it to $V^2$ and $k^2l$.  The quotient in \eqref{eq:src36} is
$C^1$ and obeys
\begin{equation}
 2\left(\int_\R k^2lP\,\dd z\right)A_u
 =\int_\R\bigl[(V_z)^2+(12Ak-5A-1)V^2
                    +6A^2k^2l^2\bigr]P\,\dd z.
 \label{eq:src38}
\end{equation}
For example the derivative of the denominator without its factor two
is
\[
 \int\bigl[A(-1+3k)k^2l+(1-3k)V^2\bigr]P.
\]
In simplifying the quotient derivative one uses
$\int(-2AV^2+4A^2k^2l)P=0$, which is precisely
\eqref{eq:src36}.

\begin{lemma}[Coefficient reduction]
\label{lem:contr-A}
Every entire limit in Lemma~\ref{lem:contr-entire} has $0\le A\le0.41$.
\end{lemma}

\begin{proof}
We first show that the right side of \eqref{eq:src38} is strictly
positive whenever $0.41\le A\le1$.  Denote $t_k=k(1-k)$.  Its
unweighted integrand is positive while $k\ge5/8$.  For a cumulative
integral ending at a slope $k_0\le5/8$, change variables from $z$ to
$k$, so the cumulative runs from $k_0$ to $1$.  Put
\[
 b=2(0.6-k),\qquad E=-0.4t_k,\qquad
 K=12Ak-5A-1-b^2+0.8(1-2k).
\]
Integrating the inequality
$V(V_k-b-E/V)^2\ge0$ gives the following lower bound for that
cumulative:
\begin{equation}
 -(bV^2+2EV)(k_0)
 +\int_{k_0}^1
       \left[2V^2+KV-2bE+(6A^2-0.16)\frac{t_k^2}{V}\right]\dd k.
 \label{eq:src39}
\end{equation}
The boundary term at $k=1$ vanishes by \eqref{eq:src32}.

The finite check \texttt{checkA} proves strict positivity of
\eqref{eq:src39}, uniformly for $A\in[0.41,1]$ and $k_0\in[0,5/8]$.
Its input is only the two explicit bounds in \eqref{eq:src32}.
It partitions $A$ and $k$, subdivides the interval between the lower
and upper allowed values of $V$, and integrates cellwise lower bounds.
For the boundary quadratic its maximum is at an endpoint when
$b\ge0$; when $b<0$ it is decreasing for $V\ge0$.  The two omitted
endpoint cells each cost at most $9$ times their width, and on the
first cell the boundary expression before its negation is less than
$0.001$.  A possible positive partial cell is subtracted rather than
included in full.  These are exactly the corrections in
\texttt{checkA}; Section~\ref{sec:certificate} verifies the arithmetic.

For completeness, positive cumulatives are preserved under the weight
$P$.  If $f$ is the unweighted integrand and
$F(z)=\int_{-\infty}^z f$, then $F>0$ at every finite point and its
terminal value is positive.  Since $P$ is positive and nonincreasing,
Stieltjes integration by parts gives
$\int fP=F(\infty)P(\infty)+\int F\,\dd(-P)>0$; limits of truncated
expressions give the same formula when a tail value is infinite.
The Gaussian/exponential bounds make the left endpoint product zero
and justify the limits.  Strictness follows either from
$P(\infty)>0$ or from the nonzero positive measure $\dd(-P)$.

Finally let $L$ be the endpoint limsup of the original rescaled $A$.
Choose translates realizing $L$ at time zero and take an entire limit.
In that limit $A(0)=L$ and $A(u)\le L$ for all $u$.  If $L\ge0.41$,
\eqref{eq:src38} makes $A_u(0)>0$, contradicting this maximum.
Thus $L<0.41$, which implies the claimed, slightly weaker bound for
every entire limit.
\end{proof}

\subsection{Stationary reference profiles}

For $a\in[0,0.41]$, let $M_a$ denote the entire solution of the first
equation of \eqref{eq:src35} with $A\equiv a$ and the value barriers
\eqref{eq:src23}, in the inherited admissible class
$-1\le M_z\le0$, $M_{zz}>0$.  This solution exists and is unique.
For existence,
evolve the heat-smoothed hinge of variance one with coefficient
$a/(1+\tau)$ in the original $m$ equation, and rescale by
$T=1+\tau$.  The resulting equation is autonomous in $\log T$; the
heat and Mills-ratio barriers, together with \eqref{eq:src32}, give
compactness of translates.  For uniqueness, subtract two entire
solutions.  Their bounded difference solves a linear equation whose
zeroth-order coefficient is $al_2-q\le a-q<0$ and whose drift is
$qz+aM_1$.  Comparison from time $v-T$ to $v$ bounds its supremum by
$Ce^{-(q-a)T}$; let $T\to\infty$.  Time translation and uniqueness
then imply stationarity.  The same compactness and uniqueness argument
proves continuity of $M_a$ in $a$ on compact spatial sets, with its
needed spatial derivatives.

The maximum principles used here and below are valid on the whole
line.  For bounded differences, localize to $[-R,R]$ and add
$\varepsilon e^{Ct}(1+z^2)$; bounded diffusivity and at most linear
drift allow $C$ independent of $R$.  Send $R\to\infty$ and then
$\varepsilon\downarrow0$.  This also proves the corresponding
comparison for bounded ratios whenever their equation has bounded
diffusivity and at most linear drift.

At a fixed reference, abbreviate $p=q-a>0$ and $W=M_a''$.  Its ODE
and its derivative give
\begin{equation}
 W=2pMl+Xk=2pM+2Bk,\qquad
 W_z=-2BW+2at_k,\qquad
 W_k=2B-2a\frac{t_k}{W}.
 \label{eq:src40}
\end{equation}
Here and subsequently $W_k$ is taken at the inverse position
$z_a(k)$.  In particular, for $x>0$,
\[
 \begin{array}{lll}
 z=x: & k\le e^{-qx^2},& M\le k/x,\\
 z=-x:& l\le e^{-px^2},& X\le l/(2px).
 \end{array}
\]
Indeed $W\ge xk$ on the positive side, and $W\ge2pxl$ on the
negative side; integrate the logarithmic derivatives of $k$ or $l$,
and then their tail integrals.  These formulas also bound the tails
of $W$ by a polynomial times the same Gaussians.

Define the reference density up to a positive factor by
\[
 \tfrac12P_a''-(B_aP_a)'=-\lambda(a)P_a,
 \qquad Q_a=(\log P_a)'.
\]
The substitution $P_a=\exp(\int_0^zB_a)\psi_a$ gives the
self-adjoint operator
\[
 \mathcal H_a=-\tfrac12\partial_z^2+\tfrac12(B_a^2+B_a').
\]
Its potential is bounded below by $c_0z^2-C_0$, uniformly in
$a\in[0,0.41]$.  Minimization of its quadratic form over unit vectors
therefore has a minimizer: the confining bound makes the mass outside
large compact intervals uniformly small, and bounded one-dimensional
Sobolev norms give compactness on each such interval.  Replace a
minimizer by its absolute value; the ODE and uniqueness for initial
values make it strictly positive.  The lowest eigenvalue is simple,
by the Wronskian identity for two positive square-integrable
solutions.  Smooth compactly supported trials give upper
semicontinuity of $\lambda(a)$; compactness of minimizing sequences
gives lower semicontinuity.  Thus $\lambda$ is continuous.

Section~\ref{sec:certificate} certifies
\[
 \lambda(0.2936533)>0.2936533,\qquad
 \lambda(0.2936535)<0.2936535.
\]
Fix a root $a_*\in(0.2936533,0.2936535)$ of $\lambda(a)=a$.
For every tabulated density reference used below, and for each root in
this bracket, the tail construction in that section proves
polynomial-order tails for $P_a$ and $Q_a\to0$ at both ends.  No table at the unknown root is needed:
at $\lambda=a_*$ the two tail coefficients in
\eqref{cert:tail-coefficients} obey
\[
 \begin{aligned}
 c_{{\rm tail},+}&=2a_*-1\in(-0.4126934,-0.4126930),\\
 c_{{\rm tail},-}&=\frac{a_*}{1/2-a_*}-1\in(0.4231063,0.4231088).
 \end{aligned}
\]
These intervals lie inside the sided hypotheses of the analytic
tail construction, and \eqref{eq:src40} supplies its coefficient
estimates uniformly for the root.  Also $Q_a\le0$: its stationary equation
is \eqref{eq:src37}, since differentiating the eigenvalue equation
removes the constant eigenvalue term.  A positive interior maximum
would have right side at most $-(q-ak)Q_a-aW<0$.  Together with the
zero tail limits this excludes positive values (also at $a=0$).

At the root the two integral identities in \eqref{eq:src36} hold,
after a common positive normalization of $P_*$.  In fact
\[
 \mathcal L(MX)=(1-a)MX-kl,
 \qquad
 \mathcal L(kl)=-a(1-2k)kl-W^2.
\]
Integration against $P_a$ replaces $\mathcal L$ by $-\lambda(a)$.
At $\lambda=a$, the first formula gives
$\int(MX-kl)P_a=0$ and the second gives
$\int W^2P_a=2a\int k^2lP_a$.  Boundary terms vanish by the reference
tail estimates.  Normalize the positive first integral to $c$.

\subsection{Comparisons with the references}
\label{subsec:contr-comparisons}

For the comparisons and integral tests, an \emph{input} consists of a
continuous coefficient $A:\R\to[0,0.41]$ and a solution $M_i(v,z)$ of
the first equation in \eqref{eq:src35}, defined for all $v,z\in\R$.
We require $0<k_i=-M_{i,z}<1$, $V_i=M_{i,zz}>0$, the scaled value
barriers \eqref{eq:src23}, and the inverse-curvature bounds
\eqref{eq:src32}, uniformly in time.  A density input also
includes a positive function $P_i(u,z)$ whose logarithmic slope
$Q_i=(\log P_i)_z$ is bounded on $\R^2$, nonpositive, and satisfies
\eqref{eq:src37}, with $u=-v$.  The regularity is that of the entire
limits above, so the equations and comparisons hold classically on
interior space--time cylinders.

Every critical entire limit is an input by
Lemmas~\ref{lem:contr-entire} and~\ref{lem:contr-A}.
Every stationary $M_a$ is an input.  The stationary density references
used below, including every root in the certified bracket, are density
inputs by the tail construction and the bound $Q_a\le0$.  The density
comparison class requires neither the full density equation in
\eqref{eq:src35} nor the integral identities \eqref{eq:src36}.  Indeed a stationary
eigenreference with $\lambda(a)\ne a$ need not solve that density
equation, but its logarithmic derivative solves \eqref{eq:src37}:
spatial differentiation removes the constant eigenvalue term.
Consequently all the comparisons and the resulting \texttt{Env}
enclosures apply to such references.  Later, the reference at
$0.293650$ enters only through direct functional evaluation and
interpolation.  Converting a functional sign into a coefficient bound
will additionally require the homogeneous consequences of
\eqref{eq:src36}; a reference satisfies them at a root.
If $a\le A\le b$, then at the same spatial position
\begin{equation}
 M_a\le M_i\le M_b.
 \label{eq:src41}
\end{equation}
For example $\delta=M_i-M_a$ solves, in $v$, a linear equation with
drift $B_i$, damping $q-Al_a\ge q-b>0$, and source
$(A-a)M_al_a\ge0$.  Its negative part vanishes by comparison from the
arbitrarily remote past.  The upper comparison is identical with the
inequality reversed.

Curvatures must instead be compared at the same slope $k$.  Fix a
reference $a$, put
\[
 T_a=1-3k-t_k\frac{W_k}{W}
     =\frac{2pMl}{W}-2k+2a\left(\frac{t_k}{W}\right)^2,
 \qquad D_a=q+aT_a,
\]
and write $R=V_i/W$ as a function of the reference position
$z=z_a(k)$.  Substitution into the inverse-curvature equation of
Section~\ref{sec:scales} gives
\begin{equation}
 R_v=\tfrac12R^2R_{zz}
    +\left[-R^2B+(aR^2+A)\frac{t_k}{W}\right]R_z
    +R\bigl[D_a(1-R^2)+(A-a)T_a\bigr].
 \label{eq:src42}
\end{equation}
Equivalently, its derivative terms in $k$ are
$\tfrac12R^2W^2R_{kk}+(R^2WW_k-At_k)R_k$.

\begin{lemma}[A bounded curvature ratio]
\label{lem:contr-heat-ratio}
For every input, at the same $k\in(0,1)$,
\begin{equation}
 \sqrt{1-0.82\cdot1.12}
 \ \le\ \frac{V_i}{W_0}
 \ \le\ \sqrt{1+0.82\cdot0.13},
 \qquad W_0(k)=\phi(\Phi^{-1}(k)).
 \label{eq:src43}
\end{equation}
Consequently the ratio of any two input curvatures is less than $3.8$.
\end{lemma}

\begin{proof}
The central reference check and the tail estimates proved below give
$-1.12\le T_0\le0.13$.  Let
$\alpha=q+0.41\cdot0.13$ and solve
$C'=C(\alpha-qC^2)$ with $C(v)\to\infty$ as $v\downarrow v_0$.
Equation~\eqref{eq:src42} says that $CW_0$ is a supersolution for
$V_i$ in the $k$ equation.  This comparison does not initially assume
a bounded ratio.  Indeed in the equation for $V_i-CW_0$ the extra
zeroth-order coefficient is
$\tfrac12(V_i+CW_0)CW_{0,kk}\le0$, because $W_{0,kk}=-1/W_0$.
The remaining positive coefficient is bounded independently of $C$.
At time $v_0+h$, the positive part of $V_i-CW_0$ has supremum tending
to zero as $h\downarrow0$: on interior $k$-intervals $C\to\infty$,
and near both endpoints the uniform upper bound
\eqref{eq:src32} tends to zero.  Comparison, followed by
$h\downarrow0$ and $v_0\to-\infty$, gives
$R\le\sqrt{\alpha/q}$.

To obtain a positive lower ratio without assuming one, set $U=1/R$
in heat-reference position.  Differentiating \eqref{eq:src42} and
dropping its nonpositive gradient-square term gives
\[
 U_v\le\tfrac12R^2U_{zz}
       +\left[-qR^2z+A\frac{t_k}{W_0}\right]U_z
       -\delta_0U+qR,
 \qquad \delta_0=q-0.41\cdot1.12>0.
\]
Here $0\le At_k/W_0\le C$ and the already proved upper bound controls
$R$.  The lower bound in \eqref{eq:src32} and elementary normal-tail
estimates give, uniformly in entire time,
\[
 U(z)\le C\bigl(1+z_+^2+e^{z_-^2/3}\bigr).
\]
The diffusion representation for this subsolution has no surviving
boundary term at infinity.  Here is a direct justification.  For
$r<q$ and an even integer $m\ge2$, use a $C^1$, piecewise $C^2$
Lyapunov function comparable to
$1+z_+^m+e^{rz_-^2}$.  On the negative half-line its generator is
bounded above by a constant: the quadratic exponential contribution
has coefficient $2r(r-q)R^2z^2\le0$, and the additional drift is
nonnegative, hence contributes nonpositively there.  On the positive
half-line its generator is at most $C(1+z_+^{m-1})$.  Stopped It\^o
and Jensen therefore bound its expectation over duration $T$ by
$C_z(1+T)^m$.  Taking $m=4$, $r=0.45$ supplies uniform integrability
for the smaller growth bound on $U$, permitting removal of stops at
each finite duration.  The representation thus yields
\[
 U(v,z)\le e^{-\delta_0T}\E[U(v-T,Z_T)]
       +q\sup R\int_0^T e^{-\delta_0s}\,\dd s.
\]
The first term tends to zero polynomially times an exponential.  We
obtain $U\le q\sup R/\delta_0$, a uniform positive lower bound for
$R$.

Now compare $R$ to the spatially constant solution of
$r'=r(\delta_0-qr^2)$ starting below that positive lower bound.
The coefficients in \eqref{eq:src42} are bounded in diffusion and
at most linear in drift, so the whole-line comparison stated above
applies.  Let its starting time tend to $-\infty$; this gives
$R\ge\sqrt{\delta_0/q}$.  The quotient of the upper and lower
constants in \eqref{eq:src43} is less than $3.8$.
\end{proof}

\subsection{Quantitative curvature and density barriers}

Suppose that $a\le A(v)\le b$ for every $v\in\R$, where $[a,b]$
is one of the intervals used in the coefficient reductions or the
tight root enclosure below.  Their finite checks are verified in
Section~\ref{sec:certificate}.  Put $g=b-a$ and define
\[
 e(z)=\frac{0.13}{\bigl(1+(2-z)_+^2/2\bigr)^{1/4}}.
\]
The following inequalities hold at equal slope values:
\begin{equation}
 V_i(k)\le W_a(k)\bigl(1+ge(z_a(k))\bigr),\qquad
 V_i(k)\ge W_b(k)\bigl(1-1.1ge(z_b(k))\bigr).
 \label{eq:src44}
\end{equation}
We spell out both the finite tests and their implication for an entire
solution, since simply testing a stationary barrier would not by
itself prove an entire comparison.

For either endpoint temporarily call the reference $a$, and let
$R_0=1+fge$, with $f=1$ at the lower reference and $f=-1.1$ at the
upper reference.  Set
$C_0=e''/2+(-B+at_k/W)e'$.  The sufficient inequalities are
\begin{equation}
 \begin{split}
 R_0^2C_0+A\frac{t_k}{W}e'
  +R_0\left[-(R_0+1)D_ae+\frac{(T_a)_+}{|f|}\right]&<0,\\
 fgC_0-R_0D_a&<-0.02.
 \end{split}
 \label{eq:src45}
\end{equation}
The first line bounds the right side of \eqref{eq:src42} after
division by $fg$; for $f<0$ the reversal gives the required
subsolution inequality.  On replacing $R_0$ by $CR_0$, the right side
divided by $C$ changes by
$(C^2-1)R_0^2(fgC_0-R_0D_a)$.
Start with $C>1$ for the upper barrier or $0<C<1$ for the lower
barrier, chosen to dominate the bounded ratios of
Lemma~\ref{lem:contr-heat-ratio}.  For a sufficiently small fixed
$c_0>0$, the solution $C'=c_0C(1-C^2)$ preserves the corresponding
super- or subsolution inequality.  Entire comparison and passage of
the starting time to $-\infty$ give $C\to1$ and prove
\eqref{eq:src44}.  The products of drift differences with barrier
derivatives are bounded, including $Be'$.  To see explicitly why
whole-line comparison is legitimate, write
\[
 \beta=-B+a\frac{t_k}{W},\qquad
 F(r)=r[D_a(1-r^2)+(A-a)T_a].
\]
For the difference between an input $R$ and a trial $S=CR_0$,
the diffusion and drift coefficients are $R^2/2$ and
$R^2\beta+At_k/W$.  The zeroth-order coefficient is
\[
 \tfrac12(R+S)S_{zz}+(R+S)\beta S_z
       +\frac{F(R)-F(S)}{R-S},
\]
where the quotient is interpreted as $F'(S)$ at equality.
Both $R$ and $S$ stay in fixed positive bounded intervals;
$D_a,T_a,t_k/W$ are bounded, $S_{zz}$ is bounded, and
$\beta S_z$ is bounded because $Be'$ is bounded.  Hence this
coefficient is bounded above uniformly, the diffusion is bounded,
and the drift grows at most linearly.  Multiplying the difference
by a decaying time exponential and using the polynomial localization
already stated proves comparison.  The barrier is $C^1$ and
piecewise $C^2$ at $z=2$; comparison applies on both sides, or follows
by smoothing inside the strict margins.

The central inequalities \eqref{eq:src45} are the test
\texttt{checkV}.  The following tail estimates complete their
whole-line verification and also supply the heat bound used above.
For $z=x\ge6$, \eqref{eq:src40} gives
\[
 \frac{t_k}{W}\le\frac1x,\qquad -2k\le T_a\le\frac1{x^2};
\]
$e$ is constant there, so both inequalities in
\eqref{eq:src45} follow directly.  For $z=-x$,
$x\ge\sqrt{13/p}$, put
$U_0=(4p^2x^2)^{-1}\le(52p)^{-1}$.  Write
$W=2pMl(1+y)$.  Then
\[
 0\le y\le U_0,\qquad k\ge1-3\cdot10^{-6},\qquad
 \frac{x}{M}\ge1-3\cdot10^{-6},
\]
and, more sharply,
$l\le e^{-13}<2.261\cdot10^{-6}$ and
$X/x\le e^{-13}/26<8.694\cdot10^{-8}$.
Thus $(kx/M)^2>1-4.7\cdot10^{-6}>1-10^{-5}$, and hence
\[
 -1-U_0+\frac{2a(1-10^{-5})U_0}{(1+U_0)^2}
 \le T_a\le1-2k+2aU_0<0.
\]
In particular $D_a>0.56p$: maximize the quadratic in $a$ in
$a[1-2a(1-10^{-5})/(1+U_0)^2]<0.185$, using $U_0<0.214$.
Also
\[
 \frac{t_k}{W}\le\frac1{2px},\quad
 0\le-B\le px,\quad
 0\le\frac{e'}e\le\frac1{2(x+2)},\quad
 \frac{|e''|}e\le\frac1{(x+2)^2}.
\]
In the first line of \eqref{eq:src45}, divided by $e$, the positive
terms are at most
\[
 1.12\left[\frac p2+\frac{0.5p}{13}+\frac{0.41}{52}\right]
       +\frac{0.41}{52}<0.8p,
\]
whereas the negative term has magnitude at least
$0.94\cdot1.94\cdot0.56p>p$.
For the second line,
$|fgC_0|<0.059(0.5+0.039+0.088)p$, while
$R_0D_a>0.94\cdot0.56p$; since $p\ge0.09$, its value is below
$-0.02$.  At $a=0$ these tail bounds and the central check give
$-1.12\le T_0\le0.13$.

For density comparison use equal spatial positions and the trial
$\widehat Q=Q_a+a\delta$, where $\delta=M_i-M_a$.  Substitution into
\eqref{eq:src37}, using the difference equation for $\delta$ with
time reversed from $v$ to $u$, gives the exact residual
\begin{equation}
 H=a(1-a)\delta
 +(A-a)\left[M_i(Q_a'-al_a)-k_iQ_a+V_i-ak_i\delta\right].
 \label{eq:src46}
\end{equation}
In particular the second spatial derivatives of $\delta$ cancel.
The difference $\widehat Q-Q_i$ solves a linear equation with drift
$Q_i-B_i$, source $H$, and damping
\[
 D_Q=q-Q_a'-ak_a-(A-a)k_i.
\]
At either endpoint the certificate proves $D_Q>0.06$ and that the
square bracket in \eqref{eq:src46} is at least $-0.12D_Q$.
Together with \eqref{eq:src41} and entire comparison in $u$, these
inequalities give
\begin{equation}
 Q_b-b(M_b-M_a)-0.12g
 \le Q_i\le
 Q_a+a(M_b-M_a)+0.12g.
 \label{eq:src47}
\end{equation}
For example at the lower reference $\delta\ge0$ and
$H\ge-0.12gD_Q$, so adding $0.12g$ to $\widehat Q$ gives a
supersolution.  At the upper reference $\delta\le0$ and $A-b\le0$,
so subtracting $0.12g$ gives a subsolution.

On $|z|\le12$ the test is \texttt{checkH}.  It uses
\eqref{eq:src41}, $V_i\ge(9/7)k_il_i^{3/2}$, and the possible
same-position slope interval obtained by inverting
\eqref{eq:src44}; the inversion is described below.  The term
$-ak_i\delta$ is discarded only at the upper reference, where it is
nonnegative.  The damping uses the worst possible positive part of
$A-a$.
For $x\ge12$, the reference tail certification gives, with
$S=\partial_x\log P_a(\pm x)$,
\[
 |S|<0.071,\qquad |xS'+S|<0.0055.
\]
On the positive side $M_i\le M_b$ is Gaussian-small, and so is $k_i$
by a secant from $z-1$; also $Q_a\le0$ and $|Q_a'|<0.0064$.
On the negative side
\[
 M_i=x+X_i,\quad X_i\le X_b<1.2\cdot10^{-6},\quad
 l_i\le xX_b(-x+1/x)<2\cdot10^{-5}.
\]
There $Q_a=-S$, $Q_a'=S'$, and the leading terms of the bracket are
$xS'+S$.  The adverse remainder terms
$X_iS'$, $-l_iS$, $-aM_il_a$, and $-ak_i\delta$ have total absolute
size below $0.0001$, using \eqref{eq:src40} and decrease of
$xe^{-px^2}$ in this range.  Since $D_Q\ge0.09-0.0064$, these
estimates prove the required tail inequalities.

\subsection{Inverse-coordinate functionals and their enclosures}
\label{subsec:contr-functional}

The comparisons now give bounds for the fields at every entire time.
We use them to evaluate a signed integral whose sign determines whether
$A$ lies above or below a tested value.  Curvature is compared at equal
slope, so we first express the other fields in that same coordinate.

The primitive reconstruction formulas are
\begin{equation}
 M(k)=\int_0^k\frac{j}{V(j)}\,\dd j,\qquad
 X(k)=\int_k^1\frac{1-j}{V(j)}\,\dd j,\qquad z(k)=X(k)-M(k).
 \label{eq:src50}
\end{equation}
Indeed $z_k=-1/V$, $M_k=k/V$, and $X_k=-(1-k)/V$; the zero
constants at the appropriate tails follow from the value barriers.
These formulas fix the spatial origin, so no translation parameter
is suppressed when comparing curvatures at equal $k$.

For a tested coefficient $\xi$ and a nonnegative number $c_+$ define
\[
 N_\xi(V,P)=\int_0^1\frac{P}{V}
       \bigl[V^2+c_+(MX-t_k)-2\xi kt_k\bigr]\,\dd k.
\]
The critical identities have the following homogeneous form:
\begin{equation}\label{eq:contr-homogeneous}
 \int_{\R}(MX-kl)P\,\dd z=0,\qquad
 \int_{\R}V^2P\,\dd z=2A\int_{\R}k^2lP\,\dd z.
\end{equation}
They hold for each critical entire limit by \eqref{eq:src36}, and for
each stationary reference at a root of $\lambda(a)=a$.  For every
input satisfying these identities, change of variables gives exactly
\[
 N_\xi(V_i,P_i)
   =2(A-\xi)\int_0^1\frac{P_i}{V_i}kt_k\,\dd k.
\]
The value comparisons give Gaussian tails for $MX$ and $kl$, and
bounded $Q_i$ gives at most exponential growth of $P_i$.  The integrals
are therefore finite; the one on the right is strictly positive.
Thus a sign enclosure for
$N_\xi$ bounds $A$ at every time; no assumption of stationarity of
the input is involved.

The remaining construction turns the value, curvature and log-slope
comparisons into enclosures of this functional.  All comparison bounds
apply to inputs before the identities \eqref{eq:contr-homogeneous} are
imposed; only the conversion of a functional sign into a coefficient
bound uses those identities.

Here we specify the mathematical meaning of the interval operations
used in the certificate.  Its middle grid covers
$[\kappa,1-\kappa]$, where $\kappa=2^{-38}$, and contains $1/2$ as
a cell boundary.  A \emph{box} is an interval containing every value
of a function throughout one closed grid cell.  A sum of boxes times
cell widths encloses the corresponding integral.
The operation \texttt{partialsum} bounds a primitive ending anywhere
in a cell: from the inclusive full-cell sum it subtracts that cell's
box multiplied by the interval from zero to its width.  The operation
\texttt{anchored} forms the analogous signed primitive from $1/2$.
Thus endpoints and partial cells are included, not replaced by point
sampling.

The object \texttt{Env} applies \eqref{eq:src44} with
\[
 r_h=1+ge_a,\quad r_l=1-1.1ge_b,\qquad
 V_h=r_hW_a,\quad V_l=r_lW_b.
\]
To bound $M$ and $X$ below, integrate $1/V_h$ in
\eqref{eq:src50}, writing it as the known $a$-reference primitive
plus the correction $(1/r_h-1)/W_a$.  Use the $b$-reference and
$1/r_l$ for upper bounds.  Since $g\le0.41$ and $0\le e\le0.13$,
\[
 |1/r_h-1|\le0.16g,\qquad |1/r_l-1|\le0.16g.
\]
Hence the omitted correction in each primitive tail costs at most
$0.16g$ times the corresponding reference tail primitive.  The
coarser direct factor bound in \texttt{adjusted} is a second valid
enclosure, and intersection improves the result.

Subtracting the primitive bounds encloses $z_i(k)$.  A spatial range
query in \eqref{eq:src47}, with this position box as input, then
bounds $Q_i(z_i(k))$.  Together with \eqref{eq:src44} and $Q_i\le0$
it yields cellwise bounds $s_l,s_h$ for
\[
 \frac{\dd}{\dd k}\log P_i(k)
       =-\frac{Q_i(z_i(k))}{V_i(k)}.
\]
Except in the final local comparison, normalize $P_i(1/2)=1$ in
inverse coordinates.  The primitive of these slope bounds is the
\texttt{logp} enclosure.  This normalization is made separately
at each fixed time, after the evolution comparisons.  It preserves
$Q_i$ and the homogeneous identities
$\int(MX-kl)P_i=0$ and
$\int V_i^2P_i=2A\int k_i^2l_iP_i$.
No evolution equation for this renormalized density is used in
the functional tests.

Conversely, the test of \eqref{eq:src46} needs possible slope values
at a given spatial position.  The lower position bounds are replaced
by their prefix minima and the upper bounds by their suffix maxima,
which preserves enclosure and makes inversion monotone.  Invert the
resulting position intervals against the given spatial cell.  At a
clipped endpoint include the corresponding entire $k$ tail.  The
checked edge positions straddle the central region: the positive edge
is above $4$, and the negative edge below $-5.5$.  A spatial cell
beyond an edge uses that extreme slope cell and its tail.  This is
\texttt{possiblek}.  Thus the density comparison
has no circular dependence on an assumed bound for $Q_i$ in its own
residual test: \texttt{checkH} uses only $M$ and $V$ enclosures to
bound its input slopes.

Several uniform tail bounds will be used in the integral tests.
Equation~\eqref{eq:src47} gives $Q_i\ge-0.4$ for $z\le-5.5$:
the central reference table checks $Q_b>-0.28$ and $X_b<0.03$
there, and the reference $S$ estimates continue the bounds outside
the table.  Also \texttt{logp} verifies $P_i(1-\kappa)\le e^5$.
The two curvature factors lie in $[0.94,1.054]$.
For either reference let $x_c$ be the absolute position of its grid
edge in the tail under consideration.  The checked bounds give
$x_c<20$, $x_c>6$ on the positive side, and $2px_c>2$ on the
negative side.  Integration of the reference slope bounds gives
\[
 \begin{array}{ll}
 k_a(x)\le\kappa e^{-q(x^2-x_c^2)},&x\ge x_c\quad\text{(positive tail)},\\
 l_a(-x)\le\kappa e^{-p(x^2-x_c^2)},&x\ge x_c\quad\text{(negative tail)}.
 \end{array}
\]
The same statements hold at $b$.  In the negative tails
$W\le(x+1)l$, by \eqref{eq:src40}.  On the positive tails the bound
$W\le(x+2)k$ suffices.

\subsection{From signed integral tests to a narrow interval}

First take $c_+=0$.  To prove an upper bound, use the monotonicity in
$V>0$ of $V-2\xi kt_k/V$ and replace $V_i$ by $V_h$.  Define $p_h$
on the middle grid by its piecewise constant logarithmic slope
$s_l$, anchored by $p_h(1/2)=1$.  Then $P_i/p_h$ is nondecreasing.
The upper \texttt{coarse} test verifies that the right cumulatives
of
\[
 f_h=p_h\left(V_h-2\xi\frac{kt_k}{V_h}\right)
\]
are negative, with margin $10^{-7}$ until the far region $k>0.95$,
where the integrand itself is checked negative.  If
$F_h(k)=\int_k^{1-\kappa}f_h$, integration by parts gives
\[
 \int_\kappa^{1-\kappa}\frac{P_i}{p_h}f_h
 =\left(\frac{P_i}{p_h}\right)(\kappa)F_h(\kappa)
  +\int_\kappa^{1-\kappa}F_h\,\dd(P_i/p_h)
 \le-10^{-7}\left(\frac{P_i}{p_h}\right)(\kappa).
\]
In the negative spatial tail the actual integrand remains negative:
\[
 \frac{V_i^2}{k^2l}
 \le1.054^2\frac{(x+1)^2l_a}{k^2},
\]
which is Gaussian-small from the edge onward.  In the positive
spatial tail $P_i$ is increasing in $k$, $p_h(\kappa)\le1$, and
the only harmful term has
$\int_0^\kappa P_iV_i\,\dd k
 <\kappa(P_i/p_h)(\kappa)$ by \eqref{eq:src32}.  The strict margin
therefore survives both tails.

For a lower bound use $V_l$ and $p_l$ with logarithmic slope $s_h$.
Then $P_i/p_l$ is nonincreasing.  The lower \texttt{coarse} test
checks positive left cumulatives with the same margin, except in
$k<0.05$, where the integrand itself is positive.  Integration by
parts with left cumulatives gives a lower bound
$10^{-7}(P_i/p_l)(1-\kappa)$ on the middle integral.
Beyond the positive edge actual positivity follows from
$V_l\ge0.94xk$ in $b$-reference position.  On the negative side the
possible negative contribution is at most
\[
 \left(\frac{P_i}{p_l}\right)(1-\kappa)
 \frac{2\xi e^5}{0.94}
 \int_{x_c}^\infty l_b(-x)e^{0.43(x-x_c)}\,\dd x
 <10^{-7}\left(\frac{P_i}{p_l}\right)(1-\kappa).
\]
Here $p_l(1-\kappa)\le e^5$, and
$x_i-x_{i,c}\le(x_b-x_{b,c})/0.94$, by the lower curvature bound.
The Gaussian rate $2px_c>2$ proves the last inequality (even the
bound $\kappa/(2-0.43)$ for the integral suffices).
This establishes the analytic validity of every \texttt{coarse}
sign test.

The sharper test uses $c_+=0.3$.  Fix the reference $r$ with
$a_r=0.293650$, which lies in every interval to which this test is
applied.  At fixed inverse coordinate $k$, set
\[
 \epsilon=\log(V_i/V_r),\qquad
 \pi=\log(P_i/P_r),\qquad
 v_\sigma=V_re^{\sigma\epsilon},\quad
 p_\sigma=P_re^{\sigma\pi},\quad 0\le\sigma\le1.
\]
This is interpolation in the functional only.  Define
$M_\sigma,X_\sigma$ by \eqref{eq:src50}; no interpolated PDE is
claimed or needed.  Both endpoints, and hence their geometric
interpolants and interpolated logarithmic density slopes, lie in
\texttt{Env}'s envelopes.  It bounds $\epsilon$ and $\pi'$ by
\texttt{eps} and \texttt{pis}, and bounds $\pi$ by the anchored
primitive of \texttt{pis}.  The interpolated primitives are obtained
by correcting the $r$ primitives.  On their omitted tails
$|e^{-\sigma\epsilon}-1|<3$ by
Lemma~\ref{lem:contr-heat-ratio}; the code uses the harmless larger
constant $3.01$.  Intersecting with the static envelopes remains
valid.

Let $v,p,M,X$ denote these interpolated profiles, and put
$S=c_+(MX-t_k)-2\xi kt_k$.  Differentiating the middle-grid
contribution gives $\int(g_p\pi+g_v\epsilon)$, up to the explicitly
bounded primitive-tail terms, where
\begin{equation}
 \begin{split}
 g_p&=\frac pv(v^2+S),\\
 g_v&=\frac pv(v^2-S)
 -c_+\frac{k}{v}\int_k^{1-\kappa}\frac pvX\,\dd j
 -c_+\frac{l}{v}\int_\kappa^k\frac pvM\,\dd j.
 \end{split}
 \label{eq:src51}
\end{equation}
The integration variable in the inner integrals is $j$, with every
profile there evaluated at $j$.  To verify the formula, differentiate
\eqref{eq:src50}:
\[
 \partial_\sigma M(k)=-\int_0^k\frac{j\epsilon(j)}{v(j)}\,\dd j,
 \qquad
 \partial_\sigma X(k)=-\int_k^1\frac{(1-j)\epsilon(j)}{v(j)}\,\dd j.
\]
The local factors $p/v$ and $v^2$ give the first terms in
\eqref{eq:src51}; Fubini for the two primitive derivatives gives its
last terms.  The omitted Fubini regions are bounded in absolute value
by
\[
 \mathcal E_{\rm prim}
 =6M_r(\kappa)\int_\kappa^{1-\kappa}c_+\frac pvX\,\dd k
 +6X_r(1-\kappa)\int_\kappa^{1-\kappa}c_+\frac pvM\,\dd k.
\]
Indeed $|\epsilon|V_r/v\le3.8\log3.8<6$.  This is
\texttt{terror}, checked below $3\cdot10^{-7}$ uniformly on each
interpolation-parameter box.

Since $\pi(1/2)=0$, the density term can instead be expressed against
$\pi'$: its coefficient is
$-\int_\kappa^k g_p$ for $k<1/2$ and
$\int_k^{1-\kappa}g_p$ for $k>1/2$.  This is just Fubini applied to
the anchored primitive; there is no endpoint term.  The code
intersects this enclosure with the direct enclosure for
$\int g_p\pi$.  It integrates the resulting derivative enclosures
in $\sigma$ and adds them to the directly enclosed reference middle
integral, called \texttt{base}.

To recover a sign for the full functional, add an absolute bound for
the actual tails.  On the positive side $P_i\le1$, $V_i\le1$;
moreover
\[
 \frac{\dd k}{V_i}\le\frac{\dd x_b}{0.94},\qquad
 M_iX_i\le\frac{M_bX_b}{0.94^2}\le1.2k.
\]
On the negative side the corresponding estimates are
$M_iX_i\le7l$ and
$P_i\le e^5e^{0.43(x_b-x_{b,c})}$.
For the term $P_iV_i\,\dd k$, use the upper reference instead:
\[
 P_iV_i\,\dd k
 \le1.054e^5(x_a+1)^2l_a^2
          e^{1.52(x_a-x_{a,c})}\,\dd x_a.
\]
The exponential factor uses the position-change bound $3.8$ from
\eqref{eq:src43} and $3.8\cdot0.4=1.52$.
Integrating with the tail rates above, and $2\xi<1.1$, bounds the
total actual tail by $3\cdot10^{-7}$.  In fact a much smaller
explicit allowance suffices.  With $s=x-x_c$ and $e^5<149$, the
positive-tail contribution is at most
\[
 \kappa+\frac{1.76}{0.94}\frac\kappa2<7.05\cdot10^{-12}.
\]
Here $|c_+(MX-t_k)-2\xi kt_k|\le1.76k$ and the tail rate is
at least $2$.  The negative-tail contribution excluding $P_iV_i$
is at most
\[
 \frac{3.5\cdot149}{0.94}\frac\kappa{2-0.43}<1.30\cdot10^{-9},
\]
since the corresponding numerator is at most $3.5l$.
Its remaining $P_iV_i$ contribution is at most
\[
 1.054\cdot149\,\kappa^2
 \int_0^\infty(21+s)^2e^{-2.48s}\,\dd s<4\cdot10^{-19}.
\]
These bounds follow from $x_c<20$ and the displayed Gaussian rates;
their sum is below $1.31\cdot10^{-9}$.
Thus the allowance $10^{-6}$
in \texttt{nonlineartest} covers both the primitive-tail error and
the actual tails with slack.  This proves the analytic validity of
that sign test as well.

The certificate applies the following interval reductions.  Entries
are the resulting allowed intervals for $A$, in units of $10^{-3}$:
\[
\begin{array}{c|l}
 \text{test}&\text{successive intervals}\\ \hline
 \texttt{coarse}&
 [167,410], [183,410], [185,400], [193,397], [197,390],\\
 & [203,386], [207,381], [211,377], [215,373], [218,369], [220,370],\\[2pt]
 \texttt{nonlineartest}&
 [228,370], [231,368], [234,363], [239,358], [244,350],\\
 & [251,341], [258,332], [266,324], [272,316], [277,311],\\
 & [281,307], [283,304], [285,303].
\end{array}
\]
The initial interval is $[0,0.410]$.  Widening the upper endpoint
from $0.369$ to $0.370$ before the last coarse step is harmless;
only the improved lower endpoint requires a sign test there.
At each step, the input satisfies \eqref{eq:contr-homogeneous} and
its coefficient lies in the incoming interval for every entire time.
The comparisons and sign test therefore apply at every time, giving
the resulting interval on the whole time axis.  That whole-time
enclosure is the hypothesis for the next round of comparisons.
The resulting conclusion is
\[
 0.285\le A(v)\le0.303\quad\text{for all entire times }v,\qquad
 \sup_v|A(v)-a_*|<0.01.
\]

\subsection{A local barrier at the root}

The signed tests have placed every input satisfying
\eqref{eq:contr-homogeneous} within distance $0.01$ of the root.  To
turn this interval bound into equality, we need comparison errors that
are proportional to the remaining deviation.  Put
$d=\sup_v|A(v)-a_*|\le0.01$.  At the root reference,
\begin{equation}
 1-dL(k)\le \frac{V_i(k)}{W_*(k)}\le1+dL(k),\qquad
 L(k)=0.26+1.30k-1.86k^2+3.25k^3.
 \label{eq:src48}
\end{equation}
For $d>0$ let $C_L=W^2L_{kk}/2+WW_kL_k$.  The two required
stationary inequalities, for $R_0=1\pm dL$, are
\[
 R_0^2C_L-At_kL_k+R_0[-(R_0+1)D_aL+|T_a|]<0,
 \qquad \pm dC_L-R_0D_a<0.
\]
Scaling and entire comparison work exactly as in
\eqref{eq:src45}.  In the linearization just given, the local trial
$S=C(1\pm dL(k(z)))$ has
$S_z=\mp CdWL_k$ and
$S_{zz}=\pm Cd(W^2L_{kk}-W_zL_k)$.
The Gaussian derivative bounds make $S_{zz}$ and $\beta S_z$
bounded here too.  The test \texttt{localcheck} verifies these
inequalities on its central $k$ grid, enclosing the root by the
reference interval $[0.293650,0.293655]$ and
\eqref{eq:src41}, \eqref{eq:src44}, \eqref{eq:src47}.
In particular, the root is an admissible stationary input for
\eqref{eq:src47}; this encloses its unknown $Q_*$ between the
two tabulated endpoint expressions without using rigidity.
The omitted positive and negative tails have respectively $z>7$ and
$x=-z>10$, as checked by those enclosures; both root edge positions
also have absolute value below $20$.  The reference Gaussian tail
bounds stated above therefore apply to the root.  On the positive side the
preceding bound for $T_a$ applies and $|C_L|<0.001$, using
$W\le(x+2)e^{-qx^2}$ and
$|WW_k|\le2|B|W+2at_k$.
On the negative side $U_0<0.059$, so
\[
 T_a\in[-1-U_0,0],\qquad D_a\ge p-aU_0>0.1889,\qquad L>2.94.
\]
Again $|C_L|<0.001$, now using $W\le(x+1)e^{-px^2}$.
In both tails $t_k|L_k|<0.001$.  These bounds give the displayed
strict inequalities, since $0<L<3$ and
$1.97\cdot0.1889\cdot2.94>1.09$.

At the same spatial position one also has
\begin{equation}
 |M_i-M_*|\le0.65d,\qquad |Q_i-Q_*|\le2d.
 \label{eq:src49}
\end{equation}
For the first, the difference equation used in
\eqref{eq:src41} is controlled by the certified bound
\[
 \sup_z\frac{M_*l_*}{q-(a_*+0.01)l_*}<0.65.
\]
The finite test covers the inverse middle grid.  Outside that grid,
$z>7$ or $x=-z>10$, while the denominator exceeds $0.196$.
On the positive tail $M_*l_*\le e^{-z^2/2}/z$; on the negative tail
$M_*l_*\le(x+1)e^{-0.2063465x^2}$, by \eqref{eq:src40} and $X_*<1$.
These decreasing envelopes make the ratio less than $10^{-6}$ on
both omitted tails.  Thus the displayed supremum bound holds on the
whole line, and the constants $\pm0.65d$ are barriers.  For the second, use
\eqref{eq:src46} in absolute value.  The central test checks positive
damping and, cell by cell,
\[
 a_*\,0.65+
 \frac{a_*(1-a_*)\,0.65+\sup|\text{bracket in \eqref{eq:src46}}|}
      {D_{Q,\min}}<2.
\]
It obtains the possible $k_i$ range by the outward inversion described
above, using the larger position boxes from \eqref{eq:src48} and
\eqref{eq:src50}; at a
clipped tail cell the extra $0.0001$ in its upper curvature bound is
justified by \eqref{eq:src32}: outside the inverse grid,
$t_k\le\kappa=2^{-38}$, so
$W_*\le(6/5)\kappa^{3/4}<3.162\cdot10^{-9}$.
Equation~\eqref{eq:src48} increases this by a factor at most $1.03$;
the overlapping grid cell is already enclosed by the range query.
For the spatial tails, $a_*+0.01<0.31$ gives $M_i\le M_{0.31}$
by \eqref{eq:src41}, and also $X_i\le X_{0.31}$.
Convex secants and \eqref{eq:src40} give, for $x\ge12$,
\[
 \begin{split}
 k_i(x)&\le M_{0.31}(x-1)
       \le\frac{e^{-(x-1)^2/2}}{x-1}<4.829\cdot10^{-28},\\
 l_i(-x)&\le xX_{0.31}(-x+1/x)
       \le\frac{x e^{-0.19(x-1/x)^2}}{0.38(x-1/x)}
       <5.075\cdot10^{-12}.
 \end{split}
\]
The bounding functions decrease on this range.  Substitution in
\eqref{eq:src32}, rather than the secants alone, yields
$V_i<4.059\cdot10^{-9}<0.0001$ on both tails.
The same $S$ bounds
and cancellation used above give bracket absolute value below
$0.08$ and $D_Q>0.18$; with $|\delta|\le0.65d$ these imply the
second inequality of \eqref{eq:src49} there also.

\subsection{Strict contraction and rigidity}

\begin{proposition}[Rigidity of the logarithmic-time limit]
\label{prop:contr-rigidity}
Let $(M_i,P_i,A)$ be a density input as defined in
Subsection~\ref{subsec:contr-comparisons}, and suppose that the two
identities \eqref{eq:contr-homogeneous} hold at every entire time.
Then $A\equiv a_*$.  In particular every entire limit of the critical
rescaled profiles has this coefficient, and
\[
 \lim_{\tau\downarrow0}\frac{\tau\gamma(\tau)}{\chi(\tau)}=a_*,
 \qquad
 \lim_{\tau\downarrow0}\frac{\log\chi(\tau)}{\log\tau}=a_*.
\]
The root of $\lambda(a)=a$ in the certified interval is unique.
\end{proposition}

\begin{proof}
Let $d=\sup|A-a_*|<0.01$.  If $d=0$ there is nothing to prove;
assume $d>0$.  Interpolate from the root to the input as above, with
$c_+=0.3$ and $\xi=a_*$.  Normalize the two densities at the \emph{same
spatial position} $z_c=z_*(1/2)$ to have value one.  This normalization
preserves both identities \eqref{eq:contr-homogeneous}.  Unlike the previous normalization,
it does not require the input density at its own median slope to be
one.

Our target is $|A(v)-a_*|\le c_{\rm loc}d$ at every entire time, with
$c_{\rm loc}<1$.  We obtain it by bounding the change in the signed
functional from its zero value at the root and dividing by the positive
denominator in the identity for $N_{a_*}$.  The estimates below control
that one quotient, including its endpoint tails, by a constant times $d$.

Write $W=W_*$.  Equation~\eqref{eq:src48} implies
\[
 \frac{|\epsilon|}{d}\le E_m:=\frac{L}{1-0.01L},\qquad
 |(V_i/W)^{-\sigma}-1|\le\sigma dE_m,
 \qquad E_m<3.1.
\]
For the second inequality, if $V_i/W<1$, convexity of
$t\mapsto(V_i/W)^{-t}$ puts its value below the chord on $[0,1]$;
if $V_i/W\ge1$, use $1-e^{-x}\le x$ and the logarithmic bound.
Integrating $E_m/W$ with weights $k$ and $l$ in
\eqref{eq:src50} gives bounds $L_M,L_X$ for primitive changes
divided by $\sigma d$.  Their sum $L_z$ bounds the position change:
\[
 |z_i(k)-z_*(k)|\le dL_z(k).
\]
The code arrays \texttt{lm}, \texttt{lx}, \texttt{lz} enclose these
functions; the global bounds $3.11M_*$ and $3.11X_*$ give convenient
caps and tail allowances.

The logarithmic density difference splits as
\[
 \begin{split}
 \pi(k)={}&\log P_i(z_*(k))-\log P_*(z_*(k))\\
 &+\log P_i(z_i(k))-\log P_i(z_*(k)).
 \end{split}
\]
The first term is zero at $k=1/2$ and has $k$-derivative bounded by
$2d/W$, by \eqref{eq:src49}.  The absolute value of the second is
bounded by $d$ times \texttt{shifts}: use the position interval
$z_*\pm dL_z$ and \eqref{eq:src49} to bound $Q_i$ there by a range
query for $Q_*$ plus $2d$.  Hence
\[
 \frac{|\pi(k)|}{d}\le2|z_*(k)-z_c|+\texttt{shifts}(k).
\]
These estimates enclose the interpolated $v,p,M,X$ on each of the
four parameter subintervals used by \texttt{localcheck}.

Apply \eqref{eq:src51} to the middle contribution, divided by $d$.
The curvature term has the enclosure \texttt{cv}; the density term
has the enclosure \texttt{cp}.  For the latter, integrate by parts
only the first, anchored, part of $\pi$ just displayed, and bound the
shift part directly.  The primitive-tail error now uses $4$ instead
of $6$: indeed
\[
 \frac{|\epsilon|}{d}\frac{W}{v}
 \le\frac{3/0.97}{0.97}<4.
\]
This holds on the full open interval $(0,1)$, including the omitted
tails, because \eqref{eq:src48} was proved there analytically.
For example the omitted lower primitive derivative obeys
\[
 \frac1d\int_0^\kappa\frac{j|\epsilon(j)|}{v(j)}\,\dd j
 \le4\int_0^\kappa\frac j{W(j)}\,\dd j=4M_*(\kappa),
\]
and the upper primitive has the corresponding bound
$4X_*(1-\kappa)$.  These are precisely the two factors in the
local \texttt{terror} bound.
The certificate again checks this error below $3\cdot10^{-7}$.
It also gives a positive lower bound, called \texttt{den}, for
\[
 D_i=2\int_0^1\frac{P_i}{V_i}kt_k\,\dd k,
\]
by integrating its lower enclosure on the middle grid only.  The last
parameter box includes $\sigma=1$, so the denominator enclosure
applies to the actual input.

The full reference functional is exactly zero by the stationary
identities at $a_*$.  Therefore its tail contribution must be
subtracted from the input tail contribution before applying the
local bound.  We bound this difference by differentiating along the
interpolation, not by an error independent of $d$.  Here are explicit
bounds for that step.  Equations~\eqref{eq:src47} and
\eqref{eq:src49}, with the certified tight reference interval, give
$|Q_i|\le1.5$.  The value barriers give
$M_*+X_*\le x+1.6$ and the certificate gives $|z_c|<2$.  Thus
\[
 \frac{|\pi|}{d}\le2(x+2)+1.5\cdot3.11(x+1.6)<7x+12.
\]
With $s=x-x_c\ge0$, the interpolated density satisfies
\[
 p\le5e^5e^{0.47s}\quad\text{on the negative spatial tail},
 \qquad p\le5e^{0.07s}\quad\text{on the positive spatial tail}.
\]
At the edge the normalization estimate above, $d<0.01$, and
$x_c<20$ permit the factor $5$.  Along a tail one may use the
reference density bound together with
$p=P_*\exp(\sigma\pi)$ and the $7s$ part of the last display.
On the negative side the reference has logarithmic growth bounded
by $0.4$; the positive-side reference density is at most its edge
value.  This gives exactly the displayed exponents.

Every primitive derivative has absolute value at most $3.1d$ times
the primitive itself: the full-interval bound
$|\epsilon(j)|<3.1d$ and positivity imply
\[
 |\partial_\sigma M_\sigma(k)|
 \le3.1d\int_0^k\frac j{v_\sigma(j)}\,\dd j
 =3.1dM_\sigma(k),
\]
and likewise for $X_\sigma$.  The reference tail bounds and
$v_\sigma\ge0.97W$ dominate these differentiated integrals, so
differentiation under their endpoint integrals is justified.
Therefore the pointwise absolute functional
derivative, divided by $d$, is at most
\[
 \frac pv\left[
 \left(\frac{|\pi|}{d}+3.1\right)
    \bigl(v^2+0.3(MX+t_k)+2a_*kt_k\bigr)
       +0.3\cdot6.2MX\right].
\]
Use $\dd k/v\le\dd x/0.97$.  On the negative side,
$MX\le2.6l$ and $v^2\le1.061(x+1)^2l^2$; on the positive side,
$MX\le1.2k$ and $v^2\le1.061(x+2)^2k^2$.
For the product bound on the negative side,
$p_*:=q-a_*\ge0.2063465$ and $x\ge10$ give explicitly
\[
 \frac{MX}{l}\le
 \frac{1/(2p_*)+\kappa/(400p_*^2)}{0.97^2}<2.576<2.6;
\]
on the positive side, $x\ge7$ gives
$MX/k\le(1+\kappa/49)/0.97^2<1.063<1.2$.
The reference tail bounds give $l$ or $k$ at most
$\kappa e^{-2s}$, and $x_c<20$.  Integrating proves a total tail
error below $10^{-6}$ after division by $d$.  More explicitly, the
non-$v^2$ part on the negative side is bounded by
\[
 \frac{5e^5}{0.97}\,\kappa
 \int_0^\infty\bigl[(156+7s)1.68+4.84\bigr]e^{-1.53s}\,\dd s.
\]
The positive-side non-$v^2$ counterpart is
\[
 \frac5{0.97}\,\kappa
 \int_0^\infty[(156+7s)1.26+2.232]e^{-1.93s}\,\dd s.
\]
The two $v^2$ contributions are bounded respectively by
\[
 \begin{split}
 &\frac{5e^5}{0.97}\,1.061\kappa^2
  \int_0^\infty(156+7s)(21+s)^2e^{-3.53s}\,\dd s,\\
 &\frac5{0.97}\,1.061\kappa^2
  \int_0^\infty(156+7s)(22+s)^2e^{-3.93s}\,\dd s.
 \end{split}
\]
Using $e^5<149$ and
$\int_0^\infty s^n e^{-bs}\,\dd s=n!/b^{n+1}$ reduces their sum
to a rational expression below $5.035\cdot10^{-7}$.
This proves the stated $10^{-6}$ allowance.
Including the primitive-tail error, the $2\cdot10^{-6}$ used by
\texttt{localcheck} is sufficient.

The certified middle derivative bound plus these errors, divided by
the positive denominator enclosure, gives a number
$c_{\rm loc}<0.95$.  Hence at each entire time
\[
 |A(v)-a_*|=\frac{|N_{a_*}(V_i,P_i)|}{D_i}
 \le c_{\rm loc}\,d.
\]
Taking the supremum yields $d\le c_{\rm loc}d$, contradicting $d>0$.
Thus every such density input has $A\equiv a_*$.  In particular this
holds for every critical entire limit.  Compactness now implies
convergence of the original $A$ to $a_*$: otherwise a sequence with a
fixed nonzero deviation would produce an entire limit with that
deviation at time zero.  Finally
$\dd\log\chi/\dd\log\tau=A$; integrate between a fixed positive
$\tau_0$ and $\tau$, divide by $\log\tau$, and use the convergence of
$A$ to obtain the logarithmic susceptibility limit.

A stationary reference at any other root in the certified interval
is a density input: its value and curvature bounds were proved above,
and its positive eigenfunction has bounded nonpositive logarithmic
slope satisfying \eqref{eq:src37}.  The stationary integral identities
are exactly \eqref{eq:contr-homogeneous}.  The implication just proved
therefore forces its constant coefficient to equal $a_*$.  This proves
uniqueness in the interval without assuming that the other reference
arises as a critical translation limit.
\end{proof}


\section{The cumulative exponents}\label{sec:conclusion}

We now deduce the cumulative exponents from the selected susceptibility
slope. This completes the main theorem using the finite inequalities
invoked in Section~\ref{sec:contraction}; their verification, including
rounding and exterior error bounds, follows in Section~\ref{sec:certificate}.

\begin{proof}[Completion of Theorem~\ref{thm:main}]
Sections~\ref{sec:universality} and~\ref{sec:threshold} establish the common threshold and mixture comparison. Theorem~\ref{crit:ordered-laws} identifies the limits in the prescribed order \eqref{eq:intro-order}, followed by removal of the gap cutoff, using a nonnegative terminal gap $Y=H_1$ and a nonnegative terminal force variable $R$ on the common Brownian space. If $c=\E R$, then
\begin{equation}\label{eq:concl-observables}
 G_J(u)=\Prob(0<Y\le u),\qquad
 F_J(s)=\Prob\left(\frac{R}{\alpha_c c}\le s\ \middle|\ R>0\right).
\end{equation}
In particular $\Prob(Y=0)=\Prob(R>0)=1/\alpha_c$. The gap convention removes exactly the first of these masses; conditioning in the second formula removes the noncontact mass of $R$. The small-force cutoff estimate and strong $L^2$ convergence give a probability law with no atom at zero, and
\[
 \E\left[\frac{R}{\alpha_c c}\,\middle|\,R>0\right]
 =\frac{c}{\alpha_c c\,\Prob(R>0)}=1.
\]

We now restore physical time as the argument of $\chi$ and $w$, and
write $\tau=1-t$. Proposition~\ref{prop:contr-rigidity} gives
\[
 \lim_{\tau\downarrow0}\frac{\log\chi(1-\tau)}{\log\tau}=a_*.
\]
Equivalently $\chi(1-\tau)=\tau^{a_*+o(1)}$. The scale estimate in Section~\ref{sec:scales} states
\[
 w(1-\tau)=\int_{1-\tau}^{1}\chi(t)^{-2}\,\dd t
             \asymp\frac{\tau}{\chi(1-\tau)^2}
             =\tau^{1-2a_*+o(1)}.
\]
Together with \eqref{eq:src31}, this yields
\begin{align*}
 \Prob(0<Y\le\sqrt\tau)&=\tau^{a_*+o(1)},\\
 \Prob(0<R\le\sqrt{w(1-\tau)})&=\tau^{a_*+o(1)}.
\end{align*}
For the gap distribution set $u=\sqrt\tau$ to obtain
$G_J(u)=u^{2a_*+o(1)}$.
For the force distribution, $w(1-\tau)$ is continuous and strictly increasing in $\tau$, and tends to zero with $\tau$. Its square root therefore covers all sufficiently small positive cutoffs. Taking logarithms of the second display gives the force exponent $2a_* /(1-2a_*)$. Conditioning and the fixed mean normalization in \eqref{eq:concl-observables} change no logarithmic exponent. Thus \eqref{eq:intro-exponents} holds.

Returning to the gap-exponent notation of Theorem~\ref{thm:main},
$\gamma=1-2a_*$, the exponent relation is exact:
\[
 2+\theta=1+\frac{2a_*}{1-2a_*}
          =\frac1{1-2a_*}=\frac1\gamma.
\]
Both maps $a\mapsto1-2a$ and $a\mapsto2a/(1-2a)-1$ are strictly monotone on the certified interval, in opposite directions. The endpoint arithmetic gives
\[
 0.4126930<\gamma<0.4126934
\]
and
\[
 0.4231063544994904\ldots<\theta<0.4231087030795289\ldots,
\]
which implies the stated rational decimal bounds. Finally, the mixture comparison identifies the fixed-temperature row marginals on an entire density neighborhood above the common threshold. Each subsequent ordered limit and each exponent is therefore the same for every fixed $\eps\in[0,\eps_0)$.
\end{proof}


\section{Verification of the finite certificate}\label{sec:certificate}

The computer-assisted part of the argument has three logically distinct
components: integer certificates for the stationary ODEs, analytic estimates
that turn small residuals into errors relative to the exact profiles, and
interval bounds for the remaining integrals.  This section describes the
interface between these components.  In particular, a successful floating-point
calculation alone is not the assertion being used: every computed box must
contain the corresponding exact mathematical quantity.
We first certify the reference profiles and eigenfunction logarithmic
derivatives using integer residuals and analytic error estimates; we then
enclose the barrier and integral tests by interval arithmetic and record
the complete execution.

Throughout this section $q=1/2$, $0\le a\le .41$, $p=q-a\ge .09$, and
\[
 B(z)=qz+aM_a(z),\qquad k=-M_a',\qquad l=1-k,\qquad X=z+M_a.
\]
The existence, uniqueness, and monotonicity of $M_a$, its ODE and tail bounds
\eqref{eq:src40}, and the confining ground-state construction of $P_a$ are
analytic inputs from Section~\ref{sec:contraction}.  All decimal constants in
inequalities below denote terminating decimals, hence rational numbers.
Machine padding constants are discussed separately in
Subsection~\ref{cert:arithmetic}.

The table maps the finite checks to their uses in
Section~\ref{sec:contraction}. The enclosure proofs and recorded execution
below establish the finite inequalities; the comparison principles and
exterior estimates of that section supply their analytic consequences.
The tail estimates used with each line have been established in
Section~\ref{sec:contraction}; a finite grid by itself would not justify
any statement over an infinite spatial domain.

\begin{center}
\begin{tabular}{p{.23\textwidth}p{.68\textwidth}}
Routine & Role in the analytic argument \\ \hline
\texttt{checkA} & Positivity of the cumulative lower bound
\eqref{eq:src39}, excluding an entire limit with $A\ge .41$. \\
\texttt{checkV} & Central strict inequalities
\eqref{eq:src45}, which validate the curvature barriers
\eqref{eq:src44}. \\
\texttt{Env.checkH} & Central positive damping and source bounds for
\eqref{eq:src46}, hence the log-slope comparison
\eqref{eq:src47}, and the density bounds needed for integration tails. \\
\texttt{Env.coarse} & Signs of the left or right signed cumulatives for
$c_+=0$, with a $10^{-7}$ allowance for the adverse tail. \\
\texttt{nonlineartest} & Signs of the full calibrated functional at
$c_+=.3$, using geometric interpolation and
\eqref{eq:src51}, with a $10^{-6}$ allowance. \\
\texttt{localcheck} & The local barriers \eqref{eq:src48},
\eqref{eq:src49}, a positive denominator, and the strict contraction
coefficient, with a $2\cdot10^{-6}$ allowance. \\
\texttt{certified} & Integer residual and shooting tests, including the
opposite signs of $\lambda(a)-a$ at the final two parameter values.
\end{tabular}
\end{center}

\subsection{Integer polynomial certificates}\label{cert:integer}

Set $D=10^{36}$ and $h=1/100$.  The program uses integer coefficients $c_n$
for the polynomial
\[
 \widetilde M(z)=D^{-1}\sum_{n=0}^{18}c_ns^n,
 \qquad z=\varepsilon(j+s)h,\quad 0\le s\le1,
 \quad \varepsilon\in\{-1,1\}.
\]
The input parameter is the rational number $a=A/\mathrm{AD}$; initially
$\mathrm{AD}=10^6$, and the two final parameter tests use
$\mathrm{AD}=10^7$.  The spatial equation is
\begin{equation}\label{cert:M-ode}
 M''+2(qz+aM)M'-2pM=0.
\end{equation}
After substitution of the polynomial, multiplication by the integer common
denominator gives exactly the coefficient expression \texttt{rhs(n)} in
\texttt{mstep}.  Successive coefficients are obtained by integer division.
Their manner of construction is only a way to find a good trial: certification
uses the residual of the completed polynomial, not the accuracy of this
recurrence as an ODE solver.

Specifically, if $r_n$ are the integer residual coefficients after the
multiplication by $\mathrm{AD}D^2$, the assertion in \texttt{mstep} is
\[
 10^{23}\sum_n|r_n|<\mathrm{AD}D^2.
\]
Because $0\le s\le1$, it proves, throughout the closed cell,
\begin{equation}\label{cert:M-residual}
 |\widetilde M''+2(qz+a\widetilde M)\widetilde M'-2p\widetilde M|
 <10^{-23}.
\end{equation}
At the next cell the initial value and the derivative in $s$ are the integer
sums $\sum c_n$ and $\sum nc_n$.  Thus joins are exactly $C^1$, including the
join at zero: the two first derivative coefficients are chosen with opposite
signs.  Second derivatives may jump at joins; the comparison arguments below
apply separately on cells and use continuity of the first derivative.

The routine \texttt{mint}, with \texttt{save=True}, additionally certifies
\begin{equation}\label{cert:M-boundary}
 |\widetilde M(24)|<10^{-23},\qquad
 |\widetilde M(-24)-24|<10^{-23}.
\end{equation}
The derivative coefficient tests in \texttt{certified} give
$-1.01\le\widetilde M'\le .01$ on every cell.  Indeed, for
$k=-\widetilde M'$, the constant coefficient is the integer
$-\varepsilon100c_1$, and the sum of the absolute values of the remaining
coefficients is at most $100\sum_{n\ge2}n|c_n|$; these are precisely
\texttt{kp} and \texttt{rad}.

\begin{lemma}[Error of the stationary profile]\label{cert:M-error}
For every certified rational parameter,
\[
 \|M_a-\widetilde M\|_{L^\infty([-24,24])}<10^{-22},\qquad
 \|M_a'-\widetilde M'\|_{L^\infty([-24,24])}<4\cdot10^{-20}.
\]
\end{lemma}
\begin{proof}
Let $E=M_a-\widetilde M$ and let $r$ denote the residual in
\eqref{cert:M-residual}.  Subtraction gives
\begin{equation}\label{cert:M-difference}
 E''+2B E'-(2p-2a\widetilde M')E=-r.
\end{equation}
The zero-order damping satisfies
$2p-2a\widetilde M'\ge .18-.0082>.17$.
The true deviations from the hinge at $\pm24$ are bounded, respectively, by
$e^{-q24^2}/24$ and $e^{-p24^2}/(2p\,24)$, by
\eqref{eq:src40}.  Combining these with \eqref{cert:M-boundary}, the maximum
principle applied to the constant barriers $\pm10^{-22}$ yields the first
claim.  At an interior join a strict positive maximum of the difference is
excluded by the same one-sided argument, since its first derivative is
continuous.

On either half-line write $x=|z|$, $\beta(x)=\varepsilon B(\varepsilon x)$,
and $F=\partial_xE(\varepsilon x)$.  Equation
\eqref{cert:M-difference} reads
\[
 F'+2\beta F=f,\qquad |f|<2.1\cdot10^{-22}.
\]
On the positive half-line $\beta\ge qx$; on the negative half-line
$\beta=px-aX\ge px-.34$, using $X\le .8$.
Consequently the total magnitude of the negative part of $2\beta$ is at most
$.34^2/p<1.29$.  On $[0,.1]$ we also have $\beta\le .38$.
The integrating-factor formula for $F$, integrated once more between zero and
$.1$, and $|E(.1)-E(0)|<2\cdot10^{-22}$ now imply
$|F(0)|<2.4\cdot10^{-21}$.  Explicitly, the coefficient of $F(0)$ in this
second integration is at least $.1e^{-.076}$, whereas the source contribution
has absolute value at most
$.1^2e^{1.29}(2.1\cdot10^{-22})/2$.
Finally, on $[0,24]$ the same formula gives
\[
 |F(x)|\le e^{1.29}\bigl(2.4\cdot10^{-21}
                      +24\cdot2.1\cdot10^{-22}\bigr)
 <4\cdot10^{-20}.
\]
\end{proof}

\subsection{The decaying eigenfunction branch at infinity}\label{cert:tail}

For the polynomial-order branch of $P_a$, put
$S_\varepsilon(x)=\partial_x\log P_a(\varepsilon x)$.
The eigenvalue equation gives the Riccati equation
\begin{equation}\label{cert:Riccati}
 S'=2\beta S-S^2+2(q-ak-\lambda).
\end{equation}
Let $p_+=q$ and $p_-=p$.  For either sign set
\begin{align}
 c&=\lambda/p_\varepsilon-1,&
 d&=\frac{c(c-1)}{2p_\varepsilon},&
 e&=\frac{(2c-3)d}{2p_\varepsilon},\notag\\
 T(x)&=\frac c x+\frac d{x^3}+\frac e{x^5},&
 K&=\frac{10|e|+3d^2+30}{p_\varepsilon}.
 \label{cert:tail-coefficients}
\end{align}
The rational assertions in \texttt{tailinit}, applied at all three tested
values of $\lambda$, verify $-.71<c<.81$ and, on the negative side,
$c>-.01$.  The same bounds hold at $\lambda=a_*$ in the final root bracket.
On the positive side $p_+=1/2$ and $|c(c-1)|<1.215$; on the
negative side $|c(c-1)|\le1/4$ and $p_-\ge .09$.  These sided
bounds imply
\begin{equation}\label{cert:tail-constants}
 |d|<1.4,\qquad |e|<24,\qquad K<3066.
\end{equation}

\begin{lemma}[Riccati tail enclosure]\label{cert:tail-enclosure}
For $x\ge12$, the polynomial-order branch is uniquely determined by
\[
 T(x)-Kx^{-7}\le S_\varepsilon(x)\le T(x)+Kx^{-7}.
\]
It depends continuously on $\lambda$ within each trial window.  Moreover,
\begin{equation}\label{cert:tail-derivatives}
 |S_\varepsilon(x)|<.071,\qquad
 |xS_\varepsilon'(x)+S_\varepsilon(x)|<.0055.
\end{equation}
\end{lemma}
\begin{proof}
First replace the coefficients in \eqref{cert:Riccati} by their limiting
values $\beta=p_\varepsilon x$ and $q-ak=p_\varepsilon$.
For the residual $\mathcal R(S)=S'-2\beta S+S^2-2(q-ak-\lambda)$,
direct expansion gives
\[
 \mathcal R(T)=\frac{(2c-5)e+d^2}{x^6}
                    +\frac{2de}{x^8}+\frac{e^2}{x^{10}}.
\]
Thus $x^6|\mathcal R(T)|\le6.62|e|+d^2+1$ for $x\ge12$.
Replacing $T$ by $T\pm Kx^{-7}$ adds the leading term
$\mp2p_\varepsilon Kx^{-6}$ and a remainder whose absolute value, after
multiplication by $x^6$, is at most
\[
 \frac{|-7+2c+2d/x^2+2e/x^4|K}{x^2}+\frac{K^2}{x^8}
 \le .7p_\varepsilon K.
\]
The true coefficients differ from their limiting values by at most
\[
 |\beta-p_\varepsilon x|
 \le\frac{.41e^{-p_\varepsilon x^2}}{2p_\varepsilon x},
 \qquad |q-ak-p_\varepsilon|\le .41e^{-p_\varepsilon x^2}.
\]
Their additional residual cost, multiplied by $x^6$, is less than $8$.
For example, the largest exponential envelope is obtained with
$p_\varepsilon=.09$ at $x=12$, since
$x^6e^{-p_\varepsilon x^2}$ is decreasing for these values; inserting
$|T\pm Kx^{-7}|<.071$ proves the stated bound.
The definition of $K$ leaves strict room:
\[
 1.3p_\varepsilon K-(6.62|e|+d^2+1+8)
 =6.38|e|+2.9d^2+30>0.
\]
Accordingly the lower barrier has positive residual and the upper barrier
negative residual.  These are the required signs for integration from larger
to smaller $x$.

Start at $T(R)$ at a finite $R>12$ and integrate inward.  Comparison keeps the
solution within the barriers.  Compactness on finite intervals supplies a
limit as $R\to\infty$.  If two such limits are compared, their difference
satisfies a homogeneous first-order equation with coefficient
$2\beta-S_1-S_2\ge p_\varepsilon x$ for $x\ge12$.  Inward integration
therefore forces their difference to vanish as the comparison point tends to
infinity.  The same formula with the additional source
$-2(\lambda_1-\lambda_2)$ proves continuity in $\lambda$ on compact
subintervals.

The first estimate in \eqref{cert:tail-derivatives} follows immediately from
\eqref{cert:tail-constants}.  For the second, substitute the enclosure into
\eqref{cert:Riccati}.  The resulting upper bound is
\begin{align*}
 \frac{2|d|}{x^3}+\frac1{x^5}\bigg(&4|e|+6.62|e|+d^2+1+8
       +2p_\varepsilon K\\
       &+\frac{(1+2\cdot.82)K}{x^2}+\frac{K^2}{x^8}\bigg).
\end{align*}
To justify the constant $.82$ in this expression, observe that $d$ and $e$
have opposite signs and $|e|/(|d|x^2)=|2c-3|/(2p_\varepsilon x^2)<1$
when $d\ne0$.  Thus $xT$ has the sign of $c$ and magnitude at most
$|c|<.81$; the case $d=0$ is immediate.
The displayed bound, using \eqref{cert:tail-constants}, is less than $.00514$.
Here one must retain the identity
$p_\varepsilon K=10|e|+3d^2+30$; replacing the two factors by separate
worst-case bounds would lose the required estimate.

Finally, integrating $S=c/x+O(x^{-3})$ gives
$P(\varepsilon x)=C x^c(1+O(x^{-2}))$ with $C>0$.
After the conjugation $P=e^{\int B}\psi$, this is the square-integrable
branch of the Schr\"odinger equation.  Reduction of order gives a second
independent solution growing by the reciprocal Gaussian factor; hence no
other square-integrable branch is available at this endpoint.
\end{proof}

\subsection{Inward shooting and eigenvalue certification}\label{cert:shooting}

The shooting routine begins at $x=20$.  The value \texttt{tailinit} is an
integer multiple of $D^{-1}$ within $D^{-1}$ of the exact rational number
$T(20)$.  The routine \texttt{revpoly} expresses the stationary polynomial in
the reversed cell coordinate.  In \texttt{qcoefs}, the resulting coefficients
represent $2\widetilde\beta$ with denominator $\mathrm{AD}D$, and
$\widetilde k$ with denominator $D$.  The inward cell step is $-1/100$;
this explains the sign in \texttt{qint}'s recurrence.

For the completed degree-$18$ Riccati trial $\widetilde S$, the integer
residual test is again a sum of absolute coefficients on $0\le s\le1$.
It proves a residual smaller than $10^{-20}$ in the spatial Riccati equation
with the approximate $\widetilde\beta,\widetilde k$.  The trial is exactly
continuous between cells.  No numerical differentiation is used.

Let $E_S=S-\widetilde S$.  With $\lambda$ equal to the trial value,
\begin{equation}\label{cert:S-difference}
 E_S'=(2\beta-2\widetilde S-E_S)E_S+f_S,
 \qquad |f_S|<3\cdot10^{-19},
\end{equation}
where the source bound includes the residual and both stationary-profile
errors from Lemma~\ref{cert:M-error}.  When $\lambda$ differs from the middle
trial value by at most $10^{-9}$, add $2\cdot10^{-9}$ to this source bound.

Under the provisional condition $|E_S|\le .001$, each integer \texttt{gs}
is a lower bound for the drift in \eqref{cert:S-difference}, with denominator
$\mathrm{AD}D$.  The subtraction $.002$ in its definition covers
$|E_S|$ and the much smaller profile error.  The tests on \texttt{gs} imply:
\begin{enumerate}
\item every inward integrating factor, including factors with endpoints
inside cells, is at most $e^2$;
\item the integrating factor from $20$ to $0$ is smaller than $e^{-20}$;
\item the factor from $20$ to any $x<12.02$ is smaller than $e^{-6}$.
\end{enumerate}
For the last assertion, the cells indexed from $1203$ upward have total
integral greater than $8$, while the sum of all negative cell contributions
is greater than $-2$.  This also covers the possible partial cell.

The initial error at $20$, including variation of $T(20)$ over the eigenvalue
window, is less than $2.41\cdot10^{-6}$.  Variation of constants gives the
following deliberately loose enclosures:
\begin{equation}\label{cert:S-errors}
 |E_S|<2.9\cdot10^{-5}\quad(0\le x\le20),\qquad
 |E_S|<4.5\cdot10^{-7}\quad(0\le x<12.02).
\end{equation}
For example, the global estimate is bounded by
$e^2[2.41\cdot10^{-6}+20(2\cdot10^{-9}+3\cdot10^{-19})]$;
the central estimate replaces the first factor $e^2$ on the initial error by
$e^{-6}$.  Both are strictly below the provisional $.001$, which closes the
bootstrap and proves that the positive branch extends to the center.
For an endpoint trial compared at its own eigenvalue, the error at zero is
bounded instead by
\[
 (2.41\cdot10^{-6})e^{-20}+20e^2(3\cdot10^{-19})<10^{-10}.
\]

The matching condition is $S_-(0)+S_+(0)=0$, because derivatives in $z$
acquire opposite signs on the two half-lines.  At the lower and upper
shooting values the integer tests give, respectively, a trial sum smaller
than $-10^{-9}$ and greater than $10^{-9}$.  The total error of the two
halves is less than $2\cdot10^{-10}$, so these signs also hold for the true
shoots.  Continuity therefore supplies a match.  It produces a positive
square-integrable eigenfunction after conjugation, and consequently the
lowest eigenvalue.  We have proved
\begin{equation}\label{cert:lambda-window}
 \frac{\texttt{lam}-10}{10^{10}}
 <\lambda(A/\mathrm{AD})<
 \frac{\texttt{lam}+10}{10^{10}}
\end{equation}
for every call to \texttt{certified}.  This also justifies the log-slope
errors assigned to the middle shoot in subsequent tables.

\subsection{Floating-point enclosures}\label{cert:arithmetic}

The interval stage uses IEEE binary64 round-to-nearest arithmetic, including
square root.  Integer polynomial residual tests precede every tabulation.
Let $u=2^{-53}$ be unit roundoff.  Away from underflow each elementary
operation has relative error at most $u$; at underflow an absolute error of
one subnormal unit suffices.  The functions \texttt{down} and \texttt{up}
subtract or add
\[
 (|x|+10^{-280})2^{-43}
\]
to a rounded endpoint.  This is more than one thousand unit roundoffs in
relative scale, and its absolute component is much larger than a subnormal
unit.  Rounding in forming and applying the padding consumes less than this
allowance.  Therefore the operations of class \texttt{I}, including the
four-product rule for multiplication, reciprocal on intervals of fixed sign,
squaring, and square root, enclose the exact operations.  Assertions exclude
division by a zero-containing interval, invalid square roots, overflow, and
invalid floating-point operations.

The class constructor itself does not enlarge a literal.  Integer and dyadic
inputs are exact.  A terminating decimal that defines a mathematical barrier
or functional is consequently entered by \texttt{rat(n,d)}, rather than by
an unpadded decimal literal.  The other non-dyadic literals in the code are
allowances or coarse caps with strict slack; they may be interpreted as their
binary64 values.  In particular the caps $.1431$ and $.1301$ in
\texttt{adjusted} are strictly weaker than the analytic bounds they represent.
Integer-coefficient conversion and node evaluation require the separate
allowance below.

No library exponential or logarithm is used for interval evaluation.
For $|x|\le100$, \texttt{iexp} evaluates a degree-$20$ Taylor polynomial at
$x/256$, adds the interval remainder $[-2^{-80},2^{-80}]$, and squares eight
times.  The remainder bound follows from
$e^{.4}.4^{21}/21!<2^{-80}$.
For $1/8\le x\le8$, \texttt{ilog} takes three interval square roots and
uses
\[
 \log x=16\left(t+\frac{t^3}3+\cdots+\frac{t^{21}}{21}\right)
             +\text{remainder},
 \qquad t=\frac{x^{1/8}-1}{x^{1/8}+1}.
\]
Here $|t|<.14$ and
$16|t|^{23}/[23(1-t^2)]<2^{-60}$, the remainder interval used by the code.
All polynomial operations in these two routines are interval operations.

For summation of at most $n$ real terms, the usual induction on rounded
addition gives
\[
 \left|\operatorname{fl}\Bigl(\sum_{j=1}^n y_j\Bigr)-\sum_{j=1}^n y_j\right|
 \le \frac{nu}{1-nu}\sum_{j=1}^n|y_j|
\]
up to the negligible absolute underflow contribution.  The quadrature arrays
have $n<2\cdot10^5$, so the coefficient is less than
$3\cdot10^{-11}$.  The routine \texttt{sums} uses the much larger allowance
$10^{-9}$ times the accumulated absolute sum, plus \texttt{tiny}.  This also
covers rounding in that accumulated sum and in the final correction.
The shorter summation in \texttt{checkA} has at most $4096$ terms of absolute
value below $9$; its error after division by $4096$ is less than
$\gamma_{4096}\,9<4.1\cdot10^{-12}$.  Its remaining scalar arithmetic
costs less than $10^{-13}$, since all terms have magnitude less than $10$.
The subtraction of $2\cdot10^{-9}$ therefore covers both, including the
unpadded bound arithmetic after the summation.

The last scalar addition or division in the nonlinear and local tests is
also unpadded.  The analytic allowances retain sufficient slack to cover it.
In the recorded run all nonlinear returned values have magnitude below
$.00204$, so the final addition of the signed $10^{-6}$ changes its exact
value by less than $10^{-18}$; the two substantive errors total less than
$6\cdot10^{-7}$, leaving nearly $4\cdot10^{-7}$ unused.
For the local test, an additional diagnostic evaluation of the unchanged
function gave
\[
 O=\texttt{out.hi}=.7290900413708653,\qquad
 D_0=\texttt{den.lo}=.7928198472818054.
\]
With $O<1$ and $.7<D_0<1$, rounding the final expression
$(O+2\cdot10^{-6})/D_0$ costs less than $6\cdot10^{-16}$.
The exterior and primitive-tail estimates cost less than $1.3\cdot10^{-6}$
in the numerator, leaving $.7\cdot10^{-6}$ unused there.  Small rounding in
the positive scalar calculation of \texttt{terror} is absorbed in the same
slack.  Thus the returned local number is an upper bound for the exact
contraction coefficient, and the returned signed nonlinear numbers retain
the required mathematical signs.

\subsection{From polynomial tables to spatial and inverse boxes}\label{cert:boxes}

Each step of length $.01$ is sampled at the $50$ fractions
$0,1/50,\ldots,49/50$.  The resulting spacing is $1/5000$.
The coefficients for $X$, $k$, and $l$ are formed by exact integer operations
from those of $M$; they are evaluated independently.  Here is a quantitative
bound for the conversion and evaluation.  CPython's integer true division
computes the correctly rounded ratio $c_n/D$: it forms an integer quotient
with guard bits, records the discarded remainder, rounds ties to even, and
then converts the rounded integer and rescales exactly.\footnote{The
implementation and its rounding argument are in CPython 3.12.13,
\href{https://github.com/python/cpython/blob/v3.12.13/Objects/longobject.c}{\texttt{long\_true\_divide}, \texttt{Objects/longobject.c}}.}
Every nonzero ratio here is at least $10^{-36}$ in magnitude, so it is normal
and has relative conversion error at most $u$.  Both an ordinary sample
$j/50$ and a reversed sample $1-j/50$ have absolute node error at most $2u$,
and the exact and rounded nodes lie in $[0,1]$.  Writing
$C=\sum|c_n|/D$, the node perturbation therefore costs at most $36uC$.
The $36$ rounded operations in degree-$18$ Horner evaluation cost at most
$\gamma_{36}(1+u)C$, where $\gamma_n=nu/(1-nu)$; coefficient conversion
costs at most $uC$.  The total is less than $74uC<8.22\cdot10^{-15}C$,
and in particular smaller than $2\cdot10^{-14}C$.
The node allowance used for each of $M,X,k,l$ is
\begin{equation}\label{cert:node-error}
 3\cdot10^{-14}\sum_n|c_n|/D+3\cdot10^{-18}.
\end{equation}
It covers this evaluation error, Lemma~\ref{cert:M-error}, and the additions
and subtractions that form node bounds.  The code checks that the coefficient
sum is less than $30D$.  The separate $Q$ allowances contain the evaluation
error and \eqref{cert:S-errors}; the global node radius is
$2.95\cdot10^{-5}$ and the central refinement uses $5\cdot10^{-7}$.

Between nodes, monotonicity gives boxes for $M,X,k,l$.  For $Q=(\log P)'$,
\texttt{zcells} first bootstraps $|Q|\le1$: the trial nodes have magnitude
below $.951$, and the Riccati equation bounds $|Q'|<30$ on a fine cell under
this provisional bound.  Its length is $1/5000$, so enlargement by $.006$
is sufficient and is strictly inside the provisional bound.
Twice applying the derivative equation then sharpens the cell box.
The cap $Q\le0$ uses the analytic maximum-principle argument for the exact
eigenprofile; the cap is not an additional numerical assumption.

Inverse lookup must also include uncertainty in node values.  Since $k$
is decreasing, suffix maxima of its node lower bounds remain valid lower
bounds, and prefix minima of its upper bounds remain valid upper bounds.
The analogous operations for increasing $l$ give its monotone guaranteed
bounds.  A binary search through these outward-monotonized arrays supplies
outer nodes for the inverse position.  The code uses $k$ or $l$ according to
which is better separated from $1$; the two schemes are mathematically
equivalent.  The inverse curvature is then evaluated through the exact
identity
\[
 W=2pMl+Xk,
\]
rather than by differentiating an inverse interpolant.

The routines \texttt{Rmq.query} and \texttt{kquery} return minima and maxima
over all cells between two indices.  For a range of length $n$, take the two
blocks of length $2^{\lfloor\log_2 n\rfloor}$ abutting its endpoints.  Their
union covers the range, and their individual extrema have been computed by
successive doubling.  Thus the returned box is valid even when these blocks
overlap.  Spatial coordinate bounds and index calculations are enlarged by
$10^{-12}$ or $10^{-10}$ before multiplication and flooring; these allowances
are larger than the binary64 errors on the bounded coordinate range.
This padding also validates the scalar negative-edge test
$2px_c>2$, even though its code uses the rounded parameter
\texttt{af}.  If $x_f$ is the negation of the padded upper position
bound used in that test, then $x_f<x_c-0.99\cdot10^{-12}$.
Since $p\ge0.09$ and $x_c<20$, replacing the rounded parameter and
product by their exact values costs less than $3\cdot10^{-14}$,
whereas the position padding leaves more than $1.7\cdot10^{-13}$
in $2px_c$.  Thus a successful scalar comparison implies the exact
tail-rate inequality.

The inversion of the envelope $z_i(k)$ in \texttt{possiblek} and
\texttt{localcheck} uses the same outward monotonicity principle.  At a clipped
edge it explicitly restores the missing $k$-tail; a point outside the central
range is never silently treated as an interior grid point.

More explicitly, let $L_j,U_j$ be the outward-monotonized bounds in inverse
cell $j$.  If an admissible inverse value belongs to the queried spatial
interval $[z_-,z_+]$, necessarily $L_j\le z_+$ and $U_j\ge z_-$.  The
binary searches give the first and last cells satisfying these necessary
inequalities.  Clipping either index to the available grid and restoring
$k=0$ or $k=1$ therefore retains all admissible slopes, including those in
the omitted tails.  For the local curvature query the retained cells obey
$V_i(k)\le1.03W(k)$ by \eqref{eq:src48}.  If an admissible $k$ is in an
omitted tail, $k(1-k)\le\kappa$ and \eqref{eq:src32} gives instead
\[
 V_i(k)\le1.2\kappa^{3/4}=1.2\,2^{-28.5}<3.2\cdot10^{-9}.
\]
Consequently the code's bound $1.031\,\texttt{wi.hi}+.0001$ contains the
curvature in both cases; the clipped query cannot discard a larger tail
curvature.

\subsection{Integral enclosures and their logical use}\label{cert:integrals}

The integration grid covers $[\kappa,1-\kappa]$ with $\kappa=2^{-38}$,
and has $2048$ subdivisions per dyadic scale.  It splits at $1/2$ and all
edges and widths are exactly representable binary numbers: the finest
required denominator is $2^{49}$.  A box on a closed cell, multiplied by
its width, encloses its integral.  The routine \texttt{partialsum} removes
an arbitrary fraction of the endpoint cell by subtracting its box multiplied
by $[0,\text{width}]$.  This is valid for either sign of the integrand.
The routine \texttt{anchored} integrates separately to the left and right
of $1/2$, avoiding unnecessary cancellation.

The opening table records the analytic uses of these enclosures.

Several implementation details are important for this correspondence.
The ratio caps in \texttt{Env} apply to the exact profiles by the analytic
barriers; primitive corrections use \eqref{eq:src50}, including the omitted
primitive tail.  In \texttt{nonlineartest}, the intermediate pairs are defined
by geometric interpolation solely inside the functional.  No assertion that
they solve a parabolic equation is made or needed.  Their log slopes lie in
the same envelope because these slopes interpolate affinely.  The two forms
of the density variation---direct integration against $\pi$, and integration
by parts against its anchored derivative---each enclose the same quantity;
taking their interval intersection is therefore justified.  The quantity
\texttt{terror} bounds the missing primitive terms in that integration by
parts, rather than the full exterior integral.  These are distinct errors:
\texttt{nonlineartest}'s $10^{-6}$ contains both, and
\texttt{localcheck}'s $2\cdot10^{-6}$ contains their local analogues after
division by $d$.

These primitive-tail allowances require derivative bounds, which can be
checked directly.  Put $r=V_i/W$ and $v_\sigma=Wr^\sigma$.  Then
\[
 \partial_\sigma(v_\sigma^{-1})
       =-(\log r)v_\sigma^{-1}.
\]
In the local regime, $|\log r|/d\le L/(1-.01L)<3.1$, so differentiating
the positive integrals in \eqref{eq:src50} gives
$|\partial_\sigma M_\sigma|\le3.1dM_\sigma$ and the identical estimate
for $X_\sigma$.  Relative to the reference integrands,
$|\log r|r^{-\sigma}/d\le3.1/.97<4$, proving the primitive-tail
coefficient $4$ after division by $d$.  For the nonlinear refinement,
$1/3.8<r<3.8$ gives instead
$|\log r|r^{-\sigma}\le3.8\log3.8<6$.  This proves its coefficient $6$.
The distinct finite-difference bound
$|r^{-\sigma}-1|<2.8<3.01$ justifies its primitive correction.  These
arguments use dominated differentiation with the reference primitive as an
integrable majorant; no derivative estimate is inferred merely from a
finite-difference estimate.

The omitted endpoint cells of \texttt{checkA} admit particularly crude
bounds.  In the notation of \eqref{eq:src39}, $|b|\le1.2$, $b\ge-.8$,
$E=-.4t_k$, $t_k\le1/4$, and $K\ge-8.24$.  Since $V<.5$ and
$6A^2-.16>0$ for $.41\le A\le1$, its integral density satisfies
\[
 2V^2+KV-2bE+(6A^2-.16)t_k^2/V
 \ge-8.24\cdot.5-.16=-4.28>-9.
\]
This holds throughout the interval, hence on both omitted cells.  At
$0\le k_0\le1/4096$, the endpoint expression obeys
\[
 bV^2+2EV\le1.2V^2\le1.728t_k^{3/2}
 \le1.728/262144<6.6\cdot10^{-6}<.001.
\]
Thus its endpoint allowances also have analytic coverage.

The driver follows the successive output intervals listed with the
signed tests in Section~\ref{sec:contraction}: eleven coarse updates
from $[0,0.410]$, then thirteen nonlinear updates, give
$[0.285,0.303]$.  Each printed checkpoint is the \emph{input} interval
for the checks about to run.  In particular, the last printed interval
is $[0.283,0.304]$; successful completion of its two sign tests proves
the final interval $[0.285,0.303]$.  The slight enlargement from
$0.369$ to $0.370$ in the coarse sequence keeps a valid
coefficient enclosure.  All sign tests include the error allowances described
above.  Their analytic consequence is the input enclosure required
by the local contraction, which then forces $A\equiv a_*$.

\subsection{Reproduction and numerical output}\label{cert:execution}

The accompanying file \texttt{certificate/original\_certificate.py} is the
unaltered code printed in the appendix.  From the manuscript directory,
run
\begin{quote}\ttfamily
python3 certificate/original\_certificate.py
\end{quote}
with NumPy installed and assertions enabled; do not use \texttt{-O}.
The certificate requires IEEE binary64 arithmetic with round-to-nearest
operations as described above.  Its interval checkpoints are printed to
standard output.  The accompanying reproducibility record identifies the
software versions, source hashes, and saved output of a complete run.

The complete unchanged driver passed on October 4, 2026, under
Python 3.12.13 and NumPy 2.3.5, with assertions enabled, in
$338.52$ seconds.  It printed all $24$ checkpoints and had empty
standard error; the actual output is retained with the certificate.
The separate return-observer record from September 24, which leaves
the source, function arguments and results unchanged, gives
\[
 \texttt{checkA()}=0.002977541856187091>0,\qquad
 \texttt{localcheck()}=0.9196188060510444<.95.
\]
The displayed decimals are a returned lower bound and upper bound,
respectively.  A fresh targeted evaluation of the unchanged local routine
reproduces the second value and the numerator and denominator reported
in Subsection~\ref{cert:arithmetic}.  The complete unchanged execution
and this diagnostic evaluation are recorded separately.
Their interpretation as rigorous bounds is supplied by the error
analysis and comparison arguments above.  The exact source hash is
\begin{quote}\ttfamily\small
7bd5dbd918d9a48fe4da62044e295239b\allowbreak
74f7a1607e2643337ad4d073a85dc7f.
\end{quote}
The integer seeds for the final two spectral tests are $2936535033$ and
$2936533047$, with denominator $10^{10}$.  Formula
\eqref{cert:lambda-window} therefore yields
\begin{align*}
 .2936535023&<\lambda(.2936533)<.2936535043,\\
 .2936533037&<\lambda(.2936535)<.2936533057.
\end{align*}
In particular the first eigenvalue is strictly larger than its parameter and
the second strictly smaller.  Continuity supplies
\[
 .2936533<a_*<.2936535.
\]
The contraction argument applies to every root in this interval and proves
its uniqueness there.  Combined with the analytic identification of the
critical exponents, these rational inequalities give the exponent bounds
stated in the conclusion.


\begin{thebibliography}{99}
\bibitem{AizenmanSimsStarr2003}
M.~Aizenman, R.~Sims and S.~L.~Starr,
\emph{Extended variational principle for the Sherrington--Kirkpatrick spin-glass model},
Physical Review B \textbf{68} (2003), 214403.
\href{https://doi.org/10.1103/PhysRevB.68.214403}{doi:10.1103/PhysRevB.68.214403}.

\bibitem{AuffingerChen2015}
A.~Auffinger and W.-K.~Chen,
\emph{The Parisi formula has a unique minimizer},
Communications in Mathematical Physics \textbf{335} (2015), 1429--1444.
\href{https://doi.org/10.1007/s00220-014-2254-z}{doi:10.1007/s00220-014-2254-z}.

\bibitem{BolthausenSznitman1998}
E.~Bolthausen and A.-S.~Sznitman,
\emph{On Ruelle's probability cascades and an abstract cavity method},
Communications in Mathematical Physics \textbf{197} (1998), 247--276.
\href{https://doi.org/10.1007/s002200050450}{doi:10.1007/s002200050450}.

\bibitem{BoueDupuis1998}
M.~Bou\'e and P.~Dupuis,
\emph{A variational representation for certain functionals of Brownian motion},
The Annals of Probability \textbf{26} (1998), 1641--1659.
\href{https://doi.org/10.1214/aop/1022855876}{doi:10.1214/aop/1022855876}.

\bibitem{CKPUZ2014}
P.~Charbonneau, J.~Kurchan, G.~Parisi, P.~Urbani and F.~Zamponi,
\emph{Exact theory of dense amorphous hard spheres in high dimension.
III. The full replica symmetry breaking solution},
Journal of Statistical Mechanics: Theory and Experiment (2014), P10009.
\href{https://doi.org/10.1088/1742-5468/2014/10/P10009}{doi:10.1088/1742-5468/2014/10/P10009}.
Preprint version: \href{https://arxiv.org/abs/1310.2549v5}{arXiv:1310.2549v5}.

\bibitem{Chatterjee2006Lindeberg}
S.~Chatterjee,
\emph{A generalization of the Lindeberg principle},
The Annals of Probability \textbf{34} (2006), no.~6, 2061--2076.
\href{https://doi.org/10.1214/009117906000000575}{doi:10.1214/009117906000000575}.

\bibitem{ChenIssaMourrat2025}
H.-B.~Chen, V.~Issa and J.-C.~Mourrat,
\emph{The convex structure of the Parisi formula for multi-species spin glasses},
preprint (2025).
\href{https://arxiv.org/abs/2508.06397}{arXiv:2508.06397}.

\bibitem{Chen2013SphericalASS}
W.-K.~Chen,
\emph{The Aizenman--Sims--Starr scheme and Parisi formula for mixed
$p$-spin spherical models},
Electronic Journal of Probability \textbf{18} (2013), no.~94, 1--14.
\href{https://doi.org/10.1214/EJP.v18-2580}{doi:10.1214/EJP.v18-2580}.

\bibitem{FranzParisi2016}
S.~Franz and G.~Parisi,
\emph{The simplest model of jamming},
Journal of Physics A: Mathematical and Theoretical \textbf{49} (2016), 145001.
\href{https://doi.org/10.1088/1751-8113/49/14/145001}{doi:10.1088/1751-8113/49/14/145001}.
\href{https://arxiv.org/abs/1501.03397}{arXiv:1501.03397}.

\bibitem{FPSUZ2017}
S.~Franz, G.~Parisi, M.~Sevelev, P.~Urbani and F.~Zamponi,
\emph{Universality of the SAT-UNSAT (jamming) threshold in non-convex
continuous constraint satisfaction problems},
SciPost Physics \textbf{2} (2017), 019.
\href{https://doi.org/10.21468/SciPostPhys.2.3.019}{doi:10.21468/SciPostPhys.2.3.019}.

\bibitem{Gardner1988}
E.~Gardner,
\emph{The space of interactions in neural network models},
Journal of Physics A: Mathematical and General \textbf{21} (1988), 257--270.
\href{https://doi.org/10.1088/0305-4470/21/1/030}{doi:10.1088/0305-4470/21/1/030}.

\bibitem{GardnerDerrida1988}
E.~Gardner and B.~Derrida,
\emph{Optimal storage properties of neural network models},
Journal of Physics A: Mathematical and General \textbf{21} (1988), 271--284.
\href{https://doi.org/10.1088/0305-4470/21/1/031}{doi:10.1088/0305-4470/21/1/031}.

\bibitem{Guerra2003}
F.~Guerra,
\emph{Broken replica symmetry bounds in the mean field spin glass model},
Communications in Mathematical Physics \textbf{233} (2003), 1--12.
\href{https://doi.org/10.1007/s00220-002-0773-5}{doi:10.1007/s00220-002-0773-5}.

\bibitem{LernerDuringWyart2013}
E.~Lerner, G.~D\"uring and M.~Wyart,
\emph{Low-energy non-linear excitations in sphere packings},
Soft Matter \textbf{9} (2013), 8252--8263.
\href{https://doi.org/10.1039/C3SM50515D}{doi:10.1039/C3SM50515D}.

\bibitem{MarcenkoPastur1967}
V.~A.~Marchenko and L.~A.~Pastur,
\emph{Distribution of eigenvalues for some sets of random matrices},
Mathematics of the USSR--Sbornik \textbf{1} (1967), 457--483.
\href{https://doi.org/10.1070/SM1967v001n04ABEH001994}{doi:10.1070/SM1967v001n04ABEH001994}.

\bibitem{MontanariZhongZhou2024}
A.~Montanari, Y.~Zhong and K.~Zhou,
\emph{Tractability from overparametrization: The example of the negative perceptron},
Probability Theory and Related Fields \textbf{188} (2024), 805--910.
\href{https://doi.org/10.1007/s00440-023-01248-y}{doi:10.1007/s00440-023-01248-y}.
\href{https://arxiv.org/abs/2110.15824}{arXiv:2110.15824}.

\bibitem{MourratUninverting2025}
J.-C.~Mourrat,
\emph{Un-inverting the Parisi formula},
Annales de l'Institut Henri Poincar\'e, Probabilit\'es et Statistiques
\textbf{61} (2025), no.~4, 2709--2720.
\href{https://doi.org/10.1214/24-AIHP1487}{doi:10.1214/24-AIHP1487}.
Preprint (2023),
\href{https://arxiv.org/abs/2308.10715v2}{arXiv:2308.10715v2}.

\bibitem{Panchenko2010Sphere}
D.~Panchenko,
\emph{A connection between the Ghirlanda--Guerra identities and ultrametricity},
The Annals of Probability \textbf{38} (2010), 327--347, Theorem~2(a).
\href{https://doi.org/10.1214/09-AOP484}{doi:10.1214/09-AOP484}.

\bibitem{Panchenko2010DS}
D.~Panchenko,
\emph{On the Dovbysh--Sudakov representation result},
Electronic Communications in Probability \textbf{15} (2010), 330--338.
\href{https://doi.org/10.1214/ECP.v15-1562}{doi:10.1214/ECP.v15-1562}.

\bibitem{Panchenko2013}
D.~Panchenko,
\emph{The Parisi ultrametricity conjecture},
Annals of Mathematics \textbf{177} (2013), 383--393.
\href{https://doi.org/10.4007/annals.2013.177.1.8}{doi:10.4007/annals.2013.177.1.8}.

\bibitem{ParisiZamponi2026}
G.~Parisi and F.~Zamponi,
\emph{A proof of an identity for the critical exponents of jamming},
Journal of Statistical Mechanics: Theory and Experiment (2026), 073301.
\href{https://doi.org/10.1088/1742-5468/ae7bd7}{doi:10.1088/1742-5468/ae7bd7}.
\href{https://arxiv.org/abs/2606.03300}{arXiv:2606.03300}.

\bibitem{Stojnic2013}
M.~Stojnic,
\emph{Another look at the Gardner problem},
preprint (2013).
\href{https://arxiv.org/abs/1306.3979}{arXiv:1306.3979}.

\bibitem{Wyart2012}
M.~Wyart,
\emph{Marginal stability constrains force and pair distributions at random close packing},
Physical Review Letters \textbf{109} (2012), 125502.
\href{https://doi.org/10.1103/PhysRevLett.109.125502}{doi:10.1103/PhysRevLett.109.125502}.

\end{thebibliography}


\appendix
\section{Executable certificate}\label{app:code}
The following listing is the complete certificate distributed with the manuscript. Run it with Python~3 and NumPy, with assertions enabled; see Section~\ref{sec:certificate} for the relation between the finite checks and the analytic estimates. An arrow at the start of a displayed line marks a continuation of the same source line.
\begin{lstlisting}
import math, numpy as np
D=10**36
HN=100 # step reciprocal
L=24
N=L*HN
DEG=18
AD=10**6 # a denominator
LD=10**10
# all stored coefficients are integers /D
def mstep(A, side, j, c0,c1, check=False):
 # z0=side*j/HN, h=side/HN, q=.5, a=A/AD
 c=[c0,c1]
 # recurrence common h² terms: 2(q-a-q*n)c_n -2q*(z0/h)*(n+1)c_{n+1}; nonlinear -2*a/h ...
 # denominator HN² AD D
 def rhs(n):
  val=0
  if n<len(c): val+=(AD*(1-n)-2*A)*D*c[n]
  if n+1<len(c): val-=j*(n+1)*AD*D*c[n+1]
  conv=sum(c[i]*(n-i+1)*c[n-i+1] for i in range(max(0,n+2-len(c)),min(n+1,len(c))) )
  val-=2*A*side*HN*conv
  return val
 den=HN**2*AD*D
 for n in range(DEG-1):
  c.append(rhs(n)//(den*(n+2)*(n+1)))
 if check:
  # equation residual in M_zz - rhs/h²; scale D² AD
  err=sum(abs(((c[n+2]*(n+2)*(n+1)*den) if n+2<len(c) else 0)-rhs(n)) for n in range(2*DEG+1))
  assert err*10**23 < AD*D**2, (side,j,err/(AD*D**2))
 return c
def mint(A,m0,k0,save=False):
 data=[]; ends=[]
 for side in [-1,1]:
  arr=[]
  c0=m0
  # to ensure initial derivative exactly same with sides define k0 here as k0/HN scaled!
  c1=-side*k0
  for j in range(N):
   c=mstep(A,side,j,c0,c1,save)
   c0=sum(c);c1=sum(n*x for n,x in enumerate(c))
   if save: arr.append(c)
  ends.append(c0-(L*D if side<0 else 0));data.append(arr)
 if save:
  assert max(abs(x) for x in ends)*10**23 < D
  return data
 return np.array([x/D for x in ends])
LQ=20
def revpoly(c):
 return [(-1)**j*sum(c[i]*math.comb(i,j) for i in range(j,len(c))) for j in range(len(c))]
def qcoefs(A,data):
 out=[]
 for side,arr in zip([-1,1],data):
  barr=[];karr=[]
  for j,m in enumerate(arr[:LQ*HN]):
   mr=revpoly(m)
   # beta stored scaled AD D
   beta=[2*A*side*x for x in mr]
   beta[0]+=AD*D*(j+1)//HN;beta[1]-=AD*D//HN
   # physical derivative z: dx reversed -1/HN
   k=[side*HN*n*mr[n] for n in range(1,len(mr))]
   barr.append(beta);karr.append(k)
  out.append((barr,karr))
 return out
from fractions import Fraction as Fr
def tailinit(A,Lam,side):
 p=Fr(1,2)-(Fr(A,AD) if side<0 else 0)
 c=Fr(Lam,LD)/p-1
 assert Fr(-71,100)<c<Fr(81,100) and (side>0 or c>Fr(-1,100))
 d=c*(c-1)/(2*p);e=(2*c-3)*d/(2*p)
 val=c/LQ+d/LQ**3+e/LQ**5
 return val.numerator*D//val.denominator
def qint(A,Lam,coefs,save=False):
 ends=[];out=[]
 for side,(barr,karr) in zip([-1,1],coefs):
  s0=tailinit(A,Lam,side); arr=[]
  # h negative; residual scaled LD AD D²
  for beta,k in zip(barr[::-1],karr[::-1]):
   c=[s0]
   def rhs(n):
    convb=sum(c[i]*beta[n-i] for i in range(max(0,n+1-len(beta)),min(n+1,len(c))))
    convs=sum(c[i]*c[n-i] for i in range(max(0,n+1-len(c)),min(n+1,len(c))))
    val=LD*(convb-AD*convs)
    if n<len(k): val-=2*LD*D*A*k[n]
    if n==0: val+=(LD-2*Lam)*AD*D**2
    return -val # sign h
   den=HN*LD*AD*D
   for n in range(DEG):
    c.append(rhs(n)//(den*(n+1)))
   s0=sum(c)
   if save:
    err=sum(abs(((c[n+1]*(n+1)*den) if n+1<len(c) else 0)-rhs(n)) for n in range(2*DEG+1))
    assert err*10**20 < LD*AD*D**2
    arr.append(c)
  ends.append(s0);out.append(arr[::-1])
 if save:return ends,out
 return sum(ends)/D
db={
0: [398942280401432677939946059934936368, 5000000000000000000000000000000001, 5000000000],
167000: [436448204647172275671154571065521908, 5111026674046595034939603490045098, 4002528709],
183000: [440536919244443400292190017104934260, 5123874426378891221881448106403949, 3886015237],
185000: [441054831040266796082023272478293862, 5125511916180277425031674550494055, 3871133409],
193000: [443141880868177733854084218630813528, 5132133382720642510533594754892904, 3810876605],
197000: [444194725820485682959362490472149280, 5135487495896646327064404732292470, 3780302564],
203000: [445785769731214940781912304479234981, 5140573657830417507240017132653823, 3733871338],
207000: [446854382388375608957334107223406674, 5144001493265635761300449570536001, 3702530072],
211000: [447929381090052885385364640445907237, 5147459297648980728746473644952589, 3670873366],
215000: [449010811041136767599350272927726311, 5150947346067397468701694365523148, 3638896171],
218000: [449826131225767018035053374695566673, 5153583395711869859552546091908423, 3614699851],
220000: [450371711347115414702287234485694676, 5155350356897422981117371212249798, 3598466110],
228000: [452570421073155581002349457455156882, 5162495706450928292015335430770643, 3532693692],
231000: [453401748336946844064210643083735018, 5165207491261130326087668688414708, 3507677435],
234000: [454236820433129662972716954382051104, 5167937063446988969961240183678679, 3482465702],
239000: [455636980863394482376305475385213303, 5172526193122040947960671147502318, 3440005122],
244000: [457047684631426696806497007443129457, 5177165590864996728236809104282291, 3396983232],
251000: [459040549972883437425970210163056860, 5183746215430131520421292962258781, 3335786897],
258000: [461054492582858640009939593690213328, 5190427871537775227086712728654635, 3273433269],
266000: [463382256258761310925112548110693165, 5198189616921800648472720382573621, 3200709619],
272000: [465146562404304407912087965480558966, 5204100048614802448626400462062651, 3145111670],
277000: [466629038869660672695849875091361589, 5209084479440533646312103130621392, 3098069677],
281000: [467823080450381224176236382390647761, 5213111054345015033023913790482294, 3059960694],
283000: [468422802833727954210264696892144102, 5215137440217787546712478602076254, 3040745124],
293650: [471646903449743964326750054546019982, 5226076475530743371555160276374328, 2936567807],
293655: [471648429278170868718179649270487911, 5226081670400701322796197274730540, 2936518150],
304000: [474830062508701471123549605498961791, 5236950205867338402676604016732928, 2832226234],
307000: [475762010965247164513736332348470605, 5240147358667012150252587393292838, 2801389817],
311000: [477011157190311061741246569619152021, 5244442225649069448736879731180273, 2759849013],
316000: [478583169388318722513619887175424457, 5249862548442940911853802351401021, 2707225932],
324000: [481123016830877402189435259911758451, 5258655641387006960242857464890572, 2621373597],
332000: [483693438251416145305364963041169198, 5267598590953712710367653201944070, 2533413436],
341000: [486622067478021139143860888527609012, 5277840475816287771505304037559031, 2431833246],
350000: [489590149217013150459035454451604544, 5288276035769017229974864077309765, 2327350370],
358000: [492261862551830464462065996735397278, 5297716154872114923625434487895606, 2231929998],
363000: [493947768520916359508306619953883636, 5303695170717058167276507568238026, 2171028129],
368000: [495646114469944087456537737737225530, 5309735201830213364679958592968569, 2109122819],
369000: [495987281117020524255538529154406594, 5310950548984813632995589189134366, 2096618989],
370000: [496328947871629374756867582623647432, 5312168346853911202387504058544337, 2084073726],
373000: [497356952751536597848406417862848528, 5315836459606920685663169284212728, 2046187111],
377000: [498734651097943921003750035530256643, 5320761672756152026272462715404391, 1995077904],
381000: [500120399939983459354541920108240900, 5325726268072088182877928567867189, 1943276576],
386000: [501863941040929695618157081757711250, 5331987493849022591294812731983611, 1877528371],
390000: [503267883372610202334955761574692406, 5337040922801031344612478693686727, 1824113384],
397000: [505744327598534384034341787833659036, 5345979603431905326263929284728143, 1728838531],
400000: [506813294676182538273397955089170882, 5349847568311569582933277586355971, 1687287204],
410000: [510409698904365400089610218268197555, 5362901520459144207633059993983021, 1545546607],
2936533: [471647910495218087096951828814828891, 5226079904138646272764052074508064, 2936535033],
2936535: [471647971528437792087151567621746964, 5226080111933846901160172613080866, 2936533047],
}
def certified(A):
 m,k,lam=db[A]; data=mint(A,m,k,save=True); co=qcoefs(A,data)
 for lr in [lam-10,lam,lam+10]:
  ends,out=qint(A,lr,co,save=True)
  for side,arr,(barr,_),sarr in zip([-1,1],data,co,out):
   for c in arr:
    kp=-side*HN*c[1]; rad=HN*sum(n*abs(x) for n,x in enumerate(c) if n>=2)
    assert (kp-rad)*100>=-D and (kp+rad)*100<=101*D
   # integer lower g scaled AD*D
   gs=[b[0]-sum(abs(x) for x in b[1:])-2*AD*(c[0]+sum(abs(x) for x in c[1:]))-AD*D//500 for b,c in zip(barr,sarr)]
   assert all(sum(abs(v) for v in c)*100<95*D for c in sarr)
   assert sum(gs[1203:])>8*HN*AD*D
   assert sum(min(g,0) for g in gs)>-2*HN*AD*D and sum(gs)>20*HN*AD*D
  if lr<lam:assert sum(ends)*10**9 < -D
  if lr>lam:assert sum(ends)*10**9 > D
  if lr==lam:qdata=out
 return data,qdata,lam

# float interval library; padding assumes IEEE
np.seterr(over="raise",invalid="raise",divide="raise")
tiny=1e-280
def down(x):return x-(abs(x)+tiny)*2**-43
def up(x):return x+(abs(x)+tiny)*2**-43
class I:
 def __init__(self,lo,hi=None):
  self.lo=lo;self.hi=lo if hi is None else hi
 def __add__(x,y):
  y=iv(y);return I(down(x.lo+y.lo),up(x.hi+y.hi))
 __radd__=__add__
 def __neg__(x):return I(-x.hi,-x.lo)
 def __sub__(x,y):return x+-iv(y)
 def __rsub__(x,y):return iv(y)+-x
 def __mul__(x,y):
  y=iv(y)
  a=x.lo*y.lo;b=x.lo*y.hi;c=x.hi*y.lo;d=x.hi*y.hi
  return I(down(np.minimum(np.minimum(a,b),np.minimum(c,d))),up(np.maximum(np.maximum(a,b),np.maximum(c,d))))
 __rmul__=__mul__
 def inv(x):
  assert np.all(x.lo*x.hi>0)
  return I(down(1/x.hi),up(1/x.lo))
 def __truediv__(x,y):return x*iv(y).inv()
 def __rtruediv__(x,y):return iv(y)*x.inv()
 def sq(x):
  a=x.lo*x.lo;b=x.hi*x.hi
  return I(np.where((x.lo<=0)&(x.hi>=0),0,down(np.minimum(a,b))),up(np.maximum(a,b)))
 def sqrt(x):
  assert np.all(x.lo>=0)
  return I(np.maximum(0,down(np.sqrt(x.lo))),up(np.sqrt(x.hi)))
 def cap(x,lo,hi):
  z=I(np.maximum(x.lo,lo),np.minimum(x.hi,hi));assert np.all(z.lo<=z.hi)
  return z
 def sel(x,s):return I(x.lo[s],x.hi[s])
 def maxabs(x):return np.maximum(abs(x.lo),abs(x.hi))
def iv(x):return x if isinstance(x,I) else I(x)
def rat(n,d):return I(n)/d
def choose(cond,x,y):
 x=iv(x);y=iv(y);return I(np.where(cond,x.lo,y.lo),np.where(cond,x.hi,y.hi))
# exp interval Taylor scale /256, require |x|<=100
def iexp(x):
 x=iv(x);assert np.all(x.maxabs()<=100)
 def fn(y,upper):
  # Taylor point using interval class
  w=I(y)/256; t=I(1.)
  for i in range(20,0,-1):t=1+w*t/i
  t=t+I(-2**-80,2**-80) # exp(.4) *.4^21/21! < 2^-80
  for i in range(8):t=t.sq()
  return t.hi if upper else t.lo
 return I(fn(x.lo,False),fn(x.hi,True))
# log for args [1/8,8] via repeated sqrt and atanh
def ilog(x):
 x=iv(x); assert np.all(x.lo>=.125) and np.all(x.hi<=8)
 for _ in range(3):x=x.sqrt()
 t=(x-1)/(x+1); ts=t.sq()
 p=rat(1,21)
 for i in range(19,0,-2):p=rat(1,i)+ts*p
 return 16*t*p+I(-2**-60,2**-60) # |t|<.14

ZF=50 # fine per HN =.0002
ZN=HN*ZF
def polytab(intarr,reverse=False):
 # table each cell at s=0,...,(ZF-1)/ZF; q reversed variable use 1-s
 co=np.array([[float(v/D) for v in c] for c in intarr])
 frac=np.arange(ZF)/ZF
 if reverse:frac=1-frac
 val=co[:,[-1]]
 for j in range(co.shape[1]-2,-1,-1): val=val*frac+co[:,[j]]
 # final endpoint omitted
 return val.reshape(-1)
class Profile:
 def __init__(self,A):
  data,qdata,lam=certified(A)
  self.A=A; self.a=rat(A,AD); self.af=A/AD
  self.lam=rat(lam,LD)+I(-1.01e-9,1.01e-9) # use pads later
  sides=[]
  for side,arr,qs in zip([-1,1],data,qdata):
   arr=arr[:LQ*HN]
   mx=[];ks=[];ls=[]
   for j,c in enumerate(arr):
    xx=c.copy();xx[0]+=side*j*D//HN;xx[1]+=side*D//HN;mx.append(xx)
    kk=[-side*HN*n*v for n,v in enumerate(c) if n>0];ks.append(kk)
    ll=[-v for v in kk];ll[0]+=D;ls.append(ll)
    # error bound for polytab rounding using sumabs
    for coeff in [c,xx,kk,ll]:
     # absolute eval error <=2e-14 sumabs; want 1e-14 when small; adaptive use? store separately
     assert sum(abs(x) for x in coeff)<30*D
   sides.append([polytab(arr),polytab(mx),polytab(ks),polytab(ls),side*polytab(qs,True)])
  # join z indices -(20-1/ZN)..(20-1/ZN)
  self.offset=LQ*ZN-1
  self.M,self.X,self.k,self.l,self.Q=[np.concatenate([x[:0:-1],y]) for x,y in zip(*sides)]
  # adaptive rounding error: for M,X,k,l near small sumabs maybe bound globally tighter by cell checking
  # compute arrays of err coefficients for each
  errsides=[]
  for side,arr in zip([-1,1],data):
   lists=[[] for _ in range(4)]
   for j,c in enumerate(arr[:LQ*HN]):
    xx=c.copy();xx[0]+=side*j*D//HN;xx[1]+=side*D//HN
    kk=[-side*HN*n*v for n,v in enumerate(c) if n>0]
    ll=[-v for v in kk];ll[0]+=D
    for li,coef in zip(lists,[c,xx,kk,ll]):
     li.append(sum(abs(x) for x in coef)/D*3e-14+3e-18)
   errsides.append([np.repeat(li,ZF) for li in lists])
  self.err=[np.concatenate([x[:0:-1],y]) for x,y in zip(*errsides)]
  # monotone search approximate k and l enforce using certified thresholds cumulative
  # exact k decreasing. guaranteed lower k and upper l prefix if node bound; take suffix max lower k?
  self.klo=np.maximum.accumulate((self.k-self.err[2])[::-1])[::-1] # suffix nodes prove lower
  self.khi=np.minimum.accumulate(self.k+self.err[2]) # prefix upper
  self.llo=np.maximum.accumulate(self.l-self.err[3]) # prefix lower (l inc)
  self.lhi=np.minimum.accumulate((self.l+self.err[3])[::-1])[::-1]
 def at(self,name,idx):
  vals=getattr(self,name)[idx]
  er=(2.95e-5 if name=='Q' else self.err[['M','X','k','l'].index(name)][idx])
  return I(vals)-0+I(-er,er) # padded
 def inverse_indices(self,k,l): # intervals arrays, find z interval outer nodes
  use=k.hi<=.6
  # lower idx: want node k_lower >= khi, choose last such; equivalently search on -klo
  l1=np.searchsorted(-self.klo,-k.hi,side='right')-1
  u1=np.searchsorted(-self.khi,-k.lo,side='left')
  # l increasing: lower node upper l <= llo
  l2=np.searchsorted(self.lhi,l.lo,side='right')-1
  u2=np.searchsorted(self.llo,l.hi,side='left')
  low=np.where(use,l1,l2);hi=np.where(use,u1,u2)
  assert np.all(low>=0) and np.all(hi<len(self.k))
  return low,hi
 def inverse(self,k,l):
  low,hi=self.inverse_indices(k,l)
  z=I((low-self.offset)/ZN-1e-12,(hi-self.offset)/ZN+1e-12)
  M=I(self.at('M',hi).lo,self.at('M',low).hi).cap(0,100)
  X=I(self.at('X',low).lo,self.at('X',hi).hi).cap(0,100)
  W=2*(.5-self.a)*M*l+X*k
  return z,M,X,W

from functools import lru_cache
profiles=lru_cache(maxsize=6)(Profile)

MGRID=2048
CUT=38
def grid():
 ed=np.concatenate([(1+np.arange(MGRID)/MGRID)*2.**(-i) for i in range(CUT,1,-1)]+[np.array([.5])])
 ed=np.concatenate([ed,1-ed[-2::-1]]) # exact dyadic? smallest subdivisions require CUT+11=49 <53 yes
 w=np.diff(ed); assert len(w)%2==0
 return I(ed[:-1],ed[1:]),I(1-ed[1:],1-ed[:-1]),w,len(w)//2
kg,lg,width,mid=grid()
def sums(x,rev=False): # intervals of weighted cumulative over cells inclusive full endpoint
 v=x*I(width)
 def side(a,sgn):
  if rev:a=a[::-1]
  cs=np.cumsum(a)
  err=np.cumsum(abs(a))*1e-9+tiny # n<1e6
  out=cs+sgn*err
  return out[::-1] if rev else out
 return I(side(v.lo,-1),side(v.hi,1))
def partialsum(x,rev=False): # integral to any point in current cell from relevant tail (within grid)
 v=x # allow either sign
 full=sums(v,rev)
 # remove optional part
 amt=x*I(0,width)
 return full-amt
def total(x):
 sm=sums(x);return I(sm.lo[-1],sm.hi[-1])
def anchored(x): # integral from .5 to any cell; avoid total cancellation
 left=choose(np.arange(mid*2)<mid,x,0)
 right=choose(np.arange(mid*2)>=mid,x,0)
 return choose(np.arange(mid*2)<mid,-partialsum(left,True),partialsum(right))
# RMQ for cell bounds on z table adjacent nodes
class Rmq:
 def __init__(self,g,values):
  self.g=g;self.los=[values.lo];self.his=[values.hi]
  while 2**len(self.los)<=len(values.lo):
   step=2**(len(self.los)-1)
   self.los.append(np.minimum(self.los[-1][:-step],self.los[-1][step:]))
   self.his.append(np.maximum(self.his[-1][:-step],self.his[-1][step:]))
 def query(self,z):
  low=np.floor((z.lo-1e-10)*ZN).astype(int)+self.g.offset
  high=np.floor((z.hi+1e-10)*ZN).astype(int)+self.g.offset
  assert np.all(low>=0) and np.all(high<len(self.los[0]))
  lens=high-low+1
  olo=np.empty_like(z.lo);ohi=np.empty_like(z.lo)
  for j,(los,his) in enumerate(zip(self.los,self.his)):
   sel=(lens>=2**j)&(lens<2**(j+1)); u=low[sel];v=high[sel]-(2**j-1)
   olo[sel]=np.minimum(los[u],los[v]);ohi[sel]=np.maximum(his[u],his[v])
  return I(olo,ohi)
def zcells(g):
 idx=np.arange(len(g.M)-1)
 def monot(name,inc):
  vals=g.at(name,np.arange(len(g.M)))
  return I(vals.lo[:-1] if inc else vals.lo[1:],vals.hi[1:] if inc else vals.hi[:-1])
 M=monot('M',False).cap(0,30);X=monot('X',True).cap(0,30)
 k=monot('k',False).cap(0,1);l=monot('l',True).cap(0,1)
 z=I((idx-g.offset)/ZN-1e-12,(idx+1-g.offset)/ZN+1e-12)
 beta=.5*z+g.a*M
 # bootstrap derivative bound: start Q node error, assume |Q|<=1 within length => |Q'|<25, enlarge, refine
 qnodes=g.at('Q',idx)
 # refine error central (cert proof)
 qnodes=qnodes.cap(g.Q[idx]-np.where(abs(z.lo)<12.01,5e-7,3e-5),g.Q[idx]+np.where(abs(z.lo)<12.01,5e-7,3e-5))
 Q=qnodes+I(-.006,.006)
 for _ in range(2):
  qp=2*beta*Q-Q.sq()+2*(.5-g.a*k-g.lam)
  Q=(qnodes+I(-1,1)*I(qp.maxabs())/ZN).cap(-1,0)
 qp=2*beta*Q-Q.sq()+2*(.5-g.a*k-g.lam)
 V=2*(.5-g.a)*M*l+X*k
 return locals()
def eloc_z(z,cp): # .25 power
 t=2-z
 t=I(np.maximum(t.lo,0),np.maximum(t.hi,0))
 den=1+t.sq()/2
 e=cp/2/den.sqrt().sqrt()
 ep=e*t/(4*den)
 epp=e*(5*t.sq()/(16*den.sq())-1/(4*den))
 epp=I(np.where(z.hi>=2,np.minimum(0,epp.lo),epp.lo),np.where(z.hi>=2,np.maximum(0,epp.hi),epp.hi))
 return e,ep,epp
def checkV(g,gapA,upper,cp):
 d=zcells(g);a=g.a
 p=.5-g.af
 sel=(d['z'].lo<6)&(d['z'].hi> -math.sqrt(13/p)-.01)
 z,M,k,l,W,B=[d[x].sel(sel) for x in ['z','M','k','l','V','beta']]
 t=k*l
 T=2*(.5-a)*M*l/W-2*k+2*a*(t/W).sq()
 if g.A==0: assert np.min(T.lo)>-1.12 and np.max(T.hi)<.13
 e,ep,epp=eloc_z(z,cp)
 fac=rat(1,1) if upper else -rat(11,10)
 R=1+fac*e*gapA
 Ar=(a+I(0,1)*gapA) if upper else (a-I(0,1)*gapA)
 core=.5*epp+(-B+a*t/W)*ep
 rhs=R.sq()*core+Ar*t/W*ep+R*(-(R+1)*(.5+a*T)*e+I(np.maximum(T.lo,0),np.maximum(T.hi,0))/(fac if upper else -fac))
 core0=fac*gapA*core-(.5+a*T)*R
 return np.max(rhs.hi),np.max(core0.hi)
CP=rat(26,100); ERR=rat(12,100)
def tail_mx_bounds(g): # cut primitives station first/last cell from inverse already include true cutoff
 z,m,x,v=g
 return m.hi[0],x.hi[-1]
class Env:
 def __init__(self,A,B,check=True):
  self.a=profiles(A);self.b=profiles(B); gap=self.gap=self.b.a-self.a.a
  if check:
   for g,u in [(self.a,1),(self.b,0)]:
    rr,co=checkV(g,gap,u,CP); assert rr<0 and co<-.02,(g.A,rr,co)
  self.va0=self.a.inverse(kg,lg); self.vb0=self.b.inverse(kg,lg)
  za,ma,xa,wa=self.va0; zb,mb,xb,wb=self.vb0
  rh=1+gap*eloc_z(za,CP)[0]; rl=1-rat(11,10)*gap*eloc_z(zb,CP)[0]
  self.vhi=rh*wa; self.vlo=rl*wb
  # envelopes primitive correction
  def adjusted(box,w,R,weight,rev):
   corr=1/R-1
   tail=I(-.16,.16)*gap*(box.sel(-1 if rev else 0))
   out=box+tail+partialsum(weight/w*corr,rev)
   direct=box/(1+I(-.1431,.1301)*gap)
   return out.cap(direct.lo,direct.hi)
  ml=adjusted(ma,wa,rh,kg,False);xl=adjusted(xa,wa,rh,lg,True)
  mh=adjusted(mb,wb,rl,kg,False);xh=adjusted(xb,wb,rl,lg,True)
  self.m=I(ml.lo,mh.hi);self.x=I(xl.lo,xh.hi)
  self.z=self.x-self.m
  # Q combo cell rmq
  da=zcells(self.a);db=zcells(self.b)
  delta=(db['M']-da['M']).cap(0,30) # may fail over tiny equal? intervals overlap
  upperq=da['Q']+self.a.a*delta
  lowerq=db['Q']-self.b.a*delta
  qhi=Rmq(self.a,upperq).query(self.z)+ERR*gap
  qlo=Rmq(self.b,lowerq).query(self.z)-ERR*gap
  qhi=I(np.minimum(0,qhi.lo),np.minimum(0,qhi.hi))
  self.sl=I(np.maximum(0,(-qhi/self.vhi).lo)); self.sh=I((-qlo/self.vlo).hi)
  # actual slope between sl.lo and sh.hi (qlo.lo negative)
  self.slope=I(self.sl.lo,self.sh.hi)
  self.logp=anchored(self.slope)
  if check:self.checkH(da,db,delta)
 def possiblek(self,z): # central away grid tails
  # high k where zhi >= z.lo; low k where zlo<= z.hi
  low=np.minimum.accumulate(self.z.lo) # enforce decreasing lower? prefix min remains lower
  hi=np.maximum.accumulate(self.z.hi[::-1])[::-1]
  i1=np.searchsorted(-low,-z.hi,side='left')
  i2=np.searchsorted(-hi,-z.lo,side='right')-1
  ok=(i1<len(kg.lo))&(i2>=0)
  i1=np.clip(i1,0,len(low)-1);i2=np.clip(i2,0,len(low)-1)
  klow=np.where(ok,kg.lo[i1],np.where(z.lo>0,0,kg.lo[-1]))
  khigh=np.where(ok,kg.hi[i2],np.where(z.lo>0,kg.hi[0],1))
  # even if ok at extreme include tail if i1==0 or i2==last
  klow=np.where(i1==0,0,klow);khigh=np.where(i2==len(low)-1,1,khigh)
  assert np.all(klow<=khigh)
  k=I(klow,khigh)
  l=1-k;l=I(np.maximum(0,l.lo),np.minimum(1,l.hi))
  # lower V via 32
  vl=rat(9,7)*k*l*l.sqrt()
  return k,vl.lo
 def checkH(self,da,db,delta):
  sel=(da['z'].lo<12)&(da['z'].hi>-12)
  z=da['z'].sel(sel)
  mi=I(da['M'].lo,db['M'].hi).sel(sel)
  ki,vl=self.possiblek(z)
  for d,g,upper in [(da,self.a,True),(db,self.b,False)]:
   q,qp,k,l=[d[x].sel(sel) for x in ['Q','qp','k','l']]
   damp=.5-qp-g.a*k-(self.gap*ki if upper else 0)
   br=mi*(qp-g.a*l)-ki*q+I(vl)
   if upper:br-=g.a*ki*delta.sel(sel)
   val=br+ERR*damp
   assert np.min(damp.lo)>.06,np.min(damp.lo)
   assert np.min(val.lo)>0,(g.A,np.min(val.lo),z.lo[np.argmin(val.lo)])
  # tail q @ z <= -5.5 numerical
  for d in [da,db]:
   s=(d['z'].lo<-5.5)
   assert np.min(d['Q'].lo[s])>-.28
  assert self.z.hi[-1]<-5.5 and self.z.lo[0]>4
  # delta tail <= X_b(-5.5) for negative z; numerical
  s=db['z'].lo<-5.5
  assert np.max(db['X'].hi[s])<.03
  assert self.logp.hi[-1]<5
  for g,arr in [(self.a,self.va0),(self.b,self.vb0)]:
   assert -arr[0].hi[-1]*2*(.5-g.af)>2 and arr[0].lo[0]>6
 def coarse(self,test,upper,detail=False):
  test=iv(test)
  slope=self.sl if upper else self.sh
  p=iexp(anchored(slope))
  v=self.vhi if upper else self.vlo
  integrand=p*(v-2*test*kg*kg*lg/v)
  # sign cumulative, allow tail error absolute 1e-7, only when away far endpoint
  if upper:
   # g negative near k=1 verify central tail say k>.95, analytical beyond cut
   far=kg.lo>.95
   if np.max(integrand.hi[far])>=0:return False
   cum=partialsum(integrand,True)
   val=np.max(cum.hi[~far])+1e-7
   # positive tail k->0 total error
   return val if detail else val<0
  else:
   far=kg.hi<.05
   if np.min(integrand.lo[far])<=0:return False
   cum=partialsum(integrand)
   val=np.min(cum.lo[~far])-1e-7
   return val if detail else val>0
@lru_cache(maxsize=3)
def refdata(g):
 z,m,x,v=g.inverse(kg,lg)
 q=Rmq(g,zcells(g)['Q']).query(z)
 slope=-q/v; lp=anchored(slope); p=iexp(lp)
 return z,m,x,v,q,slope,lp,p
def nonlineartest(env,test,upper,nt=8):
 test=iv(test);c=rat(3,10)
 gr=profiles(293650)
 z,mr,xr,vr,q,sr,lpr,pr=refdata(gr)
 tk=kg*lg
 base=total(pr/vr*(vr.sq()+c*(mr*xr-tk)-2*test*kg*tk))
 # actual slopes bounded, reference slope cell intervals
 eps=ilog(I(env.vlo.lo,env.vhi.hi)/vr)
 pis=env.slope-sr
 pi=anchored(pis) # tighter than logp - lpr
 out=base
 for it in range(nt):
  t=I((it)/nt,(it+1)/nt) # nt power2
  vf=iexp(t*eps);pf=iexp(t*pi)
  v=vr*vf;p=pr*pf
  # cap v,p actual envelope geometric? both ref within => yes slope convex for interpolated log
  v=v.cap(env.vlo.lo,env.vhi.hi)
  pbox=iexp(env.logp);p=p.cap(pbox.lo,pbox.hi)
  pv=p/v
  inv=1/vf-1
  def prim(ref,weight,rev):
   # omitted tail |1/vfac-1|<=3 from heat ratio <4
   tail=I(-3.01,3.01)*ref.sel(-1 if rev else 0)
   y=ref+tail+partialsum(weight/vr*inv,rev)
   static=(env.x if rev else env.m)
   return y.cap(static.lo,static.hi)
  m=prim(mr,kg,False);x=prim(xr,lg,True)
  signed=c*(m*x-tk)-2*test*kg*tk
  gp=pv*(v.sq()+signed)
  gv=pv*(v.sq()-signed)-c*kg/v*partialsum(pv*x,True)-c*lg/v*partialsum(pv*m)
  varv=total(gv*eps)
  # omitted M epsilon derivative tails <=6 M_ref(cut), heat/log ratio
  terror=6*(mr.hi[0]*total(c*pv*x).hi+xr.hi[-1]*total(c*pv*m).hi)
  assert terror<3e-7,terror # may pv*m divergent near tail but finite cut logs polynomial!
  coeff=choose(np.arange(len(width))<mid,-partialsum(gp),partialsum(gp,True))
  varp=total(coeff*pis)
  alt=total(gp*pi)
  varp=varp.cap(alt.lo,alt.hi)
  out+=(varv+varp)/nt
 val=out.hi+1e-6 if upper else out.lo-1e-6
 return val
def kquery(vals,i1,i2):
 # range extrema sparse using generic
 los=vals.lo;his=vals.hi
 n=i2-i1+1
 outlo=np.empty(len(n));outhi=np.empty(len(n))
 step=1
 while step<=max(n):
  sel=(n>=step)&(n<2*step)
  a=i1[sel];b=i2[sel]-step+1
  outlo[sel]=np.minimum(los[a],los[b]);outhi[sel]=np.maximum(his[a],his[b])
  los=np.minimum(los[:-step],los[step:]);his=np.maximum(his[:-step],his[step:])
  step*=2
 return I(outlo,outhi)
def localcheck():
 E=Env(293650,293655)
 a=I(E.a.a.lo,E.b.a.hi); d=rat(1,100)
 k,l=kg,lg; z=E.z;m=E.m;x=E.x
 assert -20<z.lo[-1] and z.hi[-1]<-10 and z.lo[0]>7 and z.hi[0]<20
 W=I(E.vlo.lo,E.vhi.hi)
 c0,c1,c2,c3=[rat(t,100) for t in [26,130,-186,325]]
 L=c0+k*(c1+k*(c2+k*c3));Lp=c1+2*k*c2+3*k*k*c3;Lpp=2*c2+6*k*c3
 assert np.max(L.hi)<3
 t=k*l
 T=2*(.5-a)*m*l/W-2*k+2*a*(t/W).sq()
 B=.5*z+a*m
 Wk=2*B-2*a*t/W
 R=1+I(-1,1)*d*L
 core=.5*W.sq()*Lpp+W*Wk*Lp
 rhs=R.sq()*core-(a+I(-1,1)*d)*t*Lp+R*(-(R+1)*(.5+a*T)*L+I(0,T.maxabs())) # absolute adverse source
 core0=I(-1,1)*d*core-(.5+a*T)*R
 assert np.max(rhs.hi)<0,np.max(rhs.hi)
 assert np.max(core0.hi)<0
 epsmax=L/(1-d*L)
 def lprim(weight,ref,rev):
  val=I(0,3.11)*ref.sel(-1 if rev else 0)+partialsum(weight/W*epsmax,rev)
  return I(0,np.minimum(val.hi,3.11*ref.hi))
 lm=lprim(k,m,False);lx=lprim(l,x,True);lz=lm+lx
 zi=z+I(-1,1)*d*lz
 da=zcells(E.a);db=zcells(E.b)
 delta=(db['M']-da['M']).cap(0,30)
 Qh=da['Q']+E.a.a*delta+ERR*E.gap;Ql=db['Q']-E.b.a*delta-ERR*E.gap
 Qsp=I(Ql.lo,np.minimum(0,Qh.hi))
 rq=Rmq(E.a,Qsp)
 Qstar=rq.query(z)
 shifts=lz*I(rq.query(zi).maxabs()+2*.011) # pad
 # H same z
 sel=(da['z'].lo<12)&(da['z'].hi>-12); zz=da['z'].sel(sel)
 def poss(box,useW=False):
  low=np.minimum.accumulate(box.lo);hi=np.maximum.accumulate(box.hi[::-1])[::-1]
  i1=np.searchsorted(-low,-zz.hi,side='left');i2=np.searchsorted(-hi,-zz.lo,side='right')-1
  i1=np.clip(i1,0,len(low)-1);i2=np.clip(i2,0,len(low)-1)
  kk=I(np.where(i1==0,0,k.lo[i1]),np.where(i2==len(low)-1,1,k.hi[i2]))
  assert np.all(i2>=i1)
  ww=kquery(W,i1,i2) # tail max <= .0001 below from (32)
  return kk,ww
 ks,_=poss(z);ki,wi=poss(zi)
 Ms=I(da['M'].lo[sel],db['M'].hi[sel]);ls=1-ks
 assert np.max((m*l/(.5-(a+d)*l)).hi)<.65
 dm=rat(65,100)
 Mi=Ms+I(-1,1)*dm*d
 qq=Qsp.sel(sel)
 qp=2*(.5*zz+a*Ms)*qq-qq.sq()+2*(.5-a*ks-a)
 damp=.5-qp-a*ks-d*ki
 # wi max tails use endpoint or simple .001 extra
 br=Mi*(qp-a*ls)-ki*qq+I(0,(1.031*wi.hi+.0001))+I(-1,1)*a*ki*dm*d
 assert np.min(damp.lo)>.1
 gain=a*dm+(a*(1-a)*dm+I(br.maxabs()))/damp
 assert np.max(gain.hi)<2,np.max(gain.hi)
 # functional
 cadd=rat(3,10)
 sr=-Qstar/W;lpr=anchored(sr);pr=iexp(lpr)
 zc=z.sel(mid)
 assert zc.maxabs()<2
 pimax=2*I((z-zc).maxabs())+shifts
 out=0
 for it in range(4):
  td=d*(it+1)/4
  pvf=iexp(I(-1,1)*td*pimax)
  vf=iexp(I(-1,1)*td*epsmax)
  v=W*vf;p=pr*pvf;pv=p/v
  mt=(m+I(-1,1)*td*lm).cap(.968*m.lo,1.033*m.hi)
  xt=(x+I(-1,1)*td*lx).cap(.968*x.lo,1.033*x.hi)
  signed=cadd*(mt*xt-t)-2*a*k*t
  gp=pv*(v.sq()+signed)
  gv=pv*(v.sq()-signed)-cadd*k/v*partialsum(pv*xt,True)-cadd*l/v*partialsum(pv*mt)
  cv=total(I(gv.maxabs())*epsmax)
  terror=4*(m.hi[0]*total(cadd*pv*xt).hi+x.hi[-1]*total(cadd*pv*mt).hi)
  assert terror<3e-7
  coeff=choose(np.arange(len(width))<mid,partialsum(gp),partialsum(gp,True))
  cp=total(I(coeff.maxabs())*2/W)+total(I(gp.maxabs())*shifts)
  out+=(cv+cp)/4
 den=2*total(pv*k*t)
 assert den.lo>0
 return (out.hi+2e-6)/den.lo
def checkA():
 n=4096
 ind=np.arange(1,n-1)
 k=I(ind/n,(ind+1)/n);l=1-k;t=k*l
 low=rat(9,7)*k*l*l.sqrt(); high=rat(6,5)*(t*t.sqrt()).sqrt()
 b=2*(rat(6,10)-k); B=-rat(4,10)*t
 # max endpoint term convex when b>=0, decreasing if b<0
 edge=(b*low.sq()+2*B*low).hi
 edge=np.maximum(edge,(b*high.sq()+2*B*high).hi)
 ws=[low+(high-low)*I(i/128,(i+1)/128) for i in range(128)]
 mini=100.
 for aa in range(419,1024): # covers [.41,1]
  a=I(aa/1024,(aa+1)/1024)
  K=12*a*k-5*a-1-b.sq()+rat(8,10)*(1-2*k)
  pos=(6*a.sq()-rat(16,100))*t.sq()
  mins=np.minimum.reduce([(2*w.sq()+K*w-2*b*B+pos/w).lo for w in ws])
  # last cell crude integrand >=-9; first likewise and endpoint |...| upper<.001
  assert np.max(abs(mins))<9
  cs=np.cumsum(mins[::-1])[::-1]/n-2e-9
  bound=cs-np.maximum(0,mins)/n-9/n-edge
  val=min(np.min(bound[ind<=2560]),cs[0]-18/n-.001)
  mini=min(mini,val)
  assert val>0,(aa,val,cs[0])
 return mini

def driver():
 global AD
 assert checkA()>0
 A,B=(0, 410000)
 for an,bn in [(167000, 410000), (183000, 410000), (185000, 400000), (193000, 397000), (197000, 390000), (203000, 386000), (207000, 381000), (211000, 377000), (215000, 373000), (218000, 369000), (220000, 370000)]:
  e=Env(A,B); print(A,B,flush=True)
  if an>A: assert e.coarse(rat(an,AD),False)
  if bn<B: assert e.coarse(rat(bn,AD),True)
  A,B=an,bn
 A,B=(220000, 370000)
 for an,bn in [(228000, 370000), (231000, 368000), (234000, 363000), (239000, 358000), (244000, 350000), (251000, 341000), (258000, 332000), (266000, 324000), (272000, 316000), (277000, 311000), (281000, 307000), (283000, 304000), (285000, 303000)]:
  e=Env(A,B); print(A,B,flush=True)
  if an>A: assert nonlineartest(e,rat(an,AD),False)>0
  if bn<B: assert nonlineartest(e,rat(bn,AD),True)<0
  A,B=an,bn

 assert localcheck()<.95
 AD=10**7
 for A in [2936533,2936535]:
  _,_,lam=certified(A)
  assert (lam-10>A*1000 if A==2936533 else lam+10<A*1000)
driver()
\end{lstlisting}
\end{document}
